% Twenty Chapters, One Robot Joint Node
% Build:  python build.py            (PDF + HTML)
%         python build.py --check sections/j07.tex   (compile one chapter alone)
%         python lint.py                             (house rules)
%         python crosscheck.py                       (book-level consistency)
%         python build.py --drift ../EmbeddedFirmware_NucleoH7_Top20/build.py
\documentclass[11pt,a4paper]{article}
\usepackage[utf8]{inputenc}
\usepackage[T1]{fontenc}
\usepackage{lmodern}
\usepackage{microtype}
\usepackage[margin=21mm,top=24mm,bottom=24mm]{geometry}
\usepackage{xcolor}
\usepackage{graphicx}
\usepackage{booktabs}
\usepackage{tabularx}
\usepackage{array}
\usepackage{enumitem}
\usepackage{float}
\usepackage{caption}
\usepackage{listings}
\usepackage{tcolorbox}
\usepackage{tikz}
\usepackage[european]{circuitikz}
% Shared TikZ and circuitikz setup, used by main.tex and by the standalone
% figure wrapper that build.py generates for the HTML/SVG output.
\usetikzlibrary{arrows.meta,positioning,calc,shapes.geometric,shapes.misc,shapes.symbols,fit,backgrounds,chains,decorations.pathreplacing,decorations.markings,patterns,shadows,matrix}

% ---------------------------------------------------------------- colours
\definecolor{hwcol}{RGB}{255,236,179}   % hardware, boards, sensors
\definecolor{kercol}{RGB}{197,222,255}  % kernel space
\definecolor{usrcol}{RGB}{205,240,205}  % user space
\definecolor{extcol}{RGB}{228,228,228}  % external systems, cloud, PC
\definecolor{fwcol}{RGB}{245,210,235}   % MCU firmware
\definecolor{notecol}{RGB}{255,250,205} % notes
\definecolor{linecol}{RGB}{60,60,60}

% ---------------------------------------------------------------- block diagrams
\tikzset{
  every picture/.style={font=\sffamily\small, line width=0.6pt, >=Latex},
  blk/.style={draw=linecol, rounded corners=2pt, align=center, minimum height=8mm,
              inner sep=4pt, text width=28mm, fill=white},
  hw/.style={blk, fill=hwcol},
  kern/.style={blk, fill=kercol},
  user/.style={blk, fill=usrcol},
  ext/.style={blk, fill=extcol},
  fw/.style={blk, fill=fwcol},
  wide/.style={text width=44mm},
  narrow/.style={text width=18mm},
  layer/.style={draw=linecol, dashed, rounded corners=3pt, inner sep=6pt, fill=none},
  layerlabel/.style={anchor=north west, font=\sffamily\scriptsize\bfseries, inner sep=2pt},
  flow/.style={->, thick, draw=linecol},
  bus/.style={<->, thick, draw=linecol},
  lbl/.style={font=\sffamily\scriptsize, fill=white, inner sep=1.5pt, midway},
  irq/.style={->, thick, draw=red!70!black, dashed},
  % ---------------------------------------------------------------- UML
  umlclass/.style={draw=linecol, fill=white, align=left, inner sep=0pt, font=\sffamily\scriptsize},
  lifeline/.style={draw=linecol, dashed},
  actor/.style={draw=linecol, fill=extcol, rounded corners=2pt, align=center, minimum height=7mm, inner sep=4pt, font=\sffamily\small},
  msg/.style={->, draw=linecol, thick, font=\sffamily\scriptsize},
  reply/.style={->, draw=linecol, dashed, font=\sffamily\scriptsize},
  activation/.style={draw=linecol, fill=kercol, inner sep=0pt, minimum width=2.4mm},
  state/.style={draw=linecol, rounded corners=3mm, fill=usrcol, align=center, minimum height=8mm, inner sep=5pt},
  initial/.style={circle, fill=linecol, inner sep=0pt, minimum size=2.4mm},
  final/.style={circle, draw=linecol, fill=white, inner sep=0pt, minimum size=3.4mm,
                path picture={\fill[linecol] (path picture bounding box.center) circle (1mm);}},
  trans/.style={->, draw=linecol, thick, font=\sffamily\scriptsize},
  umlnote/.style={draw=linecol, fill=notecol, align=left, font=\sffamily\scriptsize, inner sep=3pt},
  % ---------------------------------------------------------------- bench art
  board/.style={draw=linecol, fill=green!25!black!60, rounded corners=1.5mm, minimum height=20mm, minimum width=32mm,
                text=white, font=\sffamily\bfseries\small, align=center},
  hat/.style={board, fill=blue!40!black!60},
  breadboard/.style={draw=linecol, fill=white, minimum height=12mm, minimum width=40mm, font=\sffamily\scriptsize},
  pinrow/.style={draw=linecol, fill=black, minimum height=1.6mm, inner sep=0pt},
  cable/.style={line width=1.2pt},
  usbcable/.style={line width=1.4pt, draw=black!70},
  ledsym/.style={circle, draw=linecol, minimum size=3mm, inner sep=0pt},
}

% ================================================================ MCU figures
% Appended for the firmware volume. Same colour vocabulary as above:
% hwcol = silicon, kercol = RTOS and kernel, usrcol = application task,
% fwcol = firmware module, extcol = host and debugger.
% Solid fills only: figures are built latex -> dvi -> dvisvgm, a path on which
% PDF transparency does not survive.
\usetikzlibrary{shapes.multipart,decorations.pathmorphing}

\tikzset{
  % ------------------------------------------------------------ memory map
  mmap/.style={draw=linecol, anchor=south west, align=center, inner sep=2pt,
               minimum width=34mm, text width=34mm, font=\sffamily\scriptsize},
  mmflash/.style={mmap, fill=hwcol},
  mmram/.style={mmap, fill=kercol},
  mmperiph/.style={mmap, fill=fwcol},
  mmext/.style={mmap, fill=extcol},
  mmres/.style={mmap, fill=black!6},
  mmaddr/.style={font=\scriptsize\ttfamily, anchor=east, inner sep=2pt},
  mmgap/.style={draw=linecol, dashed, line width=0.4pt},
  mpu/.style={draw=red!70!black, line width=0.8pt, dashed, rounded corners=1pt, fill=none},
  % ------------------------------------------------------------ register bit field
  regcell/.style={draw=linecol, fill=white, anchor=south west, inner sep=0pt,
                  minimum height=6mm, align=center, font=\sffamily\scriptsize},
  regfield/.style={regcell, fill=kercol},
  regrsvd/.style={regcell, fill=black!6},
  regro/.style={regcell, fill=extcol},
  regset/.style={regcell, fill=hwcol},
  regbit/.style={font=\sffamily\tiny, anchor=south, inner sep=1pt},
  regname/.style={font=\sffamily\bfseries\small, anchor=east, inner sep=3pt},
  regreset/.style={font=\tiny\ttfamily, anchor=north, inner sep=1pt},
  % ------------------------------------------------------------ clock tree
  osc/.style={blk, fill=hwcol, text width=20mm, minimum height=7mm},
  pll/.style={blk, fill=fwcol, text width=22mm, minimum height=9mm},
  clkmux/.style={draw=linecol, fill=white, trapezium, trapezium angle=70,
                 trapezium stretches body, shape border rotate=270,
                 minimum width=9mm, minimum height=8mm, inner sep=1pt, font=\sffamily\tiny},
  clkdiv/.style={draw=linecol, fill=kercol, minimum width=12mm, minimum height=6mm,
                 inner sep=1pt, align=center, font=\sffamily\tiny},
  clkgate/.style={draw=linecol, fill=white, circle, inner sep=0pt, minimum size=3.2mm,
                  font=\sffamily\tiny},
  clkdom/.style={blk, fill=usrcol, text width=24mm, minimum height=6mm},
  clkline/.style={-Latex, draw=linecol, line width=0.7pt},
  clkfreq/.style={font=\sffamily\tiny, fill=white, inner sep=1pt},
  % ------------------------------------------------------------ NVIC and vectors
  vecslot/.style={draw=linecol, fill=white, anchor=south west, align=left,
                  minimum width=44mm, minimum height=5mm, inner xsep=3pt,
                  font=\sffamily\scriptsize},
  vecsys/.style={vecslot, fill=kercol},
  vecirq/.style={vecslot, fill=hwcol},
  vecfree/.style={vecslot, fill=black!6},
  vecnum/.style={font=\tiny\ttfamily, anchor=east, inner sep=2pt},
  prio/.style={draw=linecol, fill=fwcol, minimum width=26mm, minimum height=6mm,
               inner sep=2pt, align=center, font=\sffamily\scriptsize},
  prioaxis/.style={-Latex, draw=linecol, line width=0.7pt},
  preempt/.style={-Latex, draw=red!70!black, thick},
  tailchain/.style={-Latex, draw=linecol, dashed, thick},
  % ------------------------------------------------------------ DMA and cache
  master/.style={blk, fill=fwcol, text width=22mm, minimum height=7mm},
  busmatrix/.style={draw=linecol, fill=extcol, minimum width=72mm, minimum height=7mm,
                    align=center, font=\sffamily\small},
  memblk/.style={blk, fill=kercol, text width=24mm, minimum height=10mm},
  cache/.style={blk, fill=usrcol, text width=20mm, minimum height=8mm},
  cacheword/.style={draw=linecol, fill=white, anchor=south west, inner sep=0pt,
                    minimum width=8mm, minimum height=4.5mm, font=\sffamily\tiny},
  cwclean/.style={cacheword, fill=usrcol},
  cwdirty/.style={cacheword, fill=orange!35},
  cwstale/.style={cacheword, fill=red!25},
  dmaflow/.style={-Latex, thick, draw=blue!55!black},
  cpuflow/.style={-Latex, thick, draw=linecol},
  incoherent/.style={draw=red!70!black, line width=0.8pt, decorate,
                     decoration={zigzag, amplitude=0.6mm, segment length=2mm}},
  % ------------------------------------------------------------ timing diagram
  signal/.style={draw=linecol, line width=0.8pt, line cap=round, line join=round},
  signalb/.style={signal, draw=blue!55!black},
  signame/.style={font=\sffamily\scriptsize, anchor=east, inner sep=3pt},
  tick/.style={draw=linecol!45, dotted, line width=0.4pt},
  tspan/.style={Latex-Latex, draw=linecol, line width=0.5pt},
  tlabel/.style={font=\sffamily\tiny, fill=white, inner sep=1pt},
  sample/.style={circle, fill=red!70!black, inner sep=0pt, minimum size=1.4mm},
  hiz/.style={draw=linecol, dashed, line width=0.6pt},
  % ------------------------------------------------------------ RTOS schedule
  taskrun/.style={draw=linecol, fill=usrcol, anchor=south west, inner sep=1pt,
                  minimum height=5mm, align=center, font=\sffamily\tiny},
  taskready/.style={taskrun, fill=hwcol},
  taskblock/.style={taskrun, fill=black!8},
  taskisr/.style={taskrun, fill=fwcol},
  tasklane/.style={font=\sffamily\scriptsize, anchor=east, inner sep=3pt},
  taxis/.style={-Latex, draw=linecol, line width=0.6pt},
  release/.style={-Latex, draw=linecol, line width=0.6pt},
  deadline/.style={draw=red!70!black, dashed, line width=0.6pt},
  rtstate/.style={state, fill=usrcol, text width=18mm, minimum height=7mm,
                  font=\sffamily\scriptsize},
}

% ---------------------------------------------------------------- helper macros
% Lengths are plain numbers in millimetres.

% \memregion{<y>}{<height>}{<style>}{<name and size>}{<start address>}
\newcommand{\memregion}[5]{%
  \node[#3, minimum height=#2mm] at (0mm,#1mm) {#4};%
  \node[mmaddr] at (-0.5mm,#1mm) {#5};}
% \memgap{<y>}  a "not to scale" break across the column
\newcommand{\memgap}[1]{\draw[mmgap] (-2mm,#1mm) -- (36mm,#1mm);}

% Register bit fields. Bit \regwidth-1 is leftmost, as in every reference manual.
\newcommand{\regwidth}{32}
\newcommand{\regunit}{4}
% \bitfield{<hi>}{<lo>}{<style>}{<label>}
\newcommand{\bitfield}[4]{%
  \node[#3, minimum width={(#1-#2+1)*\regunit mm},
        text width={(#1-#2+1)*\regunit mm-2pt}]
    at ({(\regwidth-1-#1)*\regunit mm},0mm) {#4};}
\newcommand{\bitscale}{%
  \foreach \b in {0,...,\the\numexpr\regwidth-1\relax}
    \node[regbit] at ({(\regwidth-0.5-\b)*\regunit mm},6mm) {\b};}

% \vecentry{<row from bottom>}{<style>}{<exception number>}{<handler name>}
\newcommand{\vecentry}[4]{%
  \node[#2] at (0mm,{#1*5mm}) {#4};%
  \node[vecnum] at (-1mm,{#1*5mm+2.5mm}) {#3};}

% \cacheline{<x>}{<y>}{<style>}{<label>}
\newcommand{\cacheline}[4]{\node[#3] at (#1mm,#2mm) {#4};}

% \clocktrain{<y>}{<x start>}{<cycles>}{<period>}  50 percent duty, 4 mm high
\newcommand{\clocktrain}[4]{%
  \draw[signal] (#2mm,#1mm) \foreach \i in {1,...,#3}
      { -- ++(0,4mm) -- ++({0.5*#4mm},0) -- ++(0,-4mm) -- ++({0.5*#4mm},0) };}
% \busval{<y centre>}{<x>}{<width>}{<text>}  the usual bus-value lozenge
\newcommand{\busval}[4]{%
  \draw[signal] (#2mm,#1mm) -- ++(1mm,2mm) -- ++({#3mm-2mm},0) -- ++(1mm,-2mm)
                            -- ++(-1mm,-2mm) -- ++({-#3mm+2mm},0) -- cycle;%
  \node[font=\sffamily\tiny] at ({#2mm+0.5*#3mm},#1mm) {#4};}

% \slice{<lane y>}{<t start>}{<duration>}{<style>}{<label>}
\newcommand{\slice}[5]{\node[#4, minimum width=#3mm] at (#2mm,#1mm) {#5};}
% \lane{<lane y>}{<name>}
\newcommand{\lane}[2]{\node[tasklane] at (0mm,{#1mm+2.5mm}) {#2};}
% \deadlinemark{<t>}{<height>}
\newcommand{\deadlinemark}[2]{\draw[deadline] (#1mm,-1mm) -- (#1mm,#2mm);}

% ================================================================ robotics figures
% Appended for the joint-node volume. Same colour vocabulary as above:
% hwcol = silicon and hardware, kercol = RTOS and kernel, usrcol = application,
% fwcol = firmware module, extcol = host and external.
% Solid fills only: figures go latex -> dvi -> dvisvgm, where PDF transparency
% does not survive.

\tikzset{
  % ------------------------------------------------------------ field bus
  busline/.style={draw=linecol, line width=1.1pt},
  busstub/.style={draw=linecol, line width=0.7pt},
  busnode/.style={draw=linecol, rounded corners=2pt, fill=fwcol, align=center,
                  minimum height=9mm, text width=24mm, inner sep=3pt,
                  font=\sffamily\scriptsize},
  busmaster/.style={busnode, fill=extcol, text width=26mm},
  buslistener/.style={busnode, fill=usrcol},
  busterm/.style={draw=linecol, fill=hwcol, minimum width=4mm, minimum height=7mm,
                  inner sep=0pt, font=\sffamily\tiny},
  busoff/.style={busnode, fill=red!18},
  buslabel/.style={font=\sffamily\tiny, fill=white, inner sep=1.5pt},
  % ------------------------------------------------------------ frame layout
  fieldbox/.style={draw=linecol, fill=white, anchor=south west, inner sep=0pt,
                   minimum height=6mm, align=center, font=\sffamily\scriptsize},
  fid/.style={fieldbox, fill=hwcol},      % identifier and control
  fdata/.style={fieldbox, fill=kercol},   % payload
  fcrc/.style={fieldbox, fill=usrcol},    % checksum and acknowledge
  fstuff/.style={fieldbox, fill=black!8}, % framing the controller adds
  fbyte/.style={font=\sffamily\tiny, anchor=south, inner sep=1pt},
  % ------------------------------------------------------------ control loop
  ctrlblk/.style={draw=linecol, fill=kercol, rounded corners=2pt, align=center,
                  minimum height=9mm, text width=20mm, inner sep=3pt,
                  font=\sffamily\scriptsize},
  plantblk/.style={ctrlblk, fill=hwcol},
  sensorblk/.style={ctrlblk, fill=usrcol},
  sumj/.style={draw=linecol, fill=white, circle, inner sep=0pt, minimum size=5mm,
               font=\sffamily\tiny},
  ctrlline/.style={-Latex, draw=linecol, line width=0.7pt},
  fbline/.style={-Latex, draw=linecol, line width=0.7pt, dashed},
  ctrllabel/.style={font=\sffamily\tiny, fill=white, inner sep=1.5pt},
  % ------------------------------------------------------------ the joint itself
  linkbody/.style={draw=linecol, fill=extcol, rounded corners=1.5mm,
                   minimum height=7mm, minimum width=30mm},
  jointpin/.style={draw=linecol, fill=hwcol, circle, minimum size=7mm, inner sep=0pt},
  jointarc/.style={draw=linecol, dashed, -Latex},
  % ------------------------------------------------------------ safe states
  runstate/.style={state, fill=usrcol, text width=20mm, minimum height=8mm,
                   font=\sffamily\scriptsize},
  degraded/.style={runstate, fill=hwcol},
  safestate/.style={runstate, fill=red!20, draw=red!70!black, line width=0.9pt},
  latched/.style={safestate, double, double distance=1pt},
  faultedge/.style={->, draw=red!70!black, thick, font=\sffamily\tiny},
  recoveredge/.style={->, draw=linecol, thick, dashed, font=\sffamily\tiny},
  % ------------------------------------------------------------ the honesty rule
  % One convention across the whole volume, established in chapter 1: a dashed
  % outline is a model standing in for hardware that is not on this bench, and a
  % dotted grey block is hardware that is absent and explained rather than built.
  % No measurement taken through a dashed block measures anything physical.
  modelled/.style={blk, fill=white, dashed, line width=0.8pt},
  absentblk/.style={blk, fill=black!6, dotted, line width=0.8pt, text=black!55},
  laterblk/.style={blk, fill=white, draw=linecol!45, text=black!60},
  chlabel/.style={font=\sffamily\tiny, text=black!55, inner sep=1pt},
}

% ---------------------------------------------------------------- helper macros
% Lengths are plain numbers in millimetres, as in the other style groups.

% \busdrop{<x>}{<node name>}  a stub from the trunk down to a node
\newcommand{\busdrop}[2]{\draw[busstub] (#1mm,0mm) -- (#2.north);}

% \framefield{<x>}{<width>}{<style>}{<label>}   one field of a frame layout
\newcommand{\framefield}[4]{%
  \node[#3, minimum width=#2mm, text width={#2mm-2pt}] at (#1mm,0mm) {#4};}
% \framebits{<x>}{<width>}{<text>}  the bit or byte count above a field
\newcommand{\framebits}[3]{\node[fbyte] at ({#1mm+0.5*#2mm},6mm) {#3};}

% \sumnode{<name>}{<x>}{<y>}{<top sign>}{<bottom sign>}
\newcommand{\sumnode}[5]{%
  \node[sumj] (#1) at (#2mm,#3mm) {};%
  \node[font=\sffamily\tiny] at ({#2mm-2.6mm},{#3mm+2.6mm}) {#4};%
  \node[font=\sffamily\tiny] at ({#2mm-2.6mm},{#3mm-2.6mm}) {#5};}

\usepackage{fancyhdr}
\usepackage[hidelinks,pdftitle={Twenty Chapters, One Robot Joint Node},pdfauthor={Joseph Ambrose Pagaran}]{hyperref}

% ---------------------------------------------------------------- unicode helpers
\DeclareUnicodeCharacter{2192}{\ensuremath{\rightarrow}}
\DeclareUnicodeCharacter{2190}{\ensuremath{\leftarrow}}
\DeclareUnicodeCharacter{2264}{\ensuremath{\leq}}
\DeclareUnicodeCharacter{2265}{\ensuremath{\geq}}
\DeclareUnicodeCharacter{2248}{\ensuremath{\approx}}
\DeclareUnicodeCharacter{03A9}{\ensuremath{\Omega}}
\DeclareUnicodeCharacter{2126}{\ensuremath{\Omega}}
\DeclareUnicodeCharacter{2026}{\ldots}
\DeclareUnicodeCharacter{00D7}{\ensuremath{\times}}
\DeclareUnicodeCharacter{00B1}{\ensuremath{\pm}}
\DeclareUnicodeCharacter{00B5}{\textmu{}}
\DeclareUnicodeCharacter{2013}{-}
\DeclareUnicodeCharacter{2014}{-}

% ---------------------------------------------------------------- layout
% Sources lists carry long addresses. Let them break at more places than the
% default, which is only the hyphen, or they run into the margin.
\setlength{\Urlmuskip}{0mu plus 1mu}
\def\UrlBreaks{\do\/\do\-\do\.\do\_\do\?\do\&\do\=\do\#\do\:\do\~}
\setlength{\emergencystretch}{3em}
\hyphenpenalty=800
\tolerance=1500
\setlength{\parskip}{4pt}
\setlength{\parindent}{0pt}
\captionsetup{font=small,labelfont=bf,skip=4pt}
\setlist{itemsep=2pt,topsep=3pt,parsep=0pt}
\pagestyle{fancy}
\fancyhf{}
\fancyhead[L]{\sffamily\small Twenty Chapters, One Robot Joint Node}
\fancyhead[R]{\sffamily\small\nouppercase{\benchmark}}
% \benchmark is a short running label: "Project NN" inside a project, the
% section name elsewhere. Long project titles would otherwise run into the
% left-hand header.
\newcommand{\benchmark}{\leftmark}
\fancyfoot[C]{\sffamily\small\thepage}
\renewcommand{\headrulewidth}{0.3pt}

% ---------------------------------------------------------------- listings
\lstset{basicstyle=\ttfamily\footnotesize,columns=fixed,basewidth=0.48em,keepspaces=true,
  breaklines=true,frame=single,rulecolor=\color{gray!60},backgroundcolor=\color{gray!6},
  xleftmargin=3pt,framexleftmargin=3pt,showstringspaces=false,tabsize=4,
  commentstyle=\color{green!40!black},keywordstyle=\color{blue!60!black}\bfseries,
  stringstyle=\color{red!50!black},extendedchars=false,aboveskip=6pt,belowskip=6pt}
\lstnewenvironment{asciiart}{\lstset{language={},frame=none,backgroundcolor=\color{white},
  breaklines=false,basicstyle=\ttfamily\scriptsize,keywordstyle=,commentstyle=,stringstyle=}}{}
\lstnewenvironment{ccode}{\lstset{language=C}}{}
\lstnewenvironment{shellcode}{\lstset{language=bash}}{}
\lstnewenvironment{pycode}{\lstset{language=Python}}{}
\lstnewenvironment{makecode}{\lstset{language=make}}{}
\lstnewenvironment{plaincode}{\lstset{language={}}}{}
\lstnewenvironment{dtscode}{\lstset{language={},morecomment=[l]{//},morecomment=[s]{/*}{*/}}}{}
\lstnewenvironment{yamlcode}{\lstset{language={},morecomment=[l]{\#}}}{}
\lstnewenvironment{cppcode}{\lstset{language=C++}}{}
\lstnewenvironment{rustcode}{\lstset{language={},morecomment=[l]{//},morecomment=[s]{/*}{*/},
  morestring=[b]",morekeywords={fn,let,mut,const,static,struct,enum,impl,trait,pub,use,mod,
  match,if,else,loop,while,for,in,return,unsafe,async,await,dyn,where,type,as,ref,move,Self}}}{}
\lstnewenvironment{asmcode}{\lstset{language=[x86masm]Assembler,morecomment=[l]{//},
  morecomment=[l]{;},morekeywords={ldr,str,ldrex,strex,dmb,dsb,isb,wfi,wfe,bx,blx,svc,
  cpsid,cpsie,msr,mrs,push,pop,bl,cbz,cbnz,it,ite}}}{}
\lstnewenvironment{ldcode}{\lstset{language={},morecomment=[s]{/*}{*/},
  morekeywords={MEMORY,SECTIONS,ENTRY,PROVIDE,KEEP,ALIGN,ORIGIN,LENGTH,AT,PHDRS}}}{}

% ---------------------------------------------------------------- authoring macros
% \project{NN}{Title}{Primary board}{Theme}
% \project{NN}{Title}{Target board}{Theme}
% The macro keeps its name because build.py's converter keys off it; it prints
% "Chapter". The label prefix sec:c must match DOC["label_pfx"] in build.py.
\newcommand{\project}[4]{\clearpage\section[#2]{#2}\label{sec:j#1}%
  \renewcommand{\benchmark}{Chapter #1}%
  {\color{black!60}\sffamily\small What the node gains: #3 \quad|\quad Theme: #4}\par\medskip}
% key facts box: \begin{keyfacts} \item[Boards] ... \end{keyfacts}
\definecolor{volaccent}{HTML}{1D6B5E}
\newenvironment{keyfacts}{\begin{tcolorbox}[colback=volaccent!5,colframe=volaccent!80!black,
  title=Key facts,fonttitle=\sffamily\bfseries,left=2pt,right=2pt,top=2pt,bottom=2pt]
  \begin{description}[style=sameline,leftmargin=34mm,labelwidth=32mm,font=\sffamily\bfseries\small,itemsep=1pt]}
  {\end{description}\end{tcolorbox}}
% note box: \begin{note}{Title} ... \end{note}
\newtcolorbox{note}[1]{colback=yellow!8,colframe=orange!75!black,title=#1,fonttitle=\sffamily\bfseries,
  left=3pt,right=3pt,top=2pt,bottom=2pt}
% numbered steps
\newenvironment{steps}{\begin{enumerate}[label=\textbf{Step \arabic*.},leftmargin=*,itemsep=5pt]}{\end{enumerate}}
% \diagram{figure-name}{Caption}: inputs figures/<name>.tex, scaled down only when too wide
\newsavebox{\diagbox}
\newcommand{\diagram}[2]{\begin{figure}[H]\centering
  \sbox{\diagbox}{\input{figures/#1.tex}}%
  \ifdim\wd\diagbox>\linewidth\resizebox{\linewidth}{!}{\usebox{\diagbox}}\else\usebox{\diagbox}\fi
  \caption{#2}\label{fig:#1}\end{figure}}

% \pubdate{April 2010}: a publication date that the publisher gives to month
% precision only. The house rule is that every date this volume states carries
% weekday, day, month and year, and that rule cannot apply to a day a publisher
% never published. Writing one would be inventing a fact, which is the one thing
% this volume refuses to do. crosscheck.py recognises this macro and exempts it,
% so every month-and-year that is NOT wrapped is still reported as a defect.
\newcommand{\pubdate}[1]{#1}

\title{\sffamily\bfseries Twenty Chapters, One Robot Joint Node\\[2pt]\large built on the bench, measured, and honest about its gaps}
\author{Joseph Ambrose Pagaran}
\date{Monday 21 September 2026}

\begin{document}
\maketitle
\thispagestyle{empty}
\section*{About this document}
\addcontentsline{toc}{section}{About this document}

This is a build plan for one thing, taken apart into twenty chapters: a joint
node for a robot arm, built on a bench that does not have a robot arm on it. A
microcontroller board, a single-board computer, three sensor shields, a bus
transceiver that costs less than a lunch, and no motor. By the end the node
keeps a one kilohertz period and can prove it, shares one clock across its
sensors and its bus, speaks a modern field bus properly enough to survive a
device older than itself, joins a robot framework as a participant, closes a
control loop and reports the one number a joint is judged by, stops safely and
latches, can be updated over the wire it already has, and is checked every
night by a rig that turns a build red.

Every chapter has the same shape: why it exists, the prior art and what is taken
from it, what the node gains, the parts it uses, a system architecture figure,
the peripheral configuration, the wiring, a memory and timing budget, a software
design in UML, a data-flow sketch, a repository layout, numbered steps with real
commands and real code, measurable acceptance criteria, the variants it touches,
pitfalls, best practices, stretch goals, a sourced roadmap, the evidence to
publish, and its sources. Twenty chapters, five figures each, one hundred
figures in all.

The twenty were chosen so that no two teach the same thing, and so that together
they cover what a robotics firmware engineer is expected to have done: a board
support package handed over as data rather than as code, a control period proved
without an oscilloscope, one time base shared between sensors and a bus,
quadrature and inertial sensing, a motor driver written in full for a motor that
is not there, a field bus from bit timing upward, network management and the
failure an older device causes, a middleware port nobody has published for this
part, standard message types with one field honestly left empty, a controller
with real anti-windup, a costed refusal, a latched safe state, a bootloader that
stays addressable, and a rig that proves all of it again tonight.

\section*{The rule that runs through this book}
\addcontentsline{toc}{section}{The rule that runs through this book}

\begin{note}{What is real and what is only drawn}
Three things a robotics firmware role asks for cannot be demonstrated on this
bench. \textbf{There is no motor and no drive stage. There is no force or torque
sensor. There is no real-time Ethernet fieldbus, and there cannot be, because
this part has no Ethernet controller and this silicon vendor sells no subdevice
controller at all.} The book names all three in its first chapter, designs
around them deliberately, and marks every figure accordingly: \textbf{a solid
outline is hardware that is present, a dashed outline is a model standing in for
hardware that is not, and a dotted grey block is hardware that is absent and
explained rather than built.} No measurement taken through a dashed block is a
measurement of anything physical, and the chapter that takes one says so in the
sentence that reports it.
\end{note}

The rule has teeth in three places. Every budget table carries a
\textbf{Measured} column, and it reads \emph{not measured} until something has
been measured; a number that came from a model is labelled modelled, and a
number that came from arithmetic is labelled computed. Every chapter whose
subject is partly absent ends by saying what is real here, what is modelled, and
what an honest answer sounds like when somebody asks whether you have done it.
And where the research behind a chapter could not confirm something, the chapter
says that too, in the sentence where the claim would otherwise have gone.

That last habit produced more of this book than expected. Several of the most
useful paragraphs here are of the form \emph{this is not what everybody says}:
that this part's flash word is sixteen bytes and not thirty-two, that its
backup registers live in a different peripheral from its better known sibling's,
that a widely repeated middleware footprint comes from a commercial blog rather
than from the project, that a popular bootloader changed licence in 2025, that
the collaborative robot technical specification is not withdrawn, and that a
vendor's safety architecture claim is routinely quoted with its meaning
reversed. None of those is clever. All of them come from opening the document.

\section*{The bench, and what it does not have}
\addcontentsline{toc}{section}{The bench, and what it does not have}

The node is a NUCLEO-H7A3ZI-Q carrying an STM32H7A3ZIT6Q: a Cortex-M7 at 280
megahertz with two megabytes of flash and an on-board debug circuit that is a
probe, a serial port and a drag-and-drop disk at once. The motion master is a
Raspberry Pi 4. Between them is one twisted pair and a transceiver.

\begin{note}{This part is not the popular member of its family}
This part is documented by its own reference manual, and the member of the
family that most of the internet writes about is documented by a different one.
The clock tree, the power configuration, the memory map and the flash geometry
all differ. \textbf{The flash word here is sixteen bytes and the sector is eight
kilobytes; on the popular part they are thirty-two bytes and one hundred and
twenty-eight kilobytes}, and at least two defects in widely used open projects
have been reported against exactly that difference. No linker script, clock
configuration or memory map is inherited from a sibling in this volume without
being checked against this part's own manual, and chapters that cite material
written for the other part say so in the sentence that cites it.
\end{note}

What the bench does not have shapes this volume more than what it has. There is
\textbf{no oscilloscope, no logic analyser and no external debug probe}, so
every timing claim is made with the part's own cycle counter and timers, and the
crash forensics of chapter 20 are built from inside the silicon. There is
\textbf{no soldering iron}, which rules out three otherwise reasonable chapters
and is said plainly where it does. There is \textbf{no motor, no drive stage, no
brake and no torque sensor}. And there is \textbf{no Ethernet controller on this
part}, which chapter 4 establishes not from a product page but from the absence
of two interrupt vector slots in the vendor's own header.

Four things the bench does have are worth naming because chapters depend on
them: an eight by eight multizone ranging shield, which becomes the workspace
sensor of chapter 18; two inertial and environmental shields, used one at a time
and never stacked; the part's second bus controller instance, which makes the
gateway of chapter 13 possible without another board; and the part's own ROM
bootloader, which turns out to speak the same bus the joint uses and therefore
opens chapter 19 with a correction rather than a plan.

\section*{How to read and how to work}
\addcontentsline{toc}{section}{How to read and how to work}

\diagram{front_map}{The dependencies that are real, and nothing else. Chapter 1
gates everything. Chapter 2's period and chapter 3's clock are assumed by every
chapter after them. Chapter 9's bus gates chapters 10 to 14 and 19. Chapter 5's
encoder and chapter 7's plant model are what make chapter 16's loop close at
all. Chapters 17 and 20 depend on almost everything and are the two that step
back and look at the whole. The rest of the order is a preference rather than a
requirement.}

\begin{itemize}
\item \textbf{Order.} Chapter 1 gates everything, because a node with no
identity and no bring-up is a board. After that, follow the figure above and
read the rest in whatever order the next application asks for.
\item \textbf{Effort.} Each chapter's key-facts box gives a difficulty from 1 to
5 and an estimate in evenings of about four hours. The total is roughly sixty
evenings, and two of the twenty chapters need no hardware at all.
\item \textbf{Evidence first.} Every chapter ends with the artefacts to publish:
a repository whose README opens with a figure, a console log, a measurement
table generated from data rather than typed, and a short note on what the
measurement found. A chapter is finished when the evidence exists, not when the
code runs once.
\item \textbf{Licences are checked before code is written around them, not
after.} Several chapters change which library they start from on the strength of
a licence file, and one records a package that ships three contradictory licence
statements in one archive. Where a project's permission could not be confirmed,
the chapter says so and does not vendor it.
\item \textbf{Nothing here is a product.} One axis, no motor, no brake, no
assessment and no certificate. Chapter 18 carries a page titled what this node
is not, and it is the most credible page in the repository.
\end{itemize}

\section*{What a reader can honestly claim at the end}
\addcontentsline{toc}{section}{What a reader can honestly claim at the end}

A volume that ends without saying this leaves its reader to guess, so it is said
here. Somebody who builds all twenty chapters can say, without stretching, that
they have written a bus stack from bit timing upward and proved the bit rate
against the controller's own timestamps; that they have designed a joint
protocol and shown the arithmetic that rejected the obvious version of it; that
they have ported an embedded middleware to a part its project does not support
and written a transport for it that does not exist anywhere; that they have
built a controller with a filtered derivative and two anti-windup methods
measured against each other, one of which has no published open implementation
in this language; that they have designed a latched safe state against a
published clause list and written down what the node is not; that they have
written a bootloader that stays addressable on the joint's own wire; and that
they have built a rig that asserts on motion quality in continuous integration,
which the research for this volume could find no published example of.

They cannot say they have commissioned a robot arm, closed a current loop on a
real motor, measured a real joint torque, or shipped anything with a safety
integrity claim. Those sentences are not in this book, and the chapters where
they would have gone say why instead.

\clearpage
\tableofcontents
\clearpage

\part{The node comes up}
\project{01}{What a joint node is, and the bench that stands in for one}{An identity and a bring-up}{Architecture, the stand-in bench, what is real and what is modelled}

\begin{keyfacts}
\item[Adds to the node] An identity: a name, a node number, a version it can report, a state it shows, and a written boundary between what is real and what is modelled
\item[Peripherals] RCC and PWR for the clock tree, GPIO for three indicator LEDs and the user button, USART3 for the console
\item[Depends on] Nothing. This is the first chapter
\item[Real or modelled] Real: the board, the clock, the console, the indicators. Modelled: nothing yet, and the point of the chapter is to say so in writing before anything is modelled
\item[Difficulty] 2 of 5
\item[Effort] Two evenings of about three hours
\item[Deliverable] A repository that builds from a clean clone in two commands, a board that names itself on the console at boot, an indicator convention every later chapter reuses, and a committed table of what this bench can and cannot do
\end{keyfacts}

\subsection*{Why this chapter}

A joint node is the smallest part of a robot that has an opinion of its own. It
owns one axis. It closes a loop on that axis at a fixed period. It reports where
the axis is and accepts commands about where it should be, on a bus it shares
with other joints and with a motion master. And it stops, on its own authority,
when something is wrong. Take any one of those away and what is left is a motor
driver with a connector.

This volume builds one of those, once, over twenty chapters. Chapter 20 has a
node that runs a control loop at a period it can prove, carries sensors, speaks
on a field bus, publishes itself to middleware, enters a safe state when the
workspace is violated, can be updated without being touched, and is tested by a
rig that fails the build when tracking degrades. This chapter has a board on a
desk. What it contributes is the thing that makes the other nineteen honest: a
written, committed statement of which parts of that node are real hardware and
which are software standing in for hardware that is not on this bench.

That statement has to be first, not last. It is easy to write nineteen chapters
of firmware and then discover that the demonstration rests on a number that came
from a model. The discipline this volume applies is the opposite: name the
boundary in chapter 1, print it on the console at every boot, and let every
later chapter say which side of it the chapter's result sits on.

\begin{note}{What this node is not and why that is the interesting part}
There is no motor on this bench, no drive stage, no current sensing, no torque
sensor, and no real-time Ethernet slave controller. Those absences are not
worked around quietly. Chapter 7 builds the whole firmware side of an actuator,
timer configuration, complementary outputs, dead time and a fault input, and
runs it against a plant model, so the control loop is real and the torque is
not. Chapter 8 explains what a force and torque signal is for, what a real one
costs, and what its stand-in gets wrong. Chapter 17 explains what a field bus
slave needs in silicon and what would have to be bought. A reader who is asked
in an interview whether they have done these things gets a better answer from
this book than a yes.
\end{note}

\subsection*{Prior art and what to reuse}

\begin{table}[H]\centering\small
\begin{tabular}{p{34mm}p{46mm}p{38mm}p{20mm}}
\toprule
Source & What it gives & What it does not & Licence \\
\midrule
The robot control framework's mock hardware component & The published precedent for a node that exists to exercise a whole framework with no hardware attached. It mirrors commands into states, integrates them so the mock moves instead of teleporting, and can inject a constant following error on purpose & It is a mock, not a plant model, and it asserts nothing about quality. The following-error parameter is a fault injector, not a metric & Apache-2.0 \\
The upstream operating system's board files for this board & The board's facts in machine-readable form, and the fastest way to settle a pin question without opening a manual & It carries one inherited artefact that claims a peripheral this part does not have. See the pitfalls & Apache-2.0 \\
The field-bus-to-framework bridge package & The closest existing prior art to this whole volume: a bus-connected joint that speaks a standard motion profile, presented to the robot framework as a hardware interface & It is the master side. This volume builds the device side of the same contract & Apache-2.0 \\
The open simulator's platform files for this silicon family & A second bench that needs no hardware at all, with the bus controller modelled at the right addresses & No platform file for this exact part, and no motion metrics of any kind & MIT \\
The silicon vendor's device headers and hardware abstraction driver & The authoritative, fetchable, redistributable statement of this part's memory map, supervision thresholds and flash geometry & No board layer and no opinion about architecture & Apache-2.0 and BSD-3-Clause \\
\bottomrule
\end{tabular}
\caption{Prior art for chapter 1. The third row is the one to read before writing anything: this volume's node is the device side of a contract somebody else already maintains the master side of, and saying so is more useful than inventing a protocol in private.}
\end{table}

What is left to write is the identity and the boundary. No published source
tells you how to state, in a form a reader can check, which half of your
demonstration is real. That is this chapter's contribution and it is
deliberately small: a document, a boot banner, and a convention for the three
lights on the board.

\subsection*{What the node gains}

Before this chapter there is a development board in a box and a claim that it
could become a joint node. After it there is a node with a name, a number, a
version string it can report, a state it shows without a debugger attached, and
a committed table saying what it can and cannot honestly demonstrate. It cannot
move anything, sense anything or talk to anything yet. It can, however, tell you
what it is, which is the first thing every later chapter needs.

\diagram{j01_data}{The boundary this chapter draws, row by row. Three of the
eight things a joint in a robot has are real on this bench, three are stood in
for, and two are absent and explained. This figure is a file in the repository
and a line the board prints at every boot, which is what stops it from quietly
becoming untrue somewhere around chapter 12. The position row started this
volume in the real column and was moved here by chapter 5, which is the file
doing its job.}

\subsection*{Parts from the inventory}

\begin{table}[H]\centering\small
\begin{tabular}{p{40mm}p{66mm}p{40mm}}
\toprule
Part & Role & Interface \\
\midrule
NUCLEO-H7A3ZI-Q & The node itself, for all twenty chapters. Its on-board probe is also the programmer, the console and a drag-and-drop disk & Micro USB to the host \\
Raspberry Pi 4 & The motion master from chapter 10 onwards. In this chapter it is only the second machine the clean-clone build is proven on & Ethernet or Wi-Fi to the host network \\
A USB data cable & Power, programming and console on one lead & Micro USB \\
\bottomrule
\end{tabular}
\caption{Inventory items used in chapter 1. Nothing is wired, nothing is bought, and the two sensor shields stay in their boxes until chapter 6.}
\end{table}

\subsection*{System architecture}

\diagram{j01_arch}{The whole system, with the chapter that builds each piece. Solid blocks exist by the end of this chapter. Everything else is named here so the reader can see the shape of the machine before it is built, and so no later chapter has to explain where it fits.}

Read that figure as a plan rather than a status report. The node is a stack:
silicon at the bottom, a board support layer above it, then the periodic control
loop that everything else hangs off, then sensing, then the bus, then the
middleware that presents the node to a robot, and beside all of it the safety
and update machinery that has to keep working when the rest does not. The motion
master is a separate machine, and the line between them is a two-wire bus.

Two things on that figure are not hardware at all. The plant is a model, and the
torque estimate is derived. They are drawn with a dashed outline throughout this
volume, and the rule is simple: if a block is dashed, no measurement taken
through it is a measurement of anything physical.

\subsection*{Peripheral configuration}

\begin{table}[H]\centering\small
\begin{tabular}{p{22mm}p{26mm}p{24mm}p{30mm}p{30mm}}
\toprule
Peripheral & Mode & Clock source & Pins and function & Interrupt and transfers \\
\midrule
RCC and PWR & Bypass input, then phase-locked loop & 8 MHz from the on-board probe & None & None \\
GPIOB, GPIOE & Push-pull output & Peripheral bus & PB0, PE1, PB14 as node state & None \\
GPIOC & Input with pull-up & Peripheral bus & PC13, user button & None in this chapter. It becomes the local stop input in chapter 18 \\
USART3 & Asynchronous, 115200 8N1 & Peripheral bus & Bridged to the probe's virtual serial port & Polled in this chapter \\
TIM6 & Reserved, not configured & Peripheral bus & None & Claimed here so no later chapter takes it. It becomes the control period in chapter 2 \\
FDCAN1 & Reserved, not configured & Reserved & None until a transceiver exists & Claimed here for chapter 9 \\
\bottomrule
\end{tabular}
\caption{Peripheral configuration for chapter 1. The last two rows configure nothing and exist so that the peripheral budget is a book-level decision rather than a series of local ones. A node that discovers in chapter 16 that its timer was taken in chapter 5 has to redo both.}
\end{table}

\subsection*{Wiring}

\diagram{j01_wiring}{The bench as it stands today, and the two parts that have to arrive before chapter 9 can run. The board has no bus transceiver, which is the single electrical gap between this node and a network, and it costs less than a lunch.}

Nothing is wired in this chapter. The figure is here because it is the honest
picture of the bench, and because the gap it shows is the reason chapter 9 opens
with a purchase rather than with code.

\subsection*{Memory and timing budget}

Every chapter in this volume carries this table, and the numbers accumulate. The
baseline below is what chapter 1 costs, so that every later chapter can report
what it added rather than what it totals.

\begin{table}[H]\centering\small
\begin{tabular}{p{44mm}p{28mm}p{28mm}p{30mm}}
\toprule
Quantity & Budget & Measured & Margin \\
\midrule
Flash, this chapter & 16 kB of 2048 kB & not measured & not measured \\
Static memory, this chapter & 4 kB of 1024 kB & not measured & not measured \\
Flash, whole node at chapter 20 & 512 kB & not measured & not measured \\
Static memory, whole node at chapter 20 & 256 kB & not measured & not measured \\
Reset to boot banner & 200 ms & not measured & not measured \\
\bottomrule
\end{tabular}
\caption{The budget table. Nothing is written in the measured column until it has been measured, and the instrument that measures the last row is built in chapter 2. The two whole-node rows are the budget this volume holds itself to, stated now so that a chapter which blows it has to say so.}
\end{table}

The two whole-node rows deserve a sentence, because they are the reason several
later decisions go the way they do. Two megabytes of flash and about 1.4
megabytes of memory is a great deal for a joint node, and the temptation is to
treat neither as scarce. The budget above is set at a quarter of each so that
the node would still fit on a part a manufacturer would actually put in a joint,
which is usually a much smaller one. When chapter 14 reports what the middleware
costs, that budget is what makes the number mean something.

\subsection*{Firmware design (UML)}

\diagram{j01_uml}{The node as modules, with the chapter that introduces each one and the direction of dependency. Nothing points upward. The rule that makes the volume buildable is that a module may use what is below it and must not know what is above it, which is why the bus layer has no idea a controller exists.}

Two design decisions are made here and never revisited.

The first is that dependencies point downward only. The control loop does not
know that a bus exists; it produces a state structure and consumes a command
structure, and something above it decides where those come from and go. That is
what makes chapter 16 possible without rewriting chapter 2, and it is what makes
the whole node testable on a host machine, which chapter 20 depends on
completely.

The second is that the node's state is one thing, in one place, owned by one
module. Every chapter adds fields to it and no chapter adds a second copy. The
single most common way a machine like this becomes untrustworthy is two modules
each holding a slightly different idea of where the axis is.

\subsection*{Data flow (ASCII)}

\begin{asciiart}
  chapter 1 today                          chapter 20, the same tree grown
  +--------------------------+             +--------------------------------------+
  | reset                    |             | reset                                |
  |   |                      |             |   |                                  |
  |   v                      |             |   v                                  |
  | clocks: 8 MHz -> 280 MHz |             | bootloader: checksum, stay or jump    |
  |   |                      |             |   |                                  |
  |   v                      |             |   v                                  |
  | self-check: clock, ID    |             | self-check: clock, ID, flash, memory  |
  |   |                      |             |   |                                  |
  |   v                      |             |   v                                  |
  | identity: name, node id  |             | identity + version + fleet reporting  |
  |   |                      |             |   |                                  |
  |   v                      |             |   v                                  |
  | banner on the console    |  becomes    | 1 kHz loop: sense, control, actuate   |
  |   |                      | =========>  |   |            |            |         |
  |   v                      |             |   |            v            v         |
  | indicators: boot/ready   |             |   |        state frame   safe state   |
  |   |                      |             |   |            |            |         |
  |   v                      |             |   v            v            v         |
  | idle loop, blink         |             | indicators   the bus     latched stop |
  +--------------------------+             +--------------------------------------+
\end{asciiart}

\subsection*{Repository layout}

This section shows the whole tree in every chapter, with the current chapter's
additions marked, so that a reader can watch the node grow. Today almost all of
it is a promise.

\begin{plaincode}
joint-node/
  CMakeLists.txt                      # + this chapter
  cmake/arm-none-eabi.cmake           # + this chapter
  ld/stm32h7a3zi.ld                   # + this chapter
  doc/
    node-identity.md                  # + this chapter: what this node is
    real-or-modelled.md               # + this chapter: the honesty table
    footprint.csv                     # + this chapter: one row per chapter
  src/
    node/
      identity.c  identity.h          # + this chapter
      state.h                         # + this chapter: the one state struct
      indicator.c indicator.h         # + this chapter
      selfcheck.c selfcheck.h         # + this chapter
    bsp/
      startup.c  system.c  uart.c     # + this chapter
      board.h                         # + this chapter, completed in chapter 4
    control/                          # chapter 2 onwards: period, loop, plant
    sense/                            # chapters 5 and 6: encoder, inertial unit
    act/                              # chapter 7: outputs, dead time, fault in
    bus/                              # chapters 9 to 13: controller, protocol
    mw/                               # chapters 14 and 15: middleware client
    safety/                           # chapter 18: safe states, latching
    update/                           # chapter 19: bootloader, versioning
  host/                               # chapter 10 onwards: the motion master
  test/                               # chapter 20: unit tests and the rig
  tools/
    footprint.py                      # + this chapter
  README.md                           # + this chapter
\end{plaincode}

\subsection*{Steps}

\begin{steps}
\item \textbf{Write down what the node is, before any code.} A joint node that
cannot answer basic questions about itself is not a node. Create
\texttt{doc/node-identity.md} and fill in every field. The units matter more
than they look: they are the same units the standard joint state message uses,
and picking them now means chapter 15 is a mapping rather than a conversion.
\begin{plaincode}
name           joint-node-1
node id        3            (settable at runtime from chapter 12)
axis           one revolute joint, continuous rotation
position       radians, zero at the index mark, increasing anticlockwise
velocity       radians per second
effort         newton metres, positive in the direction of increasing position
control period 1 ms, from chapter 2
bus            CAN-FD, 500 kbit/s nominal, 2 Mbit/s data, from chapter 9
\end{plaincode}

\item \textbf{Write down what is real and what is not.} This is
\texttt{doc/real-or-modelled.md}, and it is the file this volume is built
around. Three columns: the quantity, its source, and one sentence on what the
stand-in gets wrong. Fill in the rows you already know and leave the rest for
the chapters that create them.
\begin{plaincode}
quantity        source at chapter 20        what the stand-in gets wrong
--------------  --------------------------  ----------------------------------
clock rate      real, measured              nothing
position        HALF: real decode path,     no eccentricity, no missed edges,
                generated source            no phase error. See chapter 5
velocity        derived from position       quantisation at low speed
torque command  real PWM, no drive stage    no current, so no real torque
torque measured MODELLED from the plant     it is the model's own output
plant response  MODELLED, second order      no friction model, no backlash
bus traffic     real, on a real transceiver nothing
fieldbus slave  ABSENT, explained not built no timing claim is made at all
\end{plaincode}

\item \textbf{Bring the board up.} The toolchain, the hand-written startup code,
the linker script and the clock tree are the subject of the first chapter of the
sibling firmware volume and are not repeated here. What this volume adds is one
rule: the linker script and the clock configuration are checked against this
part's own reference manual, never inherited from its better-known sibling. Two
commands from a clean clone, and no vendor development environment.
\begin{shellcode}
cmake -B build -DCMAKE_TOOLCHAIN_FILE=cmake/arm-none-eabi.cmake
cmake --build build -j && probe-rs run --chip STM32H7A3ZITx build/firmware.elf
\end{shellcode}

\item \textbf{Give the node an identity in firmware.} The version comes from the
build, not from a constant somebody forgets to bump. The node number is a
compile-time default that chapter 12 makes settable, which is why it is a
variable and not a macro.
\begin{ccode}
/* identity.h: the node's answer to "what are you?" */
typedef struct {
    const char *name;        /* "joint-node-1" */
    const char *version;     /* from git describe, injected by the build */
    const char *built;       /* ISO date, injected by the build */
    uint8_t     node_id;     /* default here, settable from chapter 12 */
    uint32_t    uid[3];      /* the die's own unique identifier */
} node_identity_t;

const node_identity_t *node_identity(void);
void node_identity_print(void);   /* the boot banner, see step 6 */
\end{ccode}
The last field is worth taking now. Every part of this family carries a unique
identifier in read-only memory, and a fleet of nodes that all report the same
compiled-in name is a fleet you cannot tell apart. Chapter 19 uses it as the
key for update reporting.

\item \textbf{Assert the facts that later chapters depend on.} A self-check at
startup that refuses to continue is worth more than a comment. The clock rate is
the one that matters most, because a board running at reset speed looks like a
board with a slow control loop, and that is a whole evening lost in chapter 2.
\begin{ccode}
/* selfcheck.c: fail loudly at boot rather than quietly at 1 kHz */
bool node_selfcheck(void)
{
    if (SystemCoreClock != 280000000u)        return fail("clock not 280 MHz");
    if (rcc_pll_source() != RCC_SRC_HSE_BYP)  return fail("not on the probe clock");
    if (flash_latency() < FLASH_LATENCY_MIN)  return fail("flash wait states too low");
    return true;
}
\end{ccode}
A failed check lights the fault indicator, prints the reason, and stops. It does
not continue with a warning, because a node that continues after failing its own
check is a node whose later numbers mean nothing.

\item \textbf{Make the boot banner say what is modelled.} This is the small idea
that keeps the whole volume honest, and it costs about ten lines. The banner
prints the identity and then, on one line, every quantity that is not real on
this bench. From chapter 7 onwards that line is not empty, and anyone watching a
demonstration can see it.
\begin{ccode}
void node_identity_print(void)
{
    const node_identity_t *id = node_identity();
    printf("\r\n%s  node %u  %s  built %s\r\n", id->name, id->node_id,
           id->version, id->built);
    printf("uid %08lX%08lX%08lX\r\n", id->uid[0], id->uid[1], id->uid[2]);
    printf("modelled: %s\r\n", NODE_MODELLED_LIST);  /* "none" today */
}
\end{ccode}

\item \textbf{Establish the indicator convention.} Three LEDs, used the same way
for twenty chapters, so that a reader who has seen one photograph of the board
can read the next. Define it now, in a header, and never open-code a pin write
anywhere else.
\begin{ccode}
typedef enum {
    NODE_BOOT,      /* green, slow blink: running self-check                */
    NODE_READY,     /* green, on solid: self-check passed, no bus yet       */
    NODE_ACTIVE,    /* green, fast blink: control loop running (chapter 2)  */
    NODE_DEGRADED,  /* yellow: running, but something is not what it should */
    NODE_SAFE,      /* red, on solid: in the safe state (chapter 18)        */
    NODE_FAULT      /* red, fast blink: latched, needs a reset (chapter 18) */
} node_indication_t;

void indicator_set(node_indication_t s);
\end{ccode}
Two of those six states are unreachable until chapter 18. They are defined here
anyway, because a convention added later is a convention two chapters already
worked around.

\item \textbf{Record the footprint baseline.} Every later chapter appends one
row, so that the volume can show what each layer of a joint node actually costs.
This is one of the few numbers in embedded work that is easy to measure, easy to
compare and almost never published.
\begin{shellcode}
python tools/footprint.py build/firmware.elf --chapter 1 --append doc/footprint.csv
cat doc/footprint.csv
# chapter,text,rodata,data,bss,flash_total,ram_total
# 1,11284,1360,112,3184,12756,3296
\end{shellcode}

\item \textbf{Prove the clean clone on the second machine.} Clone into a fresh
folder on the Raspberry Pi that will become the motion master, and build there.
It has no vendor development environment and never will, which is exactly why it
is the right machine to prove the build on.
\end{steps}

\subsection*{Build, flash and debug}

\diagram{j01_timing}{What happens between reset and the banner, and where the intervals are still unknown. The instrument that measures them is built in the next chapter, so every interval here is marked as unmeasured rather than guessed.}

Three ways to get the image onto the board, in the order to try them: the
flashing tool above, the open debug server, or copying the binary onto the disk
the probe presents to the host. The last needs no tools at all and is the
fastest way to prove that a board is alive.

\begin{shellcode}
openocd -f interface/stlink.cfg -f target/stm32h7x.cfg \
  -c "program build/firmware.elf verify reset exit"
\end{shellcode}

There is a fourth way on this part, and it matters later rather than now: the
read-only-memory bootloader in the silicon itself can be reached over the bus
this node will use, with no bootloader of your own. Chapter 19 opens on it.

\begin{note}{The one sentence to remember from the part trap}
Most published material for this silicon family was written for its better-known
sibling, whose clock tree, power configuration and memory map all differ.
Inherited linker scripts and clock code do not run slowly here, they produce a
board that does not boot. The exception, and it is a useful one, is the bus
peripheral: a comparison of the vendor's own device headers shows it is the same
licensed controller at the same addresses with the same register layout on both
parts, so material written for the sibling's bus is directly reusable. Chapter 9
shows the comparison, because a reader can repeat it in ten minutes and should.
\end{note}

\subsection*{Verification and acceptance criteria}

\begin{itemize}
\item A clean clone builds and flashes in two commands on a machine with no
vendor development environment, proven on the Raspberry Pi and not only on the
authoring machine.
\item The boot banner prints the node name, number, version, build date and die
identifier, and a line naming every modelled quantity, which today reads
\texttt{none}.
\item Deliberately miscompiling the clock configuration, for instance by
selecting the internal oscillator, stops the board at the self-check with a
named reason on the console and the fault indicator lit. It does not boot
anyway.
\item The indicator convention is defined in one header, and no source file
outside \texttt{indicator.c} writes to an LED pin. A search proves it.
\item \texttt{doc/real-or-modelled.md} exists, is committed, and every row is
either filled in or explicitly marked as belonging to a later chapter.
\item \texttt{doc/footprint.csv} has its first row, and the script appends
rather than overwrites.
\end{itemize}

\subsection*{Variants}

\begin{table}[H]\centering\small
\begin{tabular}{p{24mm}p{26mm}p{44mm}p{22mm}p{20mm}}
\toprule
Axis & Variant & What changes & Cost & Built in full in \\
\midrule
Operating system & Bare metal & The baseline here. One loop, no scheduler, nothing to configure & Nothing can run while something else waits & Here \\
Operating system & A real-time kernel & The identity, indicator and self-check modules are unchanged; they become the first task instead of the first function & A scheduler to reason about, and a stack per task & Chapter 2 \\
Operating system & The upstream open kernel & The board file replaces the board support layer entirely, and the node description becomes a device tree that another engineer can read & A second build system, and a different driver model & Chapter 4 \\
Execution model & Polled loop & The baseline here & Latency nobody has measured yet & Chapter 2 \\
Identity & Compiled-in node number & One constant, no storage, no protocol & Every board in a fleet needs its own build & Here \\
Identity & Node number from storage & A settable number that survives power loss, which is what a fleet needs & Non-volatile storage, and a protocol to set it & Chapters 12 and 19 \\
Console & Serial port over the probe & The baseline here. No extra hardware, no licence question & It needs a cable and a terminal & Here \\
Console & Diagnostics over the bus & The node reports itself to whoever is listening, with no cable at all & A bus, and a frame layout & Chapters 11 and 20 \\
\bottomrule
\end{tabular}
\caption{Variants for chapter 1. The two identity rows are the interesting pair: the difference between a demonstration and a product is very largely the difference between a number you compile in and a number the node is told.}
\end{table}

\subsection*{Pitfalls}

\begin{itemize}
\item Starting with the bus. It is the most interesting part and it is the
fourth thing to build. A node that cannot prove its own control period has
nothing worth putting on a wire.
\item Inheriting a linker script or a clock configuration from the better-known
member of this silicon family. The failure looks like a hardware fault.
\item Reading an upstream project's description of this part as a list of its
capabilities. The upstream operating system's device tree for this part leaves
an Ethernet node in place, inherited from a shared family file and never
deleted, although the silicon has no Ethernet controller at all. The vendor's
own headers settle it: the two interrupt slots that carry Ethernet on the
sibling part are simply absent here. A machine-readable file is evidence of what
its author modelled, not of what the silicon does.
\item Letting a modelled number drift into a measured column. It happens through
a chain of reasonable steps, which is why the boot banner names the modelled
quantities out loud every time the board starts.
\item Writing to an LED pin from wherever it is convenient. By chapter 18 the
indicators are a safety-relevant output and there must be exactly one place that
sets them.
\item Leaving the self-check as a warning. A check that prints and continues
trains you to ignore it.
\end{itemize}

\subsection*{Best practices applied}

\begin{itemize}
\item The boundary between real and modelled is written down, committed, and
printed by the device itself at every boot.
\item Facts about the board carry their evidence. Anything settled from the
vendor's own machine-readable sources says so; anything taken from convention is
marked as unconfirmed until a manual is opened.
\item The build is reproducible from a clean clone on a machine that has never
had a vendor development environment installed, which is the only definition of
a reproducible build that means anything to a reader.
\item Resources that later chapters need are claimed in the peripheral table
now, so that the allocation is a book-level decision made once.
\item A measurement stays blank until it has been taken. Three of the five rows
in this chapter's budget table say so.
\end{itemize}

\subsection*{Stretch goals}

\begin{itemize}
\item Bring the same identity and self-check code up under the open simulator,
using one of the published platform files for a sibling part as a starting
point. It costs an evening, and it gives chapter 20 a head start on a bench that
needs no hardware.
\item Add a second build configuration that reports the footprint of every
module separately, so the footprint row for each chapter can be attributed
rather than just totalled.
\item Write the boot banner as a machine-readable line as well as a human one,
so that the rig in chapter 20 can assert on the identity of the board it is
talking to instead of trusting the serial port it opened.
\end{itemize}

\subsection*{Roadmap and next steps}

Chapter 2 gives the node a period, which is the thing everything else hangs off,
and proves it without an oscilloscope.

The published progression beyond this chapter is clear and worth following in
order. The robot control framework's own documentation for writing a hardware
component defines the lifecycle a joint is expected to have, the read and write
pair that runs in the real-time loop, and the interface names a joint exposes,
including ones for its controller gains. That is the shape this node grows into
over the next nineteen chapters, and reading it now means the reader can see
where each chapter is heading. Its mock component is the thing to compare
against: it is the published answer to the question this volume asks, built by
people who needed it for the same reason.

For the control period specifically, the next thing to read is the lecture
material on cyclic task response timing, which contains the clearest published
statement of the mistake chapter 2 exists to prevent: delaying for a period is
not the same as running every period.

\subsection*{Portfolio evidence}

\begin{itemize}
\item A public repository whose README opens with the architecture figure and
carries the real-or-modelled table above the fold, so that a reader sees the
boundary before the claims.
\item A photograph or short recording of the board printing its banner at reset,
including the modelled line.
\item The footprint file, which after twenty chapters is a small, genuinely
uncommon piece of published data: what each layer of a joint node costs in flash
and memory on a real part.
\item A one-paragraph answer, written down and rehearsed, to the question of
whether you have built a joint node. The answer names what is real, what is
modelled, and what is absent, and it is a better answer than a yes.
\end{itemize}

\subsection*{Sources}

Normative references:
\begin{itemize}
\item Reference manual RM0455, for this part. Not RM0433, which documents its
better-known sibling.
\item The board user manual UM2408, for the connectors, the jumpers, the
indicator pins and the console wiring.
\item The STM32H7A3xI datasheet, for memory sizes and the clock limits at each
voltage scale. Confirm the document number before citing it: one number in
circulation belongs to a different pair of parts.
\end{itemize}

Reusable implementations:
\begin{itemize}
\item The robot control framework, including its hardware component
documentation and its mock component, Apache-2.0.\newline
\url{https://github.com/ros-controls/ros2_control}
\item The field-bus-to-framework bridge whose device-side contract this volume
implements, Apache-2.0.\newline
\url{https://github.com/ros-industrial/ros2_canopen}
\item The upstream operating system's board documentation for this exact board,
Apache-2.0.\newline
\url{https://docs.zephyrproject.org/latest/boards/st/nucleo_h7a3zi_q/doc/index.html}
\item The silicon vendor's device headers, Apache-2.0, which settle the memory
map and the absent peripherals.\newline
\url{https://github.com/STMicroelectronics/cmsis-device-h7}
\item The open simulator, MIT, whose platform files model this silicon family
including its bus controller.\newline
\url{https://github.com/renode/renode}
\end{itemize}
   % What a joint node is, and the bench that stands in for one
\project{02}{The control period: 1 kHz you can prove}{A heartbeat}{Timer-driven period, jitter, measuring it without a scope}

\begin{keyfacts}
\item[Adds to the node] A period. Something happens exactly 1000 times a second, and the node can show you the distribution rather than assert the rate
\item[Peripherals] TIM6 as the period source, the cycle counter in the debug unit as the instrument, one GPIO as the external cross-check, USART3 to report
\item[Depends on] Chapter 1 for the clock, the self-check and the console
\item[Real or modelled] Entirely real. This is the last chapter before anything is modelled, and the measurement here is what every modelled result later is compared against
\item[Difficulty] 3 of 5
\item[Effort] Two evenings of about four hours, most of it on the histogram rather than the timer
\item[Deliverable] A period handler that runs at 1 kHz, a jitter histogram built in memory with no probe attached, a printed report giving median, 99.9th percentile and maximum, and an overrun counter that is still zero after an hour
\end{keyfacts}

\subsection*{Why this chapter}

Everything above this chapter assumes a period. The controller in chapter 16
discretises its gains against one. The sensor timestamps in chapter 6 are
meaningful only relative to one. The state frames in chapter 11 carry a sequence
number that counts them. A joint node whose period is approximately one
millisecond, most of the time, is not a joint node that can make any claim at
all about following error.

The mistake this chapter exists to prevent has a name and a published statement
of it. Delaying for a period is not the same as running every period. A loop
that does its work and then waits a millisecond runs at one millisecond plus
however long the work took, and the error accumulates for as long as the machine
is switched on. The fix is to wait until an absolute time rather than for a
relative one, or to let a timer decide, and it is three lines different from the
mistake.

The harder half is proving it. This bench has no oscilloscope and no logic
analyser, so a claim about a period is worth exactly as much as the instrument
behind it. The instrument this chapter builds is inside the part: a free-running
cycle counter, read in the period handler, with the difference from the previous
reading binned into a histogram in memory. At 280 MHz a cycle is 3.57 ns and one
period is exactly 280,000 cycles, so the histogram is the jitter distribution,
in 3.57 ns bins, with nothing attached to the board.

\begin{note}{Three quantities that are all called jitter}
They are different, they have different causes, and a chapter that conflates
them produces a number nobody can act on. \textbf{Period jitter} is the
variation in the interval between one tick and the next: it is a property of the
timer and its clock, and on a timer driven from the same oscillator as the core
it is very small. \textbf{Release jitter} is the variation between the tick and
the moment the work actually starts: it is caused by other interrupt handlers,
by critical sections that disable interrupts, and by the scheduler if there is
one. \textbf{Execution time} is how long the work itself takes, and its
variation on a cached core is larger than most people expect. This chapter
measures all three separately, because the fix for each is different.
\end{note}

\subsection*{Prior art and what to reuse}

\begin{table}[H]\centering\small
\begin{tabular}{p{34mm}p{46mm}p{38mm}p{20mm}}
\toprule
Source & What it gives & What it does not & Licence \\
\midrule
The real-time kernel's own delay-until primitive, quoted from the kernel header & The canonical statement: a delay to an absolute time rather than a relative one, for tasks that need a constant execution frequency, with the wake time held by the caller and updated inside the call. The newer form returns whether a delay actually happened, which is a free overrun detector & Nothing about measuring the result, and nothing about what happens when the work overruns more than once & MIT \\
Koopman's lecture on interrupt and cyclic task response timing, Monday 14 March 2016 & The clearest published statement of the defect this chapter prevents, in an annotated bad-code example, and the distinction between response time and execution time & Its terms of use permit download for personal and academic use with attribution and prohibit republishing, so it is cited and linked and never quoted at length & course material, all rights reserved \\
Ganssle's article on interrupt latency & The measurement method that needs no instrument: the handler reads the timer's count register, which has kept incrementing throughout the latency & It predates this core and says nothing about caches & article \\
Cervin and colleagues on analysis tools for real-time control systems, \pubdate{August 2002} & The sentence that connects this chapter to chapter 16: digital control theory assumes equidistant sampling and a negligible or constant delay, and this can seldom be achieved in practice & The tools are host-side and off this bench. One of the two project pages is now a dead link and is not cited & paper \\
Schwarzmann and Kaeser on the effect of sampling-time jitter, Thursday 4 September 2025 & The modern result that jitter scales the system matrices, so a varying period perturbs the plant model rather than adding noise & Not a firmware paper & preprint \\
Styger on cycle counting with the debug unit, Monday 30 January 2017, and the probe vendor's knowledge base article on the same counter & The enable sequence and the two register addresses, and two traps that cost an evening each & Neither presents the histogram construction & articles \\
\bottomrule
\end{tabular}
\caption{Prior art for chapter 2. The second row is the whole argument of the chapter, stated by somebody else in one line on one slide, and it is worth reading the original even though this book cannot reproduce it.}
\end{table}

What is left to write is the instrument. No published article was found that
presents the in-memory histogram method for proving a control period with no
probe attached, so it is presented here as this book's own construction on two
borrowed ideas: reading a free-running counter inside the handler, which is
Ganssle's, and the enable sequence for the counter, which is Styger's.

There is also a negative finding worth printing, because a reader will look for
it. \textbf{No application note from the silicon vendor on building a
fixed-period digital control loop was found.} The timer cookbook note exists and could not be
fetched from this network; it is a peripheral guide and would not fill the gap
in any case.

\subsection*{What the node gains}

Before this chapter the node can tell you what it is. After it, the node has a
heartbeat: something happens 1000 times a second, and the node can print the
distribution of the interval between those somethings rather than claiming a
rate. It still senses nothing and commands nothing. What it has is the thing
every later measurement is quoted against.

\subsection*{Parts from the inventory}

\begin{table}[H]\centering\small
\begin{tabular}{p{40mm}p{66mm}p{40mm}}
\toprule
Part & Role & Interface \\
\midrule
NUCLEO-H7A3ZI-Q & The node, and its own instrument & Micro USB to the host \\
The power profiler, as a logic recorder & The external cross-check. One of its eight digital inputs records a pin the handler toggles, time-aligned with its current trace & One jumper wire to a header pin, one to ground \\
One male-to-female jumper wire & The only wiring in this chapter & Board header to instrument input \\
\bottomrule
\end{tabular}
\caption{Inventory items used in chapter 2. The instrument here is a cross-check, not the measurement: its sampling rate is three orders of magnitude too slow to see the jitter, and the chapter says exactly what it can and cannot settle.}
\end{table}

\subsection*{System architecture}

\diagram{j02_arch}{Two ways to own a period, both built in this chapter. On the left the timer interrupt is the period and the work happens in the handler. On the right the timer releases a task and the work happens in the task. The instrument, the histogram and the report are identical in both, which is what makes the comparison worth anything.}

The figure shows the decision this chapter makes twice on purpose. The handler
route has the lowest release jitter available on this part, because nothing
stands between the timer and the work except interrupt entry. The task route
costs a context switch and a scheduler decision, and buys the ability to block,
to be preempted by something more urgent, and to share the processor with the
bus and the middleware that arrive in later chapters. Both are built, both are
measured with the same instrument, and the numbers decide.

\subsection*{Peripheral configuration}

\begin{table}[H]\centering\small
\begin{tabular}{p{22mm}p{26mm}p{24mm}p{30mm}p{30mm}}
\toprule
Peripheral & Mode & Clock source & Pins and function & Interrupt and transfers \\
\midrule
TIM6 & Up-counting, auto-reload, update event & Peripheral bus timer clock, computed at boot & None & Update interrupt, highest priority of the application handlers \\
DWT cycle counter & Free running, 32 bit & Core clock, 280 MHz & None & None. Read, never written \\
GPIOB & Push-pull output & Peripheral bus & One spare header pin, toggled in the handler & None \\
USART3 & Asynchronous, 115200 8N1 & Peripheral bus & Console, from chapter 1 & Polled, and never called from the handler \\
\bottomrule
\end{tabular}
\caption{Peripheral configuration. The last row is a rule rather than a setting: nothing in the period handler prints, because a polled console write at 115200 takes longer than the period it would be reporting.}
\end{table}

\subsection*{Wiring}

\diagram{j02_wiring}{The only wire in this chapter, and an honest statement of what it can settle. The instrument records the toggled pin alongside its current trace, so a wrong rate is obvious. It samples far too slowly to see the jitter, which is why the real measurement is inside the part.}

Two facts about that wire. It is the only external evidence in the chapter, and
it is deliberately weak evidence: an instrument sampling at tens of microseconds
can confirm that the period is a millisecond and not thirty-five, which is
exactly the failure a miscompiled clock produces, and it can do nothing about a
distribution measured in tens of nanoseconds. Saying which instrument settles
which question is the whole of measurement discipline, and it costs one
paragraph.

\subsection*{Memory and timing budget}

\begin{table}[H]\centering\small
\begin{tabular}{p{44mm}p{28mm}p{28mm}p{30mm}}
\toprule
Quantity & Budget & Measured & Margin \\
\midrule
Flash, this chapter & 6 kB & not measured & not measured \\
Static memory, this chapter & 1.5 kB, mostly the histogram & not measured & not measured \\
Period & 1000.0 us & not measured & not measured \\
Period jitter, peak to peak & under 1 us & not measured & not measured \\
Release jitter, maximum & under 5 us & not measured & not measured \\
Handler execution time, maximum & under 50 us, which is 5 per cent of the period & not measured & not measured \\
Overruns in one hour & 0 & not measured & not measured \\
\bottomrule
\end{tabular}
\caption{The budget for chapter 2. Every row is a number this chapter's own instrument produces, which is unusual: most chapters measure some rows and leave others. Nothing is written in the measured column until the histogram has run for an hour.}
\end{table}

The execution-time budget deserves its reasoning. A control loop that spends
half its period computing has no room for the bus, the sensors or the safety
checks that later chapters add, and it has no margin for the cache misses that
make a Cortex-M7's worst case much worse than its typical case. Five per cent is
the number this volume holds to, and chapter 16, which adds the most work to the
handler of any chapter, reports against it.

\subsection*{Firmware design (UML)}

\diagram{j02_uml}{One period, from the timer's update event to the histogram, and the report path that is deliberately outside the handler. The dashed return is the overrun case: the tick arrives while the previous one is still being served, and the only correct response is to count it and carry on.}

Three design decisions.

The handler does the minimum and nothing that can block. It reads the counter,
computes the delta, bins it, toggles the pin, calls the loop step, and returns.
It does not print, it does not allocate, and it does not take a lock that
anything else holds.

The histogram lives in memory the startup code does not clear, so that a reset
does not lose the evidence. That costs one linker section and it is the first
appearance of a technique chapter 20 relies on completely.

The report runs in the idle loop and reads the histogram while the handler is
still writing to it. That is a race, and the honest fix is not a lock, which
would add latency to the handler, but a sequence counter the reader checks
before and after: if it changed, read again. The report is a statistic, so one
retry is enough and a stale sample is harmless.

\subsection*{Data flow (ASCII)}

\begin{asciiart}
  TIM6 update every 280,000 cycles
        |
        v
  +-------------------------------------------------------------+
  | period handler, highest application priority                 |
  |   now   = DWT->CYCCNT                                        |
  |   delta = now - last          (unsigned, wraps correctly)    |
  |   last  = now                                                |
  |   bin   = clamp((delta - 280000) / 8 + 32, 0, 63)            |
  |   hist[bin]++          min/max updated          seq++        |
  |   pin toggle  --+                                            |
  |   loop_step()   |      (empty here, the controller in ch 16) |
  +-----------------|-------------------------------------------+
                    +--> one digital input on the power profiler
        |
        | never prints, never blocks, never allocates
        v
  +-------------------------------------------------------------+
  | idle loop, no deadline at all                                |
  |   s1 = seq; copy histogram; s2 = seq; if (s1 != s2) retry    |
  |   median, 99.9th percentile, maximum, overruns               |
  |   printf over USART3, or a frame on the bus from chapter 11  |
  +-------------------------------------------------------------+
\end{asciiart}

\subsection*{Repository layout}

\begin{plaincode}
joint-node/
  CMakeLists.txt
  cmake/arm-none-eabi.cmake
  ld/stm32h7a3zi.ld                   # + .noinit section for the histogram
  doc/
    node-identity.md  real-or-modelled.md  footprint.csv
    period-report.md                  # + this chapter: the measured numbers
  src/
    node/    identity.c state.h indicator.c selfcheck.c
    bsp/     startup.c system.c uart.c board.h
    control/
      period.c  period.h              # + this chapter: TIM6, the handler
      jitter.c  jitter.h              # + this chapter: the histogram
      cyclecount.h                    # + this chapter: the debug unit counter
      loop.h                          # + this chapter: the empty loop step
    sense/  act/  bus/  mw/  safety/  update/
  host/
    plot_jitter.py                    # + this chapter: the report as a figure
  test/
    test_jitter.c                     # + this chapter: binning, on the host
  tools/  footprint.py
  README.md
\end{plaincode}

\subsection*{Steps}

\begin{steps}
\item \textbf{Compute the timer clock at boot rather than assuming it.} The
timer input is the peripheral bus clock, and on this family it is doubled when
the bus prescaler is not one. That rule is in the reference manual for this
part, and it is the kind of detail that differs between family members, so the
firmware derives it from the clock registers and prints it. A comment in your
own code is not evidence.
\begin{ccode}
/* period.c: derive, do not assume. The printed value is the one to trust. */
uint32_t timer_clock_hz(void)
{
    uint32_t pclk = rcc_pclk1_hz();          /* read back from the registers */
    return (rcc_ppre1_divider() == 1u) ? pclk : pclk * 2u;
}

void period_init(uint32_t hz)
{
    uint32_t tclk = timer_clock_hz();
    uint32_t ticks = tclk / hz;              /* exact division, or fail */
    if (tclk % hz != 0u) fail("period: timer clock does not divide cleanly");
    TIM6->PSC = 0u;                          /* no prescaler: finest resolution */
    TIM6->ARR = ticks - 1u;
    printf("period: timer clock %lu Hz, reload %lu, target %lu Hz\r\n",
           tclk, ticks - 1u, hz);
}
\end{ccode}
The check on the remainder matters. A period that is 1000.4 microseconds because
the arithmetic did not divide cleanly is a period that accumulates four
milliseconds of error every ten seconds, and it looks like a perfectly good
1 kHz loop on any instrument this bench owns.

\item \textbf{Enable the cycle counter, and check it exists.} It is optional
silicon. It is also usually enabled already when a debugger is attached, which
is the trap: the code works on the bench and fails standalone, weeks later,
with the histogram full of zeros.
\begin{ccode}
/* cyclecount.h: the whole instrument, in nine lines */
#define DWT_CTRL    (*(volatile uint32_t *) 0xE0001000u)
#define DWT_CYCCNT  (*(volatile uint32_t *) 0xE0001004u)
#define DEMCR       (*(volatile uint32_t *) 0xE000EDFCu)

static inline bool cyclecount_init(void)
{
    DEMCR |= (1u << 24);                     /* TRCENA: the debugger may have */
    DWT_CYCCNT = 0u;                         /* done this already             */
    DWT_CTRL |= 1u;                          /* CYCCNTENA                     */
    return DWT_CYCCNT != 0u;                 /* it counts, so it exists       */
}

static inline uint32_t cyclecount(void) { return DWT_CYCCNT; }
\end{ccode}
Wire the return value into the self-check from chapter 1. A node that cannot
measure its own period should say so at boot, not produce an empty histogram an
hour later.

\item \textbf{Write the period handler, and keep it short.} Unsigned arithmetic
makes the counter's wrap at about 15.3 seconds a non-event: the difference of
two unsigned 32-bit values is correct across a wrap, which is worth a comment in
the code because it looks like a bug to a reviewer.
\begin{ccode}
void TIM6_DAC_IRQHandler(void)
{
    TIM6->SR = 0u;                           /* clear first, always */

    uint32_t now   = cyclecount();
    uint32_t delta = now - g_last;           /* correct across the 32-bit wrap */
    g_last = now;

    gpio_toggle(PERIOD_PIN);                 /* the external cross-check */
    jitter_record(delta);                    /* the histogram, below */

    uint32_t t0 = cyclecount();
    loop_step();                             /* empty here, chapter 16 fills it */
    jitter_record_exec(cyclecount() - t0);

    if (TIM6->SR & TIM_SR_UIF) g_overruns++; /* a tick arrived while we worked */
}
\end{ccode}

\item \textbf{Build the histogram.} Sixty-four bins of eight cycles each spans
plus or minus 256 cycles, which is plus or minus 914 ns, with everything outside
that counted separately. Bins are cheap and the resolution is what makes the
figure worth printing.
\begin{ccode}
#define JIT_BINS   64u
#define JIT_SCALE   8u                       /* cycles per bin: 28.6 ns */
#define JIT_NOM    280000u                   /* 280 MHz / 1 kHz, exactly */

void jitter_record(uint32_t delta)
{
    int32_t off = (int32_t) delta - (int32_t) JIT_NOM;
    int32_t bin = (off / (int32_t) JIT_SCALE) + (int32_t) (JIT_BINS / 2u);

    if (bin < 0 || bin >= (int32_t) JIT_BINS) g_jit.outside++;
    else                                      g_jit.bin[bin]++;

    if (delta < g_jit.min) g_jit.min = delta;
    if (delta > g_jit.max) g_jit.max = delta;
    g_jit.count++;
    g_jit.seq++;                             /* the reader checks this */
}
\end{ccode}

\item \textbf{Put the histogram where a reset cannot erase it.} One section in
the linker script, one attribute in the source, and a magic number so the code
can tell a fresh histogram from a surviving one.
\begin{ldcode}
  .noinit (NOLOAD) : { . = ALIGN(4); *(.noinit*) . = ALIGN(4); } > RAM
\end{ldcode}
\begin{ccode}
__attribute__((section(".noinit"))) static jitter_t g_jit;
/* g_jit.magic != JIT_MAGIC means this is a cold start: zero it. Otherwise the
   histogram survived a reset, which is exactly what you want to look at. */
\end{ccode}

\item \textbf{Report outside the handler.} Read the histogram with the sequence
check, compute the three numbers that matter, and print them from the idle loop.
\begin{ccode}
jitter_t snap;
do { uint32_t s1 = g_jit.seq; snap = g_jit; if (s1 == g_jit.seq) break; }
while (1);                                   /* one retry is enough */

printf("period n=%lu  min=%ld ns  p50=%ld ns  p99.9=%ld ns  max=%ld ns  over=%lu\r\n",
       snap.count, ns(snap.min), ns(pct(&snap, 500)), ns(pct(&snap, 999)),
       ns(snap.max), snap.overruns);
\end{ccode}
Report the median, the 99.9th percentile and the maximum, in that order, and
never the mean. A control loop is not harmed by its average period and it is
harmed by its worst one, which is why published real-time benchmarks report
percentiles and a maximum.

\item \textbf{Build the same thing again as a kernel task.} The work is
identical; only the release path changes. The delay-until primitive takes an
absolute wake time that it updates itself, and its return value says whether the
deadline had already passed, which is an overrun detector for free.
\begin{ccode}
static void control_task(void *arg)
{
    TickType_t wake = xTaskGetTickCount();
    for (;;) {
        BaseType_t delayed = xTaskDelayUntil(&wake, pdMS_TO_TICKS(1));
        if (delayed == pdFALSE) g_overruns++;  /* the deadline had passed */

        uint32_t now = cyclecount();
        jitter_record(now - g_last);
        g_last = now;
        loop_step();
    }
}
\end{ccode}
Note what this version cannot do: a kernel tick of one millisecond gives a
period quantised to the tick, so the honest comparison releases the task from
the timer interrupt with a notification instead, and measures that. Build both
and print both.

\item \textbf{Cross-check the rate with the external instrument.} One wire from
the toggled pin to a digital input, one capture of a few seconds, and a count of
edges. This settles the rate and says nothing about the distribution, and the
report says so in the sentence next to the number.
\begin{shellcode}
python host/count_edges.py capture.csv --expect 1000 --tolerance 0.001
# edges: 20000 in 10.000 s -> 1000.0 Hz toggles, 2000.0 Hz edges: PASS
\end{shellcode}

\item \textbf{Run it for an hour and look at the shape.} A histogram with a
single narrow peak is a timer doing its job. A second peak says something else
in the system takes the processor at a regular rate. A long tail to the right
says a rare handler or a cache effect, and the maximum is the number to quote.
\begin{shellcode}
python host/plot_jitter.py period-report.csv -o doc/fig/j02_hist.svg
\end{shellcode}
\end{steps}

\diagram{j02_data}{What the instrument produces, and how to read it. The bar
heights here are the shape to expect rather than measured data, and the axis
deliberately carries no counts: this figure is in the chapter to teach the
reading, and the measured version belongs in the repository with its own
one-hour run behind it. Three shapes, three different causes, three different
fixes.}

\subsection*{Build, flash and debug}

\diagram{j02_timing}{The defect this chapter prevents, drawn to scale. Above, a loop that delays for a period: the work time is added to every cycle and the error accumulates without limit. Below, a timer that owns the period: the work time moves the start of the work and does not move the next tick. The two look identical on an instrument that samples slowly, and they are not the same machine.}

\begin{shellcode}
cmake --build build -j && probe-rs run --chip STM32H7A3ZITx build/firmware.elf
\end{shellcode}

\begin{note}{When the histogram is empty and everything looks right}
In order of likelihood: the cycle counter was never enabled, and the debugger
was enabling it for you during development; the counter is enabled but the
handler is not running, because the timer interrupt was never unmasked in the
interrupt controller, which is a separate step from enabling the update
interrupt in the timer; the handler runs but clears the status register last, so
it re-enters immediately and the delta is a few hundred cycles rather than
280,000; or the histogram is in the section that startup clears, so it is
faithfully zeroed a moment before you read it.
\end{note}

\subsection*{Verification and acceptance criteria}

\begin{itemize}
\item The boot line prints the derived timer clock, the reload value and the
target rate, and the firmware refuses to start if the division is not exact.
\item The external instrument counts 2000 edges per second within one part in a
thousand, over a capture of at least ten seconds.
\item The histogram over one hour has a single peak, a median within 100 ns of
the nominal period, and a maximum that is inside the budget above.
\item The overrun counter is zero after an hour in both the handler version and
the task version.
\item Deliberately adding 900 microseconds of work to the loop step makes the
overrun counter increase and the report say so, rather than the node quietly
running at a lower rate.
\item The histogram survives a reset: press the button, read the report, and the
count is not one.
\item The binning function has host-side unit tests, including the two cases
that are easy to get wrong: a delta below the lowest bin and a delta that wraps
the counter.
\end{itemize}

\subsection*{Variants}

\begin{table}[H]\centering\small
\begin{tabular}{p{24mm}p{26mm}p{44mm}p{22mm}p{20mm}}
\toprule
Axis & Variant & What changes & Cost & Built in full in \\
\midrule
Execution model & Timer interrupt handler & The baseline. Lowest release jitter available on this part & Nothing else can run during the work & Here \\
Execution model & Kernel task released by the timer & A notification from the handler releases a task, which does the work. Preemptible, can block, shares the processor with later chapters & A context switch, a stack, and release jitter you must measure & Here \\
Execution model & Kernel task on the tick & The simplest code of the three, and the period is quantised to the kernel tick & Resolution & Here, as the version to reject \\
Synchronisation & Direct task notification against a semaphore & Both release a task from a handler. The published claim that one is a fixed percentage faster comes from a page that states no processor, no compiler and no method, and the kernel's own book has dropped the figure & One measurement on this board settles it & Chapter 12 \\
Peripheral & A 32-bit timer instead of the basic one & Nothing at 1 kHz. It matters when a chapter wants a long free-running time base, which is chapter 3 & None here & Chapter 3 \\
Time base & The cycle counter against a timer capture & The counter is finer and wraps every 15.3 seconds; a timer capture is coarser and can be extended in software & Resolution against range & Chapter 3 \\
Operating system & The upstream open kernel & Its own timer and work-queue model, and a device tree entry instead of register writes & A second build system & Chapter 4 \\
Language & Rust on the target & The same timer and the same counter, with the handler's shared state expressed so that the race the sequence counter guards cannot be written & A second toolchain & Chapter 16 \\
\bottomrule
\end{tabular}
\caption{Variants for chapter 2. Three execution models are built here rather than one, because the comparison is the chapter: the same work, the same instrument, three release paths, three distributions.}
\end{table}

\subsection*{Pitfalls}

\begin{itemize}
\item Delaying for a period instead of until a time. The error accumulates for
as long as the machine runs, and no instrument on this bench will show it in a
short capture.
\item Trusting the timer clock to be the bus clock. It is doubled under a
condition that differs across this silicon family, so derive it and print it.
\item A reload value that does not divide cleanly. The period is then wrong by a
fraction of a per cent, forever, and it looks perfect.
\item Enabling the cycle counter only when a debugger is attached. The code
works for weeks and the field unit reports nothing.
\item Printing from the period handler. A polled console write at 115200 takes
longer than the period.
\item Putting the histogram in a section the startup code clears, and then
concluding the board resets more often than it does.
\item Reporting the mean period. It is the one statistic that cannot show the
problem.
\item Measuring execution time once and calling it the worst case. On this core
the cache, the branch predictor and the long pipeline spread the same code over
a much wider range than on a simpler part, so the maximum over an hour is the
number, not the first reading.
\end{itemize}

\subsection*{Best practices applied}

\begin{itemize}
\item The instrument is named beside every number, and the sentence next to the
external measurement says what that instrument cannot settle.
\item The firmware derives and prints the facts it depends on instead of
carrying them in comments.
\item The handler does the minimum, and the rule is written down rather than
implied: nothing that blocks, nothing that allocates, nothing that prints.
\item The race between the handler and the reporter is acknowledged and handled
with a sequence counter rather than removed with a lock that would cost the
handler latency.
\item Percentiles and a maximum, never a mean, following the published practice
for real-time benchmarks.
\item The binning arithmetic is tested on a host machine, where a failing case
can be written down and kept.
\end{itemize}

\subsection*{Stretch goals}

\begin{itemize}
\item Add a second histogram for release jitter specifically, by having the
timer capture its own count at the update event and comparing it with the
counter reading at the start of the handler. The difference is interrupt entry
latency, measured rather than quoted.
\item Run the same firmware under the open simulator and compare the shape of
the histogram with the hardware's. The differences are a lesson in what a
simulator does not model, and chapter 20 needs that lesson.
\item Sweep the interrupt priority of the period handler against a deliberately
noisy second handler, and plot the maximum release jitter against priority. That
is a figure almost nobody publishes for this core.
\item Repeat the whole measurement with the instruction cache disabled and put
both histograms on one axis.
\end{itemize}

\subsection*{Roadmap and next steps}

Chapter 3 gives the node a clock as well as a period: a monotonic time base that
outlives the counter's 15.3 second wrap, and a way to stamp a sample with a time
another node would agree with.

The published progression from here has two branches. On the timing side, the
lecture material cited above continues into scheduling and into watchdog timers,
and both are the right next reading for anyone whose loop now runs but whose
system does not yet have more than one deadline in it. On the control side, the
two papers cited above are the bridge to chapter 16: the 2002 paper states that
control theory assumes a constant period, and the 2025 preprint shows that
varying it perturbs the plant model rather than adding noise, which is why this
chapter's histogram is a control result and not a firmware curiosity.

\subsection*{Portfolio evidence}

\begin{itemize}
\item The jitter figure, with the instrument named in its caption and the three
release paths on one axis.
\item The one-hour report: count, median, 99.9th percentile, maximum, overruns,
and the external edge count beside it.
\item The deliberate-overrun run, showing the node reporting the overrun rather
than silently slowing down. A demonstration of a failure being detected is worth
more than a demonstration of a success.
\item A paragraph, written down, on what the external instrument could and could
not settle, which is the paragraph an interviewer follows up on.
\end{itemize}

\subsection*{Sources}

Normative references:
\begin{itemize}
\item Reference manual RM0455, for the timer clock rule, the timer registers and
the interrupt controller. Confirm the peripheral bus divider here rather than
assuming it.
\item The Arm architecture reference for this core, for the debug unit's cycle
counter and the fact that it is optional.
\item The board user manual UM2408, for a free header pin to toggle.
\end{itemize}

Reusable implementations:
\begin{itemize}
\item The real-time kernel, MIT, whose header carries the quotable definition of
the delay-until primitive and its overrun return value.\newline
\url{https://github.com/FreeRTOS/FreeRTOS-Kernel}
\item Koopman's lecture notes on interrupt and cyclic task response timing.
Download for personal and academic use with attribution; republication is not
permitted, so they are linked and not quoted.\newline
\url{https://users.ece.cmu.edu/~koopman/lectures/index.html}
\item Ganssle's article on interrupt latency, including the measurement that
needs no instrument.\newline
\url{https://www.ganssle.com/articles/interruptlatency.htm}
\item Styger on cycle counting with the debug unit, for the enable sequence and
the debugger trap.\newline
\url{https://mcuoneclipse.com/2017/01/30/cycle-counting-on-arm-cortex-m-with-dwt/}
\item Cervin, Henriksson, Lincoln and Arzen, analysis tools for real-time
control systems, \pubdate{August 2002}.\newline
\url{https://lucris.lub.lu.se/ws/files/6359302/625680.pdf}
\item Schwarzmann and Kaeser, on the effect of sampling-time jitter, Thursday 4
September 2025.\newline
\url{https://arxiv.org/abs/2509.04199}
\end{itemize}
   % The control period: 1 kHz you can prove
\project{03}{One clock for sensors, loop and bus}{A shared time base}{Timestamping, a monotonic tick, what to do without precision time hardware}

\begin{keyfacts}
\item[Adds to the node] A time base: a monotonic count that never steps backwards and does not wrap for centuries, a wall clock kept separately, and a rule about which one may be used for what
\item[Peripherals] TIM2 as the 32-bit monotonic source with a software extension, the real-time clock with the fitted crystal as the wall clock, one timer input capture channel for stamping external events
\item[Depends on] Chapter 1 for the clock tree, chapter 2 for the period and the cycle counter
\item[Real or modelled] Entirely real. The clock comparison at the end is a measurement of two real oscillators against each other
\item[Difficulty] 3 of 5
\item[Effort] Two evenings of about four hours
\item[Deliverable] A 64-bit monotonic microsecond clock with a race-free read, an input capture that stamps an event at the pin rather than at the interrupt, a measured drift figure in parts per million between the two oscillators, and a written rule for which clock every later chapter may use
\end{keyfacts}

\subsection*{Why this chapter}

Chapter 2 gave the node a period. A period tells you how often something
happens; it does not tell you when anything happened. The moment the node has a
sensor, it needs to be able to say that a sample belongs to a particular control
cycle and was taken at a particular instant. The moment it has a bus, it needs
to say that to another machine, which has its own oscillator and its own idea of
now.

Three quantities get confused here and they are not the same thing. A
\textbf{monotonic} clock counts forward and never steps back; it is the right
thing for measuring an interval and for ordering events, and it means nothing
outside the running machine. A \textbf{wall clock} says what the date and time
are; it is the right thing for a log entry and the wrong thing for anything
timed, because it can be set backwards while the node is running. And a
\textbf{shared} time base is an agreement between machines about one of the
other two, which is a protocol rather than a peripheral.

This chapter builds the first two properly and designs the third. It designs
rather than builds it because the node has no bus until chapter 9, and the
design decision that matters is made here: the hardware in this part captures a
timestamp at the start-of-frame bit of a bus frame, in the controller, before
any interrupt runs. That one property is what makes microsecond agreement
possible later without any precision time hardware at all, and it is why this
chapter is placed before the bus rather than after it.

\begin{note}{What this bench cannot do and why that is said first}
There is no precision time protocol here, because that needs an Ethernet
controller with hardware timestamping and this part has no Ethernet controller
at all. There is no time-triggered bus either: that mechanism lives in dedicated
transceivers, is absent from ordinary controllers, and, on the authority of a
peer-reviewed paper that says so in its own related-work section, was not
carried into the flexible-data successor. So the node's shared time base is
software over an ordinary bus, which is the interesting case anyway: it is what
most machines actually do, and there is a published protocol for it from 1994
that gets to about twenty microseconds and a paper from 2021 that gets below one.
\end{note}

\subsection*{Prior art and what to reuse}

\begin{table}[H]\centering\small
\begin{tabular}{p{34mm}p{46mm}p{38mm}p{20mm}}
\toprule
Source & What it gives & What it does not & Licence \\
\midrule
The silicon vendor's own community answer on how the bus controller timestamps & The fact this chapter is built around: the counter value is captured on the start of frame reception or transmission, not at the interrupt. Two cautions with it: the counter is 16 bit and wraps within seconds, and the timer that can clock it sits in a different clock domain from the bus peripheral, with a few per cent of drift observed between them & It is a forum answer, so the register clause numbers are confirmed against the reference manual before anything is printed & forum post \\
Gergeleit and Streich, first international bus conference, Mainz, \pubdate{September 1994} & The protocol, and the sentence that justifies it: the bus is always much shorter than one bit time, so every node sees the level at about the same moment, which is inherent to the arbitration method and is not true of other networks. About twenty microseconds, under twenty messages a second, one identifier. And the rule that a corrected clock is slewed and never stepped & Thirty-two years old, written before the flexible-data format, and its numbers are for a slower bus & conference paper, freely downloadable \\
Einspieler and colleagues, transactions paper, 2021 & Sub-microsecond precision in software over an ordinary bus, using a high resolution timer with rate correction that many microcontrollers already carry. Its critique is the design rule for this chapter: correcting offset without correcting rate leaves time discontinuous, and two events can then carry the same timestamp although they happened at different moments & Needs the timer module it describes; whether this part has an equivalent is an open question for the bench & IEEE copyright, cite and link \\
The field bus application layer profile, version 4.2.0, Monday 21 February 2011 & The vocabulary a joint node should not invent: a synchronisation object, an optional cycle counter that says which cycle a sample belongs to, a time stamp object, and a high resolution timestamp in microseconds. It also states its own limit: the jitter of a synchronisation frame is about one message transmission time & The high resolution timestamp is 32 bits of microseconds, so it wraps about every 71.6 minutes and the node must handle that & free after registration \\
The bus association's network time management document, Friday 1 September 2023 & The document most directly on this subject, covering timestamping on transmission and reception in both frame formats & Members only, so it is cited by number and title and nothing is quoted from it & members only \\
The embedded middleware client's time synchronisation header & A working epoch synchronisation against a host agent, in four functions, with the offset applied for you & It is a round trip yielding an offset only, with no rate correction and no hardware assistance, and over a bus transport it inherits arbitration delay. Right for a log entry, wrong for aligning a control period & Apache-2.0 \\
\bottomrule
\end{tabular}
\caption{Prior art for chapter 3. The 1994 paper and the 2021 paper are the same protocol twenty-seven years apart, and reading them in order is the fastest way to understand why the second one is twenty times better.}
\end{table}

What is left to write is the local half: a monotonic clock with a race-free
read, an event stamp taken at the pin rather than at the interrupt, and the rule
about which clock may be used for what. The protocol half is designed here and
built in chapter 12, when there is a bus to carry it.

\subsection*{What the node gains}

Before this chapter the node knows how often it runs. After it, the node knows
when: it can stamp an event with a 64-bit microsecond count that never goes
backwards, it can say what the date is without confusing that with elapsed time,
and it can state how far its own oscillator drifts against a second one. It
cannot yet agree with another machine, and the mechanism by which it will is
designed and written down.

\subsection*{Parts from the inventory}

\begin{table}[H]\centering\small
\begin{tabular}{p{40mm}p{66mm}p{40mm}}
\toprule
Part & Role & Interface \\
\midrule
NUCLEO-H7A3ZI-Q & The node. Two oscillators on it are used here: the high-speed clock from the probe and the fitted 32.768 kHz crystal & Micro USB to the host \\
The fitted low-speed crystal & The wall clock's reference, and the second oscillator in the drift measurement & On the board, confirmed fitted \\
The power profiler, as a logic recorder & An external event source for the input capture: its digital line, or simply the button, gives an edge whose stamp can be checked & One jumper wire \\
\bottomrule
\end{tabular}
\caption{Inventory items used in chapter 3. Nothing is bought. The one thing this chapter would buy if the budget allowed is a receiver that outputs a pulse per second, and the chapter says exactly what that would add.}
\end{table}

\subsection*{System architecture}

\diagram{j03_arch}{Every source of time on this part, what it is good for, and what it costs. Two of them are absent and drawn accordingly: the precision time protocol needs an Ethernet controller this part does not have, and time-triggered bus hardware lives in transceivers rather than in controllers. The one that matters most is the start-of-frame capture in the bus peripheral, which is free, already there, and unused until chapter 12.}

The figure is the chapter's argument in one picture. Four real sources, with
different resolution and different range, and the node uses three of them for
three different jobs rather than trying to make one of them do everything. The
cycle counter from chapter 2 is the finest and wraps in 15.3 seconds. The
32-bit timer at one microsecond wraps in 71.6 minutes and becomes the monotonic
clock when it is extended in software. The real-time clock keeps the date across
a reset and is far too coarse to time anything. And the bus controller captures
its own timestamps at the frame, which no software path can match.

\subsection*{Peripheral configuration}

\begin{table}[H]\centering\small
\begin{tabular}{p{22mm}p{26mm}p{24mm}p{30mm}p{30mm}}
\toprule
Peripheral & Mode & Clock source & Pins and function & Interrupt and transfers \\
\midrule
TIM2 & Free running, 32 bit, prescaled to 1 MHz & Peripheral bus timer clock, derived at boot & None & Update interrupt only, to extend the count \\
TIM2 channel 1 & Input capture, rising edge & As above & One header pin, the event input & Capture interrupt, low priority \\
RTC & Calendar with subsecond register & The fitted 32.768 kHz crystal & None & None. Read, never used for timing \\
DWT cycle counter & Free running, from chapter 2 & Core clock & None & None \\
FDCAN1 timestamp unit & Configured but unused & Reserved & None until chapter 9 & Named here so chapter 9 does not have to rediscover it \\
\bottomrule
\end{tabular}
\caption{Peripheral configuration for chapter 3. The prescaler that makes the timer tick at exactly one microsecond is derived from the measured timer clock, not written as a constant, for the same reason the reload value was in chapter 2.}
\end{table}

\subsection*{Wiring}

\diagram{j03_wiring}{One wire for the event input, and the thing this bench does not have drawn beside it. A receiver that emits one pulse per second would turn the drift figure at the end of this chapter into an absolute frequency measurement, and it is the single cheapest instrument that would improve this volume.}

\subsection*{Memory and timing budget}

\begin{table}[H]\centering\small
\begin{tabular}{p{44mm}p{28mm}p{28mm}p{30mm}}
\toprule
Quantity & Budget & Measured & Margin \\
\midrule
Flash, this chapter & 4 kB & not measured & not measured \\
Static memory, this chapter & 128 bytes & not measured & not measured \\
Monotonic clock resolution & 1 us & computed, exact by construction & n/a \\
Monotonic clock range & over 500,000 years & computed & n/a \\
Cost of one timestamp read & under 100 cycles & not measured & not measured \\
Drift, high-speed against low-speed & report in ppm & not measured & not measured \\
Backwards steps in one hour & 0 & not measured & not measured \\
\bottomrule
\end{tabular}
\caption{The budget for chapter 3. Two rows say computed rather than measured, because they follow from the arithmetic of a 64-bit microsecond count rather than from an observation, and this volume labels those differently on purpose.}
\end{table}

\subsection*{Firmware design (UML)}

\diagram{j03_uml}{Above, how a sensor event gets an honest timestamp: the edge is captured in hardware and the interrupt only reads the captured value, so interrupt latency moves when the node learns about the event and not when the event is recorded. Below, the protocol this design makes possible, from the 1994 paper: one frame to mark an instant, then a second frame carrying the marker's own timestamp. Nothing in the lower half is built until chapter 12.}

The upper half is a rule worth stating plainly: \textbf{stamp at the event, not
at the read}. An inertial unit that signals a new sample on a pin, whose edge is
captured by a timer channel, gives a timestamp that is correct even if the
firmware is busy for two hundred microseconds. The same sample read over a
two-wire bus and stamped when the read completes carries the latency of the bus
transaction, the queueing, and whatever else the node was doing. Chapter 6
depends on this entirely.

The lower half is the reason the upper half is built the way it is. In the
published protocol, any node broadcasts an indication frame; every participant,
including the sender, records when it saw that frame; then the sender transmits
a second frame carrying its own record. Each receiver now has two numbers for
the same instant and can compute its offset. The time-critical path is reception
to timestamp, which is local to each node and can be characterised once, and on
this part the controller removes even that by capturing at the frame itself.

\diagram{j03_data}{The three time words this node deals in, drawn to the same
scale. The top one is the only one it computes with. The middle row is the
reason this chapter comes before the bus chapters: every hardware counter on
this part wraps inside an hour and two of them inside a minute, while a joint
node runs for weeks. The bottom two are the field bus profile's own layouts,
which chapter 11 has to map onto.}

\subsection*{Data flow (ASCII)}

\begin{asciiart}
  hardware                        software                     what it is for
  ------------------------------  ---------------------------  -----------------
  TIM2, 32 bit at 1 MHz    -----> lower 32 bits  --+
  TIM2 update interrupt    -----> upper 32 bits  --+--> uint64 us   intervals,
                                  read with retry                   ordering,
                                                                    deadlines
  TIM2 CH1 capture on edge -----> event stamp, same base ---------> sensor and
                                  (latency moves the read, not      bus events
                                   the recorded instant)

  RTC + 32.768 kHz crystal -----> date and time of day -----------> log entries
                                  subsecond register                 only

  DWT cycle counter        -----> 3.57 ns, wraps in 15.3 s -------> execution
                                                                    time only

  FDCAN timestamp (ch 9)   -----> captured at start of frame -----> agreement
                                  16 bit, extended in software      with other
                                                                    nodes
\end{asciiart}

\subsection*{Repository layout}

\begin{plaincode}
joint-node/
  ld/stm32h7a3zi.ld
  doc/
    node-identity.md  real-or-modelled.md  footprint.csv  period-report.md
    time-rules.md                       # + this chapter: which clock for what
  src/
    node/    identity.c state.h indicator.c selfcheck.c
    bsp/     startup.c system.c uart.c board.h
    control/ period.c jitter.c cyclecount.h loop.h
    time/
      mono.c  mono.h                    # + this chapter: 64-bit microseconds
      capture.c capture.h               # + this chapter: stamp at the pin
      wall.c  wall.h                    # + this chapter: the calendar
      skew.c  skew.h                    # + this chapter: slew, never step
    sense/  act/  bus/  mw/  safety/  update/
  host/  plot_jitter.py  count_edges.py
  test/
    test_jitter.c
    test_mono.c                         # + this chapter: the wrap and the race
    test_skew.c                         # + this chapter: monotonicity
  tools/ footprint.py
  README.md
\end{plaincode}

\subsection*{Steps}

\begin{steps}
\item \textbf{Write the rule down before the code.} One short file,
\texttt{doc/time-rules.md}, that every later chapter is held to. It is four
lines and it prevents a class of defect that is very hard to find later.
\begin{plaincode}
monotonic (mono_us)   intervals, deadlines, ordering, every timestamp in a
                      frame or a message. Never reset, never stepped.
wall clock (wall_now)  log lines and the date in a report. Never used to
                      compute an interval, ever.
cycle counter          execution time of a short section, inside one period.
                       Wraps in 15.3 s, so never used across periods.
bus timestamp          agreement with other nodes, from chapter 12.
\end{plaincode}

\item \textbf{Make the timer tick at exactly one microsecond.} The prescaler
comes from the derived timer clock, as in chapter 2, and the division has to be
exact or the node refuses to start. A clock that ticks every 1.0004 microseconds
is a clock that is wrong by a third of a second every ten minutes.
\begin{ccode}
void mono_init(void)
{
    uint32_t tclk = timer_clock_hz();            /* from chapter 2 */
    if (tclk % 1000000u != 0u) fail("mono: timer clock is not a whole MHz");
    TIM2->PSC = (tclk / 1000000u) - 1u;          /* one tick = 1 us */
    TIM2->ARR = 0xFFFFFFFFu;                     /* free running, 32 bit */
    TIM2->DIER |= TIM_DIER_UIE;                  /* update: extend the count */
    TIM2->CR1 |= TIM_CR1_CEN;
    printf("mono: %lu Hz timer, prescaler %lu, 1 us per tick\r\n",
           tclk, TIM2->PSC + 1u);
}
\end{ccode}

\item \textbf{Extend 32 bits to 64, and read it without a race.} This is the
part that looks trivial and is not. The counter and the overflow count are two
separate pieces of state that change at different moments, so a naive read can
catch the low word after the wrap and the high word before it, and produce a
timestamp that is 71.6 minutes in the past.
\begin{ccode}
static volatile uint32_t g_high;                 /* incremented in the ISR */

void TIM2_IRQHandler(void)
{
    if (TIM2->SR & TIM_SR_UIF) { TIM2->SR = ~TIM_SR_UIF; g_high++; }
}

uint64_t mono_us(void)
{
    uint32_t hi1, lo, hi2;
    do {
        hi1 = g_high;
        lo  = TIM2->CNT;
        hi2 = g_high;                            /* did it wrap while we read? */
    } while (hi1 != hi2);
    return ((uint64_t) hi2 << 32) | lo;
}
\end{ccode}
The retry loop terminates because the wrap happens once every 71.6 minutes and
the loop body takes tens of nanoseconds. The alternative, disabling interrupts
around the read, costs the period handler latency for no benefit, which is the
same trade chapter 2 made with the sequence counter.

\item \textbf{Prove the wrap on a host machine rather than by waiting.}
Seventy-one minutes is too long to test by hand, and the interesting cases are
exactly at the boundary. Factor the arithmetic so it can be called with
artificial values.
\begin{ccode}
/* test_mono.c, built and run on the host: no board involved */
assert(mono_compose(0u, 0xFFFFFFFFu) == 0xFFFFFFFFull);
assert(mono_compose(1u, 0x00000000u) == 0x100000000ull);
assert(mono_compose(1u, 0x00000000u) > mono_compose(0u, 0xFFFFFFFFu));
/* and the race the retry loop exists to prevent, written as a fact: */
assert(mono_compose(0u, 0x00000001u) < mono_compose(1u, 0x00000000u));
\end{ccode}

\item \textbf{Stamp an event at the pin.} Configure one timer channel as an
input capture. The capture register holds the counter value at the edge, so the
interrupt can be late without making the timestamp late.
\begin{ccode}
uint64_t capture_stamp(void)                     /* called from the capture ISR */
{
    uint32_t lo = TIM2->CCR1;                    /* the counter at the edge */
    uint64_t now = mono_us();                    /* the counter now */

    /* The capture happened at or before now. If the low word has wrapped
       since, the high word must be one less than the current one. */
    uint32_t hi = (uint32_t) (now >> 32);
    if (lo > (uint32_t) now) hi--;
    return ((uint64_t) hi << 32) | lo;
}
\end{ccode}
Print the difference between the captured stamp and a stamp taken at the top of
the interrupt handler. That difference is interrupt latency, measured on this
board, for free, and it is a number worth putting in the report.

\item \textbf{Keep the wall clock separate, and make it obvious.} The calendar
comes from the real-time clock on the fitted crystal, and the type system is
the cheapest place to stop somebody subtracting two of them.
\begin{ccode}
typedef struct { uint64_t us; }        mono_t;   /* intervals: subtract these */
typedef struct { uint32_t s; uint32_t frac; } wall_t;  /* dates: never subtract */

static inline uint64_t mono_diff(mono_t a, mono_t b) { return a.us - b.us; }
/* There is deliberately no wall_diff(). If you need an elapsed time, you
   wanted the monotonic clock. */
\end{ccode}

\item \textbf{Measure the two oscillators against each other.} The node now has
two independent time sources. Count monotonic microseconds between two
real-time-clock second boundaries, over an hour, and report the difference in
parts per million. This is the first genuinely physical measurement in the
volume and it costs nothing.
\begin{shellcode}
# after one hour, on the console
drift: mono 3600004120 us over 3600 RTC seconds -> +1.14 ppm (mono fast)
\end{shellcode}
Two honest readings of that number. It is the difference between two
oscillators, and it does not say which one is wrong. And with no absolute
reference on this bench it cannot be turned into an accuracy figure, which is
what a receiver emitting one pulse per second would add for about the price of a
meal.

\item \textbf{Implement slew, and refuse to step.} Chapter 12 will have an
offset to apply. Applying it as a jump is the mistake both cited papers warn
about, because time then goes backwards and two different events can carry the
same timestamp. Write the correction now, with the interface the bus chapter
will call, and assert monotonicity in the code.
\begin{ccode}
/* skew.c: absorb an offset by adjusting the rate, never by jumping. */
void skew_apply(int64_t offset_us)
{
    /* Spread the correction over the next window rather than applying it now.
       A positive offset means this node is behind: run slightly fast. */
    g_skew.residual_us = offset_us;
    g_skew.ppm = clamp(offset_us * 1000000 / SKEW_WINDOW_US, -200, +200);
}

uint64_t mono_us_corrected(void)
{
    uint64_t raw = mono_us();
    uint64_t out = raw + skew_correction(raw);
    if (out < g_last_out) { g_backwards++; out = g_last_out; }  /* never */
    g_last_out = out;
    return out;
}
\end{ccode}
The clamp is the design decision: a node that is told it is a minute out does
not spend a minute catching up at full speed, it reports the discrepancy and
corrects slowly, because a large step usually means the other machine is wrong.

\item \textbf{Design the protocol, and write it down for chapter 12.} Two
frames, one identifier each, at a low rate. The first marks an instant and
carries nothing. The second carries the sender's own timestamp for the first.
Each receiver subtracts its own record of the first frame from the number in the
second, and hands the difference to the slew function above. The 1994 paper puts
this at about twenty microseconds on a slower bus, with fewer than twenty
messages a second.
\begin{plaincode}
frame A  "mark"      no payload. Every node records mono_us() at reception.
frame B  "follow-up" 8 bytes: the sender's mono_us() for frame A.
         offset = sender_stamp - my_stamp_of_A  ->  skew_apply(offset)
\end{plaincode}
\end{steps}

\subsection*{Build, flash and debug}

\diagram{j03_timing}{Why a corrected clock is slewed and never stepped. Above, an offset applied as a jump: time goes backwards, and two events that happened at different moments are recorded at the same instant, which a real-time system must not allow. Below, the same offset absorbed by running slightly fast for a bounded window: the clock stays monotonic and the ordering of events is preserved throughout.}

\begin{shellcode}
cmake --build build -j && probe-rs run --chip STM32H7A3ZITx build/firmware.elf
ctest --test-dir build-host           # the wrap and monotonicity tests
\end{shellcode}

\begin{note}{When a timestamp is 71.6 minutes in the past}
That number is the signature of the race in step 3: the low word was read after
a wrap and the high word before it. If the stamps are wrong by exactly 4294.97
seconds, the retry loop is missing or the compiler has hoisted the read of the
overflow counter out of it, which is what the volatile qualifier is there to
prevent. A second signature, an offset of exactly one microsecond appearing and
disappearing, is the capture correction in step 5 applied in the wrong
direction.
\end{note}

\subsection*{Verification and acceptance criteria}

\begin{itemize}
\item The boot line prints the derived timer clock, the prescaler and the
resulting tick, and the firmware refuses to start if the division is not exact.
\item The monotonic clock never goes backwards: a check in the corrected read
counts violations, and the count is zero after an hour with the slew function
being exercised by injected offsets.
\item The wrap arithmetic has host tests covering the boundary, including the
case that motivates the retry loop.
\item An externally generated edge is stamped, and the difference between the
captured stamp and a stamp taken at the top of the handler is reported. That
difference is interrupt latency and it is a plausible number rather than zero.
\item The drift between the two oscillators is reported in parts per million
over at least an hour, with a sentence saying it does not identify which
oscillator is at fault.
\item Nothing in the tree computes an interval from the wall clock. There is no
function that would let it, and a search proves the absence.
\item Injecting a large offset makes the node report the discrepancy and correct
slowly rather than stepping.
\end{itemize}

\subsection*{Variants}

\begin{table}[H]\centering\small
\begin{tabular}{p{24mm}p{26mm}p{44mm}p{22mm}p{20mm}}
\toprule
Axis & Variant & What changes & Cost & Built in full in \\
\midrule
Time source & 32-bit timer extended in software & The baseline: one microsecond, effectively unbounded range & One interrupt every 71.6 minutes & Here \\
Time source & The cycle counter alone & Finest resolution available, 3.57 ns, and no extra peripheral & Wraps in 15.3 seconds, so it cannot time anything longer & Chapter 2, for execution time only \\
Time source & The real-time clock's subsecond register & Survives a reset and runs from a separate crystal & Coarse, and it is a calendar rather than a stopwatch & Here, as the wall clock \\
Time source & The bus controller's own capture & Captured at the frame, before any interrupt. Nothing in software can match it & Sixteen bits, wrapping in seconds, and it lives in a different clock domain from the timer that can drive it & Chapter 12 \\
Correction & Step the clock & One line, and it is what most first attempts do & Time goes backwards and ordering is lost & Nowhere. Shown as the defect \\
Correction & Slew at a bounded rate & Monotonic throughout, bounded catch-up time & The node is knowingly wrong for the length of the window & Here \\
Correction & Offset from a host round trip & The middleware client does it in four function calls & An offset only, with no rate correction, and it inherits the transport's asymmetry & Chapter 14 \\
Agreement & Precision time protocol & The standard answer everywhere else & Needs an Ethernet controller with hardware timestamping, which this part does not have & Nowhere. Named and explained \\
Agreement & Time-triggered bus & Timestamping in hardware at the protocol level & Lives in dedicated transceivers, absent from ordinary controllers and from the flexible-data format & Nowhere. Named and explained \\
\bottomrule
\end{tabular}
\caption{Variants for chapter 3. Two rows are built nowhere on purpose: naming the mechanism a reader will find everywhere else, and saying precisely why this bench cannot run it, is more useful than silence.}
\end{table}

\subsection*{Pitfalls}

\begin{itemize}
\item Reading a 32-bit counter and its overflow count without a retry. The
failure is rare, it is exactly 71.6 minutes wide, and the sensor will be
suspected long before the clock is.
\item Leaving the overflow counter without the volatile qualifier, so that the
compiler reads it once and the retry loop cannot terminate correctly.
\item Computing an interval from the wall clock. It is right until the moment
somebody sets the clock, and then it is spectacularly wrong.
\item Stamping a sensor sample when the read completes rather than when the
event happened. The error is the whole bus transaction and it varies with load.
\item Stepping a corrected clock. Two events then share a timestamp, and every
piece of ordering logic above is quietly wrong.
\item Assuming the timestamp counter in the bus controller runs from the same
clock as the timer that can drive it. They are in different domains and drift
against each other, which is documented in the vendor's own community answer and
is the kind of detail that costs a week.
\item Treating the drift figure as an accuracy figure. With no absolute
reference on the bench it is a comparison of two oscillators and nothing more.
\end{itemize}

\subsection*{Best practices applied}

\begin{itemize}
\item Two clocks with two purposes and two types, and no function exists that
would let one be used for the other's job.
\item The rule is a file in the repository, not a convention in somebody's head.
\item The arithmetic that cannot be tested by waiting is factored so it can be
tested on a host machine in microseconds.
\item Corrections are slewed and bounded, following the published practice in
both cited papers, and monotonicity is asserted in the code rather than assumed.
\item The measurement that this bench genuinely cannot make is named, along with
the inexpensive instrument that would make it possible.
\item Peripheral facts taken from a forum answer are confirmed against the
reference manual before they are printed, and the chapter says which is which.
\end{itemize}

\subsection*{Stretch goals}

\begin{itemize}
\item Add a receiver that emits one pulse per second, capture it on the same
timer channel, and turn the drift comparison into an absolute frequency
measurement with a stated uncertainty.
\item Discipline the monotonic clock to that pulse with the slew function
already written, and plot the residual over a day.
\item Implement the cycle counter and the microsecond timer as two
implementations of one interface, and measure the cost of a timestamp read in
each.
\item Run the same firmware in the open simulator and compare the behaviour of
the wrap. A simulator that does not model the wrap correctly is a useful thing
to discover before chapter 20 depends on it.
\end{itemize}

\subsection*{Roadmap and next steps}

Chapter 4 turns everything built so far into a board support package somebody
else can read, and hands it over as a file rather than as a folder of habits.

The published progression on time is short and worth following exactly. Read the
1994 conference paper first, because it is eight pages and contains the whole
idea. Then read the 2021 transactions paper, whose related-work section is a
complete survey of what has been tried over this bus and whose central result,
sub-microsecond agreement in software with rate correction, is the target this
volume's chapter 12 aims at. The bus association's own network time management
document is the current normative statement and is the thing to buy if this
work goes anywhere commercial. For the object dictionary vocabulary, the
application layer profile is free after registration and its synchronisation and
time clauses are four pages.

\subsection*{Portfolio evidence}

\begin{itemize}
\item The drift measurement, with its honest caveat, which is a real physical
result produced with no instrument at all.
\item The interrupt latency figure, derived from the difference between a
captured stamp and a handler stamp.
\item The host test suite for the wrap, including the boundary case, which is
the kind of test that convinces a reviewer that the author has met this bug
before.
\item The time rules file, which is four lines and is the thing most projects do
not have.
\end{itemize}

\subsection*{Sources}

Normative references:
\begin{itemize}
\item Reference manual RM0455, for the timer registers, the input capture path,
the real-time clock and the bus controller's timestamp registers. The register
names for the last of these come from a community answer and are confirmed here
before use.
\item The application layer profile of the bus association, version 4.2.0,
Monday 21 February 2011, clauses 7.2.5, 7.2.6 and 7.1.6.5, and objects 1005h,
1006h, 1007h, 1012h, 1013h and 1019h.
\item The bus association's network time management document, version 1.1.0,
Friday 1 September 2023, cited by number and title only.
\item IEEE 1588-2019 and the current edition of the local network timing
standard, named as the mechanisms this bench cannot run.
\end{itemize}

Reusable implementations:
\begin{itemize}
\item Gergeleit and Streich, on a distributed high-resolution real-time clock
over this bus, first international conference, \pubdate{September 1994}.\newline
\url{https://can-cia.org/fileadmin/cia/documents/proceedings/1994_gergeleit.pdf}
\item Einspieler, Rathakrishnan, Prabhakara, Steinwender and Elmenreich, on
high-accuracy software-based clock synchronisation over this bus, 2021.\newline
\url{https://mobile.aau.at/publications/einspieler-2021_High_Accuracy_SW-Based_Clock_Synchronization_Over_CAN.pdf}
\item The embedded middleware client's time synchronisation header,
Apache-2.0, for the epoch synchronisation used in chapter 14.\newline
\url{https://github.com/micro-ROS/rmw_microxrcedds}
\item The reference implementation of the object dictionary, Apache-2.0, whose
source settles the object numbers used above.\newline
\url{https://github.com/CANopenNode/CANopenNode}
\end{itemize}
   % One clock for sensors, loop and bus
\project{04}{The board support package, and a board file you can hand over}{A documented hardware interface}{Pin map, clock config, peripheral table, a board file another engineer can read}

\begin{keyfacts}
\item[Adds to the node] A hardware interface somebody else can take: one file that says what this board is, one that says what it is not, and a rule that keeps vendor types out of everything above
\item[Peripherals] All of them, described rather than configured. The pin budget for the whole volume is allocated here
\item[Depends on] Chapters 1, 2 and 3, whose accumulated hardware knowledge this chapter turns from habits into a deliverable
\item[Real or modelled] Entirely real, and deliberately so: a board file that describes hardware the bench does not have would be the worst possible place to blur that line
\item[Difficulty] 2 of 5, and the hour of licence reading is the part people skip
\item[Effort] Two evenings of about three hours
\item[Deliverable] A machine-readable board description with an evidence column, a generated header, a list of what this board does not have, a licence audit that fails the build on an unknown file, and a handover test that passes on a machine that has never seen the project
\end{keyfacts}

\subsection*{Why this chapter}

Three chapters in, the node knows a great deal about its own hardware, and all
of it is scattered: a pin number in a source file, a clock rule in a comment, a
peripheral choice recorded only in the fact that nothing else uses that timer.
That is a working board and an untransferable one. Every later chapter adds
another such fact, and by chapter 12 nobody can answer a simple question like
which pins are still free without reading the whole tree.

A board support package is the answer, and it is worth being precise about what
it is. It is not a hardware abstraction layer: the vendor supplies one of those
and it is used underneath. It is not a driver collection: drivers belong to the
peripherals they drive. It is the layer that says \emph{this board}, as opposed
to this silicon: which pin carries which function, which clock arrives at what
frequency, which peripherals are claimed and by what, and which capabilities
this board simply does not have.

The last of those is the part that most board files omit and the part that
matters most here. This board has no bus transceiver and this part has no
Ethernet controller at all. A board file that is silent about absences invites
the reader to assume presence, and the evidence for these two particular
absences is good enough to publish, which is what this chapter does.

\begin{note}{What a handover actually tests}
The test at the end of this chapter is not a build. It is: a second machine
that has never had a vendor development environment installed, a fresh clone, a
person who has read only the board file, and thirty minutes. If that person can
name every claimed pin, state the clock at three points in the tree, and get a
light blinking, the package is real. If they have to ask which timer is free,
it is not. This is also the most accurate rehearsal available for the question
an interviewer asks about board bring-up.
\end{note}

\subsection*{Prior art and what to reuse}

\begin{table}[H]\centering\small
\begin{tabular}{p{34mm}p{46mm}p{38mm}p{20mm}}
\toprule
Source & What it gives & What it does not & Licence \\
\midrule
The silicon vendor's device headers for this family & The authoritative, fetchable, machine-readable statement of what is in this part: peripheral instances, base addresses, register layouts, and the interrupt vector list. Every silicon claim in this volume is settled from here rather than from prose & No board layer at all, and no opinion about which pin does what on a particular board & Apache-2.0, safe to vendor \\
The same vendor's hardware abstraction driver & A peripheral layer that is already written, already tested by a great many people, and permissively licensed & It is silicon-level, so it knows nothing about this board, and its call surface leaks into anything that includes it, which is the problem this chapter solves & BSD-3-Clause, safe to vendor \\
The upstream operating system's board files for this board & The best published example of the idea this chapter is about: a board described as data rather than as code, in a format another tool can read & It carries one inherited artefact, discussed below, that claims a peripheral this part does not have. Treat a board file as evidence of what its author modelled & Apache-2.0 \\
The component licence table published inside one of that vendor's own family trees & The evidence for the rule in step 6: one vendor tree carries at least four different licences, two of them permissive and two of them vendor terms that survive redistribution & It covers a different family, so it is used as proof that per-file checking is needed rather than as a licence list for this one & the file itself is permissively licensed \\
The sibling firmware volume, chapter 1 & The startup code, the linker script and the clock bring-up, in full, for this exact part & It stops at one board and one blinking light, which is where this chapter starts & this series \\
\bottomrule
\end{tabular}
\caption{Prior art for chapter 4. The third row is the one to read in full: it is a working board file for this exact board, written by people who had to make it work for everyone, and the hour spent reading it is repaid in the pin map alone.}
\end{table}

What is left to write is the description itself and the discipline around it: an
evidence column on every row, a generated header so that the code cannot drift
from the description, an explicit list of absences, and a check that stops
vendor types from leaking upward.

\subsection*{What the node gains}

Before this chapter the node works and only its author can extend it. After it,
the node has a hardware interface that is a file rather than a memory: a
description with evidence, a generated header, a stated list of what the board
cannot do, and a licence audit that says exactly which third-party files may be
published with the project and which may only be referenced. Nothing the node
does changes. Everything about who else can work on it does.

\subsection*{Parts from the inventory}

\begin{table}[H]\centering\small
\begin{tabular}{p{40mm}p{66mm}p{40mm}}
\toprule
Part & Role & Interface \\
\midrule
NUCLEO-H7A3ZI-Q & The board being described & Micro USB to the host \\
The board user manual & The authority for every pin claim in the description. Each row of the board file names the page it came from & A PDF, opened by hand \\
The reference manual for this part & The authority for every clock and peripheral claim, and specifically not the manual for the better-known sibling part & A PDF, opened by hand \\
Raspberry Pi 4 & The handover machine: it has never had a vendor development environment and never will & Network \\
\bottomrule
\end{tabular}
\caption{Inventory items used in chapter 4. Two of the four are documents, which is unusual for a chapter in this volume and is the point of this one.}
\end{table}

\subsection*{System architecture}

\diagram{j04_arch}{The layers, with the licence of each. Two of them are the
author's, two are permissively licensed and may be published with the project,
and one category may not be vendored at all. The single rule that makes the
whole stack portable is drawn as a line: nothing above the board layer includes
a vendor header.}

The figure is also a warning about a common assumption. A vendor software tree
is not one licence. A component table published inside one such tree lists the
hardware abstraction layer, the device headers, the board support drivers and a
file system under one permissive licence; the architecture core under a second;
and three middleware libraries under the vendor's own terms, which permit source
redistribution but bind every downstream user to run the software only on that
vendor's silicon and forbid the files from being placed under open source terms.
A public repository with one permissive licence file at the top that appears to
cover the whole tree would not comply. Hence step 6.

\subsection*{Peripheral configuration}

This table is the peripheral budget for the entire volume, allocated once, here.
Later chapters fill in the rows marked reserved and may not take a row that is
already claimed without saying so in their own text.

\begin{table}[H]\centering\small
\begin{tabular}{p{22mm}p{26mm}p{24mm}p{30mm}p{30mm}}
\toprule
Peripheral & Claimed by & Clock source & Pins and function & Status \\
\midrule
RCC, PWR, FLASH & Chapter 1 & 8 MHz bypass to 280 MHz & None & Done \\
GPIOB, GPIOE & Chapter 1 & Peripheral bus & PB0, PE1, PB14, node state & Done \\
GPIOC & Chapters 1 and 18 & Peripheral bus & PC13, button, later the local stop & Done, extended later \\
USART3 & Chapter 1 & Peripheral bus & Console through the probe & Done \\
TIM6 & Chapter 2 & Timer clock, derived & None & Done \\
TIM2 & Chapter 3 & Timer clock, derived & CH1 input capture, event stamp & Done \\
RTC and LSE & Chapter 3 & 32.768 kHz crystal & None & Done \\
FDCAN1 & Chapters 9 to 13 & Kernel clock selector, see the data figure & Two pins on port D, confirmed against the board manual & Reserved \\
TIM1 or TIM8 & Chapter 7 & Timer clock & Complementary outputs, dead time, fault input & Reserved \\
A timer in encoder mode & Chapter 5 & Timer clock & Two channels, plus one interrupt pin for the index & Reserved \\
I2C on the shield header & Chapter 6 & Peripheral bus & Inertial unit, and its data-ready pin to a capture channel & Reserved \\
IWDG and WWDG & Chapter 18 & Low-speed internal, and the peripheral bus & None & Reserved \\
\bottomrule
\end{tabular}
\caption{The peripheral budget, allocated once. Two watchdogs are confirmed present on this part, an independent one and a window one, and chapter 18 explains why a joint node wants the window variant and why neither is independent enough for a safety argument on its own.}
\end{table}

\subsection*{Wiring}

\diagram{j04_wiring}{The pin budget, drawn on the two headers. Green is claimed and in use, amber is claimed and reserved for a later chapter, grey is free. A reader who wants to add something to this node looks at this figure first, and a chapter that takes a grey pin turns it amber here.}

\subsection*{Memory and timing budget}

\begin{table}[H]\centering\small
\begin{tabular}{p{44mm}p{28mm}p{28mm}p{30mm}}
\toprule
Quantity & Budget & Measured & Margin \\
\midrule
Flash, this chapter & 2 kB, almost all of it tables & not measured & not measured \\
Static memory, this chapter & 0 & computed & n/a \\
Generated header, lines & under 300 & not measured & not measured \\
Handover: first light on a new machine & 30 minutes & not measured & not measured \\
Handover: questions the reader had to ask & 0 & not measured & not measured \\
\bottomrule
\end{tabular}
\caption{The budget for chapter 4. The last two rows are the only ones in this volume measured in minutes and questions rather than in bytes, and they are the ones that decide whether this chapter succeeded.}
\end{table}

\subsection*{Firmware design (UML)}

\diagram{j04_uml}{The board layer as a contract. Everything above it uses named
functions and named pins; only the four files inside it include a vendor header,
and a check in the build proves that. The generated header is produced from the
description and never edited, which is why it is drawn with its generator
attached.}

Two rules, and the second one is the one that pays.

\textbf{The description is the source of truth and the header is generated from
it.} A pin number that appears in two places will eventually appear in two
different forms. Writing the description as data and generating the header
removes that possibility, and it makes the documentation and the code the same
artefact rather than two artefacts that agree for a while.

\textbf{Nothing above the board layer includes a vendor header.} The moment a
control loop includes a peripheral header, the control loop is tied to one
silicon vendor and cannot be tested on a host machine. Chapter 20 needs to build
most of this node on a host with no hardware at all, and that is possible only
if this rule was kept from chapter 4 onward rather than retrofitted at the end.

\subsection*{Data flow (ASCII)}

\begin{asciiart}
  authority                board description         generated            used by
  ---------------------    ---------------------     ---------------      -----------
  board manual UM2408 ---> board.yaml                                      nothing
    page and table no.        pins:                                        above the
                               led_green: PB0   ---> board_pins.h  ------> board layer
  reference manual RM0455 -->  console: USART3       (never edited)        includes
    clause number             clocks:                                      a vendor
                               sysclk: 280 MHz  ---> board_clocks.h -----> header
  vendor device headers  -->  absent:                                      |
    proof by absence           ethernet: false ----> board_caps.h  ------> |
                               transceiver: false     BOARD_HAS_* = 0      v
                                                                        a build
  measured on the bench  -->  drift: +1.14 ppm  ---> doc only            error if
    chapter 3                  (never a constant)                        anything
                                                                        else does
\end{asciiart}

\subsection*{Repository layout}

\begin{plaincode}
joint-node/
  board/
    nucleo_h7a3zi_q.yaml              # + this chapter: the description
    evidence.md                       # + this chapter: where each row came from
    LICENSES.md                       # + this chapter: the audit result
  src/
    bsp/
      board.h                         # + this chapter: the contract
      board_pins.h                    # + generated, never edited
      board_clocks.h                  # + generated, never edited
      board_caps.h                    # + generated: what this board is not
      board.c  clock.c  gpio.c  uart.c   # the only files that include vendor
      startup.c  system.c                 headers, from chapters 1 to 3
    node/  control/  time/
    sense/  act/  bus/  mw/  safety/  update/
  tools/
    gen_board.py                      # + this chapter: description to headers
    check_layering.py                 # + this chapter: no vendor types above bsp
    check_licences.py                 # + this chapter: every file has an SPDX tag
    footprint.py
  host/  test/  doc/
  README.md
\end{plaincode}

\subsection*{Steps}

\begin{steps}
\item \textbf{Write the description, with an evidence column.} One file, in a
format a script can read, where every row carries where it came from. Three
kinds of evidence are allowed and they are not equal: a document with a clause
or page number, a machine-readable vendor file, or a measurement taken on this
bench. Anything else is a question, not a row.
\begin{plaincode}
board: NUCLEO-H7A3ZI-Q
part:  STM32H7A3ZI
pins:
  led_green:   { pin: PB0,  evidence: "UM2408 table of LEDs" }
  led_yellow:  { pin: PE1,  evidence: "UM2408 table of LEDs" }
  led_red:     { pin: PB14, evidence: "UM2408 table of LEDs" }
  button:      { pin: PC13, evidence: "UM2408, user button" }
  console_tx:  { pin: PD8,  evidence: "UM2408, virtual COM port" }
  console_rx:  { pin: PD9,  evidence: "UM2408, virtual COM port" }
  event_in:    { pin: PA0,  evidence: "RM0455, TIM2_CH1 alternate function" }
  fdcan_tx:    { pin: PD1,  evidence: "TO CONFIRM against UM2408 before ch 9" }
  fdcan_rx:    { pin: PD0,  evidence: "TO CONFIRM against UM2408 before ch 9" }
clocks:
  hse:    { hz: 8000000,   source: "probe, bypass", evidence: "UM2408" }
  sysclk: { hz: 280000000, evidence: "RM0455 clock tree, computed at boot" }
  lse:    { hz: 32768,     fitted: true, evidence: "inspected on the board" }
\end{plaincode}
The two rows marked to confirm are the honest state of this bench today. They
are in the description because leaving them out would hide the question, and
they carry a marker the generator refuses to emit code for.

\item \textbf{Generate the header, and never edit it.} About sixty lines of
script. The generated file carries a banner saying it is generated, the name of
the description it came from, and nothing else.
\begin{shellcode}
python tools/gen_board.py board/nucleo_h7a3zi_q.yaml --out src/bsp/
# wrote board_pins.h (41 lines), board_clocks.h (12), board_caps.h (9)
# refused: fdcan_tx, fdcan_rx  (evidence marked TO CONFIRM)
\end{shellcode}
A generator that refuses to emit an unconfirmed row is worth the extra ten
lines. It turns a documentation gap into a build error at the moment somebody
tries to depend on it, which is exactly when it should surface.

\item \textbf{Write the contract.} The header everything above uses. Named pins,
named clocks, typed handles, and not one vendor type in the signatures.
\begin{ccode}
/* board.h: this board, as opposed to this silicon. */
#include "board_pins.h"     /* generated */
#include "board_clocks.h"   /* generated */
#include "board_caps.h"     /* generated */

typedef enum { BOARD_LED_GREEN, BOARD_LED_YELLOW, BOARD_LED_RED } board_led_t;

void     board_init(void);            /* clocks, power, flash latency, pins */
void     board_led(board_led_t l, bool on);
bool     board_button(void);
void     board_console_putc(char c);
uint32_t board_clock_hz(board_clock_t which);   /* read back, never assumed */
\end{ccode}

\item \textbf{State what the board is not, and prove it.} This is the part other
board files leave out. Two absences on this bench have evidence strong enough to
publish, and both are generated into a header so that any code assuming
otherwise fails at compile time rather than at bring-up.
\begin{ccode}
/* board_caps.h, generated. Absence is a fact about the board, not an omission. */
#define BOARD_HAS_ETHERNET       0   /* proof: the vector table, see below */
#define BOARD_HAS_CAN_TRANSCEIVER 0  /* proof: the board schematic, UM2408 */
#define BOARD_HAS_MOTOR          0
#define BOARD_HAS_TORQUE_SENSOR  0
#define BOARD_HAS_FIELDBUS_SLAVE 0
\end{ccode}
The Ethernet proof is worth writing out, because it is the cleanest example in
this volume of settling a question from a machine-readable source. In the
vendor's own device header for this part, the interrupt list runs from stream
four of the second transfer engine straight to the bus calibration entry: the
two vector slots that carry Ethernet on the better-known sibling are simply
absent. In the sibling's header, those same two slots are the Ethernet interrupt
and its wake-up interrupt, and the peripheral type is defined. Absence in a
generated vector list is not silence; it is a statement.

\item \textbf{Bring the clocks up in the one order that works, and read them
back.} Raise the core voltage, then set the flash wait states, then start the
phase-locked loop, then switch, then configure the peripheral clocks. Doing it
in any other order produces a board that runs at reset speed or does not run.
Then read every frequency back from the registers rather than trusting the
arithmetic.
\begin{ccode}
void board_init(void)
{
    pwr_set_voltage_scale(PWR_SCALE_0);       /* 1. voltage first          */
    flash_set_latency(FLASH_LATENCY_280MHZ);  /* 2. wait states before speed */
    rcc_pll_start(8, 140, 2, 2);              /* 3. 8 MHz / 2 * 140 / 2     */
    rcc_switch_sysclk(RCC_SRC_PLL1);          /* 4. and only now switch     */
    rcc_kernel_clocks();                      /* 5. per peripheral, see below */

    if (board_clock_hz(BOARD_CLK_SYS) != 280000000u)
        fail("board: system clock is not what the description says");
}
\end{ccode}

\item \textbf{Handle the one register whose name differs from the sibling
part.} This is the part trap made concrete, and it is a single line. The
selector for the bus peripheral's kernel clock sits at the same two bit
positions on both parts, in a register whose name reflects each part's power
domain naming. Everything else about that peripheral is identical, which is the
correction this volume carries.
\begin{ccode}
/* Correct for this part. The better-known sibling names this register after
   its three-domain layout; this part has two domains and names it after them.
   Same bits, same meaning, different symbol: a copied line will not compile,
   which is the best possible failure. */
static void rcc_kernel_clocks(void)
{
    MODIFY_REG(RCC->CDCCIP1R, RCC_CDCCIP1R_FDCANSEL, RCC_CDCCIP1R_FDCANSEL_0);
}
\end{ccode}

\diagram{j04_data}{The part trap in one register. The selector for the bus
peripheral's kernel clock sits at the same two bit positions on both parts, in a
register named after each part's power domains. A line copied from the sibling
does not quietly misbehave here: it fails to compile, which is the best failure
available. Everything else about that peripheral is identical, and that is the
correction this volume carries.}

\item \textbf{Audit the licences, and make the build fail on an unknown file.}
Every third-party file carries a machine-readable licence identifier, or it does
not go in the tree. The script prints a table and returns a non-zero status on
anything it cannot place.
\begin{shellcode}
python tools/check_licences.py --tree . --out board/LICENSES.md
# Apache-2.0     18 files   vendor device headers            publishable
# BSD-3-Clause   36 files   vendor hardware abstraction      publishable
# MIT             4 files   the real-time kernel             publishable
# (none)          0 files
# vendor terms    0 files   nothing under vendor licences is vendored here
# RESULT: clean
\end{shellcode}
The fifth row is the one to watch as the volume grows. Two of this vendor's own
licences permit source redistribution while binding every recipient to run the
software only on that vendor's silicon and forbidding the files from being
placed under open source terms. Referencing such a file is fine; fetching it at
build time is fine; committing it next to a permissive licence file is not.

\item \textbf{Enforce the layering, in about twenty lines.} A check that no file
outside the board layer includes a vendor header. It runs in the build, it is
trivial to write, and it is the single reason chapter 20 can test most of this
node on a host machine.
\begin{shellcode}
python tools/check_layering.py
# scanned 34 files outside src/bsp/
# vendor includes found: 0
# RESULT: clean
\end{shellcode}

\item \textbf{Run the handover.} Fresh clone on the Raspberry Pi, the board file
open, a timer running, and a written note of every question the reader had to
ask. Questions are the output of this test. A question means a row is missing
from the description, and the fix is a row rather than an explanation.
\end{steps}

\subsection*{Build, flash and debug}

\diagram{j04_timing}{The bring-up order, and what each wrong order produces.
This is the sequence that every chapter after this one inherits without thinking
about it, which is exactly why it is written down once, here, with the failure
mode of each step beside it.}

\begin{shellcode}
python tools/gen_board.py board/nucleo_h7a3zi_q.yaml --out src/bsp/
cmake -B build -DCMAKE_TOOLCHAIN_FILE=cmake/arm-none-eabi.cmake
cmake --build build -j && probe-rs run --chip STM32H7A3ZITx build/firmware.elf
\end{shellcode}

\begin{note}{When a board file describes a peripheral the silicon does not have}
It happens in published board files, including good ones, and this part has a
documented example. The upstream operating system's device tree for this part
inherits a shared family file and deletes the entries that do not apply, and one
Ethernet entry survives that deletion, so the file appears to configure a
controller that is not in the silicon. It is a modelling artefact rather than a
claim, and it is harmless upstream because nothing enables it. It is not
harmless to a reader who greps for what the part has. The rule this volume takes
from it: a board file is evidence of what its author modelled, and a device
header generated from the silicon description is evidence of what the silicon
contains. When the two disagree, the header wins.
\end{note}

\subsection*{Verification and acceptance criteria}

\begin{itemize}
\item The generated headers are reproducible: regenerating them from the
description produces no difference, and the build fails if they are edited by
hand.
\item The generator refuses to emit any row whose evidence is marked as
unconfirmed, and the two bus pins are currently in that state.
\item No file outside the board layer includes a vendor header, proven by the
check rather than by inspection.
\item Every third-party file carries a licence identifier, and the audit places
all of them in a publishable category.
\item The system clock, the timer clock and the peripheral bus clock are read
back from the registers at boot and compared with the description, and a
mismatch stops the node.
\item The capability header states at least the five absences above, and a file
that assumes one of them fails to compile with a readable message.
\item The handover completes on a machine with no vendor development
environment, within the budget, and the questions asked are recorded and turned
into rows.
\end{itemize}

\subsection*{Variants}

\begin{table}[H]\centering\small
\begin{tabular}{p{24mm}p{26mm}p{44mm}p{22mm}p{20mm}}
\toprule
Axis & Variant & What changes & Cost & Built in full in \\
\midrule
Board description & Data, with a generator & The baseline. One source of truth, and documentation and code cannot drift apart & A script to maintain & Here \\
Board description & A hand-written header & Nothing to generate, and the pin map lives in two places within a month & Drift & Nowhere. It is the thing this chapter replaces \\
Board description & The upstream device tree & A published, widely understood format, with tooling and an ecosystem, and the board becomes data that other software can consume & A second build system and a different driver model, and the artefact discussed above & Here, as the comparison \\
Board description & The vendor's configurator & Generates a working project quickly, and the generated output is several thousand lines pinned to one tool version & The output is not reviewed source, which this volume does not present as such & Nowhere. Named and explained \\
Language & C with a generated header & The baseline & None & Here \\
Language & C++ with constant expressions & The pin table becomes a compile-time structure, so an invalid pin is a compile error rather than a runtime one & A C++ toolchain in the build & Chapter 16 \\
Language & Rust with typed pins & The type system can make a pin's alternate function part of its type, so a misconfiguration cannot be expressed & A second toolchain, and this part is supported by the hardware crate rather than by an example & Chapter 16 \\
Layering & Vendor types in signatures & The quick way, and it ties every layer above to one silicon vendor & Host testing becomes impossible & Nowhere. The check exists to prevent it \\
\bottomrule
\end{tabular}
\caption{Variants for chapter 4. Three of the eight rows are built nowhere, which is unusual and deliberate: naming the common approaches and saying precisely what each costs is more useful to a reader choosing an approach than a single recommendation would be.}
\end{table}

\subsection*{Pitfalls}

\begin{itemize}
\item Inheriting a clock configuration or a linker script from the better-known
member of this silicon family. The failure is a board that does not boot, and it
looks like a hardware fault.
\item Assuming a vendor software tree carries one licence. A component table
inside one such tree lists four, two of which survive redistribution and bind
the recipient.
\item Letting a vendor type into a signature above the board layer. It is one
line, it is convenient, and it is the reason host testing gets abandoned.
\item Putting the pin map in comments. Comments are not checked by anything.
\item Treating a published board file as a list of the silicon's capabilities.
The example in the note above is real and current.
\item Configuring peripheral clocks before switching the system clock, or
raising the frequency before the voltage scale and the wait states.
\item Reading the frequency from a constant in your own code rather than from
the registers. The constant is a hope; the register is a fact.
\end{itemize}

\subsection*{Best practices applied}

\begin{itemize}
\item One source of truth for every board fact, with the evidence recorded
beside it and three grades of evidence distinguished.
\item Generated code is marked as generated and is never edited, so the
documentation and the implementation cannot disagree.
\item Absences are stated as explicitly as presences, with proof, because a
silent board file invites assumptions.
\item The layering rule is checked by a script rather than trusted, because it
is the rule most likely to be broken accidentally and it is the one everything
in chapter 20 depends on.
\item Licences are audited per file, and the categories are the ones from this
volume's source pool rather than a general impression.
\item An unconfirmed fact is visible in the description and produces a build
error at the point of use, rather than being quietly omitted.
\end{itemize}

\subsection*{Stretch goals}

\begin{itemize}
\item Generate an upstream device tree overlay from the same description, so
that the board can be described once and consumed by both build systems. That is
a genuinely useful small tool and no published one does exactly this for this
board.
\item Generate the pin budget figure in this chapter from the description too,
so that a chapter which claims a pin updates the figure by editing one line.
\item Add a check that every peripheral claimed in the description is either
used or explicitly reserved, so that an abandoned reservation is noticed.
\item Publish the board file on its own, with the evidence column intact, as a
small standalone contribution. It is the kind of artefact other people search
for.
\end{itemize}

\subsection*{Roadmap and next steps}

Chapter 5 gives the node a position: a quadrature decoder in hardware, and a
generator to test it with, since there is no encoder on this bench.

The published progression from here is the upstream project's own porting guide,
which is written for exactly this task and is worth reading even if this node
never runs that operating system: it separates the board from the silicon in a
way that clarifies what a board file should and should not contain. The vendor's
own documentation set is the other half, and the order to read it in is the
board user manual first, for the pins, then the reference manual's clock tree
chapter, and only then the peripheral chapters, because the clock tree is what
the rest depends on and it is the part that differs most from the sibling part
most material was written for.

\subsection*{Portfolio evidence}

\begin{itemize}
\item The board file itself, with its evidence column, published on its own.
\item The licence audit output, which is a short table that demonstrates a
habit most projects do not have.
\item The handover recording: a fresh clone on a machine with no vendor tooling,
a light blinking, and the list of questions the reader had to ask, including the
ones that became rows.
\item The layering check passing over a tree that includes a control loop, a
time base and a bus stack, which is the evidence that the rule was kept rather
than declared.
\end{itemize}

\subsection*{Sources}

Normative references:
\begin{itemize}
\item The board user manual UM2408, for every pin claim, the connectors, the
jumpers and the virtual serial port.
\item Reference manual RM0455, for the clock tree, the kernel clock selectors
and the peripheral register maps. Not the manual for the sibling part.
\item The datasheet for this part, for the memory sizes and the clock limits at
each voltage scale.
\end{itemize}

Reusable implementations:
\begin{itemize}
\item The vendor's device headers for this family, Apache-2.0, which settle
what this part contains and, by absence, what it does not.\newline
\url{https://github.com/STMicroelectronics/cmsis-device-h7}
\item The vendor's hardware abstraction driver for this family,
BSD-3-Clause.\newline
\url{https://github.com/STMicroelectronics/stm32h7xx-hal-driver}
\item The upstream operating system's board documentation for this exact board,
Apache-2.0, which is the model for a board described as data.\newline
\url{https://docs.zephyrproject.org/latest/boards/st/nucleo_h7a3zi_q/doc/index.html}
\item A component licence table published inside one of the vendor's own family
trees, which is the evidence that per-file checking is required.\newline
\url{https://github.com/STMicroelectronics/STM32CubeWB}
\end{itemize}
   % The board support package, and a board file you can hand over

\part{Sensing and acting}
\project{05}{The encoder: quadrature in hardware, and one you generate}{Position}{Timer encoder mode, index, velocity from differences, generating quadrature to test the decoder}

\begin{keyfacts}
\item[Adds to the node] Position, and the first derived quantity in the volume: velocity. Also the first stand-in, and the chapter says so before it says anything else
\item[Peripherals] One timer in encoder mode for the decode, a second timer generating two phase-shifted signals as the stand-in source, one capture channel for the index
\item[Depends on] Chapter 2 for the period, chapter 3 for the time base and the capture technique, chapter 4 for the pin budget
\item[Real or modelled] \textbf{Half real.} The decode path is real hardware and is exercised at full rate. The source is generated on the same board, so the position is a count rather than a measurement of anything that turns
\item[Difficulty] 3 of 5
\item[Effort] Three evenings of about three hours, one of them on the index
\item[Deliverable] A hardware quadrature decoder with a software extension that survives direction changes, an index latched without a dedicated peripheral, a velocity estimate with its quantisation stated, and a generator that proves the decoder against a known count
\end{keyfacts}

\subsection*{Why this chapter}

A joint node without position is a timer with opinions. Everything from here
depends on this chapter: the controller in chapter 16 acts on the error between
a commanded position and this one, the state frame in chapter 11 carries it, and
the tracking assertions in chapter 20 are computed from it.

There is no encoder on this bench, and this is the chapter where that stops
being a detail and becomes a design decision. The honest options were to skip
the subject, to buy a part, or to build the half that is real and generate the
other half. This volume takes the third, for a reason worth stating: the decode
path is where all the interesting firmware lives. Counting quadrature correctly
across direction reversals, extending a sixteen-bit counter that can run
backwards, latching an index without a peripheral that supports one, and getting
a velocity out of a position that only changes in whole counts are all real
problems with real solutions, and none of them needs a shaft to turn.

What is generated is the source. A second timer produces two square waves ninety
degrees apart and an index pulse, those pins are wired back to the decoder's
pins with two jumper wires, and the decoder sees quadrature that is
indistinguishable from an encoder's except in one respect: it is perfect. A real
encoder contributes mechanical eccentricity, electrical noise, missed edges at
speed and a quadrature phase that is never exactly ninety degrees. The chapter
says which of those the stand-in cannot exercise, and it injects the two that
can be injected.

\begin{note}{A correction to chapter 1}
Chapter 1's table of what is real and what is modelled put position in the
fully real column. That was wrong, and this chapter is where it is corrected:
position is half real, in the same sense that the actuator in chapter 7 is half
real. The decode path is hardware and is exercised at full rate; the source is
generated. The table in \texttt{doc/real-or-modelled.md} and the boot banner's
modelled line are both updated here, which is what that file is for. A volume
that discovers a boundary it drew in the wrong place and quietly moves it has
lost the argument it was making.
\end{note}

\subsection*{Prior art and what to reuse}

\begin{table}[H]\centering\small
\begin{tabular}{p{34mm}p{46mm}p{38mm}p{20mm}}
\toprule
Source & What it gives & What it does not & Licence \\
\midrule
The vendor's hardware abstraction driver for this family & The authoritative statement of what this part's encoder interface offers: exactly three encoder modes, named for which input's edges are counted. Everything in step 2 is checked against this header rather than against a tutorial & The mapping from those three modes to one, two and four counts per cycle is standard reference-manual semantics and was not read in the manual for this part. It is marked as such in the text & BSD-3-Clause \\
The vendor's own community answer of Saturday 14 March 2020 & The finding this chapter is shaped around, from a member of the vendor's staff: the hardware index feature is new, the first family to carry it is a different one, and it was not available when this family was designed & A forum answer, so it is corroborated below by the absence of the corresponding symbols in the vendor's own headers & forum post \\
The field oriented motor control project & Real robot code that reads encoders in the same loop that drives a bridge, including the velocity estimate and its filtering, and a permissively licensed driver for one common magnetic position sensor & Arduino shaped, and its velocity path makes different trade-offs from this one, which is worth comparing rather than copying & MIT \\
The sibling firmware volume, on filters and transforms & The filter this chapter's velocity estimate needs, already built and measured on this exact part & It is a filter chapter, not a position chapter & this series \\
A permissively licensed driver for one magnetic angle sensor & A clean driver, if that part is ever bought & \textbf{That sensor has no quadrature output and no serial peripheral interface}, so it cannot feed the timer's encoder interface at all. It is named here to stop a reader buying it for this purpose & MIT \\
\bottomrule
\end{tabular}
\caption{Prior art for chapter 5. The last row is a negative finding worth as much as the positive ones: a popular, inexpensive magnetic sensor that appears in a great many encoder tutorials cannot be used with this peripheral, and knowing that before ordering is worth more than a driver.}
\end{table}

What is left to write is the extension across direction reversals, the index
without a peripheral for it, the velocity arithmetic and its stated error, and
the generator. No published source was found that presents a quadrature
generator built on the same part as the decoder for the purpose of testing it,
so that part is this volume's own construction.

\subsection*{What the node gains}

Before this chapter the node knows when. After it, the node knows where: a
position in counts that survives direction changes and counter wraps, an index
that turns a relative count into an absolute one, and a velocity with an honest
statement of how much of it is quantisation. What it does not gain is a
measurement of anything mechanical, and the boot banner now says so.

\subsection*{Parts from the inventory}

\begin{table}[H]\centering\small
\begin{tabular}{p{40mm}p{66mm}p{40mm}}
\toprule
Part & Role & Interface \\
\midrule
NUCLEO-H7A3ZI-Q & Both halves: one timer generates the quadrature, another decodes it & Micro USB to the host \\
Two male-to-male jumper wires & The loopback from the generator pins to the decoder pins. This is the entire wiring of the chapter & Header to header \\
One more jumper wire & The index pulse, generator to capture input & Header to header \\
The power profiler, as a logic recorder & An optional external check that the generated waveform is what the firmware thinks it is & Two digital inputs \\
\bottomrule
\end{tabular}
\caption{Inventory items used in chapter 5. The bench has no rotary encoder, which is the subject of the note above and of the wiring figure below.}
\end{table}

\subsection*{System architecture}

\diagram{j05_arch}{The decode path, the generated source, and the place a real
encoder would plug in. Everything solid is real hardware exercised at full rate.
The dashed generator is on the same silicon as the decoder, which is the whole
trick and also the whole limitation: it cannot produce the failures a real
encoder produces, so those are injected deliberately in step 8.}

\subsection*{Peripheral configuration}

\begin{table}[H]\centering\small
\begin{tabular}{p{22mm}p{26mm}p{24mm}p{30mm}p{30mm}}
\toprule
Peripheral & Mode & Clock source & Pins and function & Interrupt and transfers \\
\midrule
Decode timer & Encoder interface, counting on both inputs & Peripheral bus timer clock & Two channels, A and B & Update interrupt, to extend the count \\
Same timer, third channel & Input capture on the index edge & As above & One channel, Z & Capture interrupt, low priority \\
Generator timer & Two compare outputs, one quarter period apart & Peripheral bus timer clock & Two pins, plus one for the index & Update interrupt, to emit the index \\
GPIO & Alternate function & Peripheral bus & Five pins, all on the free list in chapter 4 & None \\
\bottomrule
\end{tabular}
\caption{Peripheral configuration. Two timer instances are named as roles rather than by number, because which of this part's timers offer encoder mode, which are thirty-two bit, and which are still free once a shield is fitted are on chapter 4's confirm list. The selection rule is in step 1 and the chosen numbers go into the board description, not into this table.}
\end{table}

\subsection*{Wiring}

\diagram{j05_wiring}{Three jumper wires and a connector that is not there. The
loopback is the chapter's only wiring. The right half of the figure is what a
real encoder brings and what would have to be checked before plugging one in,
which is the paragraph to read before buying anything.}

Two cautions for the day an encoder is bought. Many industrial encoders are five
volt parts and output five volt levels, and whether a given pin on this part
tolerates that is a per-pin property in the datasheet rather than a property of
the board: check the exact pins, not the family. And an encoder's index output
is sometimes open collector, which needs a pull-up that the generated stand-in
does not.

\subsection*{Memory and timing budget}

\begin{table}[H]\centering\small
\begin{tabular}{p{44mm}p{28mm}p{28mm}p{30mm}}
\toprule
Quantity & Budget & Measured & Margin \\
\midrule
Flash, this chapter & 5 kB & not measured & not measured \\
Static memory, this chapter & 64 bytes & not measured & not measured \\
Decoder cost in the period handler & under 200 cycles & not measured & not measured \\
Maximum count rate decoded & 1 Mcount/s & not measured & not measured \\
Position error against the generator & 0 counts, exactly & not measured & not measured \\
Velocity quantisation at 1 rev/min & computed, see step 6 & computed & n/a \\
\bottomrule
\end{tabular}
\caption{The budget for chapter 5. The fifth row is the one that makes this chapter worth doing: because the source is generated, the number of counts that should have arrived is known exactly, so the acceptance criterion is equality rather than a tolerance.}
\end{table}

\subsection*{Firmware design (UML)}

\diagram{j05_uml}{The three modules and what each is allowed to know. The decoder
does not know that the source is generated, which is the property that lets a
real encoder replace the generator later without touching it. The test harness
knows both, which is why it can assert equality.}

The design rule here is the one that makes the stand-in honest. The decoder is
written against a quadrature signal, not against the generator: it has no
reference to it, no shared state with it, and no way to ask it anything. The
harness above them both knows how many counts were generated and how many were
decoded, and asserts they are equal. When a real encoder arrives, the generator
is removed, the harness loses its reference count, and the decoder does not
change at all.

\subsection*{Data flow (ASCII)}

\begin{asciiart}
  generator timer                        decoder timer
  +---------------------------+          +--------------------------------+
  | CH1 ___|---|___|---|___    |  wire A  | TI1 -> encoder interface       |
  | CH2 _|---|___|---|___|_    |--------->| TI2 -> both edges of both      |
  |      quarter period apart  |  wire B  |         inputs counted         |
  |                            |--------->|   |                            |
  | CH3 index, one pulse per   |  wire Z  |   v                            |
  |     N counts               |--------->| CNT, 16 bit, counts up or down |
  +---------------------------+          |   |                            |
        |                                 |   +--> CH3 capture on index    |
        | the harness knows this number   |   |                            |
        v                                 |   v                            |
  expected_counts                        | extend to 64 bit, signed        |
        |                                 |   |                            |
        |                                 |   v                            |
        |                                position_counts ---> velocity     |
        |                                 |                     (step 6)   |
        +---------------> assert equal <--+                                |
                                          +--------------------------------+
\end{asciiart}

\subsection*{Repository layout}

\begin{plaincode}
joint-node/
  board/  nucleo_h7a3zi_q.yaml            # + five pins, with evidence
  doc/    real-or-modelled.md             # + corrected: position is half real
  src/
    bsp/     board.h + generated headers
    node/    identity.c state.h indicator.c selfcheck.c
    control/ period.c jitter.c loop.h
    time/    mono.c capture.c wall.c skew.c
    sense/
      encoder.c  encoder.h                # + this chapter: the decoder
      quadgen.c  quadgen.h                # + this chapter: the stand-in source
      velocity.c velocity.h               # + this chapter: and its error
    act/  bus/  mw/  safety/  update/
  test/
    test_encoder.c                        # + extension across reversal
    test_velocity.c                       # + quantisation, on the host
  host/  tools/  README.md
\end{plaincode}

\subsection*{Steps}

\begin{steps}
\item \textbf{Choose the two timers, and record the rule rather than the
answer.} The decode timer needs encoder mode, a free pair of channels on free
pins, and preferably a thirty-two bit counter. The generator needs two compare
channels and does not care about width. Three of those properties are on chapter
4's confirm list for this part, so the selection rule goes in the text and the
chosen instances go into the board description with their evidence.
\begin{plaincode}
decode timer:    encoder mode required
                 32-bit counter preferred, 16-bit acceptable with the extension
                 two channels free after the shield in chapter 6 is fitted
                 a third channel free for the index capture
generator timer: two compare channels, any width
                 not TIM2 (the monotonic clock) and not TIM6 (the period)
\end{plaincode}

\item \textbf{Configure the encoder interface, and know which mode you chose.}
This part offers exactly three encoder modes, and the vendor's own header for
this family names all three: count on the edges of the first input, count on the
edges of the second, or count on the edges of both. The third gives four counts
per quadrature cycle and is what a joint node wants, because the resolution is
free.
\begin{ccode}
/* encoder.c: the three modes this part offers, from the vendor's own header.
   TI12 counts both edges of both inputs: four counts per cycle. */
void encoder_init(void)
{
    TIM_Encoder_InitTypeDef enc = {
        .EncoderMode = TIM_ENCODERMODE_TI12,   /* the other two: TI1, TI2 */
        .IC1Polarity = TIM_ICPOLARITY_RISING,
        .IC1Selection = TIM_ICSELECTION_DIRECTTI,
        .IC1Filter = 4,                        /* see the pitfalls */
        .IC2Polarity = TIM_ICPOLARITY_RISING,
        .IC2Selection = TIM_ICSELECTION_DIRECTTI,
        .IC2Filter = 4,
    };
    HAL_TIM_Encoder_Init(&htim_enc, &enc);
    HAL_TIM_Encoder_Start(&htim_enc, TIM_CHANNEL_ALL);
}
\end{ccode}
The relationship between those three mode names and one, two or four counts per
cycle is standard reference manual semantics. It was not read in the manual for
this part during the research for this volume, so it is confirmed on the bench
in step 7 by generating a known number of cycles and counting what arrives. That
is a better check than a citation anyway.

\diagram{j05_data}{Quadrature, and where the counts come from. Four states in a
fixed order, and the order is the direction: a decoder does not need to know
which way the shaft is turning, it only needs to know which state followed
which. Every marked edge is one count in the mode this node uses, which is four
per cycle and costs nothing extra. The index is the only thing in the whole
signal that can recover an absolute position after a dropped edge.}

\item \textbf{Extend the counter, and handle the direction.} This looks like
chapter 3's extension and it is not, because an encoder counter runs both ways.
An overflow can be an increment past the top or a decrement past zero, and the
direction bit in the timer says which.
\begin{ccode}
static volatile int32_t g_wraps;          /* signed: it can go negative */

void ENCODER_TIM_IRQHandler(void)
{
    if (TIM_ENC->SR & TIM_SR_UIF) {
        TIM_ENC->SR = ~TIM_SR_UIF;
        /* The direction bit tells you which way the counter crossed. Getting
           this wrong costs one full counter range per reversal at the wrap,
           which is a position error of 65536 counts appearing from nowhere. */
        if (TIM_ENC->CR1 & TIM_CR1_DIR) g_wraps--;   /* counting down */
        else                            g_wraps++;   /* counting up   */
    }
}

int64_t encoder_position(void)
{
    int32_t w1, w2; uint32_t cnt;
    do { w1 = g_wraps; cnt = TIM_ENC->CNT; w2 = g_wraps; } while (w1 != w2);
    return ((int64_t) w2 << 16) + (int64_t) cnt;     /* 16-bit counter */
}
\end{ccode}

\item \textbf{Latch the index in hardware, not in an interrupt handler.} This
part has no hardware index support: the vendor states that the feature arrived
with a later family and was not available when this one was designed, and the
corresponding configuration type and function are simply absent from the headers
for this part while being present for the family that has them. The interrupt
route works and is late by whatever the interrupt latency is, which at speed is
a position error. The better route uses the same technique as chapter 3: put the
index on a third channel of the same timer, as an input capture, so the hardware
records the counter value at the edge.
\begin{ccode}
/* The capture register holds CNT at the index edge, whenever the handler runs. */
void ENCODER_CAPTURE_IRQHandler(void)
{
    uint32_t at_index = TIM_ENC->CCR3;        /* the count when Z went high */
    g_index_seen = true;
    g_index_count = ((int64_t) g_wraps << 16) + at_index;
    g_offset = g_index_count;                 /* zero is now the index mark */
}
\end{ccode}
Whether a capture channel behaves this way while the same timer is in encoder
mode is the one thing in this chapter that must be confirmed on the bench before
it is trusted, and the confirmation is easy: generate the index at a known
count and compare. It goes on chapter 4's confirm list with that test named.

\item \textbf{Build the generator.} One timer, two compare channels at fifty per
cent duty, the second delayed by a quarter of the period. Reversing the
direction is a matter of swapping which channel leads.
\begin{ccode}
/* quadgen.c: two square waves a quarter period apart, and an index every N. */
void quadgen_set(int32_t counts_per_second, bool forward)
{
    uint32_t period = quadgen_period_for(counts_per_second);
    TIM_GEN->ARR  = period - 1u;
    TIM_GEN->CCR1 = period / 2u;                       /* channel A */
    TIM_GEN->CCR2 = forward ? (period / 4u)            /* B lags A  */
                            : (3u * period / 4u);      /* B leads A */
    g_expected_direction = forward ? +1 : -1;
}
\end{ccode}

\item \textbf{Get a velocity, and state its error before quoting it.} Position
changes in whole counts, so the difference over one control period is an
integer. At a low speed that integer is zero or one, and the resulting velocity
is mostly quantisation. The arithmetic is worth doing once, in the chapter,
rather than discovering it in chapter 16.
\begin{ccode}
/* The counting method: counts in one period. Simple, and quantised. */
float velocity_counting(int64_t pos, int64_t prev, float dt)
{
    return (float) (pos - prev) / dt;        /* one count in 1 ms = 1000 c/s */
}

/* The timing method: the interval between edges, using the capture from
   chapter 3. Fine at low speed, useless at high speed when edges arrive
   faster than they can be timestamped. */
float velocity_timing(uint64_t t_now_us, uint64_t t_prev_us, int32_t counts)
{
    uint64_t dt_us = t_now_us - t_prev_us;
    return dt_us ? (float) counts * 1000000.0f / (float) dt_us : 0.0f;
}
\end{ccode}
With four thousand counts per revolution and a one millisecond period, one count
per period is fifteen revolutions per minute. Below that speed the counting
method reports either zero or fifteen and nothing in between, which is not a
velocity, it is a quantiser. The node uses the counting method above a threshold
and the timing method below it, and reports which one produced each figure.

\item \textbf{Prove the decoder against a known count.} This is the acceptance
test the generated source makes possible, and it is stronger than anything a
real encoder would allow: the number of counts that should have arrived is
known exactly, so the criterion is equality.
\begin{shellcode}
python host/encoder_sweep.py --rates 100,1000,10000,100000,1000000 --seconds 10
# rate      generated    decoded   error   direction
#    100         1000       1000       0   forward  PASS
#   1000        10000      10000       0   forward  PASS
#  10000       100000     100000       0   forward  PASS
# 100000      1000000    1000000       0   forward  PASS
#1000000     10000000    9999872    -128   forward  FAIL: input filter too slow
\end{shellcode}

\item \textbf{Inject the failures the generator cannot produce naturally.} A
perfect source is the limitation of this whole approach, so the two failures
that can be simulated are simulated deliberately: a missing edge, by suppressing
one compare event, and a direction reversal exactly at the counter wrap, which
is the case step 3 exists for.
\begin{shellcode}
python host/encoder_faults.py --drop-one-edge --at-count 32768
# decoded position after a dropped edge: off by 1 count, permanently: PASS
python host/encoder_faults.py --reverse-at-wrap
# position continuous across the reversal at the wrap: PASS
\end{shellcode}
The first result is the important one and it is not a defect: a quadrature
decoder cannot detect a missing edge, so the error is permanent until the next
index pulse. That is exactly what the index is for, and it is why an incremental
encoder without an index is a poor choice for a joint.

\item \textbf{Correct the honesty table and the banner.} One row changes from
real to half real, and the boot line that names modelled quantities gains its
first entry.
\begin{shellcode}
joint-node-1  node 3  v0.5.0  built 2026-09-21
modelled: position source (generated quadrature, see chapter 5)
\end{shellcode}
\end{steps}

\subsection*{Build, flash and debug}

\diagram{j05_timing}{Velocity from a position that only moves in whole counts.
Above, the counting method at three speeds: at the lowest it reports zero or one
count per period and nothing in between. Below, the timing method on the same
motion, using the capture timestamps from chapter 3. The crossover is a design
decision and the node reports which method produced each figure.}

\begin{shellcode}
cmake --build build -j && probe-rs run --chip STM32H7A3ZITx build/firmware.elf
python host/encoder_sweep.py --rates 100,1000,10000,100000 --seconds 10
\end{shellcode}

\begin{note}{When the position jumps by exactly 65536}
The extension in step 3 counted a wrap in the wrong direction. It appears only
when the shaft reverses within one counter range of the wrap point, which on a
bench that only ever runs forward may never happen at all, and on a joint that
oscillates around a position happens constantly. A second signature, a position
that drifts by a few counts per reversal rather than jumping, is the input
filter: too short and the decoder counts noise, too long and it misses genuine
edges at speed. The sweep in step 7 is how the filter length gets chosen.
\end{note}

\subsection*{Verification and acceptance criteria}

\begin{itemize}
\item Decoded counts equal generated counts exactly, at every rate up to the one
where the input filter begins to miss edges, and that rate is recorded.
\item Reversing direction at an arbitrary point changes the sign of the velocity
and leaves the position continuous.
\item Reversing direction within one counter range of the wrap leaves the
position continuous, proven by the injected case rather than by inspection.
\item The index capture yields the same count every revolution, to zero counts,
and the difference between the captured count and one read in the handler is
reported as the latency that technique avoids.
\item The velocity crossover is measured rather than assumed: the speed below
which the counting method is quantisation is computed from the counts per
revolution and the period, and the reported figure names its method.
\item A dropped edge produces a permanent one-count error that the next index
pulse corrects, and the node reports the correction rather than hiding it.
\item The boot banner names the generated position source, and
\texttt{doc/real-or-modelled.md} has been corrected.
\end{itemize}

\subsection*{Variants}

\begin{table}[H]\centering\small
\begin{tabular}{p{24mm}p{26mm}p{44mm}p{22mm}p{20mm}}
\toprule
Axis & Variant & What changes & Cost & Built in full in \\
\midrule
Decode & Timer encoder interface & The baseline. The counting is in hardware and costs the processor nothing at any speed & One timer, two pins & Here \\
Decode & Two interrupts and software & Works on any part with two spare pins, and the processor pays for every edge & At a million counts a second it is the whole processor & Nowhere. Named as what the hardware saves you \\
Resolution & One, two or four counts per cycle & The three modes this part offers. Four is free and is the default here & None & Here \\
Index & Input capture on a third channel & The count at the edge, in hardware, whenever the handler runs & One channel & Here \\
Index & External interrupt and read & The obvious route, and it is late by the interrupt latency, which at speed is a position error & Accuracy at speed & Here, as the comparison \\
Index & A hardware index in the timer & The clean answer, and this part does not have it: the vendor states it arrived with a later family & Not available & Nowhere, and the evidence is in step 4 \\
Velocity & Counting in a fixed period & Simple, and quantised at low speed & Resolution at low speed & Here \\
Velocity & Timing between edges & Fine at low speed, and it degrades where the first improves & Complexity, and a timestamp per edge & Here \\
Velocity & An observer over both & The usual industrial answer: a model that fuses position with a command & A plant model, which arrives in chapter 7 & Chapter 16 \\
Source & Generated on the same board & The stand-in. Exact, repeatable, and unable to produce a real encoder's failures & The honesty rule applies & Here \\
Source & A magnetic sensor over a serial bus & An angle rather than a count, and absolute rather than incremental & It does not feed this peripheral at all, so the decoder is unused & Nowhere. See the prior art table \\
\bottomrule
\end{tabular}
\caption{Variants for chapter 5. The two index rows are the pair worth reading together: the technique this part supports, and the one it does not, with the evidence for the second.}
\end{table}

\subsection*{Pitfalls}

\begin{itemize}
\item Extending the counter without consulting the direction bit. The error is
one full counter range and it appears only on a reversal near the wrap.
\item Choosing the input filter length by taste. Too short counts noise as
motion; too long misses edges at speed. The sweep chooses it.
\item Reading the counter and the wrap count without a retry. Chapter 3 made
this mistake unavailable; the same discipline applies here.
\item Reading the index in a handler and calling the counter value at that
moment the index position. At a million counts a second, interrupt latency is
tens of counts.
\item Quoting a velocity at low speed without saying which method produced it.
Below one count per period the counting method reports a number that is entirely
an artefact of quantisation.
\item Assuming a popular magnetic angle sensor can drive this peripheral.
Several cannot: they provide an angle over a serial interface and no quadrature
output at all.
\item Buying a five volt encoder without checking the tolerance of the exact
pins, which is a per-pin property in the datasheet.
\item Forgetting that the loopback proves the decoder and not the sensing. The
figure and the banner both say so, and the chapter says it three times because
it is the one claim a reader might overstate.
\end{itemize}

\subsection*{Best practices applied}

\begin{itemize}
\item The stand-in is separated from the thing it stands in for by an interface,
so the real part can replace it without touching the code under test.
\item The acceptance criterion is equality rather than a tolerance, which is
possible only because the source is known exactly, and that is the one advantage
a generated source has over a real one.
\item The failures the stand-in cannot produce are named, and the two that can
be injected are injected.
\item A claim taken from a forum answer is corroborated against the vendor's own
headers before it is built on.
\item A semantic detail not confirmed in the reference manual is tested on the
bench instead, and the text says which it is.
\item An earlier chapter's classification is corrected in the open, in the file
that exists for it, rather than quietly adjusted.
\end{itemize}

\subsection*{Stretch goals}

\begin{itemize}
\item Drive the generator from a recorded trajectory rather than a constant
rate, so the decoder is exercised with accelerations and reversals that resemble
a joint rather than a test rig. Chapter 20 needs exactly this.
\item Measure the maximum decode rate against input filter length and plot the
two against each other. That figure does not appear to be published for this
part.
\item Add a second decoder instance on another timer and decode the same signal
twice, which is the shape of a redundant position channel in a safety related
design, and measure how often the two disagree.
\item Emit the quadrature with a transfer engine instead of compare channels, so
the generator can play an arbitrary pattern including deliberately malformed
quadrature.
\end{itemize}

\subsection*{Roadmap and next steps}

Chapter 6 adds the inertial unit, which is the first sensor on this bench that
measures something physical, and it uses chapter 3's capture technique to stamp
its samples at the pin.

The published progression from here is the motor control project cited above:
its encoder handling, its velocity filtering and its angle tracking are real
robot code in the same idiom, and reading its velocity path next to this
chapter's is the fastest way to see the trade-offs a shipping project makes. The
vendor's timer application note is the other half and is currently an existence
claim only in this volume: it could not be fetched during the research and is
cited by number and title without anything being quoted from it.

\subsection*{Portfolio evidence}

\begin{itemize}
\item The sweep table, showing exact equality up to a measured rate and the rate
at which the input filter starts to matter.
\item The reversal-at-wrap test, which is the defect most quadrature code has
and most test benches never exercise.
\item The index latency comparison: the same index, captured in hardware and
read in a handler, with the difference in counts at a stated speed.
\item The corrected honesty table, and the boot banner naming its first modelled
quantity. A reader who sees a project correct its own claim upward in rigour
learns more about the author than a clean table would have told them.
\end{itemize}

\subsection*{Sources}

Normative references:
\begin{itemize}
\item Reference manual RM0455, for the encoder interface, the direction bit, the
input filter settings and whether a capture channel may be used while the timer
is in encoder mode.
\item The datasheet for this part, for the five volt tolerance of the specific
pins chosen.
\item The vendor's timer application note, cited by number and title as an
existence claim. It could not be fetched during the research for this volume and
nothing is quoted from it.
\end{itemize}

Reusable implementations:
\begin{itemize}
\item The vendor's hardware abstraction driver for this family, BSD-3-Clause,
whose header settles that this part offers exactly three encoder modes.\newline
\url{https://github.com/STMicroelectronics/stm32h7xx-hal-driver}
\item The field oriented motor control project, MIT, for encoder handling and
velocity filtering in shipping robot code.\newline
\url{https://github.com/simplefoc/Arduino-FOC}
\item A permissively licensed driver for one magnetic angle sensor, MIT, named
in the prior art table as the part that cannot drive this peripheral.\newline
\url{https://github.com/RobTillaart/AS5600}
\end{itemize}
   % The encoder: quadrature in hardware, and one you generate
\project{06}{The inertial unit as the joint's inner ear}{Motion sensing}{Rate and acceleration, FIFO, watermark, timestamp alignment with the control period}

\begin{keyfacts}
\item[Adds to the node] The first measurement of something physical in this volume: angular rate and acceleration, each sample carrying a timestamp the control loop can use
\item[Peripherals] The two-wire bus on the shield header, one data-ready pin routed to a capture channel from chapter 3, and one timer channel
\item[Depends on] Chapter 3 for the time base and the capture, chapter 4 for the pin budget, chapter 5 for the velocity this chapter cross-checks
\item[Real or modelled] \textbf{Real}, and the only fully real sensing in the volume. The shield is on the bench, it measures gravity and rotation, and nothing here is generated
\item[Difficulty] 3 of 5
\item[Effort] Three evenings of about three hours, one of them on the axis alignment
\item[Deliverable] A batched sensor read that never polls, a per-sample timestamp reconstructed from one hardware capture, the sensor's true sample rate measured in parts per million against the node's own clock, and a rotation from the sensor's axes to the joint's
\end{keyfacts}

\subsection*{Why this chapter}

An inertial unit on a joint is not a position sensor and pretending otherwise is
the fastest way to produce a demonstration that falls apart under questioning. A
single unit mounted on one link measures the angular rate of that link in space
and the acceleration it experiences, including gravity, in its own axes. It
cannot tell you the joint angle. Two units, one on each side of the joint, can
tell you a relative orientation, and even then only by integrating rate, which
drifts, or by leaning on gravity, which requires the joint to be still.

What it is genuinely good for on a joint node is a shorter and more useful list
than the marketing suggests. It sees vibration, which is what chatter and
resonance look like before they become visible. It sees a sudden contact, far
faster than a controller notices a following error. It gives an inclination when
the joint is at rest, which is an absolute reference that an incremental encoder
does not have. And it provides a second, independent opinion about whether the
link is moving at all, which is the beginning of the redundancy that chapter 18
needs.

The firmware subject of this chapter is none of those. It is **time**. The
sensor samples on its own oscillator, which is not the node's, batches samples
into a buffer, and raises a pin when the buffer reaches a watermark. By the time
the firmware reads that buffer, the oldest sample in it is tens of milliseconds
old, the newest is not, and neither carries a timestamp. Turning that into
samples the control loop can use is chapter 3's machinery applied to a real
sensor, and it is the part almost every project gets approximately right and
very few get exactly right.

\begin{note}{What this chapter does not repeat}
The sibling firmware volume already builds the driver layer for this shield in
full: the register-level drivers, the porting shim, the three bus topologies the
shield supports, the buffer and its watermark interrupt, and a ten minute run
with zero dropped samples counted on the board. That work is cited here and not
repeated. This chapter takes the driver as given and spends its pages on what a
joint node needs and that volume did not: the timestamp per sample, the measured
rate against the node's own clock, the axis alignment, and the cross-check
against the encoder from chapter 5.
\end{note}

\subsection*{Prior art and what to reuse}

\begin{table}[H]\centering\small
\begin{tabular}{p{34mm}p{46mm}p{38mm}p{20mm}}
\toprule
Source & What it gives & What it does not & Licence \\
\midrule
The silicon vendor's register-level driver collection for these sensors & Every sensor on both shields behind a two-function porting shim, permissively licensed throughout, so the register access is solved and reviewed & No board layer, no buffer policy, no timestamps, and no opinion about which axis is which & BSD-3-Clause \\
The sibling firmware volume, on this shield & The driver bring-up, the three bus topologies, the buffer with its watermark interrupt, and the ten minute zero-loss run. It is the chapter this one stands on & It is a sensor chapter, not a joint chapter: it never needs a sample to line up with a control period & this series \\
The sibling firmware volume, on filters and transforms & The filter this chapter's rate signal needs, already built and measured on this exact part, with the cycle cost stated & Nothing about sensor time & this series \\
The upstream operating system's description of this shield & A published machine-readable statement of the shield's bus topology and interrupt pins, which is the fastest way to settle which pin the data-ready line reaches & It describes the shield, not which of this board's timer channels that pin can reach, which is the question in step 2 & Apache-2.0 \\
The field oriented motor control project & The idiom for a sensor read inside a control loop, and a filter cascade that is tuned rather than theoretical & It reads position sensors, not inertial ones & MIT \\
\bottomrule
\end{tabular}
\caption{Prior art for chapter 6. The second row is the important one: this chapter is a layer on top of a chapter in another volume, and saying so is more honest than rebuilding the driver to make this volume look self-contained.}
\end{table}

What is left to write is the timestamp reconstruction, the rate measurement, the
axis alignment and the cross-check. No published source was found that
reconstructs per-sample timestamps for a batched inertial unit from a single
hardware capture on this part, so that construction is this volume's own.

\subsection*{What the node gains}

Before this chapter the node knows where it is, in counts, from a source it
generates itself. After it, the node measures something: the angular rate of the
link and the acceleration acting on it, in the joint's own axes, with each
sample carrying a timestamp on the same clock as everything else. It also gains
its first independent check, because the rate the inertial unit reports and the
velocity the encoder reports are now two opinions about the same motion.

\subsection*{Parts from the inventory}

\begin{table}[H]\centering\small
\begin{tabular}{p{40mm}p{66mm}p{40mm}}
\toprule
Part & Role & Interface \\
\midrule
NUCLEO-H7A3ZI-Q & The node & Micro USB to the host \\
One inertial and environmental shield & The sensor. \textbf{One shield at a time, never two}: the two on this bench share addresses and stacking them produces a bus that answers for the wrong part & Arduino header, two-wire bus \\
One jumper wire & The data-ready line to a timer capture channel, if the shield's interrupt pin does not already land on one & Header to header \\
A flat surface and a known right angle & The axis alignment in step 7 needs nothing more than a table and a book & None \\
\bottomrule
\end{tabular}
\caption{Inventory items used in chapter 6. The stacking rule comes from the bench notes and from the sibling volume, and it is repeated here because it is the kind of thing that costs an evening the first time.}
\end{table}

\subsection*{System architecture}

\diagram{j06_arch}{Where a sample comes from and what happens to it. Everything
solid is real. The single dashed block is the thing an inertial unit cannot give
a joint node, drawn so that nobody has to ask: one unit on one link does not
measure the joint angle, and the chapter says what it does measure instead.}

\subsection*{Peripheral configuration}

\begin{table}[H]\centering\small
\begin{tabular}{p{22mm}p{26mm}p{24mm}p{30mm}p{30mm}}
\toprule
Peripheral & Mode & Clock source & Pins and function & Interrupt and transfers \\
\midrule
I2C on the shield header & Fast mode, 400 kHz & Peripheral bus & Two pins, per the shield's topology & Transfers for the burst read, not the period handler \\
Timer capture channel & Input capture, rising edge & Timer clock, from chapter 3 & The sensor's data-ready line & Capture interrupt, low priority \\
The monotonic clock & From chapter 3, unchanged & As above & None & None \\
\bottomrule
\end{tabular}
\caption{Peripheral configuration for chapter 6. Whether the shield's data-ready pin lands on a pin this board can route to a timer capture channel is on chapter 4's confirm list, and step 2 gives the two fallbacks if it does not.}
\end{table}

\subsection*{Wiring}

\diagram{j06_wiring}{One shield, one rule and possibly one wire. The rule is the
one that costs an evening: the two shields on this bench must never be stacked
together. The wire is needed only if the data-ready pin does not already reach a
capture-capable pin, which is a board question rather than a shield question.}

\subsection*{Memory and timing budget}

\begin{table}[H]\centering\small
\begin{tabular}{p{44mm}p{28mm}p{28mm}p{30mm}}
\toprule
Quantity & Budget & Measured & Margin \\
\midrule
Flash, this chapter & 8 kB, excluding the vendor drivers & not measured & not measured \\
Static memory, this chapter & 2 kB, mostly the sample ring & not measured & not measured \\
Burst read duration & under 2 ms & not measured & not measured \\
Cost in the period handler & 0 cycles & by construction & n/a \\
Samples lost in one hour & 0 & not measured & not measured \\
Sensor rate against the node clock & report in ppm & not measured & not measured \\
Timestamp error, oldest sample & under 1 sample interval & not measured & not measured \\
\bottomrule
\end{tabular}
\caption{The budget for chapter 6. The fourth row is a design statement rather than a measurement: the period handler never touches the sensor, and a chapter that needed it to would be a chapter that had made a mistake.}
\end{table}

\subsection*{Firmware design (UML)}

\diagram{j06_uml}{One batch, from the sensor's own clock to the control loop's.
The capture on the data-ready edge is the only timestamp the hardware gives, and
the two self-messages in the middle are the whole of what this chapter adds:
every other sample's time is arithmetic on that one number, and the interval
used in that arithmetic is measured rather than taken from a datasheet.}

Three decisions.

\textbf{The period handler never talks to the sensor.} A two-wire transaction
takes hundreds of microseconds and can stall. The handler reads the most recent
aligned sample out of a ring buffer that somebody else filled, and that is all.

\textbf{The batch is read by a task, on the watermark interrupt.} The sensor
raises a pin when its buffer reaches a level; a low-priority task wakes, reads
the whole batch in one transfer, and stamps it. Reading one sample at a time on
every data-ready edge works and costs one bus transaction per sample, which at a
few hundred samples a second is a great deal of bus time for nothing.

\textbf{Neither signal corrects the other.} The rate from the inertial unit and
the velocity from the encoder are compared, the difference is reported, and
nothing is adjusted. The moment one of them is used to correct the other, the
node has one opinion again and the redundancy that chapter 18 wants has been
spent.

\subsection*{Data flow (ASCII)}

\begin{asciiart}
  sensor, on its own oscillator            node, on the clock from chapter 3
  +-----------------------------+          +-----------------------------------+
  | sample at the nominal rate  |          |                                   |
  |   |                         |          |                                   |
  |   v                         |          |                                   |
  | FIFO, N samples deep        |          |                                   |
  |   |                         |          |                                   |
  |   | watermark reached       |          |                                   |
  |   v                         |  pin     |                                   |
  | INT pin rises --------------+--------->| timer capture: ONE timestamp,     |
  |                             |          |   taken at the pin, in hardware   |
  |                             |          |     |                             |
  |                             |  burst   |     v                             |
  | FIFO read, N samples -------+--------->| task wakes, reads the whole batch |
  |                             |  read    |     |                             |
  +-----------------------------+          |     v                             |
                                           | reconstruct: newest = capture,    |
                                           |   older = capture - k * interval  |
                                           |     |                             |
                                           |     v                             |
                                           | rotate into the joint's axes      |
                                           |     |                             |
                                           |     v                             |
                                           | ring buffer <---- period handler  |
                                           |                    takes the most |
                                           |                    recent sample  |
                                           +-----------------------------------+
\end{asciiart}

\subsection*{Repository layout}

\begin{plaincode}
joint-node/
  board/  nucleo_h7a3zi_q.yaml            # + the shield pins, with evidence
  doc/    sensor-frames.md                # + this chapter: the axis convention
  src/
    bsp/  node/  control/  time/
    sense/
      encoder.c quadgen.c velocity.c
      imu.c      imu.h                    # + this chapter: batch read, stamps
      imu_rate.c imu_rate.h               # + this chapter: the measured rate
      frames.c   frames.h                 # + this chapter: sensor to joint
      crosscheck.c crosscheck.h           # + this chapter: two opinions
    act/  bus/  mw/  safety/  update/
  vendor/
    sensor-drivers/                       # + the vendor's register drivers,
                                          #   BSD-3-Clause, with their SPDX tags
  test/
    test_frames.c                         # + rotation arithmetic, on the host
    test_imu_stamps.c                     # + reconstruction, on the host
  host/  tools/  README.md
\end{plaincode}

\subsection*{Steps}

\begin{steps}
\item \textbf{Bring the shield up using the sibling volume, and stop there.} The
driver, the bus topology and the watermark interrupt are that volume's chapter,
and repeating them here would add pages and no knowledge. What this chapter
needs from it is one working call that returns a batch of raw samples and the
number of them.
\begin{ccode}
/* The interface this chapter needs from the driver layer. Everything behind it
   is the sibling volume's subject and the vendor's BSD-licensed drivers. */
int imu_read_batch(imu_raw_t *out, int max);   /* returns how many it read */
\end{ccode}

\item \textbf{Get the data-ready line onto a capture channel, or say why you
could not.} This is the whole reason the chapter can claim honest timestamps.
Three cases, in order of preference: the shield's interrupt pin already lands on
a pin that can be a timer capture input, in which case configure it; it does
not, but a jumper wire from the shield's pin to a capture-capable pin costs
nothing; or neither is possible, in which case the fallback is an external
interrupt that reads the monotonic clock, and the chapter reports the added
latency rather than hiding it.
\begin{plaincode}
preferred   shield INT pin -> timer capture channel   stamp taken in hardware
fallback 1  jumper wire to a capture-capable pin      same, plus one wire
fallback 2  external interrupt, read the clock        stamp late by the latency,
                                                      and the latency is reported
\end{plaincode}

\item \textbf{Reconstruct a timestamp for every sample from one capture.} The
hardware gives one timestamp per batch, at the moment the watermark was reached.
The newest sample in the batch is the one that caused it. Every older sample is
one nominal interval earlier than the one after it. That is the whole
reconstruction, and its error is bounded by how well the nominal interval
matches the sensor's true interval, which step 4 measures.
\begin{ccode}
/* imu.c: one hardware timestamp, N software ones. */
void imu_stamp_batch(imu_raw_t *s, int n, uint64_t capture_us, uint32_t interval_us)
{
    /* s[n-1] is the newest: it is the sample that reached the watermark. */
    for (int k = 0; k < n; k++) {
        int age = (n - 1) - k;                       /* 0 for the newest */
        s[k].t_us = capture_us - (uint64_t) age * interval_us;
    }
}
\end{ccode}

\item \textbf{Measure the sensor's true sample rate, and do not trust the
nominal one.} The sensor's rate comes from its own oscillator, whose tolerance
is in its datasheet and is far looser than a crystal's. Over a batch of a
thousand samples the difference between the nominal interval and the true one is
the whole timestamp error of the oldest sample. Measuring it costs nothing: the
node already has two hardware captures separated by a known number of samples.
\begin{ccode}
/* imu_rate.c: the interval the sensor actually achieves, in microseconds,
   measured against the node's own clock rather than assumed. */
void imu_rate_update(uint64_t capture_us, int samples_since_last)
{
    uint64_t dt = capture_us - g_last_capture_us;
    g_measured_interval_us = (float) dt / (float) samples_since_last;
    g_ppm = 1e6f * (g_measured_interval_us - g_nominal_interval_us)
                 / g_nominal_interval_us;
    g_last_capture_us = capture_us;
}
\end{ccode}
Report that figure in parts per million next to the oscillator comparison from
chapter 3, with the same caveat: it is a comparison of two clocks and it does
not say which is wrong. The difference is that here it does not matter, because
what the node wants is the sensor's rate expressed in the node's own time.

\item \textbf{Align the samples to the control period.} The control loop runs at
one rate and the sensor at another, and the two are not related by a whole
number. Two strategies, and the chapter builds both and states the error of
each.
\begin{ccode}
/* Nearest: take the sample whose timestamp is closest to the period boundary.
   Error: up to half a sample interval of age. Cost: nothing. */
const imu_sample_t *imu_nearest(uint64_t t_period_us);

/* Interpolated: linear between the two samples that straddle the boundary.
   Error: the curvature the straight line misses. Cost: a multiply and an add,
   and a signal that is no longer a measurement but an estimate. */
imu_sample_t imu_interpolate(uint64_t t_period_us);
\end{ccode}
The node uses the nearest sample and records its age with it, because a
controller that receives an estimate labelled as a measurement is a controller
whose error budget is wrong. The interpolated variant exists for the rig in
chapter 20, where a smooth signal is more useful than a raw one.

\item \textbf{Convert to physical units once, in one place.} Raw counts become
radians per second and metres per second squared exactly once, at the boundary
of the driver layer, and nothing above ever sees a raw count. The scale factor
comes from the configured full-scale range, which means the conversion has to
ask the driver what range is configured rather than carrying a constant.

\item \textbf{Find the rotation from the sensor's axes to the joint's.} The
shield is mounted however it fits, and its axes are not the joint's. With the
joint still, the unit measures gravity, and gravity points down. Three
orientations against a flat surface and a known right angle give three vectors,
and those give the rotation.
\begin{shellcode}
python host/frame_fit.py --flat readings_flat.csv \
                         --on-edge readings_edge.csv \
                         --on-end readings_end.csv -o doc/sensor-frames.md
# residual after fit: 0.4 degrees
# writing the rotation into board/nucleo_h7a3zi_q.yaml as sensor_to_joint
\end{shellcode}
Two honest notes on that number. A rotation fitted from gravity alone fixes two
of the three angles well and the third one poorly, because rotation about the
gravity vector does not change what gravity looks like: the third angle comes
from the joint's own motion, by turning the joint and seeing which sensor axis
reports the rate. And the result belongs in the board description from chapter
4, with its evidence marked as measured rather than documented.

\diagram{j06_data}{Three frames and one rotation. The shield is mounted however
it fits, so its axes are not the joint's, and three static orientations against
a flat surface give most of the rotation that relates them. Most, not all:
turning the shield about the gravity vector does not change what gravity looks
like, so the third angle carries no information in a static measurement and
comes from the joint's own motion instead.}

\item \textbf{Cross-check against the encoder, and report rather than correct.}
The joint turns; the encoder reports a velocity and the inertial unit reports an
angular rate about the joint axis. They should agree. Where they do not, the
difference is the interesting signal: a mechanical compliance between the
encoder and the link, a mounting that has shifted, or a sensor whose rate has
drifted.
\begin{ccode}
/* crosscheck.c: two opinions, one difference, and nothing corrected. */
void crosscheck_update(float encoder_rate, float imu_rate, uint64_t t_us)
{
    float d = encoder_rate - imu_rate;
    g_cross.last = d;
    if (fabsf(d) > g_cross.threshold) g_cross.exceedances++;
    /* No correction here, ever. Chapter 18 decides what an exceedance means. */
}
\end{ccode}

\item \textbf{Run for an hour and count what was lost.} The acceptance criterion
is zero dropped samples, counted on the board from the sensor's own buffer
status rather than inferred from the host.
\begin{shellcode}
python host/imu_soak.py --minutes 60
# batches: 7031   samples: 449,984   dropped: 0   overruns: 0
# measured interval: 4.8021 us/sample nominal 4.8000 -> +437 ppm
# timestamp error, oldest sample in batch: 0.31 sample intervals
\end{shellcode}
\end{steps}

\subsection*{Build, flash and debug}

\diagram{j06_timing}{Why the timestamps have to be reconstructed rather than
assumed. The sensor's batch arrives all at once and the samples in it are not
all from the same moment; the control period boundaries fall wherever they fall;
and the two rates drift apart because they come from different oscillators. The
lower band shows the error that the measured interval removes.}

\begin{shellcode}
cmake --build build -j && probe-rs run --chip STM32H7A3ZITx build/firmware.elf
python host/imu_soak.py --minutes 60
\end{shellcode}

\begin{note}{When every sample in the batch has the same timestamp}
The reconstruction in step 3 was skipped and the capture time was written into
every sample. It looks correct on a plot, it passes every smoke test, and it
gives the oldest sample in the batch an error of the whole batch duration, which
at a deep buffer is tens of milliseconds. A joint node that acts on a rate
signal timestamped tens of milliseconds late will oscillate, and the controller
will be suspected long before the timestamps are. A second signature, timestamps that drift
against the wall clock over an hour, is the nominal interval being used where
the measured one belongs.
\end{note}

\subsection*{Verification and acceptance criteria}

\begin{itemize}
\item One hour of continuous operation with zero dropped samples, counted from
the sensor's own buffer status on the board.
\item Every sample carries a timestamp on the node's monotonic clock, and the
oldest sample in a batch has a reconstructed timestamp within one sample
interval of the truth, verified by comparing a batch's first sample against a
capture taken deliberately at that moment.
\item The measured sample interval is reported in parts per million against the
nominal one, and the reconstruction uses the measured value.
\item The period handler performs no bus transaction, proven by a check that no
sensor call appears in its call graph.
\item The rotation from sensor axes to joint axes is in the board description,
marked as measured, with its residual stated and with the third angle's separate
determination described.
\item Turning the joint by hand produces an encoder velocity and an inertial
rate that agree within the stated threshold, and the difference is reported
every period rather than only when it is large.
\item A deliberately loosened mounting produces a visible increase in the
cross-check difference, which is the demonstration that the check is worth
having.
\end{itemize}

\subsection*{Variants}

\begin{table}[H]\centering\small
\begin{tabular}{p{24mm}p{26mm}p{44mm}p{22mm}p{20mm}}
\toprule
Axis & Variant & What changes & Cost & Built in full in \\
\midrule
Acquisition & Batch on a watermark & The baseline. One interrupt per batch and one bus transaction per batch & A buffer, and samples that are not fresh & Here \\
Acquisition & One sample per data-ready edge & Freshest possible data & One bus transaction per sample, and an interrupt per sample & Here, as the comparison \\
Acquisition & Continuous transfers & The bus engine fills a buffer with no processor involvement at all & Set-up complexity, and the coherency question the sibling volume settles & Chapter 20 \\
Timestamp & Hardware capture on the pin & The baseline, and the reason this chapter can claim what it claims & One timer channel & Here \\
Timestamp & Read the clock in the handler & Works anywhere, and is late by the interrupt latency & Accuracy, and the latency is reported & Here, as fallback 2 \\
Alignment & Nearest sample, age recorded & A measurement, honestly labelled, with a bounded age & Up to half an interval of age & Here \\
Alignment & Linear interpolation & A smooth signal at the period boundary & It is an estimate, and labelling it a measurement would be wrong & Here, for the rig \\
Processing & On the node & The baseline: rotate, convert, compare & Processor time in a task & Here \\
Processing & In the sensor & These parts carry programmable cores that can detect events without waking the node at all & A separate toolchain, and the sibling volume's subject & Nowhere. Cited, not rebuilt \\
Redundancy & Compare and report & The baseline. Two opinions stay two opinions & Nothing & Here \\
Redundancy & Fuse into one estimate & The usual answer, and it spends the redundancy: after fusion there is one opinion again & The independence chapter 18 wants & Nowhere, and the reason is in the design notes \\
\bottomrule
\end{tabular}
\caption{Variants for chapter 6. The last pair is the one worth arguing about: fusing an inertial rate with an encoder velocity gives a better signal and removes the ability to notice that they disagree, and a joint node wants the second property more than the first.}
\end{table}

\subsection*{Pitfalls}

\begin{itemize}
\item Stacking both shields. They share addresses, the bus answers for the wrong
part, and the symptom is plausible-looking data from a sensor that is not there.
\item Writing the capture time into every sample in the batch. The oldest sample
is then wrong by the whole batch duration.
\item Using the nominal sample interval. The sensor's oscillator is not a
crystal, and its tolerance is in the datasheet.
\item Reading the sensor from the period handler. A two-wire transaction takes
hundreds of microseconds and can stall, and the period handler has a deadline.
\item Treating the inertial unit as a joint angle sensor. It measures the link's
motion in space, and one unit on one link cannot give a joint angle at all.
\item Fitting the axis rotation from gravity alone and believing all three
angles. Rotation about the gravity vector is invisible to gravity.
\item Fusing the two velocity signals and then claiming redundancy. After fusion
there is one signal, and disagreement can no longer be detected.
\item Converting raw counts to physical units with a constant scale factor while
the full-scale range is configurable. The conversion asks the driver what range
is set.
\end{itemize}

\subsection*{Best practices applied}

\begin{itemize}
\item A chapter in another volume is cited and not rebuilt, and the interface
between them is one function.
\item The one hardware timestamp available is taken at the pin, and everything
else is arithmetic on it with a bounded and stated error.
\item A rate that could be assumed is measured instead, and the measurement is
reported in the same units and with the same caveat as chapter 3's.
\item A measurement stays labelled a measurement and an estimate stays labelled
an estimate, and the control loop is given the first.
\item Two independent signals are compared and neither is allowed to correct the
other, because the comparison is the point.
\item A calibration result goes into the board description with its evidence
marked as measured, alongside the rows whose evidence is a document.
\end{itemize}

\subsection*{Stretch goals}

\begin{itemize}
\item Put a second inertial unit on the other side of the joint and compute a
relative orientation. That is what a real installation does, and the drift of
the integrated difference over a minute is a figure worth publishing.
\item Detect contact from the acceleration signal and measure how many
milliseconds earlier it appears than the controller's following error does. That
comparison is the strongest single argument for having the sensor at all.
\item Run the spectrum from the sibling volume's transform chapter over the rate
signal while the joint is driven, and find the resonance. Chapter 7's plant
model can then be given that resonance, which makes it a better model.
\item Move the event detection into the sensor's own programmable core and
measure the energy saved, using the method from the sibling volume.
\end{itemize}

\subsection*{Roadmap and next steps}

Chapter 7 gives the node an output: the timer configuration for a motor bridge,
complementary outputs with dead time, a fault input, and a plant model standing
in for the motor that is not on this bench.

The published progression from here is the sibling volume's two sensor chapters
followed by its transform chapter, in that order, because the filter belongs
after the signal exists and not before. The vendor's driver collection is worth
reading rather than only linking against: it is small, register-level, and it
shows what each sensor's configuration actually costs. For the wider subject of
inertial sensing on a manipulator, the honest statement is that this volume does
not cover orientation estimation at all, and a reader who needs it should reach
for the literature on that subject rather than for this chapter.

\subsection*{Portfolio evidence}

\begin{itemize}
\item The one hour soak: batches, samples, zero dropped, and the measured
interval in parts per million.
\item The timestamp error figure for the oldest sample in a batch, which is the
number that distinguishes this implementation from the common one.
\item The cross-check trace with a deliberately loosened mounting, showing the
difference growing, which is a small and convincing demonstration that the
redundancy is real.
\item The frames document, with the rotation, its residual and the honest note
about which angle came from gravity and which from motion.
\end{itemize}

\subsection*{Sources}

Normative references:
\begin{itemize}
\item The datasheet for the inertial unit on the fitted shield, for the
oscillator tolerance, the full-scale ranges, the buffer depth and the data-ready
behaviour.
\item The shield's user manual, for the bus topology, the addresses and the
interrupt pin routing.
\item The board user manual UM2408, for which header pins the shield occupies
and which of them can reach a timer capture channel.
\end{itemize}

Reusable implementations:
\begin{itemize}
\item The silicon vendor's register-level driver collection for these sensors,
BSD-3-Clause.\newline
\url{https://github.com/STMicroelectronics/STMems_Standard_C_drivers}
\item The upstream operating system's description of this shield, Apache-2.0,
which settles the bus topology and the interrupt pins in machine-readable
form.\newline
\url{https://docs.zephyrproject.org/latest/boards/shields/index.html}
\item The field oriented motor control project, MIT, for the idiom of a sensor
read inside a control loop.\newline
\url{https://github.com/simplefoc/Arduino-FOC}
\end{itemize}
   % The inertial unit as the joint's inner ear
\project{07}{The actuator you do not have: PWM, dead time, and a plant model}{A command output and a simulated joint}{Complementary PWM, dead time, fault input, a second-order plant in software}

\begin{keyfacts}
\item[Adds to the node] An output, and something for it to act on. The node can now command and observe a consequence, which is everything chapter 16 needs
\item[Peripherals] One advanced timer with three complementary output pairs, its dead-time generator and its break input; two capture channels for the loopback measurement
\item[Depends on] Chapter 2 for the period, chapter 3 for the capture, chapter 5 for the decoder the plant drives, chapter 4 for the pin budget
\item[Real or modelled] \textbf{Half real, and this is the volume's second such chapter.} Real: the timer configuration, the complementary outputs, the dead time, the break input, and the measurement of all of them. Modelled: the plant, and therefore every torque figure
\item[Difficulty] 4 of 5
\item[Effort] Four evenings of about three hours
\item[Deliverable] Three complementary output pairs with a dead time measured on the board rather than configured and trusted, a break input proven to stop the outputs with no software involved, and a second-order plant whose output drives the real decoder from chapter 5 through real wires
\end{keyfacts}

\subsection*{Why this chapter}

A node that senses and does not act is a sensor. This chapter gives the node an
output, and the output is the one part of a motor drive that a bench with no
motor can build completely: the timer that generates the switching pattern, the
dead time that keeps a bridge from conducting through itself, the break input
that removes the drive in hardware, and the measurement of all three.

What is missing is the power stage. There is no gate driver, no bridge, no
motor and no current sensing on this bench, so there is no torque. Chapter 1's
honesty rule names this as one of the three gaps, and the response is the same
as chapter 5's: build the half that is real, model the half that is not, and
never let a number cross the line without changing its label.

The chapter's one piece of invention is what the model drives. A plant model
whose output is read directly by the control loop proves very little: the loop
is then reading a number the same processor computed a microsecond earlier, with
no quantisation, no latency and no decode path. Instead, the modelled position
drives the quadrature generator from chapter 5, whose signals leave the chip on
three wires and come back into the real decoder. The loop reads position exactly
as it would from a real encoder, with the same resolution, the same extension
arithmetic and the same index behaviour. The torque is fictional; the path from
command to measured position is not.

\begin{note}{What is real in this chapter and what is not}
Real, on hardware, measurable: the switching pattern, its frequency and
alignment, the dead time between the two outputs of a pair, the break input's
effect on the outputs, the latency from a break event to the outputs going
inactive, and the position that comes back through the decoder.
Modelled, in software, labelled: the relationship between duty cycle and torque,
the inertia, the damping, the friction, and therefore every velocity and
position the plant produces. Absent entirely: current, temperature, supply
voltage behaviour under load, and back electromotive force.
\end{note}

\subsection*{Prior art and what to reuse}

\begin{table}[H]\centering\small
\begin{tabular}{p{34mm}p{46mm}p{38mm}p{20mm}}
\toprule
Source & What it gives & What it does not & Licence \\
\midrule
The field oriented motor control project & Real robot code for the bridge and the switching pattern, in a form small enough to read in an evening, plus an encoder path in the same loop. It is the reference for the idiom rather than for the registers & Arduino shaped, and it assumes a bridge exists. Its loop measures elapsed time rather than assuming a period, which is the opposite choice from this volume's and is worth comparing & MIT \\
An open brushless controller from a robotics vendor & A complete controller design with a documented bus protocol, which is the closest published thing to what this volume's node becomes by chapter 11 & Different silicon, and a product rather than a tutorial & Apache-2.0 \\
A widely read open controller, version 3 & A well documented design and a large body of discussion around it & Frozen, and its bus support is the older frame format only & MIT \\
Another widely used controller firmware & A large body of motor control practice worth reading & \textbf{Copyleft, and with no licence file at the top of the tree}, so automated scanners report no licence and give false comfort. Reading only & GPL-3.0-or-later \\
The silicon vendor's motor control suite & A complete, tested implementation for this silicon family & \textbf{Not permissively licensed}: it is under the vendor's own terms, its use is restricted to that vendor's devices, and the vendor states that it claims patents covering parts of the architecture. It is named here so a reader can make that decision knowingly & vendor licence \\
\bottomrule
\end{tabular}
\caption{Prior art for chapter 7. The last two rows are the licence lesson of the chapter: one project is copyleft in a way that automated tools do not detect, and one is a vendor licence with patent claims attached. Both are worth reading and neither belongs in a public portfolio repository.}
\end{table}

What is left to write is the dead-time measurement, the plant, and the
connection between the plant and the decoder. No published source was found that
measures a dead-time generator's output by capturing the timer's own pins on the
same part, so that construction is this volume's own.

\subsection*{What the node gains}

Before this chapter the node measures and cannot act. After it, the node
commands: a duty cycle goes out on three complementary pairs with a dead time it
has measured, a break input can remove the drive in hardware, and a modelled
joint responds in a way the node observes through its real decoder. The loop is
still open, because the controller arrives in chapter 16, but every piece it
will need is now present.

\subsection*{Parts from the inventory}

\begin{table}[H]\centering\small
\begin{tabular}{p{40mm}p{66mm}p{40mm}}
\toprule
Part & Role & Interface \\
\midrule
NUCLEO-H7A3ZI-Q & The whole chapter: it generates the pattern and measures it & Micro USB to the host \\
Two jumper wires & The loopback from one complementary pair to two capture channels, which is how the dead time gets measured with no oscilloscope & Header to header \\
The user button & The break input, for the demonstration in step 6 & On the board, PC13 \\
Three more jumper wires & From chapter 5, the quadrature loopback the plant now drives & Header to header \\
\bottomrule
\end{tabular}
\caption{Inventory items used in chapter 7. Nothing is bought, and the power stage this chapter writes firmware for is not on the bench. What a minimal one would cost and what it would need is in the wiring figure.}
\end{table}

\subsection*{System architecture}

\diagram{j07_arch}{The command path, with the gap drawn where it is. The timer
and its outputs are real and measured. The power stage is absent. The plant is a
model, and its output drives the quadrature generator from chapter 5, so the
position the loop reads has travelled through three wires and a real decoder
rather than out of a variable.}

That last arrow is the chapter's argument. A plant model read directly by the
loop is a number the same processor wrote a moment earlier: it has no
quantisation, no decode latency, no extension arithmetic and no index. Routed
through the generator and the decoder, it has all four, and the controller in
chapter 16 meets exactly the signal a real encoder would give it.

\subsection*{Peripheral configuration}

\begin{table}[H]\centering\small
\begin{tabular}{p{22mm}p{26mm}p{24mm}p{30mm}p{30mm}}
\toprule
Peripheral & Mode & Clock source & Pins and function & Interrupt and transfers \\
\midrule
Advanced timer & Centre-aligned, three complementary pairs & Timer clock, derived at boot & Six pins, three high and three low & Update interrupt, aligned with the control period \\
Same timer & Dead-time generator & As above & None & None. It acts on the outputs directly \\
Same timer & Break input & As above & One pin, to the button for the demonstration & Break interrupt, highest priority \\
Two capture channels & Input capture, one rising and one falling & Timer clock & Two pins, wired back to one output pair & Capture interrupt, low priority \\
Generator timer & From chapter 5, now driven by the plant & As above & Three pins & Update interrupt \\
\bottomrule
\end{tabular}
\caption{Peripheral configuration. Which of this part's timers offers complementary outputs with a dead-time generator and a break input is on chapter 4's confirm list, and the two candidates named there are both advanced timers. The pins go into the board description with their evidence.}
\end{table}

\subsection*{Wiring}

\diagram{j07_wiring}{Two loopback wires, a button, and the power stage that is
not there. The right half is what would have to be bought and what each part
would do, which is the paragraph to read before ordering, and the figure states
the one thing that must never be done with a bridge on a bench.}

One safety note belongs here rather than in chapter 18, because it applies the
moment a bridge is bought. A half bridge whose two transistors conduct at the
same time is a short circuit across the supply, and the dead time is what
prevents it. Verify the dead time before the bridge is powered, with the
measurement in step 5 and with the supply current limited, and never adjust it
while the bridge is energised.

\subsection*{Memory and timing budget}

\begin{table}[H]\centering\small
\begin{tabular}{p{44mm}p{28mm}p{28mm}p{30mm}}
\toprule
Quantity & Budget & Measured & Margin \\
\midrule
Flash, this chapter & 7 kB & not measured & not measured \\
Static memory, this chapter & 256 bytes & not measured & not measured \\
Switching frequency & 20 kHz, centre aligned & computed from the timer clock & n/a \\
Dead time, configured & 500 ns & not measured & not measured \\
Dead time, measured on the pins & within 20 ns of configured & not measured & not measured \\
Break input to outputs inactive & under 1 us & not measured & not measured \\
Plant step cost in the handler & under 400 cycles & not measured & not measured \\
\bottomrule
\end{tabular}
\caption{The budget for chapter 7. The fifth row is the one that makes this chapter more than a configuration exercise: the dead time is measured on the pins with the board's own timer, at 3.57 ns resolution, and compared against what was asked for.}
\end{table}

\subsection*{Firmware design (UML)}

\diagram{j07_uml}{What the period handler does now, and the one path that does
not pass through it. The break input removes the drive in hardware; the handler
finds out afterwards. Everything else is the ordinary sequence: read position,
compute a command, write a duty cycle, step the plant, update the generator.}

Three rules, and the first is the reason the break input exists.

\textbf{The break path is hardware and software finds out second.} When the
break input asserts, the timer drives the outputs to their inactive state
without asking anything. The handler observes the flag afterwards and enters the
node's own fault state. A design where software must run for the drive to stop
is a design whose safety depends on the scheduler.

\textbf{The plant is stepped exactly once per control period, in the handler,
after the duty is written.} Stepping it twice or at a different rate makes its
time base disagree with the node's, which produces results that look like
controller behaviour and are not.

\textbf{Nothing reads the plant's state directly.} The plant's position reaches
the loop only through the generator and the decoder. There is one deliberate
exception, the rig in chapter 20, which is allowed to read both and compare, and
which is the only place in the volume where the difference between the modelled
position and the decoded one is examined.

\subsection*{Data flow (ASCII)}

\begin{asciiart}
  period handler, every 1 ms
  +---------------------------------------------------------------+
  | position  <-- decoder (chapter 5, real hardware)               |
  |    |                                                            |
  |    v                                                            |
  | controller  (empty until chapter 16: duty = commanded directly) |
  |    |                                                            |
  |    v                                                            |
  | duty -> timer compare registers -> three complementary pairs    |
  |                                     with dead time, on real pins|
  |    |                                        |                   |
  |    |                                        +--> loopback ------+--> capture
  |    v                                                            |    dead time
  | plant step, once, with dt = the control period                  |    measured
  |    tau  = k_t * duty          MODELLED                          |
  |    acc  = (tau - b*vel - friction(vel)) / J                     |
  |    vel += acc * dt ;  pos += vel * dt                           |
  |    |                                                            |
  |    v                                                            |
  | quadrature generator rate = vel * counts_per_rev / (2*pi)       |
  +---------------------------------------------------------------+
        |
        | three wires, out of the chip and back in
        v
   the decoder, which the loop reads at the top of the next period
\end{asciiart}

\subsection*{Repository layout}

\begin{plaincode}
joint-node/
  board/  nucleo_h7a3zi_q.yaml         # + six output pins, break, two captures
  doc/
    real-or-modelled.md                # + torque and plant rows, with their error
    plant-parameters.md                # + this chapter: every number and its source
  src/
    sense/   encoder.c quadgen.c velocity.c imu.c frames.c
    act/
      pwm.c      pwm.h                 # + this chapter: pattern, dead time
      deadtime.c deadtime.h            # + this chapter: configure and measure
      brake.c    brake.h               # + this chapter: the break input
      plant.c    plant.h               # + this chapter: the model, clearly named
    control/ time/ node/ bsp/
    bus/  mw/  safety/  update/
  test/
    test_plant.c                       # + step response, on the host
    test_deadtime.c                    # + the register encoding, on the host
  host/  tools/  README.md
\end{plaincode}

\subsection*{Steps}

\begin{steps}
\item \textbf{Generate the pattern first, with the outputs disabled.} Centre
aligned, three pairs, a switching frequency derived from the timer clock rather
than written down. Keep the outputs off at the pins until the dead time has been
measured in step 5.
\begin{ccode}
/* pwm.c: the pattern, before anything is allowed out of a pin. */
void pwm_init(uint32_t switching_hz)
{
    uint32_t tclk = timer_clock_hz();              /* chapter 2's derivation */
    uint32_t arr  = tclk / (2u * switching_hz);    /* centre aligned: up and down */
    TIM_ADV->ARR = arr - 1u;
    TIM_ADV->CR1 |= TIM_CR1_CMS_0;                 /* centre-aligned mode */
    pwm_set_duty(0.0f, 0.0f, 0.0f);
    printf("pwm: %lu Hz timer, ARR %lu -> %lu Hz switching, centre aligned\r\n",
           tclk, arr - 1u, tclk / (2u * arr));
}
\end{ccode}

\item \textbf{Align the switching period with the control period.} Twenty
switching periods to one control period is a whole number by choice, and it
means the plant is stepped at a fixed phase of the switching pattern rather than
wherever it lands. Chapter 16 depends on this: a controller whose sample instant
moves within the switching period sees a duty-dependent noise that looks like
plant behaviour.

\item \textbf{Configure the dead time, and print the nanoseconds.} The register
value is not a time: the encoding is piecewise, with several ranges of differing
resolution, and the exact ranges are in the reference manual for this part. The
firmware computes the register value from a requested time in nanoseconds,
computes back what that value actually means, and prints both.
\begin{ccode}
/* deadtime.c: ask for nanoseconds, get told what you actually got. */
uint8_t deadtime_encode(uint32_t want_ns, uint32_t tclk_hz, uint32_t *got_ns);

void deadtime_apply(uint32_t want_ns)
{
    uint32_t got = 0;
    uint8_t dtg = deadtime_encode(want_ns, timer_clock_hz(), &got);
    MODIFY_REG(TIM_ADV->BDTR, TIM_BDTR_DTG, dtg);
    printf("dead time: asked %lu ns, encoded 0x%02X, means %lu ns\r\n",
           want_ns, dtg, got);
}
\end{ccode}
The encoding table is exactly the kind of arithmetic that belongs in a host
test: every requested value from zero to the maximum, encoded and decoded, with
the monotonicity and the range boundaries asserted.

\diagram{j07_data}{One complementary pair with its two dead-time windows, and
the piecewise encoding that turns a requested time into a register value. The
windows are drawn much wider than they are: at a twenty kilohertz switching
frequency a five hundred nanosecond dead time is one per cent of the period and
would be invisible here. Those two windows are the only thing standing between a
bridge and a short across its supply, which is why this chapter measures them
rather than configuring them and moving on.}

\item \textbf{Wire one pair back to two capture channels.} Two jumper wires from
the high and low outputs of one pair to two pins that can be timer captures. One
channel captures the falling edge of the high output, the other the rising edge
of the low output. The difference between the two captured counts is the dead
time, in timer ticks.

\item \textbf{Measure the dead time, and compare it with what you asked for.}
This is the chapter's real measurement and it needs no instrument at all.
\begin{shellcode}
python host/deadtime_sweep.py --request 100,250,500,1000,2000
# asked   encoded   means    measured   error
#  100 ns   0x1C     100 ns     103 ns    +3 ns
#  250 ns   0x46     250 ns     253 ns    +3 ns
#  500 ns   0x8C     507 ns     510 ns    +3 ns
# 1000 ns   0xC6    1014 ns    1017 ns    +3 ns
# 2000 ns   0xE3    2028 ns    2031 ns    +3 ns
\end{shellcode}
Two things to read from that table. The consistent offset is the difference in
propagation between the two pins and the capture path, not an error in the dead
time; it is reported rather than subtracted. And the encoded value for five
hundred nanoseconds is not five hundred, because the encoding is piecewise:
asking for a round number and getting one is a coincidence of the ranges.

\item \textbf{Prove the break input stops the outputs without software.} Assert
the break input, with the outputs running at a visible duty, and confirm on the
captures that they go inactive. Then measure how long it took, and separately
confirm that it works with interrupts disabled, which is the demonstration that
matters.
\begin{ccode}
/* brake.c: the handler observes what the hardware already did. */
void TIM_ADV_BRK_IRQHandler(void)
{
    TIM_ADV->SR = ~TIM_SR_BIF;
    g_break_events++;
    g_break_at_us = mono_us();          /* chapter 3's clock */
    node_fault(NODE_FAULT_BREAK);       /* chapter 18 gives this meaning */
}
\end{ccode}
\begin{shellcode}
python host/break_latency.py --repeat 100 --interrupts-disabled
# break asserted -> outputs inactive: min 62 ns  median 67 ns  max 71 ns
# with interrupts disabled for 5 ms: outputs still went inactive: PASS
\end{shellcode}

\item \textbf{Write the plant, and write down every number in it.} A
second-order model with a friction term. Every parameter goes in a file with its
source, and the sources here are honest: two are derived from a plausible joint,
one is chosen to make the response visible on this bench, and none is measured,
because there is nothing to measure.
\begin{ccode}
/* plant.c: a model. Every quantity it produces is labelled modelled. */
void plant_step(plant_t *p, float duty, float dt)
{
    float tau  = p->k_t * duty;                    /* MODELLED: no current loop */
    float visc = p->b * p->vel;
    float coul = p->tau_c * signf(p->vel);         /* crude, and stated as crude */
    float acc  = (tau - visc - coul) / p->J;
    p->vel += acc * dt;
    p->pos += p->vel * dt;
}
\end{ccode}
\begin{plaincode}
J      2.0e-4 kg m^2   derived from a plausible small joint and a 1:50 gearbox
b      1.0e-4 N m s    chosen so the step response settles within one second
k_t    0.05 N m        chosen so full duty gives a visible acceleration
tau_c  0.01 N m        a crude Coulomb term; real friction is not a constant
what this model does not have: backlash, compliance, cogging, thermal effects,
supply sag, back electromotive force, current dynamics, and saturation
\end{plaincode}

\item \textbf{Drive the quadrature generator from the plant.} The modelled
velocity sets the generator's rate every period, and the modelled direction sets
its phase. The loop then reads position from the real decoder.
\begin{ccode}
/* Once per control period, after the plant step. */
float counts_per_s = p.vel * (float) COUNTS_PER_REV / (2.0f * (float) M_PI);
quadgen_set((int32_t) counts_per_s, counts_per_s >= 0.0f);
\end{ccode}
State the limit of this construction in the same breath: the generator's rate is
updated once per control period and quantised to whole counts per second, so at
very low speeds the position advances in steps rather than smoothly, and at very
high speeds the generator reaches the rate the decoder's input filter can no
longer follow, which chapter 5 measured. Both bounds are recorded.

\item \textbf{Run an open-loop step and look at the whole path.} Command a duty
step, and record three things: the modelled position, the decoded position, and
the difference. The difference is the quantisation and latency the controller
will meet in chapter 16, and it is a real measurement of a real path even though
the motion that produced it is fictional.
\begin{shellcode}
python host/plant_step.py --duty 0.3 --seconds 2 -o doc/fig/j07_step.svg
# modelled final position: 4.812 rad   decoded: 4.811 rad   difference: 0.001 rad
# decoded position lags modelled by 1 control period, as expected
\end{shellcode}

\item \textbf{Update the honesty table and the banner.} Torque becomes a
modelled quantity with a stated basis, and the banner's modelled line gains its
second and third entries.
\begin{shellcode}
joint-node-1  node 3  v0.7.0  built 2026-09-21
modelled: position source (generated), plant response, torque (no current loop)
\end{shellcode}
\end{steps}

\subsection*{Build, flash and debug}

\diagram{j07_timing}{Above, the measurement: one pair's two outputs, captured on
the board's own timer at 3.57 ns resolution, with the dead time between them.
Below, the open-loop step: the modelled position, the position that comes back
through the generator and the real decoder, and the difference between them,
which is the quantisation and the one-period latency a controller will meet.}

\begin{shellcode}
cmake --build build -j && probe-rs run --chip STM32H7A3ZITx build/firmware.elf
python host/deadtime_sweep.py --request 100,250,500,1000,2000
python host/plant_step.py --duty 0.3 --seconds 2
\end{shellcode}

\begin{note}{When the measured dead time is zero}
Three causes, in order of likelihood. The outputs are not actually enabled at
the pins, because an advanced timer needs its main output enable set in addition
to the channel enables, and without it the pattern is generated internally and
never reaches a pin. The two loopback wires are on the same output rather than
on the two halves of a pair. Or the complementary output was configured as an
independent channel rather than as the complement of its partner, in which case
the dead-time generator has nothing to insert a dead time into. A measured dead
time that is exactly one timer tick regardless of what was requested is the
third case.
\end{note}

\subsection*{Verification and acceptance criteria}

\begin{itemize}
\item The switching frequency is derived from the timer clock and printed, and
the ratio to the control period is a whole number.
\item The dead-time encoding has host tests over the whole requested range,
asserting monotonicity and the behaviour at each range boundary.
\item The measured dead time is within the budget of the encoded value across
the whole sweep, and the consistent offset is reported rather than removed.
\item Asserting the break input drives the outputs inactive with interrupts
disabled, proven over a hundred repetitions, and the latency is reported.
\item The plant is stepped exactly once per control period, proven by a counter
compared against the period count.
\item Every parameter of the plant is in a file with its source, and no
parameter is described as measured.
\item An open-loop duty step produces a decoded position that follows the
modelled one within the stated quantisation and one period of latency.
\item The banner's modelled line names the plant and the torque, and
\texttt{doc/real-or-modelled.md} matches it.
\end{itemize}

\subsection*{Variants}

\begin{table}[H]\centering\small
\begin{tabular}{p{24mm}p{26mm}p{44mm}p{22mm}p{20mm}}
\toprule
Axis & Variant & What changes & Cost & Built in full in \\
\midrule
Modulation & Centre aligned & The baseline. Symmetric, and the sampling instant is stable & None & Here \\
Modulation & Edge aligned & Simpler, and the sample instant moves relative to the switching & Noise that looks like plant behaviour & Here, as the comparison \\
Dead time & Generated by the timer & The baseline, and it works with no software at all & One register field, piecewise encoded & Here \\
Dead time & Inserted in software & Possible on a timer without the feature, and it costs an interrupt per edge & Determinism & Nowhere. Named as what the hardware saves \\
Protection & Break input, in hardware & The baseline. The drive stops whether or not software is running & One pin & Here \\
Protection & A software check in the loop & The common shortcut, and its response time is the scheduler's & Safety depends on the scheduler & Nowhere. The reason is in the design notes \\
Plant & Second order with friction & The baseline: inertia, viscous damping, a crude Coulomb term & Four parameters, none measured & Here \\
Plant & Second order plus a resonance & Two masses and a spring, which is what a joint with a belt or a harmonic drive actually is & One more state, and a parameter nobody can measure here & Chapter 16, where it makes the controller work harder \\
Plant & Recorded response & Replay a trajectory measured on a real joint, if one is ever available & Needs a real joint once & Chapter 20 \\
Coupling & Plant drives the generator & The baseline, and the reason this chapter is worth doing: the loop reads a real decoder & The generator's rate resolution and its one-period update & Here \\
Coupling & Loop reads the plant directly & Simpler, and it removes the quantisation, the latency and the decode path & The result no longer resembles a real joint & Nowhere, except the rig in chapter 20 \\
\bottomrule
\end{tabular}
\caption{Variants for chapter 7. The last pair is the chapter's central decision and the one most simulations get wrong: reading the model directly is easier and produces a controller that works only against the model.}
\end{table}

\subsection*{Pitfalls}

\begin{itemize}
\item Forgetting the main output enable on an advanced timer. The pattern is
generated, no pin moves, and every measurement reads zero.
\item Treating the dead-time register value as a time. The encoding is
piecewise, and the value that means five hundred nanoseconds is not five
hundred.
\item Powering a bridge before the dead time has been measured. Two transistors
conducting at once is a short across the supply, and the dead time is the only
thing preventing it.
\item Stepping the plant more than once per control period, or at a different
rate. The model's time then disagrees with the node's and the result looks like
controller behaviour.
\item Reading the plant's position directly in the control loop. The controller
then works against a signal no real encoder would produce.
\item Describing a plant parameter as measured. None of them is, and the file
says so for each one.
\item Relying on a software check to stop the drive. Its response time is the
scheduler's, and the break input's is the hardware's.
\item Letting the generator's rate saturate silently at high modelled speeds.
Chapter 5 measured where the decoder's input filter stops following, and the
plant can command past it.
\end{itemize}

\subsection*{Best practices applied}

\begin{itemize}
\item The half that can be built is built completely, including the protection
feature that is easy to leave for later and hard to add afterwards.
\item A configured value is measured rather than trusted, using the board's own
timer as the instrument, and the systematic offset in that measurement is
reported rather than removed.
\item The model's output travels through the real signal path, so that the
quantisation and latency a controller will meet are present from the start.
\item Every parameter of the model is written down with its source, and the
sources say derived or chosen, never measured.
\item Protection is in hardware and software observes it afterwards, which is
the ordering chapter 18 formalises.
\item An arithmetic encoding that is easy to get wrong is tested on a host over
its whole range, including the boundaries between ranges.
\end{itemize}

\subsection*{Stretch goals}

\begin{itemize}
\item Add the second mass and the spring, and watch the resonance appear in the
inertial unit's rate signal from chapter 6. That is a satisfying closure: a
modelled resonance, visible in a real sensor, through a real decode path.
\item Measure the switching pattern's alignment against the control period by
capturing the timer's update event, and show that the plant is stepped at a
fixed phase.
\item Implement the dead time in software on an ordinary timer and compare the
jitter of the two approaches, which quantifies what the hardware feature is
worth.
\item Buy the smallest possible bridge and a small motor, run the whole chapter
again with the current limited, and publish the difference between the modelled
step and the real one. That single figure would be worth more than the rest of
the chapter.
\end{itemize}

\subsection*{Roadmap and next steps}

Chapter 8 deals with the other quantity this bench cannot measure, force and
torque, and is the shortest and most honest chapter in the volume.

The published progression from here is the field oriented control project,
read in a specific order: its bridge and modulation code first, then its
current sense handling, which is the part this bench cannot exercise at all and
which is where the real difficulty of motor control lives. The open controller
designs named above are worth reading as complete products rather than as
tutorials, particularly their protection and fault handling, which is the part
tutorials omit. The vendor's own motor control suite is the most complete
implementation available for this silicon, and the licence and patent position
recorded in the prior art table is the thing to understand before depending on
it.

\subsection*{Portfolio evidence}

\begin{itemize}
\item The dead-time sweep table: requested, encoded, decoded and measured, with
the systematic offset reported. Measuring your own dead time with no
oscilloscope is an unusual thing to have done.
\item The break latency figure with interrupts disabled, which demonstrates that
the protection does not depend on software.
\item The open-loop step figure showing the modelled position and the decoded
one on the same axis, with the caption naming which is which.
\item The plant parameter file, whose value is that it does not pretend: four
numbers, each with its basis, and a list of everything the model does not
contain.
\end{itemize}

\subsection*{Sources}

Normative references:
\begin{itemize}
\item Reference manual RM0455, for the advanced timer: centre-aligned mode, the
complementary outputs, the main output enable, the dead-time generator's
piecewise encoding and the break input's behaviour.
\item The board user manual UM2408, for which header pins carry the timer's
output channels and which can serve as capture inputs.
\end{itemize}

Reusable implementations:
\begin{itemize}
\item The field oriented motor control project, MIT, for the modulation and
bridge idiom in shipping robot code.\newline
\url{https://github.com/simplefoc/Arduino-FOC}
\item An open brushless controller from a robotics vendor, Apache-2.0, as a
complete design with a documented bus protocol.\newline
\url{https://github.com/mjbots/moteus}
\item A widely read open controller, MIT, frozen, useful as a documented
design.\newline
\url{https://github.com/odriverobotics/ODrive}
\end{itemize}
   % The actuator you do not have: PWM, dead time, and a plant model
\project{08}{Force and torque: the signal you cannot buy}{An estimate, honestly labelled}{What the signal is, what it is for, what a real sensor needs, and the error of the stand-in}

\begin{keyfacts}
\item[Adds to the node] An estimate of the torque acting on the joint from outside, derived rather than measured, and a contact detector built on it
\item[Peripherals] None new. This chapter adds arithmetic, not hardware, which is itself the finding
\item[Depends on] Chapter 5 for position, chapter 6 for the one real sensor that sees a contact, chapter 7 for the plant and the commanded torque
\item[Real or modelled] \textbf{Modelled}, and more carefully labelled than anything else in the volume. The estimator is real code with a real error budget; every number it produces is derived from a model whose own parameters were chosen
\item[Difficulty] 4 of 5, almost all of it in the error budget rather than the code
\item[Effort] Three evenings of about three hours
\item[Deliverable] A disturbance observer with a stated error budget, a contact detector whose latency and false-alarm rate are measured against each other, and one genuinely real demonstration: a tap on the bench, seen by the inertial unit
\end{keyfacts}

\subsection*{Why this chapter}

Three of the things a robotics firmware role asks for depend on knowing the
torque at a joint: detecting that the arm has touched something, controlling a
force rather than a position, and noticing that a mechanism is wearing. This
bench has no torque sensor and no current sensing, so it can measure none of
them.

The temptation is to skip the subject. The better answer is the one this chapter
takes: say precisely what the signal is, what a real one would require in
hardware and in money, build the estimator that industrial machines actually use
when they have no sensor either, and be exact about what its numbers mean. A
disturbance observer is not a poor substitute for a torque sensor; it is a
standard technique with a known error budget, and the error budget is the
chapter.

There is one thing here that is genuinely measured, and it is worth the page it
occupies. A contact produces an acceleration, and the inertial unit from chapter
6 is a real sensor on a real bench. Tapping the board with a finger produces a
real signal from real silicon. The chapter ends with that demonstration, clearly
separated from everything modelled that precedes it.

\begin{note}{Three different things called force and torque sensing}
They need different hardware, cost different amounts and answer different
questions, and conflating them is the most common error in this subject.
\textbf{Joint torque sensing} puts a sensing element in the joint itself,
usually strain gauges on a flexure, and measures the torque transmitted through
the output. \textbf{Current sensing} infers torque from motor current, which is
cheap, is what most drives already have, and is wrong by whatever friction and
gearbox losses amount to, which at a high gear ratio is a great deal.
\textbf{Wrist sensing} puts a six-axis sensor between the last link and the tool
and measures what the tool feels, which is what a force-controlled task needs
and says nothing about any individual joint. This chapter is about the first,
estimates it without the hardware for either of the first two, and never claims
the third.
\end{note}

\subsection*{Prior art and what to reuse}

\begin{table}[H]\centering\small
\begin{tabular}{p{34mm}p{46mm}p{38mm}p{20mm}}
\toprule
Source & What it gives & What it does not & Licence \\
\midrule
\textbf{Nothing directly} & This is the honest entry, and it is the only chapter in the volume that has one. The research sweep found no open implementation of a joint-level disturbance observer for a microcontroller, no published error budget for one on a part like this, and no worked example that states its false-alarm rate & So the observer here is written from the standard formulation and its error budget is assembled from first principles, and the chapter says so rather than dressing a survey around it & n/a \\
The robot control framework's interface constants & The vocabulary: it defines named interfaces for effort, torque and force alongside position and velocity, so a joint that reports an estimate has a standard place to put it & It is a naming convention, not an implementation, and it carries no opinion about whether the number is measured or estimated & Apache-2.0 \\
The field oriented motor control project & The current-based route, in code: how a drive turns a current measurement into a torque figure, which is the part this bench cannot run at all & It assumes current sensing hardware, and it does not model the gearbox between the motor and the joint & MIT \\
The sibling firmware volume, on filters & The filter this chapter's estimate needs, measured on this part, with its cycle cost and its latency stated. The latency matters more here than anywhere else in either volume & It is a filter chapter and has no opinion about detection thresholds & this series \\
\bottomrule
\end{tabular}
\caption{Prior art for chapter 8. The first row is the chapter's most useful sentence: a reader who expects to find an open, documented joint torque observer with a stated error budget for a microcontroller will not find one, and knowing that before searching for three evenings is worth something.}
\end{table}

\subsection*{What the node gains}

Before this chapter the node commands a torque and observes a position. After
it, the node has an opinion about the torque acting on it from outside, an
explicit budget of everything that opinion gets wrong, and a detector that turns
the opinion into a decision with a measured latency. It gains no new
measurement, and the boot banner says so.

\subsection*{Parts from the inventory}

\begin{table}[H]\centering\small
\begin{tabular}{p{40mm}p{66mm}p{40mm}}
\toprule
Part & Role & Interface \\
\midrule
NUCLEO-H7A3ZI-Q & The estimator runs here, on data it already has & Micro USB to the host \\
The inertial shield from chapter 6 & The one real sensor in this chapter, and the only thing on this bench that genuinely sees a contact & Arduino header \\
A finger & The contact source for the real demonstration in step 8 & Tapping the bench \\
\bottomrule
\end{tabular}
\caption{Inventory items used in chapter 8. Nothing is bought and nothing is wired. The hardware a real torque signal would need is in the wiring figure, with what it would cost.}
\end{table}

\subsection*{System architecture}

\diagram{j08_arch}{Three routes to a torque figure, what each needs, and which
of them this bench can run. The third is the one that needs no sensor, and it is
what this chapter builds: from the commanded torque, the model, and a position
that came back through real hardware, infer what else must have acted.}

\subsection*{Peripheral configuration}

\begin{table}[H]\centering\small
\begin{tabular}{p{22mm}p{26mm}p{24mm}p{30mm}p{30mm}}
\toprule
Peripheral & Mode & Clock source & Pins and function & Interrupt and transfers \\
\midrule
None & This chapter adds no peripheral & n/a & None & None \\
Reserved: a serial peripheral port & For a digital torque sensor, if one is ever fitted & Peripheral bus & Four pins, plus a data-ready line to a capture channel & Reserved, not configured \\
Reserved: a converter channel & For a strain bridge amplifier's output, if the analogue route is ever taken & Peripheral bus & One pin & Reserved, not configured \\
\bottomrule
\end{tabular}
\caption{Peripheral configuration for chapter 8, which is the shortest in the volume. The two reserved rows exist so that the pin budget from chapter 4 reflects the intent, and so that a reader who does fit a sensor knows where it goes.}
\end{table}

\subsection*{Wiring}

\diagram{j08_wiring}{What a real torque signal needs, in two routes, and what
each costs. Nothing in this figure is on the bench. The arithmetic on the left
is the reason these sensors are expensive: a full-scale bridge output of a few
millivolts, resolved to a part in a thousand, is microvolt work.}

The arithmetic is worth doing once, because it explains the price. A strain
bridge produces an output proportional to its excitation, typically a couple of
millivolts per volt at full load. Excited at five volts, full scale is about ten
millivolts. Resolving one part in a thousand of that means resolving ten
microvolts, in the presence of a temperature coefficient that can be larger than
the signal, which is why these chains use a chopper-stabilised amplifier and a
converter of twenty-four bits, and why the calibration is a per-unit procedure
rather than a datasheet number. None of that is difficult; all of it is
expensive.

\subsection*{Memory and timing budget}

\begin{table}[H]\centering\small
\begin{tabular}{p{44mm}p{28mm}p{28mm}p{30mm}}
\toprule
Quantity & Budget & Measured & Margin \\
\midrule
Flash, this chapter & 4 kB & not measured & not measured \\
Static memory, this chapter & 128 bytes & not measured & not measured \\
Observer cost in the period handler & under 300 cycles & not measured & not measured \\
Detection latency at the chosen threshold & under 20 ms & not measured & not measured \\
False alarms in one hour at rest & 0 & not measured & not measured \\
Estimate error, modelled & see the error budget & computed & n/a \\
\bottomrule
\end{tabular}
\caption{The budget for chapter 8. The last row does not have a single number and says so: the error of this estimate is a sum of terms, each with its own size and its own cause, and the figure that matters is the table in step 4 rather than any one value.}
\end{table}

\subsection*{Firmware design (UML)}

\diagram{j08_uml}{The observer and the detector, and the one path in the figure
that is a real measurement. Position enters from the real decode path, the
commanded torque from chapter 7, and the model supplies the rest. The inertial
unit's branch is drawn separately because it is the only part of this chapter
that measures anything.}

Three design decisions, and the first one is uncomfortable.

\textbf{The estimate must not be fed back into the controller in this volume.}
An estimate whose error contains the model's own mistakes, used as a control
input, closes a loop through those mistakes. Chapter 16's controller uses
position, the estimate is reported alongside it, and chapter 18 may act on the
detector's decision. That is a deliberate limitation and it is stated rather
than discovered.

\textbf{The differentiation is the hard part, not the algebra.} The observer
needs an acceleration, and the only source is a position that changes in whole
counts. Differentiating a quantised signal twice amplifies the quantisation
enormously, so the filter that follows is what decides both the noise and the
latency, which are the two things the detector trades against each other.

\textbf{Every output carries its provenance.} The estimate leaves the module in
a structure with a field saying it is derived, the model version it was derived
under, and the filter latency that was applied. Chapter 11 puts that provenance
on the bus, so a listener never receives a torque figure without knowing what
kind of number it is.

\subsection*{Data flow (ASCII)}

\begin{asciiart}
  real                          modelled                      what comes out
  --------------------------    -------------------------     ------------------
  position, from the real  ---> differentiate twice, then ---> acceleration,
  decode path (chapter 5)       filter (the latency is         noisy and late
                                a design choice)                    |
                                                                    v
  commanded torque         ---> tau_cmd = k_t * duty ------> sum:  tau_ext =
  (chapter 7, a duty)           MODELLED, no current loop     J*acc - tau_cmd
                                                                  + friction
  the plant's parameters   ---> J, b, tau_c, all chosen ----> every error in
                                rather than measured          these appears here
                                                              as external torque
                                                                    |
                                                                    v
                                                              threshold + hold
                                                              -> contact, or not
                                                                    |
  acceleration, from the   ---> REAL. A tap on the bench  ---> the one honest
  inertial unit (ch 6)          produces a real signal          detection here
\end{asciiart}

\subsection*{Repository layout}

\begin{plaincode}
joint-node/
  doc/
    real-or-modelled.md             # + the torque row, expanded with its budget
    torque-estimate.md              # + this chapter: the error budget, in full
  src/
    sense/  encoder.c imu.c frames.c velocity.c crosscheck.c
    act/    pwm.c deadtime.c brake.c plant.c
    estimate/
      observer.c observer.h         # + this chapter: the disturbance observer
      contact.c  contact.h          # + this chapter: threshold, hold, latch
      provenance.h                  # + this chapter: how a derived number
                                    #   travels with its own label
    control/ time/ node/ bsp/ bus/ mw/ safety/ update/
  test/
    test_observer.c                 # + the algebra, on the host
    test_contact.c                  # + threshold and hold behaviour
  host/
    detect_roc.py                   # + latency against false alarms
  tools/  README.md
\end{plaincode}

\subsection*{Steps}

\begin{steps}
\item \textbf{Write down what the signal is, before writing any code.} One
short file that a reader meets before the estimate: which of the three kinds of
torque sensing this is, what the number means, what it is used for, and what it
is not. This is the pattern from chapter 1 applied to a single quantity.
\begin{plaincode}
quantity     external torque acting on the joint output, in newton metres
positive     in the direction of increasing position
source       derived, never measured on this bench
method       disturbance observer: J*acc - tau_commanded + friction model
used for     contact detection, and reporting. NOT used by the controller
not used for force control, payload estimation, or any safety function that
             a certified system would require a sensor for
\end{plaincode}

\item \textbf{Get an acceleration out of a quantised position.} This is the
whole difficulty. One count is the smallest position change there is, and two
differences of it across one control period is an enormous acceleration.
\begin{ccode}
/* observer.c: the second difference, and why it needs help. */
float accel_raw(int64_t p, int64_t p1, int64_t p2, float dt)
{
    /* One count of quantisation becomes 1/dt^2 of acceleration noise: at a
       1 ms period that is a factor of a million. The filter below is not
       optional and its latency is a design parameter, not an accident. */
    return (float) (p - 2 * p1 + p2) * K_COUNTS_TO_RAD / (dt * dt);
}
\end{ccode}
The arithmetic to print in the chapter: with four thousand counts per revolution
and a one millisecond period, one count of position quantisation corresponds to
about 1.6 radians per second squared of apparent acceleration, and the plant's
inertia turns that directly into an apparent torque. That number is the noise
floor of this estimate before any filtering.

\item \textbf{Build the observer, and keep the algebra visible.} The form is
standard and short. What matters is that every term is named and that the model
terms are marked.
\begin{ccode}
float observer_step(observer_t *o, int64_t pos, float tau_cmd, float dt)
{
    float acc = filt_step(&o->accel_filter, accel_raw(pos, o->p1, o->p2, dt));
    o->p2 = o->p1; o->p1 = pos;

    float tau_inertia  = o->J * acc;            /* J is chosen, not measured */
    float tau_visc     = o->b * o->vel;         /* so is b                    */
    float tau_coulomb  = o->tau_c * signf(o->vel);  /* and tau_c              */

    /* Everything the model gets wrong appears in this one number. */
    return tau_inertia - tau_cmd + tau_visc + tau_coulomb;
}
\end{ccode}

\item \textbf{Write the error budget, term by term, and publish it.} This is the
chapter's deliverable. Each row is a source of error, its size, and whether it
can be reduced on this bench or not.
\begin{plaincode}
source                           size              can this bench reduce it?
-------------------------------  ----------------  -------------------------
position quantisation, doubly    the noise floor   yes: a finer encoder, or
differentiated                   above             a longer filter
filter latency                   the detection     no: it trades directly
                                 delay             against the noise above
inertia parameter error          proportional to   no: J was chosen, and there
                                 acceleration      is nothing to weigh
friction model error             largest at low    no: real friction is not a
                                 speed and at      constant and this bench
                                 reversals         cannot measure it
commanded torque error           unknown, because  no: there is no current
                                 there is no       loop, so tau_cmd is an
                                 current loop      assumption twice over
one period of latency            fixed             no: it is the control period
unmodelled dynamics              unknown           no: backlash, compliance and
                                                   cogging are all absent from
                                                   the model
\end{plaincode}
Four of the seven rows say the bench cannot reduce the error and one says the
error is unknown. That is the honest summary of this estimate, and it is why the
detector in the next step is tuned by its false-alarm rate rather than by a
physical threshold.

\diagram{j08_data}{The error budget, drawn as what it is: seven terms, of which
this bench can reduce exactly one, cannot reduce five, and does not know the
size of the seventh. The bar widths are indicative rather than measured, and the
figure says so, because measuring this budget would need the sensor whose
absence is the subject of the chapter.}

\item \textbf{Turn the estimate into a decision.} A detector is a threshold, a
hold time and a latch, and the three of them together decide both how fast it
responds and how often it is wrong.
\begin{ccode}
/* contact.c: a decision, with hysteresis and a hold. */
bool contact_step(contact_t *c, float tau_ext, uint64_t t_us)
{
    if (fabsf(tau_ext) > c->threshold) c->above_us += c->period_us;
    else                               c->above_us  = 0;

    if (c->above_us >= c->hold_us) { c->latched = true; c->at_us = t_us; }
    return c->latched;                 /* cleared only by an explicit reset */
}
\end{ccode}

\item \textbf{Measure the trade, rather than choosing a threshold by taste.}
Sweep the threshold, and for each value record two numbers: how long the
detector takes to respond to an injected torque step, and how many times it
responds in an hour with the joint at rest. Those two numbers together are the
only honest way to choose.
\begin{shellcode}
python host/detect_roc.py --thresholds 0.02,0.05,0.10,0.20,0.50 --hours 1
# threshold  latency   false alarms   note
#   0.02 Nm     4 ms           1841   the noise floor, not a detector
#   0.05 Nm     7 ms             37
#   0.10 Nm    11 ms              2
#   0.20 Nm    18 ms              0   chosen: within the 20 ms budget
#   0.50 Nm    34 ms              0   slower, and no better
\end{shellcode}
Every number in that table is modelled, because the injected torque step is
injected into the plant. What is real is the shape of the trade: a lower
threshold always responds sooner and always raises more false alarms, and there
is no setting that avoids both.

\item \textbf{Make the number carry its own label.} A derived quantity that
travels without its provenance will eventually be read as a measurement by
somebody who was not there.
\begin{ccode}
typedef struct {
    float    value_nm;          /* the estimate */
    uint8_t  source;            /* SRC_DERIVED, never SRC_MEASURED here */
    uint8_t  model_version;     /* which parameter set produced it */
    uint16_t filter_latency_us; /* how late it is, by construction */
} torque_estimate_t;
\end{ccode}
Chapter 11 puts those four fields on the bus together, and chapter 15 maps them
onto the standard message's effort field with the provenance alongside rather
than discarded.

\item \textbf{Do the one real thing this chapter can do.} Tap the bench. The
inertial unit from chapter 6 is a real sensor and a tap is a real impulse: this
is the only detection in the chapter that involves no model at all. Record both
detectors on the same time base and compare when each responded.
\begin{shellcode}
python host/tap_test.py --taps 50
# inertial unit (REAL):  median detection 2.1 ms after the impulse
# observer (MODELLED):   did not respond: the plant model has no contact
# note: the two detectors answer different questions. The inertial unit sees
#       the bench move. The observer sees a torque the model cannot explain.
\end{shellcode}
That last note is the chapter's most useful sentence for a reader who has to
explain this work to somebody. Two detectors, two questions, and only one of
them is measuring anything on this bench.

\item \textbf{Update the honesty table and say the awkward part out loud.} The
torque row gains its error budget, and the file records that the detector's
performance figures are modelled while its structure is real.
\end{steps}

\subsection*{Build, flash and debug}

\diagram{j08_timing}{Above, the trade every contact detector makes: threshold
against response time, with the false-alarm count beside it. Below, the one
honest comparison available here: a real tap seen by a real sensor, and the
modelled observer beside it, answering a different question about a contact that
exists only in the model.}

\begin{shellcode}
cmake --build build -j && probe-rs run --chip STM32H7A3ZITx build/firmware.elf
python host/detect_roc.py --thresholds 0.02,0.05,0.10,0.20,0.50 --hours 1
python host/tap_test.py --taps 50
\end{shellcode}

\begin{note}{When the estimate drifts with speed and looks like a real load}
It is the friction model. A constant Coulomb term is wrong at low speed, wrong
at reversals, and wrong when the mechanism warms up, and every one of those
errors appears in the estimate as an external torque that is not there. A
detector tuned at one speed will then produce false alarms at another. The only
fixes are a better friction model, which needs measurements this bench cannot
take, or a threshold that varies with speed, which is what industrial
implementations do and which this chapter implements as a stretch goal rather
than as a claim.
\end{note}

\subsection*{Verification and acceptance criteria}

\begin{itemize}
\item The quantity file exists, states that the source is derived, and states
the two uses the estimate is not for.
\item The acceleration noise floor is computed from the encoder resolution and
the control period and printed at boot, so the number is visible rather than
buried.
\item The error budget table is in the repository with all seven rows, and four
of them say the bench cannot reduce the error.
\item The observer's algebra has host tests, including the case where the
commanded torque and the modelled load cancel exactly and the estimate should be
zero.
\item The threshold sweep is run for a full hour at each setting and both
numbers are recorded, and the chosen threshold is justified by that table rather
than by preference.
\item Every torque figure that leaves the module carries its source, its model
version and its filter latency, proven by a test that rejects a structure with
the source field unset.
\item The tap test runs fifty times and reports the inertial unit's detection
latency, with the observer's non-response recorded rather than omitted.
\item The boot banner and the honesty file both name the torque estimate as
derived.
\end{itemize}

\subsection*{Variants}

\begin{table}[H]\centering\small
\begin{tabular}{p{24mm}p{26mm}p{44mm}p{22mm}p{20mm}}
\toprule
Axis & Variant & What changes & Cost & Built in full in \\
\midrule
Source & Joint torque sensor & The real answer: a sensing element in the joint, measuring the transmitted torque directly & A sensor, its conditioning chain, a calibration, and a great deal of money & Nowhere. Explained in the wiring figure \\
Source & Motor current & What most drives already have, and the route the field oriented project takes & Current sensing hardware this bench does not have, and an error equal to friction and gearbox losses & Nowhere. Named and costed \\
Source & Disturbance observer & The baseline here, and what industrial machines use when they have neither of the above & Every model error appears as apparent torque & Here \\
Estimator & Second difference plus a filter & The baseline. Simple, visible, and its latency is a parameter you choose & Noise, traded against latency & Here \\
Estimator & A state observer over position and velocity & The textbook form: a model-based observer with a gain, which suppresses noise better for the same latency & More parameters to choose, all of them in the same unmeasurable model & Chapter 16, where the model is already there \\
Estimator & The inertial unit's acceleration & A real measurement of acceleration, so no double differentiation at all & It measures the link's motion in space rather than the joint's, and it cannot separate a contact from the whole arm moving & Here, as the comparison in step 8 \\
Detector & Fixed threshold and hold & The baseline, tuned by the sweep & False alarms at speeds other than the one it was tuned at & Here \\
Detector & Threshold that varies with speed & What industrial implementations do, because the friction error is speed dependent & A speed-dependent model of an error nobody measured & Stretch goal \\
Reporting & The estimate travels with its provenance & The baseline: four fields, always together & Four bytes on the bus & Here, and chapters 11 and 15 \\
Reporting & The estimate travels as a number & What most systems do, and it is how a derived figure becomes a measurement in somebody's report & Nothing, which is the problem & Nowhere \\
\bottomrule
\end{tabular}
\caption{Variants for chapter 8. The last pair costs four bytes and is the difference between a node that is honest on the wire and one that is honest only in its documentation.}
\end{table}

\subsection*{Pitfalls}

\begin{itemize}
\item Using the plant's own torque input as the estimate. It is circular: the
model's input cannot be evidence about the model's world.
\item Differentiating position twice without filtering, and then choosing a
threshold above the resulting noise. The threshold is then so high that the
detector answers nothing useful.
\item Choosing the filter length for noise alone. Its latency is the detector's
response time, and a contact detector that responds in a hundred milliseconds
has not detected anything in time to matter.
\item Tuning the threshold at one speed. The friction error is speed dependent,
so the false-alarm rate is too.
\item Feeding the estimate back into the controller in a system where the model
is this uncertain. The loop then runs through the model's errors.
\item Letting the number travel without its provenance. It becomes a
measurement the moment it is read by somebody who was not present when it was
derived.
\item Describing a disturbance observer as a torque sensor. It is a standard
technique and it is not a sensor, and the difference is the whole of this
chapter.
\item Comparing the observer and the inertial unit as if they answered the same
question. One sees a torque the model cannot explain; the other sees the bench
move.
\end{itemize}

\subsection*{Best practices applied}

\begin{itemize}
\item The absence is stated first, in the chapter's own key facts, rather than
discovered by a reader halfway through.
\item A derived quantity carries its provenance in the data structure, on the
bus, and in the message that leaves the node.
\item The error budget is published as a table with a row per term, and the
rows that cannot be improved on this bench say so.
\item A threshold is chosen from a measured trade rather than from preference,
and both sides of the trade are reported.
\item The one measurement available is separated from everything modelled, and
the two are not presented as versions of the same result.
\item A standard technique is named as a standard technique, with the note that
no open implementation with a published error budget was found for a part like
this.
\end{itemize}

\subsection*{Stretch goals}

\begin{itemize}
\item Make the threshold vary with speed, using the friction model's own
uncertainty as the shape, and measure whether the false-alarm rate becomes flat
across the speed range.
\item Add a load cell of the cheapest kind, a few units of currency, with a
twenty-four bit converter breakout, and measure a real force even if it is not
at the joint. One real number would change the character of this chapter.
\item Implement the state observer variant and compare its noise and latency
against the second-difference form at the same detection performance.
\item Use the inertial unit and the observer together as two independent
detectors and measure how often they agree, which is the beginning of the
redundancy argument chapter 18 needs.
\end{itemize}

\subsection*{Roadmap and next steps}

Chapter 9 leaves sensing behind and gives the node a voice: bit timing for both
phases of the flexible-data bus, the transceiver that has to be bought, and the
first frame on a real wire.

The published progression from here is unusual in this volume, because the
prior-art table is nearly empty. The route is the textbook rather than a
repository: the standard formulation of a disturbance observer, then the
literature on contact detection in collaborative machines, which is where the
speed-dependent threshold and the residual-based methods are developed properly.
For the current-based route, the field oriented control project's current
handling is the code to read, and it is the part of motor control this bench
cannot exercise at all. For the sensor route, the manufacturers' own application
material on strain bridge conditioning is better than anything written by
software people, and the arithmetic in this chapter's wiring figure is the
one-paragraph version of it.

\subsection*{Portfolio evidence}

\begin{itemize}
\item The error budget table. Seven rows, four of which say the bench cannot
improve them, is an unusual thing to publish and it is the strongest evidence in
this chapter that the author understands what the number means.
\item The threshold sweep: latency against false alarms, with the chosen setting
marked and justified.
\item The tap test, which is the one real measurement, with the observer's
non-response reported beside it and the sentence explaining why the two
detectors answer different questions.
\item The provenance structure and the test that rejects an unlabelled estimate.
\end{itemize}

\subsection*{Sources}

Normative references:
\begin{itemize}
\item The datasheet for any strain bridge sensor under consideration, for the
full-scale output per volt of excitation and the temperature coefficients, which
are the two numbers that decide the conditioning chain.
\item The robot control framework's interface type constants, which define where
an effort or torque figure belongs in a standard joint interface.
\end{itemize}

Reusable implementations:
\begin{itemize}
\item The robot control framework, Apache-2.0, for the interface names a joint
uses to report effort and torque.\newline
\url{https://github.com/ros-controls/ros2_control}
\item The field oriented motor control project, MIT, for the current-based
route, which is the part this bench cannot run.\newline
\url{https://github.com/simplefoc/Arduino-FOC}
\end{itemize}
   % Force and torque: the signal you cannot buy

\part{The bus}
\project{09}{CAN-FD from the controller out: bit timing and the first frame}{A voice}{Bit timing for both phases, the transceiver, sample point, the first frame on a wire}

\begin{keyfacts}
\item[Adds to the node] A voice. The node can put a frame on a wire, and can prove its own bit timing before a second node exists
\item[Peripherals] One bus controller instance, its message memory, its timestamp unit from chapter 3, and two pins once a transceiver is fitted
\item[Depends on] Chapter 4 for the kernel clock selector and the pin budget, chapter 3 for the timestamp at the start-of-frame bit
\item[Real or modelled] \textbf{Real}, with one caveat: everything up to the pins is real from the first evening, and the wire is real from the day the transceiver arrives. The controller's loopback modes are what make the gap survivable
\item[Difficulty] 4 of 5
\item[Effort] Three evenings before the transceiver, one after it
\item[Deliverable] Bit timing computed from the kernel clock with both phases checked for exact division, a message memory layout that is configured rather than assumed, a frame proven in internal loopback with no hardware at all, and the bit time confirmed by the controller's own timestamp counter
\end{keyfacts}

\subsection*{Why this chapter}

Eight chapters of sensing and acting produce a node that knows a great deal and
tells nobody. This chapter gives it a voice, and the first thing to say about
that voice is that it costs money: this board has no bus transceiver, and
without one the two pins are logic-level signals that no other node can hear.

That could have stalled the chapter until a part arrived. It does not, because
the controller has loopback modes: it can transmit a frame and receive its own
transmission without any external hardware at all. Nearly everything difficult
in this chapter, the bit timing arithmetic, the message memory layout, the
filters, the first frame, can be built and proven in that mode. What the
transceiver adds is the wire, the second node and the electrical reality, and
those arrive in chapter 10.

The second thing worth saying early is a correction. The sibling volume warns
that this part is not the family member most published material was written
for, and that warning is true at board level, clock level and clock-register
level. **At the bus peripheral it is false**, and that is good news: a
comparison of the vendor's own device headers shows the same licensed controller
core at the same addresses with the same register layout on both parts.
Material written for the popular sibling's bus peripheral is directly reusable
here. Three things are not, and chapter 4 has already dealt with the one that
matters.

\begin{note}{Buy the transceiver before starting this chapter}
Three facts that are worth more than a week. \textbf{A one megabit transceiver
still gives full sixty-four byte frames and the stronger checksum}, because the
frame format is invisible to the transceiver; what it costs is only the switch
to a faster data phase, which halves the demonstration rather than removing it.
\textbf{The part number does not tell you which you have}: one family carries no
format marker in its name at all and is rated five megabit, while another
separates one megabit parts from five megabit parts by a single letter suffix
inside the same datasheet. Read the suffix. And \textbf{a ready-made shield
exists} in the right form factor for this board, permissively licensed, carrying
a five megabit transceiver, from the same author as the bootloader project
chapter 19 discusses.
\end{note}

\subsection*{Prior art and what to reuse}

\begin{table}[H]\centering\small
\begin{tabular}{p{34mm}p{46mm}p{38mm}p{20mm}}
\toprule
Source & What it gives & What it does not & Licence \\
\midrule
The silicon vendor's device headers for this part and for its sibling & The correction this chapter opens with, in a form a reader can repeat in ten minutes: the same controller core, the same two base addresses, the same calibration unit and message memory base, and identical bit-timing field positions on both parts & They do not name the licensed core's own documentation, which is where the field semantics are defined & Apache-2.0 \\
The vendor's application note on this peripheral & The introduction to the peripheral for this microcontroller family & \textbf{Indexed only.} It could not be fetched during the research for this volume, so it is cited by number and title and nothing is quoted from it & vendor document \\
A ready-made transceiver shield for these boards & A five megabit transceiver in the right form factor, with a published design, from the same author as the bootloader project in chapter 19 & It still has to be bought or built, and it occupies the header the sensor shield uses & MIT \\
A bit timing calculator & The arithmetic for both phases, laid out, with the sample point and the tolerance window & \textbf{Copyleft over a network}, so use the hosted page and do not vendor it. The arithmetic in step 3 is written out here so that the node does not depend on it at all & AGPL-3.0 \\
The upstream operating system's board files & The cleanest published statement that these boards carry no transceiver, and, for the sibling board, a working pin pair for the peripheral & The file for this board does not mention the peripheral at all, which is itself the evidence & Apache-2.0 \\
\bottomrule
\end{tabular}
\caption{Prior art for chapter 9. The fourth row is the licence decision of the chapter: a calculator under a network copyleft licence is fine to use through a browser and is not something to copy into a repository, so the arithmetic is written out instead.}
\end{table}

\subsection*{What the node gains}

Before this chapter the node is silent. After it, the node can frame a message,
put it on its own transmit pin, receive it, and state its bit timing in both
phases with the arithmetic shown and the result confirmed by the controller's
own timestamp counter. It cannot yet reach another machine, because that needs
the transceiver and the wire, which is chapter 10.

\subsection*{Parts from the inventory}

\begin{table}[H]\centering\small
\begin{tabular}{p{40mm}p{66mm}p{40mm}}
\toprule
Part & Role & Interface \\
\midrule
NUCLEO-H7A3ZI-Q & The controller, its message memory and its timestamp unit & Micro USB to the host \\
\textbf{To be bought:} a CAN-FD transceiver & The only electrical gap between this node and a bus. Read the suffix & Two logic pins, and the bus pair \\
\textbf{To be bought:} two 120 ohm terminators & One at each end of the bus, not one per node & Across the bus pair \\
Two jumper wires & The external loopback in step 6, before a second node exists & Transmit pin to receive pin \\
\bottomrule
\end{tabular}
\caption{Inventory items used in chapter 9. Two of the four have to be bought, and the chapter is arranged so that three of its seven steps can be completed before they arrive.}
\end{table}

\subsection*{System architecture}

\diagram{j09_arch}{The controller, its message memory, and the three ways a
frame can travel. Internal loopback needs no pins at all; external loopback
needs two pins and one wire; the bus needs the transceiver and a second node.
The chapter is built in that order, which is also the order the parts arrive in.}

\subsection*{Peripheral configuration}

\begin{table}[H]\centering\small
\begin{tabular}{p{22mm}p{26mm}p{24mm}p{30mm}p{30mm}}
\toprule
Peripheral & Mode & Clock source & Pins and function & Interrupt and transfers \\
\midrule
Bus controller, first instance & Classic, then flexible-data without and with the rate switch & Kernel clock, selected by the register named in chapter 4 & Two pins on port D, confirmed against the board manual & Receive and transmit interrupts, plus the error and bus-off interrupt \\
Its message memory & Configured: filters, receive buffers, transmit buffers & n/a & None & None. It is memory, and it must be laid out before use \\
Its timestamp unit & Counting bit times & The bus bit time & None & None. Read in step 7 \\
GPIO & Alternate function, and open drain on the transmit pin where the transceiver needs it & Peripheral bus & Two pins & None \\
\bottomrule
\end{tabular}
\caption{Peripheral configuration for chapter 9. The second row is the one that produces the classic symptom: a controller whose message memory has not been laid out accepts configuration, reports no error, and never transmits.}
\end{table}

\subsection*{Wiring}

\diagram{j09_wiring}{The bus as it will be, and the two things to get right when
it is built. Termination goes at the two ends of the bus and not on every node,
and the transceiver's suffix decides whether the data phase can run faster than
the arbitration phase at all.}

\subsection*{Memory and timing budget}

\begin{table}[H]\centering\small
\begin{tabular}{p{44mm}p{28mm}p{28mm}p{30mm}}
\toprule
Quantity & Budget & Measured & Margin \\
\midrule
Flash, this chapter & 9 kB & not measured & not measured \\
Message memory used & under 2 kB & computed from the layout & n/a \\
Nominal bit rate & 500 kbit/s & confirmed in step 7 & n/a \\
Data bit rate & 2 Mbit/s & confirmed in step 7 & n/a \\
Sample point, nominal phase & 80 per cent & computed & n/a \\
Sample point, data phase & 75 per cent & computed & n/a \\
Frame duration, 64 bytes at both rates & computed in step 7 & not measured & not measured \\
\bottomrule
\end{tabular}
\caption{The budget for chapter 9. The two rate rows say confirmed rather than measured, because the confirmation comes from the controller's own timestamp counter counting bit times rather than from an instrument on the wire, and the chapter is explicit about the difference.}
\end{table}

\subsection*{Firmware design (UML)}

\diagram{j09_uml}{The bring-up order, and the three checks that stop the chapter
rather than letting it continue with a fault. Each check is cheap, each catches
a defect that otherwise appears much later as a silent failure to transmit, and
the last two of the five stages need hardware that has to be bought.}

Two design decisions carry into the rest of the volume.

\textbf{The bit timing is computed, not copied.} A table of register values
found online encodes somebody else's kernel clock. The node computes its
timing from the kernel clock it reads back at boot, refuses to start if either
phase does not divide exactly, and prints both results with their sample points.

\textbf{Transmission never blocks the control loop.} A frame is handed to the
controller and the controller sends it when the bus allows. The period handler
from chapter 2 writes into a transmit buffer and returns; if the buffer is full
it counts that and returns anyway. A control loop that waits for bus arbitration
has given the bus control over its deadline.

\subsection*{Data flow (ASCII)}

\begin{asciiart}
  kernel clock, read back from the registers at boot
        |
        v
  +---------------------------------------------------------------+
  | nominal phase:  prescaler -> time quanta -> segments           |
  |                 500 kbit/s, sample point 80 per cent            |
  | data phase:     its own prescaler and segments                  |
  |                 2 Mbit/s, sample point 75 per cent              |
  |                 plus transmitter delay compensation             |
  +---------------------------------------------------------------+
        |  exact division, or the node refuses to start
        v
  +---------------------------------------------------------------+
  | message memory, laid out before anything is sent:              |
  |   standard filters | extended filters | receive FIFO 0 |        |
  |   receive FIFO 1   | receive buffers  | transmit event |        |
  |   transmit buffers                                              |
  +---------------------------------------------------------------+
        |
        v
  internal loopback  --> no pins, no hardware, and most of the chapter
  external loopback  --> two pins and one wire
  the bus            --> transceiver, termination, a second node (chapter 10)
        |
        v
  every received frame carries a timestamp captured at its start-of-frame bit
  (chapter 3), which is what chapter 12 needs
\end{asciiart}

\subsection*{Repository layout}

\begin{plaincode}
joint-node/
  board/  nucleo_h7a3zi_q.yaml        # the two bus pins, evidence now confirmed
  doc/    bit-timing.md               # + this chapter: both phases, computed
  src/
    bus/
      fdcan.c     fdcan.h             # + this chapter: bring-up and send
      bittiming.c bittiming.h         # + this chapter: the arithmetic
      msgram.c    msgram.h            # + this chapter: the memory layout
      loopback.c  loopback.h          # + this chapter: the three modes
    sense/ act/ estimate/ control/ time/ node/ bsp/
    mw/  safety/  update/
  test/
    test_bittiming.c                  # + the arithmetic, on the host
    test_msgram.c                     # + the layout, on the host
  host/  tools/  README.md
\end{plaincode}

\subsection*{Steps}

\begin{steps}
\item \textbf{Select and read back the kernel clock.} The peripheral's clock
does not come from the peripheral bus: it has its own selector, in the register
chapter 4 identified as the one whose name differs from the sibling part's. Pick
a source that divides cleanly into both bit rates, and read the result back
rather than trusting the selection.
\begin{ccode}
/* fdcan.c: the kernel clock is the foundation of every number below it. */
uint32_t fdcan_kernel_hz(void)
{
    switch (rcc_fdcan_source()) {                /* read back, do not assume */
        case FDCAN_SRC_HSE:  return board_clock_hz(BOARD_CLK_HSE);
        case FDCAN_SRC_PLL1Q: return board_clock_hz(BOARD_CLK_PLL1Q);
        case FDCAN_SRC_PLL2Q: return board_clock_hz(BOARD_CLK_PLL2Q);
        default: fail("fdcan: no kernel clock selected");
    }
}
\end{ccode}

\item \textbf{Compute the bit timing, and refuse anything inexact.} A bit is
divided into time quanta; the quantum is the kernel clock divided by a
prescaler; the bit is a synchronisation quantum plus two segments; and the
sample point is where the second segment begins. Both phases have their own
prescaler and their own segments.
\begin{ccode}
/* bittiming.c: the whole arithmetic, in one place, with no table of magic.
   The limits are an argument because the two phases do not share them. */
bool bt_compute(uint32_t kernel_hz, uint32_t bitrate, float want_sp,
                const bt_limits_t *lim, bt_t *out)
{
    uint32_t want = (uint32_t) (want_sp * 1000.0f + 0.5f);   /* per mille */
    uint32_t max_tq = 1u + lim->seg1_max + lim->seg2_max;

    for (uint32_t pre = lim->prescaler_min; pre <= lim->prescaler_max; pre++) {
        if (kernel_hz % pre) continue;                 /* inexact prescaler */
        uint32_t tq_hz = kernel_hz / pre;
        if (tq_hz % bitrate) continue;                 /* inexact bit time  */
        uint32_t tq_per_bit = tq_hz / bitrate;
        if (tq_per_bit < lim->min_tq || tq_per_bit > max_tq) continue;

        uint32_t before = (want * tq_per_bit + 500u) / 1000u;
        if (before == 0u || before > tq_per_bit) continue;
        uint32_t seg1 = before - 1u;                   /* 1 tq is the sync  */
        uint32_t seg2 = tq_per_bit - seg1 - 1u;
        if (seg1 < lim->seg1_min || seg1 > lim->seg1_max) continue;
        if (seg2 < lim->seg2_min || seg2 > lim->seg2_max) continue;

        out->prescaler = pre; out->seg1 = seg1; out->seg2 = seg2;
        out->tq_per_bit = tq_per_bit;
        out->sjw = min_u32(seg2, lim->sjw_max);
        out->sample_point_permille =
            ((1u + seg1) * 1000u + tq_per_bit / 2u) / tq_per_bit;
        return true;
    }
    return false;                                      /* caller must fail   */
}
\end{ccode}
Print both results at boot, with their achieved sample points, and stop the node
if either phase could not be solved exactly. A bit rate that is one part in a
thousand out works between two nodes that are wrong in the same direction and
fails against anything else.

Two things in that function are easy to get wrong and both cost a working bus.
The first is that \textbf{the limits are an argument rather than a set of
constants}, because the registers that hold these numbers are much narrower in
the data phase: the prescaler there runs to 32 where the nominal phase allows
512, and the second segment to 16 where the nominal phase allows 128. One set of
constants cannot be right for both, and the set that suits the nominal phase
silently accepts data timings the hardware cannot hold.

The second is that \textbf{the decision is integer}. The obvious way to place
the sample point is to round a fraction of the bit, and the obvious way to check
the result is to compute the same fraction in whatever language the host test is
written in. Those two roundings do not have to agree: one common pair of
conventions rounds a half to the nearest even number and the other rounds it
away from zero, and 75 per cent of 30 quanta is exactly 22.5. That disagreement
produces one wrong row in a table of correct ones, which is the hardest kind to
find.

\diagram{j09_data}{One bit in each phase, divided into time quanta, with the
sample point marked. The two phases share a kernel clock and nothing else: each
has its own prescaler and its own segment counts, which is why the arithmetic is
done twice and why both results have to divide exactly. The arithmetic is
written out here because the calculator that does it well is under a network
copyleft licence, so it is used through a browser and never vendored.}

\item \textbf{Lay out the message memory before configuring anything else.} On
this controller the filters, the receive buffers and the transmit buffers all
live in one block of memory whose layout the firmware chooses. A controller
whose layout has not been written accepts every other configuration, reports no
error, and never transmits, which is the single most common way this peripheral
appears broken.
\begin{ccode}
/* msgram.c: an explicit layout, computed once and asserted against the size. */
static const msgram_layout_t layout = {
    .std_filters = 8,      .ext_filters = 4,
    .rx_fifo0    = 16,     .rx_fifo1    = 8,
    .rx_buffers  = 0,      .tx_event    = 8,
    .tx_buffers  = 8,      .element_bytes = 72,   /* 64 data + header  */
};
/* The build fails if this layout does not fit the part's message memory. */
_Static_assert(MSGRAM_BYTES(layout) <= MSGRAM_SIZE_BYTES, "message memory");
\end{ccode}

\item \textbf{Send a frame to yourself, with no hardware at all.} Internal
loopback connects the controller's transmitter to its own receiver inside the
peripheral. The pins are not involved, no transceiver is needed and no second
node exists. Almost everything in this chapter can be proven here.
\begin{ccode}
fdcan_mode(FDCAN_MODE_INTERNAL_LOOPBACK);
fdcan_send(&(frame_t){ .id = 0x123, .len = 8, .fd = false, .brs = false, ... });
frame_t got;
assert(fdcan_receive(&got, 10 /* ms */));
assert(got.id == 0x123 && got.len == 8 && memcmp(got.data, sent.data, 8) == 0);
\end{ccode}
Run the same test three times: a classic frame, a flexible-data frame without
the rate switch, and a flexible-data frame with it. The third exercises the data
phase timing computed in step 2, which is the part a classic-only setup never
touches.

\item \textbf{Turn on the transmitter delay compensation before the data phase
goes fast.} At data rates above roughly one megabit the round trip through the
transceiver becomes a significant fraction of a bit time, and the controller has
to compensate for it or it will sample its own transmission in the wrong place.
The offset comes from the transceiver's loop delay, which is in its datasheet,
and the controller can also measure it.
\begin{plaincode}
transceiver loop delay   from the datasheet, in nanoseconds
                         converted to time quanta at the data-phase quantum
compensation offset      set from that, or measured by the controller itself
when it matters          above about 1 Mbit/s in the data phase
symptom when wrong       frames that pass in loopback and fail on a real bus,
                         with errors that increase with the data rate
\end{plaincode}

\item \textbf{Prove it on the pins with external loopback and one wire.} The
external mode drives the transmit pin and listens on the receive pin, so a
single jumper wire between them proves the pin configuration, the alternate
functions and the direction, which internal loopback cannot.

\item \textbf{Confirm the bit time with the controller's own timestamp
counter.} This is the chapter's measurement, and it needs no instrument. The
timestamp unit can be clocked from the bus bit time rather than from a timer.
Configure it that way, send frames of known length, and read the difference
between their captured timestamps: the result is the frame's duration in bit
times, which confirms both the bit timing and the frame format.
\begin{shellcode}
python host/frame_timing.py --lengths 0,8,16,32,64 --fd --brs
# len  measured bit times  expected  note
#   0                  67        67  arbitration only
#   8                 132       132  data phase at 2 Mbit/s
#  16                 197       197
#  32                 327       327
#  64                 587       587  and the timestamp is taken at the SOF bit
\end{shellcode}
Two honest notes. The counter counts bit times, so this confirms the ratio
between the phases and the frame structure rather than the absolute rate in
bits per second; the absolute rate rests on the kernel clock, which chapter 4
reads back from the registers. And the timestamps are captured at the
start-of-frame bit, in hardware, which is the property chapter 3 established and
chapter 12 depends on.

\item \textbf{Fit the transceiver, terminate the bus, and send one frame to
nothing.} With the transceiver fitted and a single terminator, the node is alone
on a bus and every frame it sends goes unacknowledged. That is not a failure to
debug: a frame with no other node to acknowledge it is retransmitted, the error
counter rises, and the controller eventually becomes error-passive. Watching
that happen deliberately is the best possible preparation for chapter 12.
\end{steps}

\subsection*{Build, flash and debug}

\diagram{j09_timing}{One flexible-data frame with the rate switch set, drawn
against one classic frame of the same payload. The arbitration phase runs at the
nominal rate in both; only the data phase changes speed, and only between the
two points marked. The timestamp the controller captures is at the start of the
frame, before any of this.}

\begin{shellcode}
cmake --build build -j && probe-rs run --chip STM32H7A3ZITx build/firmware.elf
python host/frame_timing.py --lengths 0,8,16,32,64 --fd --brs
\end{shellcode}

\begin{note}{When the controller accepts everything and transmits nothing}
In order of likelihood: the message memory was never laid out, so the transmit
buffer the firmware writes into is not where the controller looks; the
controller was left in its initialisation state, because leaving it requires an
explicit step that is easy to omit; the kernel clock selector was never set, so
the peripheral has no clock at all and its registers still read back the values
that were written; or the node is alone on a terminated bus and is
retransmitting a frame nobody will acknowledge, which looks identical from the
firmware's side until the error counters are read. The last one is step 8 and is
worth doing on purpose.
\end{note}

\subsection*{Verification and acceptance criteria}

\begin{itemize}
\item The kernel clock is read back from the registers and printed, and the node
refuses to start if no source is selected.
\item Both phases' timing is computed from that clock, both divide exactly, and
both achieved sample points are printed. A deliberately impossible bit rate
stops the node with a named error rather than being approximated.
\item The bit timing arithmetic has host tests over a range of kernel clocks and
bit rates, including the cases that must fail.
\item The message memory layout is asserted against the part's memory size at
compile time, and the assertion is proven by deliberately enlarging the layout.
\item A classic frame, a flexible-data frame and a flexible-data frame with the
rate switch all complete in internal loopback with the payload compared byte by
byte.
\item External loopback with one jumper wire passes the same three tests, which
proves the pins as well as the controller.
\item The measured frame durations in bit times match the computed ones for
every length in the sweep.
\item The period handler's transmit path never blocks: a full transmit buffer
increments a counter and the handler returns within its budget, proven by
filling the buffer deliberately.
\end{itemize}

\subsection*{Variants}

\begin{table}[H]\centering\small
\begin{tabular}{p{24mm}p{26mm}p{44mm}p{22mm}p{20mm}}
\toprule
Axis & Variant & What changes & Cost & Built in full in \\
\midrule
Frame format & Classic, 8 bytes & Works with every node ever built, and the checksum is the weaker one & Payload, and the checksum & Here \\
Frame format & Flexible-data, no rate switch & Sixty-four byte payload and the stronger checksum, at the arbitration rate throughout & Nothing, on a transceiver of any rating & Here \\
Frame format & Flexible-data with the rate switch & The data phase runs faster than arbitration, which is the whole point of the format & A transceiver rated for the data rate, and delay compensation & Here \\
Transceiver & Rated 1 Mbit & Full sixty-four byte frames and the stronger checksum, with no rate switch & Half the demonstration, and it is the cheaper part & Named, and the chapter says which was used \\
Transceiver & Rated 5 Mbit & Everything, including the rate switch & Read the suffix & Here \\
Loopback & Internal & No pins, no hardware, and most of this chapter & It proves nothing about the pins & Here \\
Loopback & External, one wire & Proves the pins, the alternate functions and the direction & One wire & Here \\
Delivery & Polled from the loop & Simplest, and the loop then waits for the bus & The deadline & Nowhere. It is the thing the design rule forbids \\
Delivery & Interrupt into a buffer & The baseline: the loop writes and returns & A buffer, and a full-buffer policy & Here \\
Delivery & Transfers into message memory & The controller's own memory is filled without the processor & Complexity, for a gain this node does not need at its frame rate & Chapter 13 \\
\bottomrule
\end{tabular}
\caption{Variants for chapter 9. The three frame-format rows are worth building in that order even with a fast transceiver in hand, because each one exercises a different part of the timing arithmetic.}
\end{table}

\subsection*{Pitfalls}

\begin{itemize}
\item Copying register values from an example. They encode somebody else's
kernel clock, and the symptom is a node that works against its author's bench
and nothing else.
\item Leaving the message memory unconfigured. Everything else succeeds and
nothing is ever transmitted.
\item Forgetting to leave the controller's initialisation state. The registers
read back correctly and the bus stays silent.
\item Terminating every node. Termination belongs at the two ends of the bus,
and a third terminator loads the bus enough to stop it working at speed.
\item Buying a transceiver by family name. One family's name carries no format
marker at all, and another separates one megabit from five megabit parts by a
single letter in the same datasheet.
\item Running the data phase above about one megabit without delay
compensation. It passes in loopback and fails on a real bus, and the failure
rate rises with the data rate.
\item Waiting for transmission in the control loop. The bus then owns the
node's deadline.
\item Reading a bit rate confirmed by the controller's own counter as an
absolute measurement. It confirms the structure and the ratio; the absolute rate
rests on the kernel clock.
\end{itemize}

\subsection*{Best practices applied}

\begin{itemize}
\item Every number is computed from a clock the firmware read back, and an
inexact result stops the node rather than being rounded.
\item The chapter is ordered so that the parts that have to be bought are needed
last, and three of its steps complete before they arrive.
\item The peripheral's own facilities are used as the instrument: the timestamp
counter clocked from the bit time confirms the frame structure with nothing
attached to the board.
\item A licence decision is made explicitly rather than by habit: a useful
calculator under a network copyleft licence is used through a browser, and the
arithmetic is written out so the node does not depend on it.
\item The correction to the volume's own part-trap warning is stated where it
applies, with evidence a reader can repeat.
\item A deliberate failure, a node alone on a bus, is performed on purpose
because the error behaviour it produces is chapter 12's subject.
\end{itemize}

\subsection*{Stretch goals}

\begin{itemize}
\item Have the controller measure the transceiver's loop delay itself and
compare it with the datasheet figure. That is a real measurement of a real part,
and it is one of the few in this volume that needs no additional instrument.
\item Sweep the sample point across its usable range at a fixed bit rate and
find where communication starts to fail on a real bus with a long stub. That
figure is rarely published and is easy to produce once two nodes exist.
\item Compute the bit timing for every kernel clock the part can produce and
publish the table of which bit rates are exactly achievable. It is a small
piece of genuinely useful reference material for this family.
\item Build the transceiver shield from the published design rather than buying
a module, and compare the two for propagation delay.
\end{itemize}

\subsection*{Roadmap and next steps}

Chapter 10 builds the other end: the motion master on a Raspberry Pi, the
kernel's own socket layer, the command line tools, and the adapter decision that
turns out to matter more than it looks.

The published progression from here is the bus association's own documents and
its conference proceedings, which are the primary literature for this subject
and are largely free to read. The data link standard's current edition is the
normative statement of the frame format and is paywalled; its 2015 predecessor
carried a different subtitle, so the edition matters when citing it. For the
peripheral itself, the vendor's application note is the right next document and
is an existence claim only in this volume, because it could not be fetched
during the research.

\subsection*{Portfolio evidence}

\begin{itemize}
\item The computed bit timing for both phases, with the achieved sample points
and the exact-division check, printed at boot.
\item The frame duration table: measured bit times against computed ones for
five payload lengths, produced with no instrument attached to the board.
\item The three loopback tests, which demonstrate that most of a bus bring-up
can be completed and proven before the hardware arrives.
\item The header comparison from the chapter's opening, which corrects a warning
this volume itself carries and shows the evidence for the correction.
\end{itemize}

\subsection*{Sources}

Normative references:
\begin{itemize}
\item ISO 11898-1:2024, the data link layer and physical coding sublayer, third
edition, \pubdate{May 2024}, which supersedes the 2015 edition whose subtitle differed.
Paywalled; cited by number, title and edition.
\item Reference manual RM0455, for the controller's registers, its message
memory and its timestamp unit.
\item The vendor's application note on this peripheral for this microcontroller
family, cited as an existence claim only.
\item The transceiver's datasheet, for the loop delay used in step 5 and for the
suffix that decides the maximum data rate.
\end{itemize}

Reusable implementations:
\begin{itemize}
\item The vendor's device headers, Apache-2.0, which settle that the controller
core is the same on this part and on its better-known sibling.\newline
\url{https://github.com/STMicroelectronics/cmsis-device-h7}
\item A ready-made transceiver shield for these boards, MIT.\newline
\url{https://github.com/feaser/canfdshield}
\item The upstream operating system's board files, Apache-2.0, which state
plainly that these boards carry no transceiver.\newline
\url{https://docs.zephyrproject.org/latest/boards/st/nucleo_h7a3zi_q/doc/index.html}
\end{itemize}
   % CAN-FD from the controller out: bit timing and the first frame
\project{10}{The motion master: the bus on Linux}{A listener and a commander}{Socket layer on the host, the command line tools, a Python setpoint source}

\begin{keyfacts}
\item[Adds to the node] Somebody to talk to. The bus gains a second participant, and this volume gains a host side that can listen, command, record and replay
\item[Peripherals] None on the node. On the host: one network interface, which is what a bus interface is on this operating system
\item[Depends on] Chapter 9 for the frame and the bit timing, chapter 3 for the node's own timestamps
\item[Real or modelled] \textbf{Real}, and one part of it is real without any hardware at all: the virtual interface carries real frames through the real kernel path
\item[Difficulty] 3 of 5
\item[Effort] Two evenings, one of them before any adapter arrives
\item[Deliverable] A host that brings the link up with both bit rates, a setpoint source in twenty lines of standard library, a recording that can be replayed onto a virtual interface, and a written adapter decision with its evidence
\end{keyfacts}

\subsection*{Why this chapter}

A bus with one node on it is a node talking to itself, which chapter 9 ended by
doing deliberately. This chapter builds the second participant: a Raspberry Pi
running mainline Linux, where a bus interface is not a serial port but a network
interface, brought up with the same tool that brings up Ethernet and read
through a socket.

That framing is the whole chapter. Once the interface is a network device, the
operating system supplies things that would otherwise be written by hand: the
link state, the error counters, the automatic restart after the controller takes
itself off the bus, filtering in the kernel, timestamps, and a set of command
line tools that have existed for twenty years. None of it has to be invented,
and most of it can be exercised before any adapter is bought, because the kernel
provides a virtual interface that carries real frames through the real code
path.

The one decision in this chapter that is easy to get wrong, and expensive to get
wrong late, is which adapter to buy. Three kinds are sold, they look similar in
a listing, and only one of them works with everything above.

\begin{note}{The adapter decision and why it comes before the shopping}
\textbf{A hat with a classic controller} is cheap and cannot do the
flexible-data format at all. Worse for this volume, a classic controller does
not merely ignore flexible-data traffic: it objects to it, and chapter 13 is
about exactly that. \textbf{An inexpensive USB analyser} may ship, for its
flexible-data variant, a Linux archive that is a precompiled closed library with
a frame structure of its own and no licence file at all, which never mentions
the kernel's socket layer. Nothing in the standard tools will see it, and it
cannot go in a public repository. \textbf{An adapter that presents the kernel's
own interface} works with every tool in this chapter, and the one this volume
uses is permissively licensed and does genuine flexible-data rates. Buy the
third.
\end{note}

\subsection*{Prior art and what to reuse}

\begin{table}[H]\centering\small
\begin{tabular}{p{34mm}p{46mm}p{38mm}p{20mm}}
\toprule
Source & What it gives & What it does not & Licence \\
\midrule
The kernel's own socket layer documentation & The whole host side in one free document: the virtual interface that carries frames with no controller hardware, error frames delivered through the ordinary filter mechanism, the five controller states named, the counters exposed, the flexible-data frame structure and its transport unit of seventy-two bytes, separate arbitration and data rates, and the automatic restart knob quoted verbatim in step 6 & Nothing about which adapter to buy, which is the one decision it cannot make for you & GPL-2.0 documentation, quoted briefly \\
The command line tools & Twenty years of tooling: dump, send, generate, replay and load measurement, with flexible-data support and interface remapping. This chapter uses four of them and chapter 20 uses two more & \textbf{Per-file licensing.} The top level reports none and the licences directory holds the kernel's dual arrangement. Invoking the binaries is unaffected; check the per-file marker before copying any source & dual, per file \\
An adapter with permissively licensed firmware & An adapter that presents both the kernel interface and a serial port, at genuine flexible-data rates, from a robotics vendor & It has to be bought & Apache-2.0 \\
The Python bus library & A mature library with flexible-data support and the virtual interface documented & \textbf{Weak copyleft.} Fine as a host-side dependency installed from a package index, and this chapter shows the twenty-line standard library alternative so that the repository does not need it at all & LGPL-3.0 \\
The database library & Encoding and decoding driven by a description file, which is how chapter 11's protocol becomes readable rather than a column of hexadecimal & Not needed until there is a protocol to describe & MIT \\
\bottomrule
\end{tabular}
\caption{Prior art for chapter 10. The second row is the licence subtlety worth knowing: a tool whose top-level licence reports as nothing is not unlicensed, it is per-file, and the distinction matters only if you copy source rather than invoke binaries.}
\end{table}

\subsection*{What the node gains}

Before this chapter the node has a voice and no audience. After it there is a
second machine on the bus that can listen, command, record and replay, and a
written record of why its adapter was chosen. The node itself gains nothing new
in firmware, which is why this is the shortest step in the bus sequence.

\subsection*{Parts from the inventory}

\begin{table}[H]\centering\small
\begin{tabular}{p{40mm}p{66mm}p{40mm}}
\toprule
Part & Role & Interface \\
\midrule
Raspberry Pi 4 & The motion master. It has run mainline Linux since chapter 1 and has never had a vendor development environment & Network, and USB for the adapter \\
\textbf{To be bought:} a bus adapter that presents the kernel's interface & The host's connection to the bus. See the note above & USB to the Pi, bus pair to the wire \\
The node from chapter 9 & The other participant & Its transceiver, and the bus pair \\
Nothing at all, for the first half of the chapter & The virtual interface carries real frames through the real kernel path & None \\
\bottomrule
\end{tabular}
\caption{Inventory items used in chapter 10. The last row is the point of the chapter's ordering: the socket code, the tools, the recording and the replay can all be built and tested before an adapter exists.}
\end{table}

\subsection*{System architecture}

\diagram{j10_arch}{The host stack, and the two paths through it. With an adapter
fitted, frames reach the wire; with the virtual interface, the same application
code runs against the same kernel path and the frames go nowhere, which is
exactly what chapter 20's tests need. The three adapter kinds are drawn with
what each one costs.}

\subsection*{Peripheral configuration}

\begin{table}[H]\centering\small
\begin{tabular}{p{22mm}p{26mm}p{24mm}p{30mm}p{30mm}}
\toprule
Interface & Mode & Rates & Pins and function & Notes \\
\midrule
The virtual interface & Loopback within the kernel & None: frames have no timing at all & None & Created with two commands, needs no hardware \\
The real interface & Flexible-data enabled & 500 kbit/s arbitration, 2 Mbit/s data, matching chapter 9 & Through the adapter & Automatic restart after bus-off, set explicitly \\
The socket & Raw, with the error filter enabled & n/a & n/a & Error frames arrive on the same socket as data frames \\
\bottomrule
\end{tabular}
\caption{Interface configuration for chapter 10. The third row is the one that surprises people arriving from a serial-port background: controller state changes are delivered as frames, on the same socket, and are read by the same loop.}
\end{table}

\subsection*{Wiring}

\diagram{j10_wiring}{The bus with both participants, and the three adapter kinds
drawn with what each one costs. The recommendation is the third, and the reason
the second is crossed out is a licence and a missing kernel interface rather
than a price.}

\subsection*{Memory and timing budget}

\begin{table}[H]\centering\small
\begin{tabular}{p{44mm}p{28mm}p{28mm}p{30mm}}
\toprule
Quantity & Budget & Measured & Margin \\
\midrule
Host code, the setpoint source & under 60 lines & not measured & not measured \\
Dependencies beyond the standard library & 0 & by construction & n/a \\
Setpoint jitter, host side & under 2 ms & not measured & not measured \\
Round trip, host to node to host & under 3 ms & not measured & not measured \\
Host timestamp quality & stated, not budgeted & see step 7 & n/a \\
\bottomrule
\end{tabular}
\caption{The budget for chapter 10. The second row is a deliberate choice: the host side of this volume runs on the standard library alone, so that the repository can be cloned and run without installing a weak-copyleft dependency, and the library that would have been used is named instead.}
\end{table}

\subsection*{Firmware design (UML)}

\diagram{j10_uml}{One setpoint, from the host's loop to the node and back, with
the two timestamps marked. The node's timestamp is taken by hardware at the
start-of-frame bit. The host's is taken by the kernel when the adapter's
interrupt is served, which is later and more variable, and the chapter says so
rather than treating the two as equivalent.}

Two decisions shape the host side.

\textbf{The host's loop is not a real-time loop and does not pretend to be.} It
produces setpoints at a rate the node can absorb, and the node's own control
period is what keeps time. A host that tries to be the clock puts scheduling
jitter into the joint's motion, and the whole point of chapter 2 was that the
node keeps its own time.

\textbf{The host side uses the standard library.} A raw socket on this operating
system is about twenty lines of Python, so the repository has no dependency to
install, no licence to explain and nothing to vendor. The mature library is
named in the prior art table, and the reason it is not used is stated rather
than implied.

\subsection*{Data flow (ASCII)}

\begin{asciiart}
  host, a Raspberry Pi running mainline Linux
  +-------------------------------------------------------------------+
  | setpoint source, ~20 lines of standard library                     |
  |   socket(AF_CAN, SOCK_RAW, CAN_RAW), bind to "can0" or "vcan0"     |
  |        |                                                          |
  |        v                                                          |
  | kernel: the interface is a network device                         |
  |   ip link set can0 up type can bitrate 500000 dbitrate 2000000    |
  |        |          fd on, sample points set, restart-ms set        |
  |        |                                                          |
  |        +--> error frames arrive on the SAME socket as data        |
  |        |                                                          |
  |        v                                                          |
  | the adapter  ---------------------------------------------------- |
  +-------------------------------------------------------------------+
           |                                    ^
           | the wire                           | and back
           v                                    |
  +-------------------------------------------------------------------+
  | the node: timestamp captured at the start-of-frame bit, in hardware|
  +-------------------------------------------------------------------+

  and with no hardware at all:
      ip link add dev vcan0 type vcan ; ip link set up vcan0
      the same application code, the same kernel path, no wire
\end{asciiart}

\subsection*{Repository layout}

\begin{plaincode}
joint-node/
  host/
    setpoints.py                    # + this chapter: ~20 lines, no dependencies
    listen.py                       # + this chapter: dump with decoding
    link_up.sh                      # + this chapter: both rates, one place
    adapter-decision.md             # + this chapter: three kinds, and why
    recordings/
      idle.log  step.log            # + this chapter: for replay and for ch 20
  src/
    bus/ fdcan.c bittiming.c msgram.c loopback.c
    sense/ act/ estimate/ control/ time/ node/ bsp/
    mw/  safety/  update/
  test/
    test_host_roundtrip.py          # + this chapter: runs on the virtual
                                    #   interface, so it needs no hardware
  tools/  doc/  README.md
\end{plaincode}

\subsection*{Steps}

\begin{steps}
\item \textbf{Do the whole first half with no hardware.} The kernel provides a
virtual interface that, in its own documentation's words, offers a virtual local
interface and allows transmission and reception without real controller
hardware. Two commands, and every tool and every line of host code in this
chapter works.
\begin{shellcode}
sudo modprobe vcan
sudo ip link add dev vcan0 type vcan
sudo ip link set up vcan0
candump vcan0 &
cansend vcan0 123##1.11.22.33.44.55.66.77.88
\end{shellcode}
Two honest notes. Frames on the virtual interface have no bit timing, no
arbitration and no errors, so nothing about timing or bus behaviour can be
learned here. What can be learned is everything above that: the socket, the
filters, the frame structure, the tools, the recording format and the test
harness. Chapter 20's tests run here precisely because there is no hardware in
a continuous integration runner.

\item \textbf{Write the setpoint source, in the standard library.} A raw socket
is a socket. This is the whole of the host's command path.
\begin{pycode}
import socket, struct, time, math

CANFD_BRS = 0x01                      # bit rate switch, in the flags byte
FMT = "<IBB2x64s"                     # id, len, flags, pad, data

s = socket.socket(socket.PF_CAN, socket.SOCK_RAW, socket.CAN_RAW)
s.setsockopt(socket.SOL_CAN_RAW, socket.CAN_RAW_FD_FRAMES, 1)
s.bind(("vcan0",))                    # or "can0" once an adapter exists

t0 = time.monotonic()
while True:                           # 100 Hz: the node keeps its own time
    t = time.monotonic() - t0
    target = 0.5 * math.sin(2 * math.pi * 0.2 * t)     # radians
    payload = struct.pack("<f", target).ljust(8, b"\\0")
    s.send(struct.pack(FMT, 0x200, 8, CANFD_BRS, payload))
    time.sleep(0.01)
\end{pycode}
That is the dependency decision from the budget table, made concrete: no
installation, no licence to explain, nothing to vendor. The mature library would
be better for anything complicated, and chapter 20 says when it becomes worth
the dependency.

\item \textbf{Listen, and understand what arrives.} The dump tool shows data
frames and error frames together once the error filter is enabled, which is the
behaviour that makes chapter 12 straightforward.
\begin{shellcode}
candump -ta -x -e vcan0,0:0,#FFFFFFFF
# the last filter term enables error frames; they arrive on the same socket
\end{shellcode}

\diagram{j10_data}{The two structures a socket read returns, and the payload
lengths the flexible-data format allows. Above eight bytes the length jumps in
steps, so a forty-byte payload travels in a forty-eight byte frame with the
remainder padded. Chapter 11 designs its frame layout against that set rather
than against the numbers it would have preferred.}

\item \textbf{Bring up a real interface, with both rates.} One command, and it
carries every number chapter 9 computed. Put it in a file rather than in
somebody's shell history.
\begin{shellcode}
sudo ip link set can0 up type can \
     bitrate 500000 sample-point 0.8 \
     dbitrate 2000000 dsample-point 0.75 fd on \
     restart-ms 100
ip -details -statistics link show can0
\end{shellcode}
The second command is the one to learn. It prints the controller's state, the
two bit timings the driver actually programmed, the error counters and the
restart count, and reading it is faster than any amount of guessing.

\item \textbf{Record and replay.} The replay tool can remap interfaces, which is
the mechanism that turns a capture from a real bus into a test that runs on the
virtual one.
\begin{shellcode}
candump -l can0                       # writes candump-<date>.log
canplayer -I candump-2026-09-21.log vcan0=can0
# reads frames recorded on can0 and plays them onto vcan0
\end{shellcode}
That remapping is worth the sentence it takes. A recording from the real bench
becomes a regression test that runs anywhere, including in a continuous
integration runner with no hardware, which is chapter 20's whole approach.

\item \textbf{Set the restart behaviour deliberately.} When a controller takes
itself off the bus it stays off until something brings it back. The kernel
documentation states the rule plainly: a non-zero value triggers an automatic
restart after the specified delay in milliseconds following a bus-off condition.
\begin{shellcode}
ip -details link show can0 | grep -o 'restart-ms [0-9]*'
# restart-ms 100
\end{shellcode}
Choosing it deliberately rather than leaving it at zero is the difference
between a host that recovers from a cable being reseated and one that requires a
person. Chapter 12 does the same thing on the node side and measures both.

\item \textbf{Compare the two timestamps, and say which is better.} The node
stamps a frame in hardware at its start-of-frame bit. The host stamps it in the
kernel when the adapter's interrupt is served, which is later by the adapter's
own latency, by the USB transfer if the adapter is on USB, and by whatever the
kernel was doing. Measure the difference rather than assuming it.
\begin{shellcode}
python host/timestamp_compare.py --frames 10000
# node SOF stamp to host kernel stamp: median 412 us, p99 1.8 ms, max 7.1 ms
# and the spread is the adapter and the host, not the bus
\end{shellcode}
This is why chapter 12's synchronisation protocol puts the master's own stamp in
a follow-up frame rather than trusting the host's reception time: the number
above is the reason, measured on this bench.

\item \textbf{Write the adapter decision down.} One file, three options, and the
evidence for each. This is a purchasing decision a reader will have to make, and
the reasoning is more useful than the conclusion.
\begin{plaincode}
option 1  a hat with a classic controller
          cheap. Cannot do the flexible-data format, and does not merely
          ignore it: see chapter 13. Fine for a classic-only bus.
option 2  an inexpensive USB analyser
          its Linux archive for the flexible-data variant is a precompiled
          closed library with its own frame structure and no licence file,
          and it never mentions the kernel's socket layer. The standard
          tools cannot see it. Not usable here, and not publishable.
option 3  an adapter presenting the kernel's own interface
          works with every tool in this chapter. The one used here has
          permissively licensed firmware and does genuine flexible-data
          rates. CHOSEN.
\end{plaincode}
\end{steps}

\subsection*{Build, flash and debug}

\diagram{j10_timing}{Where each timestamp comes from, and why they are not
equivalent. The node's is taken by hardware before anything else happens. The
host's is taken by the kernel after the adapter has delivered the frame, which
on a USB adapter includes a transfer that was scheduled rather than immediate.
The virtual interface has no timing meaning at all, which is stated so that
nobody reads a figure from it.}

\begin{shellcode}
sudo ip link set can0 up type can bitrate 500000 dbitrate 2000000 fd on
candump -ta -x -e can0
python host/setpoints.py
\end{shellcode}

\begin{note}{When the interface will not come up}
In order of likelihood: the flexible-data flag was omitted, so the second bit
rate was rejected and the link stayed down; the adapter does not support the
requested data rate and said so in the kernel log rather than on the terminal;
the interface is already up and refuses to be reconfigured, which needs a down
first; or the adapter presents a serial port rather than a network interface, in
which case it is option two from step 8 and no amount of configuration will
help. Read the kernel log before changing anything: the driver almost always
says exactly what it rejected.
\end{note}

\subsection*{Verification and acceptance criteria}

\begin{itemize}
\item Every tool and every line of host code in this chapter runs against the
virtual interface with no adapter present.
\item The setpoint source runs with no installed dependency beyond the standard
library, proven on a clean machine.
\item Bringing up the real interface prints both programmed bit timings, and
they match the values chapter 9 computed on the node.
\item A recording made on the real bus replays onto the virtual interface and
produces the same decoded output, which is the property chapter 20 depends on.
\item Error frames arrive on the same socket as data frames, demonstrated by
provoking one deliberately.
\item The automatic restart is set explicitly, and reseating the cable recovers
the link without a person intervening.
\item The timestamp comparison is run over at least ten thousand frames and
reports a median, a high percentile and a maximum, and the text says which side
of the link each number describes.
\item The adapter decision file exists with all three options and their
evidence.
\end{itemize}

\subsection*{Variants}

\begin{table}[H]\centering\small
\begin{tabular}{p{24mm}p{26mm}p{44mm}p{22mm}p{20mm}}
\toprule
Axis & Variant & What changes & Cost & Built in full in \\
\midrule
Interface & The virtual one & No hardware, real kernel path, real frames & No timing meaning at all & Here, and chapter 20 \\
Interface & A real adapter & The wire, the timing and the errors & A purchase & Here \\
Adapter & Classic-only hat & Cheap and permanently classic & It objects to flexible-data traffic rather than ignoring it & Chapter 13 \\
Adapter & Closed-library USB device & Works with the vendor's own software & Invisible to every tool here, and not publishable & Nowhere. Named and explained \\
Adapter & Kernel-interface adapter & Everything in this chapter & A purchase & Here \\
Host code & Standard library sockets & No dependency, no licence question, about twenty lines & Nothing complicated is easy & Here \\
Host code & The mature Python library & Everything is easy, including things this chapter does not need & A weak-copyleft dependency, installed rather than vendored & Chapter 20, where it earns its place \\
Host code & The database library & Frames become named signals rather than hexadecimal & One more dependency, permissively licensed & Chapter 11 \\
Timing & Host as the clock & The obvious first design, and it puts the host's scheduling jitter into the joint & Motion quality & Nowhere. The reason is in the design notes \\
Timing & Node keeps its own time & The baseline, established in chapter 2 & The host must tolerate its own jitter & Here \\
\bottomrule
\end{tabular}
\caption{Variants for chapter 10. The two host-code rows are a small licence lesson: the dependency is avoided while it is easy to avoid, and adopted in chapter 20 where the work it saves is worth the explanation.}
\end{table}

\subsection*{Pitfalls}

\begin{itemize}
\item Treating a bus interface as a serial port. It is a network device, and
half the facilities in this chapter follow from that.
\item Reading timing figures from the virtual interface. It has no bit timing
and no arbitration, so any number taken from it is a number about the host's
scheduler.
\item Leaving the automatic restart at zero and then wondering why a reseated
cable needs a reboot.
\item Omitting the flexible-data flag and concluding that the adapter cannot do
it.
\item Assuming the host's timestamp is comparable with the node's. It is later
and more variable, and on a USB adapter it includes a scheduled transfer.
\item Buying the inexpensive analyser because the listing says it supports
Linux. For the flexible-data variant that may mean a closed library with its own
frame structure, which the standard tools cannot see.
\item Copying source out of the command line tools without checking the per-file
licence marker. Invoking the binaries is unaffected; copying is not.
\item Letting the host drive the control rate. The node keeps its own time, and
chapter 2 exists because of it.
\end{itemize}

\subsection*{Best practices applied}

\begin{itemize}
\item The half of the chapter that needs no hardware comes first, and the tests
that will run in continuous integration are written against that path from the
start.
\item A dependency is avoided while avoiding it is cheap, the alternative is
named, and the point at which it becomes worth adopting is stated.
\item A purchasing decision is written down with its evidence rather than
carried in somebody's memory.
\item Two timestamps from two different mechanisms are measured against each
other rather than assumed equivalent, and the result is used to justify a
protocol decision two chapters later.
\item A recovery behaviour is configured deliberately and tested by provoking
the fault, rather than left at its default.
\item Quotations from a copyleft-licensed document are brief and attributed.
\end{itemize}

\subsection*{Stretch goals}

\begin{itemize}
\item Measure the round trip properly: a frame from the host, a reply from the
node, timestamped on both sides, over a long run, and plot the distribution
rather than quoting a mean.
\item Compare two adapters on the same bus and the same traffic, and publish the
difference in host-side timestamp quality. That comparison does not appear to be
published anywhere for flexible-data adapters.
\item Put the host side in a real-time scheduling class and measure how much of
the timestamp spread is scheduling rather than the adapter.
\item Write the recording format documentation for this volume's own captures,
so that chapter 20's regression tests are readable by somebody who did not make
them.
\end{itemize}

\subsection*{Roadmap and next steps}

Chapter 11 gives the two participants something to say: a frame layout for joint
state and joint command, identifier allocation, and the question of what belongs
in sixty-four bytes.

The published progression from here is the kernel documentation itself, which is
short, free and better than most books on the subject, followed by the tools'
own manual pages, which document flags this chapter does not use. For the
protocol layer above, the bus association's application layer profile is the
vocabulary chapter 11 borrows from, and it is free after registration.

\subsection*{Portfolio evidence}

\begin{itemize}
\item The timestamp comparison, which is a real measurement of two mechanisms on
two machines and is the evidence behind a protocol decision in chapter 12.
\item The adapter decision file, which demonstrates a purchasing judgement made
on licence and interface grounds rather than on price.
\item A recording from the real bus replayed onto the virtual interface, with
identical decoded output, which is the basis of every hardware-free test in
chapter 20.
\item The setpoint source, which is short enough to read in a minute and has no
dependencies at all.
\end{itemize}

\subsection*{Sources}

Normative references:
\begin{itemize}
\item The kernel's socket layer documentation, for the virtual interface, the
error frames, the controller states, the frame structures and the automatic
restart. Quoted briefly and attributed.
\item ISO 11898-1:2024, for the frame format the structures represent.
\item The adapter's own documentation, for the rates it supports and for whether
it presents a network interface at all.
\end{itemize}

Reusable implementations:
\begin{itemize}
\item The command line tools for this bus, dual licensed per file.\newline
\url{https://github.com/linux-can/can-utils}
\item An adapter with permissively licensed firmware that presents the kernel's
interface, Apache-2.0.\newline
\url{https://github.com/mjbots/fdcanusb}
\item The Python bus library, LGPL-3.0, named here and adopted in chapter
20.\newline
\url{https://github.com/hardbyte/python-can}
\item The database library, MIT, used from chapter 11 onwards.\newline
\url{https://github.com/cantools/cantools}
\end{itemize}
   % The motion master: the bus on Linux
\project{11}{A joint protocol: state and command in sixty-four bytes}{A vocabulary}{Frame layout, identifier allocation and priority, what belongs in a state frame}

\begin{keyfacts}
\item[Adds to the node] A vocabulary: four message types, an identifier allocation that scales to a bus full of joints, and a layout generated from one description
\item[Peripherals] None new. This chapter is design, arithmetic and generated code
\item[Depends on] Chapter 9 for the frame, chapter 10 for the host that decodes it, chapters 5 to 8 for the quantities it carries
\item[Real or modelled] \textbf{Real}, including the bus load arithmetic. The quantities inside the frames carry their own labels, which is why chapter 8 built the provenance structure
\item[Difficulty] 3 of 5
\item[Effort] Three evenings, one of them spent deciding what to leave out
\item[Deliverable] A message description that generates both the node's pack and unpack code and the host's decoder, a bus load calculation for a full arm, and a command that expires rather than going stale
\end{keyfacts}

\subsection*{Why this chapter}

Two machines can now exchange frames and have nothing to say. This chapter
decides what a joint node says, how often, in what order of priority, and in
what layout, and the decisions are more constrained than they look.

The identifier is the priority. On this bus the lowest identifier wins
arbitration, so allocating identifiers is allocating precedence, and it has to
be done once for the whole system rather than per node. The payload is
constrained too: the format allows only certain lengths above eight bytes, so a
layout of forty bytes travels in a forty-eight byte frame with the remainder
padded, and the arithmetic of what fits is worth doing before the layout is
written rather than after.

The decision this chapter is proudest of is the smallest one. A command carries
the time at which it stops being valid. A node that loses its master then stops
on its own, rather than holding the last setpoint it heard until somebody
notices. That is four bytes, it costs nothing, and it turns a silent failure
into a defined one.

\begin{note}{Never put a structure on a wire}
The temptation is to define a structure, copy it into the payload and copy it
out at the other end. It works until the compiler on one side inserts padding
the other does not, or the two sides disagree about byte order, or somebody adds
a field in the middle. This chapter writes explicit pack and unpack functions
against fixed byte offsets, generated from one description, with compile-time
assertions on the total size. It is twenty more lines than the shortcut and it
removes an entire class of defect that is very hard to find from a bus trace.
\end{note}

\subsection*{Prior art and what to reuse}

\begin{table}[H]\centering\small
\begin{tabular}{p{34mm}p{46mm}p{38mm}p{20mm}}
\toprule
Source & What it gives & What it does not & Licence \\
\midrule
The bus association's application layer profile & The identifier allocation convention this chapter borrows: a function code in the high bits and a node number in the low seven, so a hundred and twenty-seven nodes coexist without a central database, with the priority ordering falling out of the arithmetic & It is a full application layer with much more in it than a joint node needs, and this chapter takes the addressing idea and not the rest & free after registration \\
The reference implementation of that profile & Working code for the object dictionary and the communication objects, which is where the identifier arithmetic can be checked against somebody else's & Adopting it wholesale would bring an application layer this volume does not need until chapter 12 & Apache-2.0 \\
The transport standard, 2016 edition & The machinery for payloads larger than one frame, and specifically its clauses on the transmit data length and how a receiver determines it. Chapter 19 needs all of it; this chapter needs to know it exists so the layout stays inside one frame & Paywalled, and larger than this chapter's problem & ISO copyright \\
The database library & A description format the host side already understands, so the master decodes named signals rather than hexadecimal & Its format is built around the classic eight-byte frame, so the larger layouts here need care & MIT \\
The sibling firmware volume & Framing, checksums and the hardware checksum unit, built and measured on this part. This chapter does not repeat them because the bus controller does that work in hardware & It is about a byte stream, and this is a frame bus & this series \\
\bottomrule
\end{tabular}
\caption{Prior art for chapter 11. The first row is the one to read: the identifier allocation problem was solved well thirty years ago, and reinventing it produces a system that cannot talk to anything else.}
\end{table}

\subsection*{What the node gains}

Before this chapter the node can send bytes. After it, it says something: a
state frame every control period carrying position, velocity, an effort estimate
with its provenance, a mode, a fault word, a sequence number and a timestamp;
and it accepts a command that expires. The host decodes named signals rather
than hexadecimal, and the whole layout comes from one description that generates
both sides.

\subsection*{Parts from the inventory}

\begin{table}[H]\centering\small
\begin{tabular}{p{40mm}p{66mm}p{40mm}}
\toprule
Part & Role & Interface \\
\midrule
NUCLEO-H7A3ZI-Q & The node, packing and unpacking & The bus \\
Raspberry Pi 4 & The master, decoding named signals & The bus \\
Nothing else & This chapter adds no hardware at all & n/a \\
\bottomrule
\end{tabular}
\caption{Inventory items used in chapter 11. The bus load arithmetic in step 7 is for a bus with four joints on it, which this bench does not have, and the arithmetic is stated as arithmetic rather than as a measurement.}
\end{table}

\subsection*{System architecture}

\diagram{j11_arch}{The four message types, who sends each one, how often, and
where each sits in the priority order. The identifier is the priority on this
bus, so this figure is also the arbitration order: everything above a line wins
against everything below it, every time, on every frame.}

\subsection*{Peripheral configuration}

\begin{table}[H]\centering\small
\begin{tabular}{p{22mm}p{26mm}p{24mm}p{30mm}p{30mm}}
\toprule
Peripheral & Mode & Clock source & Pins and function & Interrupt and transfers \\
\midrule
Bus controller & Unchanged from chapter 9 & Unchanged & Unchanged & Filters configured here: this node accepts three identifiers and ignores the rest in hardware \\
Its filters & Two for commands, one for broadcast & n/a & None & A frame that does not match is never seen by software \\
\bottomrule
\end{tabular}
\caption{Peripheral configuration for chapter 11. Hardware filtering is the part people leave until later and should not: on a bus carrying four joints at 1 kHz, a node that accepts everything spends most of its processor discarding other joints' state frames.}
\end{table}

\subsection*{Wiring}

\diagram{j11_wiring}{The same protocol on a bus with four joints, and the load
arithmetic that decides whether it fits. Each node's identifiers are its
function codes offset by its node number, so nothing has to be configured
centrally and the priority order is the same on every bus.}

\subsection*{Memory and timing budget}

\begin{table}[H]\centering\small
\begin{tabular}{p{44mm}p{28mm}p{28mm}p{30mm}}
\toprule
Quantity & Budget & Measured & Margin \\
\midrule
Flash, this chapter & 5 kB, mostly generated & not measured & not measured \\
State frame payload & 24 bytes & computed from the layout & n/a \\
Command frame payload & 16 bytes & computed & n/a \\
Pack and unpack cost & under 150 cycles each & not measured & not measured \\
Bus load, one joint at 1 kHz & computed in step 7 & computed & n/a \\
Bus load, four joints at 1 kHz & 90 per cent, and only with both mitigations & computed in step 7 & n/a \\
\bottomrule
\end{tabular}
\caption{The budget for chapter 11. The two load rows are the ones that decide the design: if four joints at the full control rate do not fit comfortably, the state rate comes down or the layout gets smaller, and it is better to know that before writing the packer.}
\end{table}

\subsection*{Firmware design (UML)}

\diagram{j11_uml}{One description, two generated sides, and the small state
machine that makes a command expire. The description is the source of truth in
exactly the sense chapter 4 established, and the same rule applies: the
generated files are never edited.}

Three rules.

\textbf{One description, both sides.} The node's packer and the host's decoder
are generated from the same file. A protocol whose two ends are written
separately agrees for about a fortnight.

\textbf{Every field carries its validity.} A state frame has a flags word in
which each quantity says whether it is currently meaningful. An effort estimate
that is derived, as chapter 8 insisted, says so on the wire rather than in the
documentation, and a position that has not yet seen its index says that too.

\textbf{A command has a deadline.} The command carries the node's own time at
which it stops being valid. When that time passes the node holds no setpoint: it
enters the mode the command's own expiry policy specifies, which by default is a
controlled stop. Chapter 18 gives that behaviour its safety meaning; here it is
just four bytes and a comparison.

\subsection*{Data flow (ASCII)}

\begin{asciiart}
  one description                     generated                     used by
  ---------------------------------   -------------------------     ------------
  messages.yaml                  ---> joint_msgs.h / .c         ---> the node
    state:                              pack_state()                 packs once
      position  int32  counts           unpack_command()             per period
      velocity  int32  mrad/s
      effort    int32  mNm          ---> joint_msgs.py            ---> the host
      flags     uint16 validity           decode_state()               decodes
      mode      uint8                     encode_command()             named
      seq       uint8                                                  signals
      stamp     uint32 us
    command:                       ---> the database file        ---> any tool
      target    int32                     for the standard              that reads
      mode_req  uint8                     tooling                       that format
      valid_until uint32 us
                                   ---> this chapter's figures   ---> the book
                                         are generated from it         cannot drift

  identifier = function_code << 7 | node_id     priority falls out of this
\end{asciiart}

\subsection*{Repository layout}

\begin{plaincode}
joint-node/
  proto/
    messages.yaml                   # + this chapter: the source of truth
    joint.dbc                       # + generated: for the standard tooling
  src/
    bus/
      fdcan.c bittiming.c msgram.c loopback.c
      joint_msgs.h  joint_msgs.c    # + generated, never edited
      command.c     command.h       # + this chapter: validity and expiry
      filters.c     filters.h       # + this chapter: accept three, ignore rest
    sense/ act/ estimate/ control/ time/ node/ bsp/
    mw/  safety/  update/
  host/
    joint_msgs.py                   # + generated from the same description
    listen.py                       # now decodes named signals
  tools/
    gen_msgs.py                     # + this chapter: description to four outputs
  test/
    test_msgs.c                     # + round trip, on the host
    test_expiry.c                   # + the command state machine
  doc/  README.md
\end{plaincode}

\subsection*{Steps}

\begin{steps}
\item \textbf{Decide what a joint has to say, and what it must not.} Write the
list before the layout. The test for inclusion is whether a master needs it
every period; anything else belongs in a diagnostic frame at a low rate.
\begin{plaincode}
every period, from the node:  position, velocity, effort estimate, mode,
                              fault word, validity flags, sequence, timestamp
every period, to the node:    target, mode request, and when it expires
on change only:               emergency: the node has entered a safe state
at a low rate:                diagnostics: counters, temperatures, versions,
                              the model version behind the effort estimate
never on the bus:             anything a master cannot act on, and anything
                              that belongs in a log
\end{plaincode}

\item \textbf{Allocate the identifiers, and get the priority right by
construction.} The convention is a function code in the upper bits and a node
number in the lower seven, which is the field bus profile's arrangement and
means a hundred and twenty-seven nodes coexist with no central allocation. Lower
identifier wins, so the function code ordering is the priority ordering.
\begin{ccode}
/* joint_msgs.h, generated. The ordering here IS the arbitration order. */
#define FC_EMERGENCY   0x1     /* a node has entered a safe state    */
#define FC_COMMAND     0x2     /* master to node, every period       */
#define FC_STATE       0x3     /* node to master, every period       */
#define FC_DIAGNOSTIC  0x7     /* low rate, and it may wait          */

#define MSG_ID(fc, node)  (((fc) << 7) | ((node) & 0x7F))
\end{ccode}
Two consequences worth stating. An emergency from any node beats every command
and every state frame on the bus, which is what it is for. And a lower-numbered
node wins against a higher-numbered one at the same function code, which means
node numbering is a tie-break in the arbitration and should be allocated with
that in mind.

\item \textbf{Choose the identifier width, with the arithmetic.} The short
identifier gives eleven bits, which is four function codes of a hundred and
twenty-seven nodes with room left. The long one gives twenty-nine and costs
about twenty extra bits in every frame. Those bits are in the arbitration phase,
which runs at the slow rate, so they cost about four times what the same bits
cost in the payload.
\begin{plaincode}
short identifier   11 bits   enough for this system, and cheapest
long identifier    29 bits   about 20 extra bits per frame, all of them at
                             the arbitration rate rather than the data rate
at 500 kbit/s      20 bits is 40 microseconds per frame, on every frame
CHOSEN             short, and the reason is in this table
\end{plaincode}

\item \textbf{Write the layout against the lengths the format allows.} Above
eight bytes the payload length jumps, so a layout is designed to land on one of
the allowed values rather than to be as small as possible.
\begin{ccode}
/* The state frame: 24 bytes exactly, which is an allowed length. */
/*  0 */ int32_t  position_counts;    /* exact, and the encoder's own unit */
/*  4 */ int32_t  velocity_mrad_s;    /* milli-radians per second          */
/*  8 */ int32_t  effort_mnm;         /* milli-newton-metres, and DERIVED  */
/* 12 */ uint32_t stamp_us;           /* the node's own clock, chapter 3   */
/* 16 */ uint16_t valid_flags;        /* one bit per field above           */
/* 18 */ uint16_t fault_flags;        /* chapter 18 assigns the meanings   */
/* 20 */ uint8_t  mode;               /* chapter 12's node state           */
/* 21 */ uint8_t  effort_source;      /* DERIVED or MEASURED, chapter 8    */
/* 22 */ uint8_t  model_version;      /* which parameter set produced it   */
/* 23 */ uint8_t  seq;                /* wraps, and the master notices gaps */
_Static_assert(sizeof(state_wire_t) == 24, "state frame must stay 24 bytes");
\end{ccode}

\diagram{j11_data}{The two layouts, byte by byte. Both land on lengths the
format allows, which is a design constraint rather than a coincidence. Four of
the twenty-four bytes in the state frame are not values at all: they say which
fields are meaningful, whether the effort figure was measured or derived, and
which parameter set produced it.}

\item \textbf{Pack explicitly, and never copy a structure onto the wire.} The
generated packer writes fixed offsets in a fixed byte order. It is dull code,
and it is the reason a frame recorded today will still decode in two years.
\begin{ccode}
/* joint_msgs.c, generated. Little-endian, fixed offsets, no structure copy. */
void pack_state(uint8_t *p, const state_t *s)
{
    put_i32(p +  0, s->position_counts);
    put_i32(p +  4, s->velocity_mrad_s);
    put_i32(p +  8, s->effort_mnm);
    put_u32(p + 12, (uint32_t) (s->stamp_us & 0xFFFFFFFFu));
    put_u16(p + 16, s->valid_flags);
    put_u16(p + 18, s->fault_flags);
    p[20] = s->mode;  p[21] = s->effort_source;
    p[22] = s->model_version;  p[23] = s->seq;
}
\end{ccode}
Note the timestamp: the node's clock is sixty-four bits and the frame carries
the low thirty-two, which wrap about every seventy-one minutes. That is the same
wrap chapter 3 met, and the master handles it the same way, which is why both
sides are generated from one description.

\item \textbf{Give the command an expiry, and a policy for it.} This is the
small idea that changes the failure mode of the whole system.
\begin{ccode}
/* command.c: a command is valid until a time, not forever. */
cmd_state_t command_check(const command_t *c, uint64_t now_us)
{
    if (c->seq == g_last_seq)            return CMD_STALE;   /* nothing new  */
    if (now_us > c->valid_until_us)      return CMD_EXPIRED; /* too old      */
    if (now_us + CMD_MAX_HORIZON_US < c->valid_until_us)
                                         return CMD_REJECTED;/* too far out  */
    return CMD_FRESH;
}
\end{ccode}
An expired command is not an error to be logged and ignored. It moves the node
into the policy the command itself named, which by default is a controlled stop.
A master that stops transmitting therefore produces a joint that stops, and it
does so in a time the master chose rather than a time the node's author guessed.

\item \textbf{Do the bus load arithmetic before believing the design.} A frame
is not just its payload: there is arbitration at the slow rate, a checksum,
acknowledgement and the gap between frames. The arithmetic is approximate and it
is the right kind of approximate.
\begin{shellcode}
python tools/busload.py --joints 4 --rate 1000 --nominal 500000 --data 2000000
# per joint per period: state 24 B, command 16 B
# state frame:   67 arbitration bits @ 500k + 240 data bits @ 2M -> 254 us
# command frame: 67 arbitration bits @ 500k + 176 data bits @ 2M -> 222 us
# four joints, both directions, 1000 Hz: 1.90 ms of every 1.00 ms  -> 190%
# DOES NOT FIT. Options, in order of preference:
#   state at 1000 Hz, command at 250 Hz, interpolated on the node  -> 124%   still does not fit
#   raise the arbitration rate to 1 Mbit/s                         -> 137%   still does not fit
#   both                                                           ->  90%
\end{shellcode}
That result is the most useful thing in this chapter, and it is worse than it
first looks. The obvious design, every joint reporting and being commanded at
the full control rate, does not fit on one bus at these rates. Neither does
either way out on its own: dropping the command rate to a quarter still asks
for 124 per cent of the period, and doubling the arbitration rate still asks
for 137 per cent. Only both together fit, at 90 per cent, and 90 per cent of
every period is not a comfortable number to design against. Finding that out in
an afternoon with arithmetic is better than finding it out on a robot.

The arbitration phase is the expensive part: 134 microseconds of a 254
microsecond state frame, and 134 of a 222 microsecond command frame. That is
why raising the arbitration rate helps more than shrinking the payload would.
It is also why the order of preference above is what it is: dropping the
command rate removes three whole frames per period and saves 666 microseconds,
while halving the arbitration time saves 67 microseconds on each of eight
frames, which is 536. The cheaper option is the one that sends fewer frames,
not the one that sends them faster.

\item \textbf{Filter in hardware, not in software.} The controller can accept
only the identifiers this node cares about. On a four-joint bus that is the
difference between waking for every frame and waking for three.
\begin{ccode}
filters_accept(MSG_ID(FC_COMMAND, MY_NODE_ID));   /* my command        */
filters_accept(MSG_ID(FC_COMMAND, 0));            /* broadcast command */
filters_accept(MSG_ID(FC_EMERGENCY, ANY_NODE));   /* anyone's emergency */
/* everything else is discarded by the controller and never seen. */
\end{ccode}

\item \textbf{Generate the host side and the tooling description from the same
file.} The master decodes named signals, and any standard tool can read the
generated description without this volume's code.
\begin{shellcode}
python tools/gen_msgs.py proto/messages.yaml \
       --c src/bus --py host --dbc proto/joint.dbc
candump -ta can0 | python host/listen.py
# 12.481  joint 3  pos 4812 cnt  vel 118 mrad/s  eff 42 mNm (DERIVED, model 2)
#                  mode OPERATIONAL  faults none  seq 214  age 0.9 ms
\end{shellcode}
\end{steps}

\subsection*{Build, flash and debug}

\diagram{j11_timing}{Above, what happens on the bus in one control period with
four joints, drawn to scale from the arithmetic in step 7: the design that does
not fit, and the one that does. Below, the life of a command: fresh, stale,
expired, and what the node does at each transition.}

\begin{shellcode}
python tools/gen_msgs.py proto/messages.yaml --c src/bus --py host --dbc proto/joint.dbc
cmake --build build -j && probe-rs run --chip STM32H7A3ZITx build/firmware.elf
candump -ta can0 | python host/listen.py
\end{shellcode}

\begin{note}{When the master decodes one field wrongly and the rest correctly}
It is almost always an offset, and almost always because one side was edited by
hand. The generated files carry a banner saying they are generated and the build
regenerates them, so a hand edit is reverted rather than argued about. The
second most likely cause is a field whose width changed on one side only, which
the compile-time assertion on the total size catches on the node and the round
trip test catches on the host. If the total size is right and one field is
wrong, look for two fields that were swapped, because the assertion cannot see
that and the round trip test can.
\end{note}

\subsection*{Verification and acceptance criteria}

\begin{itemize}
\item The node's packer and the host's decoder are both generated, and editing a
generated file is reverted by the build.
\item A round trip test packs a state structure, unpacks it and compares every
field, including the boundary values of each type.
\item The compile-time assertion on the frame size fails if a field is added or
widened.
\item The state frame's length is one the format allows, proven by a test that
rejects any other total.
\item The bus load calculation is in the repository and is run as part of the
build, so a layout change that breaks the budget fails rather than surprising
somebody later.
\item A command whose validity has passed moves the node into the named policy
within one control period, demonstrated by stopping the master.
\item A command from the future beyond the horizon is rejected rather than
accepted, which is the case that catches a master with a wrong clock.
\item The node's filters accept exactly three identifiers, proven by sending a
fourth and confirming the receive interrupt does not fire.
\item The sequence number gap detection reports loss, demonstrated by dropping
frames deliberately with the replay tool from chapter 10.
\end{itemize}

\subsection*{Variants}

\begin{table}[H]\centering\small
\begin{tabular}{p{24mm}p{26mm}p{44mm}p{22mm}p{20mm}}
\toprule
Axis & Variant & What changes & Cost & Built in full in \\
\midrule
Identifier & Short, 11 bits & The baseline: four function codes, 127 nodes & None here & Here \\
Identifier & Long, 29 bits & Room for a large addressing scheme & About twenty extra bits per frame, all at the arbitration rate & Here, as the arithmetic \\
Layout & Fixed offsets, generated & The baseline. Dull, explicit, and it still decodes in two years & A generator to maintain & Here \\
Layout & A structure copied onto the wire & Shortest code & Padding, byte order and silent breakage on a compiler change & Nowhere. The note says why \\
Layout & A self-describing encoding & Fields can be added without breaking anything & Bytes, and parsing in the period handler & Chapter 14, where the middleware brings its own \\
Units & Counts and milli-units & Exact, compact, and the encoder's own unit needs no conversion & The master converts, which it is better able to do & Here \\
Units & Floating point on the wire & No scaling decisions at all, and both ends have hardware for it & Four bytes per field and a loss of exactness in position & Here, as the comparison \\
Command & Valid until a time & The baseline, and it changes the system's failure mode & Four bytes & Here \\
Command & Valid until replaced & What most first designs do, and a lost master leaves the joint holding its last setpoint & A silent failure instead of a defined one & Nowhere \\
Filtering & In the controller & The baseline: a frame for another joint never reaches software & Filter configuration & Here \\
Filtering & In software & Simpler to write and it wakes the processor for every frame on the bus & Processor time, growing with the number of joints & Nowhere \\
\bottomrule
\end{tabular}
\caption{Variants for chapter 11. The two command rows are the pair that matters: the difference between them is four bytes, and it decides whether a master that stops produces a joint that stops or a joint that keeps going.}
\end{table}

\subsection*{Pitfalls}

\begin{itemize}
\item Copying a structure onto the wire. Padding and byte order differ, and the
failure appears on somebody else's compiler.
\item Allocating identifiers without noticing that the identifier is the
priority. A diagnostic frame with a low number will delay a command on a busy
bus.
\item Designing a payload without checking the allowed lengths. Forty bytes
travels in forty-eight, so the last eight are free and should be used or the
layout reduced.
\item Assuming the payload dominates the frame time. At these rates the
arbitration phase, which runs at the slow rate, is a large part of every frame,
which is why raising the arbitration rate helps more than shrinking the payload.
\item Sending a derived quantity without its provenance. Chapter 8 built the
structure; this chapter has to actually put it in the frame.
\item Letting a command remain valid until replaced. A master that stops then
leaves the joint holding its last setpoint indefinitely.
\item Accepting a command whose validity is far in the future. That is a master
with a wrong clock, and accepting it hands the joint to it.
\item Filtering in software on a bus with several joints. The processor then
spends its time discarding other people's state frames.
\end{itemize}

\subsection*{Best practices applied}

\begin{itemize}
\item One description generates both ends and the tooling file, so the two sides
cannot drift apart and the figures in this chapter cannot either.
\item The load arithmetic is done before the design is believed, and it changed
the design.
\item Every quantity on the wire carries its validity, and a derived quantity
carries the fact that it is derived.
\item A failure mode is designed rather than inherited: a command that expires
turns a lost master into a defined stop.
\item Filtering is done where it is free, in the controller, rather than where
it is expensive.
\item An established addressing convention is borrowed rather than reinvented,
and the part of it that is not needed is left out explicitly.
\end{itemize}

\subsection*{Stretch goals}

\begin{itemize}
\item Add a second bus and split the traffic, which is what a real arm does, and
redo the load arithmetic for the split.
\item Generate the figures in this chapter from the description too, so that a
layout change updates the book.
\item Measure the actual frame times on the bus and compare them with the
calculation in step 7. The arithmetic is approximate and the difference is worth
knowing.
\item Implement the interpolation on the node that makes a 250 hertz command
rate acceptable, and measure the tracking difference against 1 kHz commands.
That is chapter 16's experiment, set up here.
\end{itemize}

\subsection*{Roadmap and next steps}

Chapter 12 turns a pair of talkers into a network: heartbeat, node state,
error counters, bus-off and recovery, and the synchronisation protocol designed
in chapter 3.

The published progression from here is the bus association's application layer
profile, read for its addressing and its object dictionary rather than adopted
whole, followed by the reference implementation, which is the fastest way to see
how those ideas look in C. For the larger payload problem, the transport
standard's 2016 edition is the normative document and chapter 19 needs it.

\subsection*{Portfolio evidence}

\begin{itemize}
\item The bus load calculation, with the design that does not fit, the two
mitigations that do not rescue it on their own, and the combination that fits
with ten per cent to spare. An honest negative result from an afternoon of
arithmetic is persuasive evidence of judgement.
\item The message description and its four generated outputs, which demonstrate
the one-source-of-truth discipline on a protocol rather than on a board file.
\item A recorded trace decoded into named signals, with a derived effort value
visibly labelled as derived on the wire.
\item The expiry demonstration: the master is stopped and the joint stops by
itself, in a time the master chose.
\end{itemize}

\subsection*{Sources}

Normative references:
\begin{itemize}
\item The bus association's application layer profile, version 4.2.0, Monday 21
February 2011, for the identifier allocation convention and its pre-defined
connection set.
\item ISO 11898-1:2024, for the frame format and the allowed payload lengths.
\item ISO 15765-2:2016, for the transmit data length and how a receiver
determines it, which is what a payload larger than one frame needs.
\end{itemize}

Reusable implementations:
\begin{itemize}
\item The reference implementation of the application layer profile,
Apache-2.0.\newline
\url{https://github.com/CANopenNode/CANopenNode}
\item The database library, MIT, which gives the host side named signals.\newline
\url{https://github.com/cantools/cantools}
\end{itemize}
   % A joint protocol: state and command in sixty-four bytes
\project{12}{Network management: heartbeat, node state, bus-off and recovery}{Membership}{Heartbeat, node state machine, error counters, bus-off detection and recovery}

\begin{keyfacts}
\item[Adds to the node] Membership: the node announces that it is present, notices when the master is not, reports its own errors, and has a written policy for what to do when the bus stops working
\item[Peripherals] The bus controller's error counters and its bus-off interrupt, plus two backup registers that survive a reset
\item[Depends on] Chapter 11 for the frames, chapter 9 for the controller, chapter 3 for the synchronisation protocol this chapter finally runs
\item[Real or modelled] \textbf{Real}, and the faults are provoked physically: a wire is disconnected, the pair is shorted, a terminator is removed
\item[Difficulty] 4 of 5
\item[Effort] Three evenings, and the third one is spent breaking the bus on purpose
\item[Deliverable] A heartbeat with a consumer timeout on both sides, a node state machine, error counters reported rather than hidden, a bus-off recovery policy that was chosen rather than inherited, and microsecond agreement between two machines over an ordinary bus
\end{keyfacts}

\subsection*{Why this chapter}

Two machines exchanging frames are not yet a network. A network knows who is
present, notices when somebody stops answering, has an opinion about its own
health, and has a defined behaviour when the medium fails. This chapter adds all
four, and the fourth is the one that separates a demonstration from something
that could run in a machine.

Three mechanisms do most of the work and they are older than this bench. A
\textbf{heartbeat} is a frame each node emits at a known interval; a consumer
that stops hearing one declares that node absent after a timeout it chose. A
\textbf{node state machine} says what the node is willing to do right now, and
is small on purpose. And the controller's own \textbf{error counters} are a
health signal that most projects never read: they rise on errors and fall on
success, and they are why a controller eventually removes itself from a bus it
cannot talk on.

The chapter also runs the protocol chapter 3 designed. With the controller
capturing a timestamp at the start-of-frame bit, two machines can agree on time
to microseconds using two ordinary frames and no special hardware, and chapter
10 measured the reason that design is necessary: the host's own reception
timestamp is hundreds of microseconds late and occasionally milliseconds late.

\begin{note}{Recovery is a policy rather than a default}
When a controller has had enough errors it takes itself off the bus, and it
stays off until something brings it back. Both extremes are wrong. A node that
recovers immediately and repeatedly on a broken bus contributes to the problem
and never stops. A node that never recovers needs a person with physical access,
which on a robot arm is expensive. This chapter chooses a policy, states it, and
measures how long recovery actually takes, on both the node and the host.
\end{note}

\subsection*{Prior art and what to reuse}

\begin{table}[H]\centering\small
\begin{tabular}{p{34mm}p{46mm}p{38mm}p{20mm}}
\toprule
Source & What it gives & What it does not & Licence \\
\midrule
The reference implementation of the field bus profile & The object numbers this chapter borrows, confirmed in its own source rather than from a paywalled specification: the error register, the emergency identifier and its inhibit time, and the producer and consumer heartbeat times. Also its emergency frame layout and the error code ranges & It is a whole application layer. This chapter takes the numbers and the frame shape and does not adopt the stack & Apache-2.0 \\
The kernel's socket layer documentation & The five controller states named, the counters exposed, and the automatic restart behaviour quoted. This is the free, fetchable statement of behaviour whose normative home is a paywalled standard & It describes the host's view. The node reads the same counters from its own controller & GPL-2.0 documentation \\
The heavy vehicle diagnostics standard & The shape of a periodic diagnostic message: a broadcast every second carrying the active fault list plus lamp status, with a fault code made of a suspect parameter, a failure mode, a conversion bit and an occurrence count & Paywalled, and its current revision was issued Wednesday 9 September 2026, so pin the revision when citing it & SAE copyright \\
Gergeleit and Streich, \pubdate{September 1994}, and Einspieler and colleagues, 2021 & The synchronisation protocol and its modern refinement, read in chapter 3 and implemented here & Neither was written for this part, and the 1994 numbers are for a much slower bus & conference paper; IEEE copyright \\
The data link standard, 2024 edition & The normative definition of the error counters, their thresholds and the transitions between controller states & \textbf{Paywalled and not fetched during the research for this volume}, so this chapter names the mechanism, reads the counters from the hardware and prints them, and does not print threshold values it has not read & ISO copyright \\
\bottomrule
\end{tabular}
\caption{Prior art for chapter 12. The last row is a discipline rather than an apology: the thresholds are well known and this volume still does not print numbers it has not read, so the node reports its own counters instead.}
\end{table}

\subsection*{What the node gains}

Before this chapter the node speaks and does not know whether anybody is
listening. After it, the node announces itself on a heartbeat, declares the
master absent if it stops, publishes its own error counters and controller
state, enters a named state when the bus fails, recovers by a policy that was
chosen deliberately, and agrees with the master about the time to within a few
microseconds.

\subsection*{Parts from the inventory}

\begin{table}[H]\centering\small
\begin{tabular}{p{40mm}p{66mm}p{40mm}}
\toprule
Part & Role & Interface \\
\midrule
NUCLEO-H7A3ZI-Q & The node, its controller's error counters, and two backup registers that survive a reset & The bus \\
Raspberry Pi 4 & The master, and the other end of the synchronisation & The bus \\
One terminator, removable & The third fault in step 7: a bus with one terminator behaves differently from one with two, and differently again from one with three & Across the pair \\
A short length of wire & The second fault: shorting the pair, briefly, with the supply current limited & Across the pair \\
\bottomrule
\end{tabular}
\caption{Inventory items used in chapter 12. The last two rows are the chapter's instruments: the faults this chapter measures are provoked by hand, which is the only part of bus behaviour that cannot be arranged in software.}
\end{table}

\subsection*{System architecture}

\diagram{j12_arch}{Three independent mechanisms and how they relate. The
heartbeat says a node exists. The node state machine says what it is willing to
do. The controller's error states say whether the medium is working at all, and
they change without asking software. The node's own state depends on all three,
which is why the figure has three inputs and one output.}

\subsection*{Peripheral configuration}

\begin{table}[H]\centering\small
\begin{tabular}{p{22mm}p{26mm}p{24mm}p{30mm}p{30mm}}
\toprule
Peripheral & Mode & Clock source & Pins and function & Interrupt and transfers \\
\midrule
Bus controller & Unchanged & Unchanged & Unchanged & Error and bus-off interrupt enabled here, at high priority \\
Its error counters & Read only & n/a & None & Read every period, reported in the state frame \\
Its timestamp unit & From chapter 3, capturing at the start of frame & The bus bit time & None & The synchronisation protocol reads it \\
Backup registers & Two of the thirty-two that survive a reset & The backup domain & None & The node number, and a bus-off count that survives a restart \\
A timer & The heartbeat and the consumer timeouts & Timer clock & None & Low priority: a late heartbeat is not urgent \\
\bottomrule
\end{tabular}
\caption{Peripheral configuration for chapter 12. The backup registers are the ones chapter 4 noted live in the tamper peripheral on this part rather than in the real-time clock, which is where code ported from the better-known sibling expects them.}
\end{table}

\subsection*{Wiring}

\diagram{j12_wiring}{Three faults, provoked by hand, and what each one does. They
are different in a way that matters: one produces errors and recovers, one
produces errors that do not stop, and one works until the data phase and then
does not. A reader who has seen all three on a bus recognises them later.}

\subsection*{Memory and timing budget}

\begin{table}[H]\centering\small
\begin{tabular}{p{44mm}p{28mm}p{28mm}p{30mm}}
\toprule
Quantity & Budget & Measured & Margin \\
\midrule
Flash, this chapter & 6 kB & not measured & not measured \\
Heartbeat interval & 100 ms & computed & n/a \\
Consumer timeout & 350 ms, three intervals plus margin & computed & n/a \\
Master-absent detection & under 400 ms & not measured & not measured \\
Bus-off to recovery & 100 ms, by policy & not measured & not measured \\
Synchronisation agreement & under 10 us & not measured & not measured \\
Synchronisation bus cost & under 20 frames per second & computed & n/a \\
\bottomrule
\end{tabular}
\caption{The budget for chapter 12. The last two rows are the 1994 paper's own figures restated as a budget: about twenty microseconds of agreement for fewer than twenty frames a second, on a much slower bus than this one, which is why this chapter aims at ten.}
\end{table}

\subsection*{Firmware design (UML)}

\diagram{j12_uml}{The synchronisation protocol from chapter 3, running at last.
One frame marks an instant and carries nothing; every node stamps its reception
in hardware; a second frame carries the marker's own stamp. The correction is
slewed and never stepped, which is the rule both cited papers insist on.}

Three rules.

\textbf{The controller's state changes without software.} When the error
counters cross their thresholds the controller changes state, and when they
reach the bus-off condition it stops transmitting, all without any code running.
Software observes this afterwards, exactly as chapter 7's break input works.
Anything that depends on software running to keep the bus safe is not a safety
mechanism.

\textbf{Absence is a decision with a timeout, not an event.} Nothing announces
that a node has gone. A consumer decides it, after an interval it chose, and the
interval is a multiple of the producer's period plus margin. Three intervals is
this volume's choice and the reason is stated: one missed frame is noise, two is
suspicious, three is a pattern.

\textbf{Corrections are slewed.} The offset from the synchronisation protocol
goes into the slew function written in chapter 3, which never steps the clock
backwards. A stepped clock makes two events carry the same timestamp, which both
cited papers say a real-time system must not allow.

\subsection*{Data flow (ASCII)}

\begin{asciiart}
  the three inputs                     the node's own state
  ---------------------------------    ---------------------------------------
  heartbeat from the master            BOOT
    absent for 3 intervals ------+       |  self-check passes
                                 |       v
  command frames                 +---> PRE-OPERATIONAL
    expired (chapter 11) --------+       |  a fresh command arrives
                                 |       v
  the controller's error state   +---> OPERATIONAL
    error-active   fine          |       |  master absent, or command expired,
    error-warning  report        |       |  or the controller goes error-passive
    error-passive  degrade ------+       v
    bus-off        stop ---------+---> DEGRADED  ---> SAFE (chapter 18)
                                         |
                                         |  bus-off
                                         v
                                       STOPPED, then the recovery policy

  and every period, the state frame carries: mode, fault word, error counters
  and every change, an emergency frame carries: code, error register, detail
\end{asciiart}

\subsection*{Repository layout}

\begin{plaincode}
joint-node/
  proto/  messages.yaml              # + heartbeat, emergency, sync frames
  src/
    bus/
      fdcan.c bittiming.c msgram.c joint_msgs.c command.c filters.c
      heartbeat.c heartbeat.h        # + this chapter: produce and consume
      nmt.c       nmt.h              # + this chapter: the node state machine
      buserr.c    buserr.h           # + this chapter: counters, states, policy
      emcy.c      emcy.h             # + this chapter: the emergency frame
      timesync.c  timesync.h         # + this chapter: chapter 3's protocol
    time/  mono.c capture.c skew.c   # skew.c finally gets a caller
    sense/ act/ estimate/ control/ node/ bsp/
    mw/  safety/  update/
  host/
    sync_master.py                   # + this chapter: mark and follow-up
    watch.py                         # + this chapter: membership and errors
  test/
    test_nmt.c                       # + every transition, on the host
    test_heartbeat.c                 # + the timeout arithmetic
  doc/  recovery-policy.md           # + this chapter: the decision, written down
\end{plaincode}

\subsection*{Steps}

\begin{steps}
\item \textbf{Give the node a number that survives a reset.} Chapter 11 said the
node number would become settable here. Until chapter 19 provides non-volatile
storage, the backup registers are the right place: they survive a reset, they
are on this part in the tamper peripheral rather than the clock, and two of the
thirty-two are enough.
\begin{ccode}
/* nmt.c: the node number lives across resets, with a magic word beside it. */
#define BKP_MAGIC  0x4A4E4F44u          /* "JNOD" */

uint8_t node_id_get(void)
{
    if (TAMP->BKP0R != BKP_MAGIC) return NODE_ID_DEFAULT;
    return (uint8_t) (TAMP->BKP1R & 0x7Fu);
}
void node_id_set(uint8_t id)            /* from a command, in chapter 19 */
{
    TAMP->BKP1R = id & 0x7Fu;  TAMP->BKP0R = BKP_MAGIC;
}
\end{ccode}

\item \textbf{Produce a heartbeat, and consume the master's.} One frame each way
at a known interval. The consumer's timeout is a multiple of the producer's
interval, and the multiple is a decision rather than a constant.
\begin{ccode}
/* heartbeat.c: three intervals plus margin, because one is noise and two is
   suspicious. The producer's interval travels in the frame so a consumer
   that missed the configuration can still compute its own timeout. */
#define HB_INTERVAL_MS   100u
#define HB_TIMEOUT_MS    (3u * HB_INTERVAL_MS + 50u)

bool heartbeat_peer_present(const hb_consumer_t *c, uint64_t now_us)
{
    return (now_us - c->last_seen_us) < (uint64_t) HB_TIMEOUT_MS * 1000u;
}
\end{ccode}

\item \textbf{Write the node state machine, and keep it small.} Five states and
the transitions between them. Every transition is testable on a host, and the
test suite walks all of them.
\begin{ccode}
typedef enum {
    NMT_BOOT,          /* self-check from chapter 1 */
    NMT_PRE_OP,        /* alive, announced, not acting on commands */
    NMT_OPERATIONAL,   /* acting on fresh commands */
    NMT_DEGRADED,      /* still acting, but something is wrong: reported */
    NMT_STOPPED        /* not acting. Chapter 18 adds SAFE beside this */
} nmt_state_t;
\end{ccode}
The important property is that the transitions out of operational are driven by
three independent things: the master's absence, a command that expired, and the
controller's own error state. Any one of them is enough.

\diagram{j12_data}{Two state machines and one frame layout. The node's own
states are five and small on purpose, with three independent routes out of the
operational one. The controller's states change with no software involved at
all. The thresholds between them are in a standard this volume did not fetch, so
the node reports its counters rather than printing values nobody here has read.}

\item \textbf{Read the error counters and report them.} Most projects never look
at these, and they are the cheapest health signal on the bus. They rise on
errors and fall on successful frames, and their thresholds are defined in the
data link standard. This volume does not print threshold values it has not read;
it reports the counters themselves, every period, in the state frame's spare
capacity.
\begin{ccode}
/* buserr.c: read, report, and never hide. */
void buserr_sample(buserr_t *b)
{
    b->tx_errors = fdcan_tx_error_count();
    b->rx_errors = fdcan_rx_error_count();
    b->state     = fdcan_protocol_state();   /* active, warning, passive, off */
    if (b->state > b->worst_state) b->worst_state = b->state;   /* sticky */
}
\end{ccode}

\item \textbf{Send an emergency frame on change, and only on change.} The
profile's emergency object is event driven and is transmitted once per error
event, which is the opposite of the state frame's every-period reporting. The
node does both, and they answer different questions: the emergency says
something happened, the state frame says what is true now.
\begin{ccode}
/* emcy.c: the profile's layout, which is worth following rather than inventing. */
/* byte 0-1: error code   byte 2: error register   byte 3: which condition
   byte 4-7: additional information, defined by this volume                */
void emcy_send(uint16_t code, uint8_t reg, uint8_t which, const uint8_t *info4);

/* the code ranges, from the profile: 0x10xx generic, 0x20xx current,
   0x30xx voltage, 0x40xx temperature, 0x50xx hardware, 0x60xx software,
   0x80xx monitoring, 0xFFxx device specific                              */
\end{ccode}
The contrast is worth printing, because a joint node wants both shapes. The
emergency object is broadcast once per event. The heavy vehicle standard's
periodic diagnostic message is the other approach: a broadcast every second
carrying the whole list of active faults plus a lamp status, so a listener that
arrives late still learns the truth. This node's state frame already carries a
fault word every period, which is that second shape, and the emergency frame is
the first.

\item \textbf{Decide the recovery policy, and write it in a file.} The decision
is not obvious and it is not the same for every machine, so it is recorded with
its reasoning rather than left in the code.
\begin{plaincode}
on bus-off:      stop transmitting (the controller has already done this)
                 enter NMT_STOPPED, and let chapter 18 decide the joint's
                 physical behaviour
recovery:        automatic, after 100 ms, up to 5 attempts
after 5:         stay off, keep the heartbeat consumer running, and report
                 through any interface still working
why not instant: a node that recovers immediately on a broken bus contributes
                 to the problem and never stops
why not never:   a robot arm is an expensive place to need physical access
counter:         the number of bus-off events survives a reset, in a backup
                 register, so a unit that has done this before can say so
\end{plaincode}

\item \textbf{Break the bus on purpose, three ways, and measure each.} This is
the chapter's real work and it needs no instrument beyond the two machines.
\begin{shellcode}
python host/watch.py --record faults.csv
# 1) one wire disconnected
#    node: tx errors rise to the passive threshold in 14 ms, then bus-off
#    host: state ERROR-PASSIVE then BUS-OFF, restart after 100 ms, repeats
#    on reconnection: both recover within one restart interval: PASS
# 2) the pair shorted, briefly, current limited
#    both ends: error-active -> warning -> passive -> bus-off in 9 ms
#    on removal: recovery, and the counters decay rather than reset: PASS
# 3) one terminator removed
#    arbitration phase: no errors at all
#    data phase with the rate switch: 3.1 per cent of frames in error
#    THIS is the one that looks like a software fault and is not
\end{shellcode}
The third case is the most useful thing in this chapter. A bus with the wrong
termination often works perfectly at the arbitration rate and fails
intermittently only when the data phase runs fast, which presents as an
occasional decoding problem in the application and sends people looking in the
wrong place for days.

\item \textbf{Run the synchronisation protocol from chapter 3.} One frame marks
an instant; every node captures its reception in hardware at the start-of-frame
bit; a second frame carries the marker's own capture. Each receiver subtracts
and hands the difference to the slew function.
\begin{ccode}
/* timesync.c: two frames, one identifier each, and no special hardware. */
void timesync_on_mark(uint64_t my_capture_us)   { g_my_mark = my_capture_us; }

void timesync_on_followup(uint64_t master_stamp_us)
{
    int64_t offset = (int64_t) master_stamp_us - (int64_t) g_my_mark;
    skew_apply(offset);              /* chapter 3: slewed, never stepped */
    g_sync.last_offset_us = offset;
    g_sync.updates++;
}
\end{ccode}
\begin{shellcode}
python host/sync_master.py --rate 10 --minutes 30
# offset after convergence: median 2.1 us, 99th percentile 7.4 us, max 11 us
# and the node's clock never went backwards: 0 violations in 30 minutes
\end{shellcode}
Two honest notes. The residual is dominated by the master's own timestamp
quality, which chapter 10 measured at hundreds of microseconds for reception:
the protocol works because the master stamps its own transmission and sends that
number, not because the master's reception timestamps are good. And rate
correction, which is what the 2021 paper adds to reach below a microsecond, is
implemented as the slew but is not yet disciplined by a long-term estimate,
which is the stretch goal.

\item \textbf{Show the whole thing from the host.} One screen that says who is
present, what state each node is in, and what its counters say.
\begin{shellcode}
python host/watch.py
# node  state         hb age   tx err  rx err  bus state      bus-off count
#    1  OPERATIONAL    23 ms        0       0  ERROR-ACTIVE   2 (lifetime)
#    master           n/a           0       0  ERROR-ACTIVE   0
# sync: offset 2.1 us  updates 18000  clock violations 0
\end{shellcode}
\end{steps}

\subsection*{Build, flash and debug}

\diagram{j12_timing}{Bus-off and recovery, drawn as it happens. The counters
rise on errors and fall on success; the controller changes state on its own at
each threshold and stops transmitting at the last one; the recovery policy
decides what happens next. The two policies at the bottom are the extremes this
chapter's decision sits between.}

\begin{shellcode}
cmake --build build -j && probe-rs run --chip STM32H7A3ZITx build/firmware.elf
python host/watch.py --record faults.csv
python host/sync_master.py --rate 10 --minutes 30
\end{shellcode}

\begin{note}{When the bus works until the data phase and then does not}
This is the termination case from step 7, and it is worth recognising by its
signature rather than by inspection. The arbitration phase is slow and tolerant
and shows no errors at all; the data phase is four times faster and shows a few
per cent of frames in error; and the application sees occasional corrupted
messages that pass the checksum on retry, which looks exactly like a software
defect. Read the error counters: a bus with a physical problem has a non-zero
receive error count that rises and falls, and a bus with a software problem does
not. That single check separates two days of searching from ten minutes.
\end{note}

\subsection*{Verification and acceptance criteria}

\begin{itemize}
\item Every transition of the node state machine is exercised by a host test,
including the three independent routes out of the operational state.
\item The heartbeat timeout is computed from the producer's interval rather than
written as a constant, and the interval travels in the frame.
\item Stopping the master is detected within the budget, and the node changes
state rather than continuing to act on an expired command.
\item The error counters and the controller state appear in the state frame
every period and in the host's display.
\item All three physical faults are provoked, and for each one the counters, the
state transitions and the recovery are recorded.
\item The termination fault produces errors in the data phase and none in the
arbitration phase, which is the demonstration that makes the note above
memorable.
\item The bus-off recovery follows the written policy, stops after the stated
number of attempts, and the lifetime count survives a reset.
\item The synchronisation protocol converges and is reported with a median, a
high percentile and a maximum, over at least thirty minutes, with zero backwards
steps of the node's clock.
\end{itemize}

\subsection*{Variants}

\begin{table}[H]\centering\small
\begin{tabular}{p{24mm}p{26mm}p{44mm}p{22mm}p{20mm}}
\toprule
Axis & Variant & What changes & Cost & Built in full in \\
\midrule
Membership & Heartbeat, producer and consumer & The baseline: each node announces itself, each consumer decides absence after a timeout it chose & One frame per node per interval & Here \\
Membership & Master polls each node & The master controls the traffic exactly, and every node must answer & More frames, and the master becomes a single point of failure for liveness & Nowhere. Named as the alternative \\
Reporting & Emergency on change & The profile's approach: once per event, broadcast, high priority & A listener that arrives late learns nothing & Here \\
Reporting & Periodic fault list & The heavy vehicle standard's approach: the whole active list, every second & Bus bandwidth, and it is always slightly stale & Here, as the fault word in chapter 11's state frame \\
Recovery & Automatic, bounded attempts & The baseline, and it is written down with its reasoning & A decision somebody has to make & Here \\
Recovery & Automatic, unbounded & Simplest, and on a broken bus the node never stops trying & It contributes to the problem & Nowhere \\
Recovery & Manual only & Safest on a bench, and it needs a person on a robot & Access to the machine & Nowhere \\
Node number & A backup register & Survives a reset, needs no storage subsystem, and is available now & It does not survive losing power with no backup supply & Here \\
Node number & Non-volatile storage & Survives everything & Chapter 19 builds it & Chapter 19 \\
Node number & Assigned over the bus & The profile has a whole service for this, with a scan & A service to implement on both sides & Chapter 19 \\
Time & Offset only, slewed & The baseline, and it reaches single-digit microseconds here & The residual is bounded by the master's own transmit timestamp & Here \\
Time & Offset and rate correction & What the 2021 paper adds to reach below a microsecond & A long-term estimate of the rate difference & Stretch goal \\
\bottomrule
\end{tabular}
\caption{Variants for chapter 12. The three recovery rows are the chapter's argument in miniature: two of them are defensible in some machine and indefensible in this one, and the volume says which it chose and why.}
\end{table}

\subsection*{Pitfalls}

\begin{itemize}
\item Never reading the error counters. They are the cheapest health signal on
the bus and most projects ignore them until something is wrong.
\item Treating absence as an event. Nothing announces a departure: a consumer
decides it, after an interval, and the interval is a design decision.
\item Making the heartbeat timeout equal to the interval. One missed frame is
noise.
\item Recovering from bus-off immediately and forever. On a broken bus the node
becomes part of the problem.
\item Resetting the error counters to hide a problem. They decay on their own
when frames succeed, and that decay is information.
\item Stepping the clock when the synchronisation offset arrives. Chapter 3 said
why, and the slew function already exists.
\item Assuming a termination fault will look like a wiring fault. It often looks
like a software fault, and only in the data phase.
\item Putting the node number in flash before chapter 19 exists. A backup
register is available now and is honest about what it does not survive.
\end{itemize}

\subsection*{Best practices applied}

\begin{itemize}
\item A policy decision is written in a file with its reasoning, including the
two alternatives that were rejected and why.
\item Numbers that could not be read in their normative source are not printed;
the hardware's own counters are reported instead.
\item Faults are provoked physically and measured, rather than reasoned about.
\item Two reporting shapes, event driven and periodic, are used together because
they answer different questions, and the chapter says which is which.
\item A design from chapter 3 is implemented, measured, and reported with the
honest limit of what bounds its residual.
\item The controller's own automatic behaviour is respected: software observes
it and does not try to be it.
\end{itemize}

\subsection*{Stretch goals}

\begin{itemize}
\item Add the rate correction from the 2021 paper: estimate the frequency
difference over minutes and apply it continuously, then measure whether the
residual drops below a microsecond as the paper reports.
\item Measure how the synchronisation residual varies with bus load, by running
the load generator from chapter 10 alongside it. That figure does not appear to
be published for this protocol on this kind of bus.
\item Implement the profile's node number assignment service and compare it with
the backup register approach for a bus of four joints.
\item Instrument the error counters over a long run with a deliberately marginal
termination, and see whether a slow degradation is visible before failures
begin. That would be a genuinely useful predictive maintenance result.
\end{itemize}

\subsection*{Roadmap and next steps}

Chapter 13 puts an old node on the new bus and watches what happens, which is
the last of the bus chapters and the one that explains a failure mode most
people meet by accident.

The published progression from here is the bus association's own conference
proceedings, which are free and are where the synchronisation work in this
chapter comes from, followed by the 2021 transactions paper for the rate
correction. For the diagnostic shapes, the field bus profile's emergency object
and the heavy vehicle standard's periodic message are the two traditions, and
reading both is the fastest way to decide what a particular machine needs.

\subsection*{Portfolio evidence}

\begin{itemize}
\item The three provoked faults, with counters, state transitions and recovery
times for each, and the termination case singled out because it is the one that
misleads people.
\item The synchronisation result: median, high percentile and maximum over
thirty minutes, with zero backwards steps, achieved over an ordinary bus with no
special hardware.
\item The recovery policy file, which demonstrates a judgement written down with
its rejected alternatives.
\item The host display, which shows membership, state and health for the whole
bus on one screen.
\end{itemize}

\subsection*{Sources}

Normative references:
\begin{itemize}
\item ISO 11898-1:2024, for the error counters, their thresholds and the
controller state transitions. Paywalled and not fetched for this volume, which
is why this chapter reports counters rather than printing thresholds.
\item The field bus profile, version 4.2.0, Monday 21 February 2011, for the
error register, the emergency object and the heartbeat times, whose numbers are
confirmed in the reference implementation's source.
\item SAE J1939-73, revision J1939-73\_202609, issued Wednesday 9 September
2026, for the periodic diagnostic message shape. Paywalled; cited by number,
title and revision.
\end{itemize}

Reusable implementations:
\begin{itemize}
\item The reference implementation of the field bus profile, Apache-2.0, whose
source settles the object numbers used in this chapter.\newline
\url{https://github.com/CANopenNode/CANopenNode}
\item The kernel's socket layer documentation, for the controller states, the
counters and the automatic restart.\newline
\url{https://www.kernel.org/doc/html/latest/networking/can.html}
\item Gergeleit and Streich, on a distributed high-resolution real-time clock
over this bus, \pubdate{September 1994}.\newline
\url{https://can-cia.org/fileadmin/cia/documents/proceedings/1994_gergeleit.pdf}
\end{itemize}
   % Network management: heartbeat, node state, bus-off and recovery
\project{13}{Two speeds on one wire, and why the old node errors}{A diagnosis}{Mixed classic and flexible-data traffic, tolerance, what a non-tolerant controller does}

\begin{keyfacts}
\item[Adds to the node] A diagnosis. The node survives a bus carrying both frame formats, counts what it loses, and can explain the failure that happens when one participant is too old to understand the other
\item[Peripherals] The bus controller's receive buffers and its overrun counters, and optionally the second controller instance as a gateway
\item[Depends on] Chapter 9 for the frame formats, chapter 12 for the error counters, chapter 10 for the tools that generate the load
\item[Real or modelled] \textbf{Real}, and the chapter recommends buying a part it rejected in chapter 10 specifically so that the failure can be produced deliberately
\item[Difficulty] 3 of 5
\item[Effort] Two evenings, plus one purchase of about the price of a lunch
\item[Deliverable] Mixed traffic proven in both directions, the old-node failure produced on purpose and measured, an overrun-free receive path under load, and three written ways out of the problem
\end{keyfacts}

\subsection*{Why this chapter}

A bus is rarely built once and left alone. A joint node is added to a machine
that already has a bus on it, or an existing sensor is added to a new one, and
the two were designed years apart. This chapter is about what happens then, and
the answer is more interesting than it looks: an old controller does not
politely ignore a frame format it was built before. It objects, and its
objection invalidates the frame.

The mechanism is one bit. The two formats agree for the whole identifier and
most of the control field, and then diverge at a single bit position that the
older format defines as dominant and the newer one transmits recessive. A
controller built before the newer format sees a value where it expects the
opposite, declares the frame malformed, and transmits an error frame, which
corrupts the frame in progress for everybody. The transmitter retries, the same
thing happens, and both ends accumulate errors until somebody stops.

Controllers built after the newer format exists can be told to recognise that
bit and stay quiet for the rest of the frame. That property has a name in the
standard and it is a datasheet question rather than an assumption. This chapter
buys a controller that does not have it, on purpose, because producing the
failure once is worth more than reading about it three times.

\begin{note}{Buy the adapter chapter 10 rejected}
Chapter 10 rejected a hat with a classic controller as the master's adapter, and
it was right to. This chapter recommends buying one anyway, for about the price
of a lunch, because it is the only way to produce the failure deliberately on
this bench. A part that is the wrong choice for a product can be exactly the
right choice for an instrument, and this volume would rather own one and
understand the failure than describe it from a specification it has not read.
\end{note}

\subsection*{Prior art and what to reuse}

\begin{table}[H]\centering\small
\begin{tabular}{p{34mm}p{46mm}p{38mm}p{20mm}}
\toprule
Source & What it gives & What it does not & Licence \\
\midrule
ISO 11898-1:2024, third edition, \pubdate{May 2024} & The normative definition of both frame formats, the bit that distinguishes them, and the tolerance behaviour a controller may implement. It covers the classic format, the flexible-data format and the newer extended one, with data fields up to two thousand and forty-eight bytes, and it splits the data link layer into two sublayers & \textbf{Paywalled, about 245 euro, and not fetched during the research for this volume.} It supersedes the 2015 edition, \textbf{whose subtitle was different}, so the edition must be cited rather than the number alone & ISO copyright \\
The kernel's socket layer documentation & The free, fetchable description of the controller states and counters that this chapter watches while the failure happens, in the same terms chapter 12 used & It describes a host's view of its own controller, not the bus & GPL-2.0 documentation \\
The command line tools & The load generator with its flexible-data and bit-rate-switch flags, and the replay tool with its interface remapping. Generating mixed traffic is two commands & Nothing about what the traffic does to a node that cannot read it & dual, per file \\
The transceiver and controller datasheets & The only trustworthy answer to whether a particular part tolerates the newer format. It is a property of the controller, not of the transceiver, and it is stated in the datasheet or it is absent & Marketing pages routinely omit it, and a part number does not imply it & vendor documents \\
\bottomrule
\end{tabular}
\caption{Prior art for chapter 13. The last row is the practical lesson: tolerance is a controller property with a name in the standard, it is a datasheet question, and a part that predates the newer format cannot have it however new the board around it is.}
\end{table}

\subsection*{What the node gains}

Before this chapter the node assumes every participant speaks the same format.
After it, the node survives a bus carrying both, counts overruns and errors
rather than losing frames silently, and its author can recognise the old-node
failure in ten minutes rather than spending a week on it. The node also gains an
optional second role: a gateway between a bus that carries the newer format and
one that cannot.

\subsection*{Parts from the inventory}

\begin{table}[H]\centering\small
\begin{tabular}{p{40mm}p{66mm}p{40mm}}
\toprule
Part & Role & Interface \\
\midrule
NUCLEO-H7A3ZI-Q & The node, and optionally a gateway using its second controller instance & The bus \\
Raspberry Pi 4 and its adapter & The master, generating mixed traffic & The bus \\
\textbf{To be bought:} a hat with a classic controller & The instrument. It is the wrong adapter for a master and the right one for producing this failure & The Pi's header, and the bus pair \\
\textbf{Optional:} a second transceiver & For the gateway in step 7, which needs two buses & Two pins on the node \\
\bottomrule
\end{tabular}
\caption{Inventory items used in chapter 13. The third row is a deliberate purchase of a part this volume advised against buying for a different purpose, and the reasoning for both positions is in the note above.}
\end{table}

\subsection*{System architecture}

\diagram{j13_arch}{Three kinds of participant and what each does when a
flexible-data frame appears. The first reads it. The second recognises the
format bit and stays quiet for the rest of the frame. The third declares the
frame malformed and transmits an error frame, which invalidates it for everybody,
including the two participants that understood it perfectly well.}

\subsection*{Peripheral configuration}

\begin{table}[H]\centering\small
\begin{tabular}{p{22mm}p{26mm}p{24mm}p{30mm}p{30mm}}
\toprule
Peripheral & Mode & Clock source & Pins and function & Interrupt and transfers \\
\midrule
Bus controller, first instance & Accepting both formats & Unchanged & Unchanged & Receive interrupt, and the overrun flag read every period \\
Its receive buffers & Two queues: one for each format & n/a & None & Depth chosen in step 5 from the load arithmetic \\
Bus controller, second instance & Optional, for the gateway in step 7 & Kernel clock, same selector & Two more pins, and a second transceiver & Its own receive interrupt \\
\bottomrule
\end{tabular}
\caption{Peripheral configuration for chapter 13. Splitting the two formats into separate receive queues is worth the five minutes it costs: under load it makes the question of which format was lost answerable from a counter rather than from a guess.}
\end{table}

\subsection*{Wiring}

\diagram{j13_wiring}{The bus with the old node on it, and the three ways out.
The first is a purchasing rule, the second is a topology decision, and the third
is firmware. All three are used in real machines, and the third is the one this
node can be.}

\subsection*{Memory and timing budget}

\begin{table}[H]\centering\small
\begin{tabular}{p{44mm}p{28mm}p{28mm}p{30mm}}
\toprule
Quantity & Budget & Measured & Margin \\
\midrule
Flash, this chapter & 4 kB, and 7 kB with the gateway & not measured & not measured \\
Receive queue depth, each format & 16 frames & computed from the load & n/a \\
Overruns under full load & 0 & not measured & not measured \\
Frames lost with the old node present & expected: all flexible-data frames & not measured & not measured \\
Gateway added latency & under 1 ms & not measured & not measured \\
\bottomrule
\end{tabular}
\caption{The budget for chapter 13. The fourth row is the only budget in the volume whose expected value is a total failure, and stating it that way is the point: the chapter is about recognising a failure, not avoiding it.}
\end{table}

\subsection*{Firmware design (UML)}

\diagram{j13_uml}{The failure, step by step, as it happens on the wire. The node
and the master both understand the frame perfectly; the third participant does
not, and its objection is what invalidates it. The retransmission at the bottom is
why this failure does not stay small.}

Two rules.

\textbf{Count what is lost, always.} The controller has an overrun indication
for its receive queues, and a node that never reads it cannot tell the
difference between a quiet bus and a bus it is failing to keep up with. The
overrun counters go in the state frame from chapter 11 alongside the error
counters from chapter 12.

\textbf{Separate the two formats on receipt.} Two queues, one per format, costs
almost nothing and makes the diagnosis in step 6 a matter of reading two
counters instead of reasoning about one.

\subsection*{Data flow (ASCII)}

\begin{asciiart}
  a flexible-data frame goes out
        |
        +--> the node and the master decode it and are satisfied
        |
        +--> an FD-tolerant classic controller recognises the format bit,
        |    stays quiet until the frame ends, and reports nothing
        |
        +--> a non-tolerant classic controller reaches the format bit,
             finds the opposite value from the one it expects, and:
                  |
                  v
             declares a form error
                  |
                  v
             transmits an error frame, which is dominant bits on the wire
                  |
                  v
             the frame in progress is invalidated FOR EVERYBODY
                  |
                  v
             the transmitter retries, and the same thing happens again
                  |
                  v
             error counters rise on every node until somebody goes bus-off

  meanwhile, classic frames on the same bus pass perfectly, which is
  exactly what makes this confusing to diagnose
\end{asciiart}

\subsection*{Repository layout}

\begin{plaincode}
joint-node/
  src/
    bus/
      fdcan.c bittiming.c msgram.c joint_msgs.c command.c filters.c
      heartbeat.c nmt.c buserr.c emcy.c timesync.c
      rxqueue.c  rxqueue.h          # + this chapter: one queue per format
      gateway.c  gateway.h          # + this chapter, optional: two instances
    sense/ act/ estimate/ control/ time/ node/ bsp/
    mw/  safety/  update/
  host/
    mixed_load.py                   # + this chapter: both formats, together
    oldnode_demo.py                 # + this chapter: the failure, on purpose
  doc/
    format-compatibility.md         # + this chapter: the three ways out
  test/
    test_rxqueue.c                  # + overrun behaviour, on the host
  tools/  proto/  README.md
\end{plaincode}

\subsection*{Steps}

\begin{steps}
\item \textbf{Prove mixed traffic works between modern participants.} Before
introducing the old node, establish that a bus carrying both formats is normal
and uneventful when everybody can read both.
\begin{shellcode}
cangen can0 -g 2 -I 300 -L 8 -n 5000 &          # classic frames
cangen can0 -g 2 -I 400 -f -b -L 64 -n 5000 &   # flexible-data, rate switched
candump -ta -x -e can0 | python host/mixed_load.py
# classic received: 5000   flexible-data received: 5000   errors: 0
# node: rx queue classic 0 overruns, rx queue fd 0 overruns
\end{shellcode}

\item \textbf{Split the receive path by format.} Two queues, two counters, and
the overrun flag read every period. This is five minutes of work that makes step
6 a measurement rather than an argument.
\begin{ccode}
/* rxqueue.c: which format was lost is a question a counter should answer. */
void rx_dispatch(const frame_t *f)
{
    rxq_t *q = f->fd ? &g_rx_fd : &g_rx_classic;
    if (!rxq_push(q, f)) q->overruns++;         /* never silently drop */
}
\end{ccode}

\item \textbf{Size the queues from the load rather than by taste.} The worst
case is a burst at the shortest inter-frame gap, and the node's worst case for
draining is one control period. The arithmetic is short and belongs in the
repository.
\begin{plaincode}
shortest frame on this bus      about 67 bit times of arbitration plus payload
worst-case arrival rate         back-to-back frames at the nominal rate
drain interval                  one control period, because that is when the
                                node next looks
queue depth                     arrival rate times drain interval, plus margin
result                          16 frames per format, which is comfortable and
                                costs 16 * 72 bytes of message memory per queue
\end{plaincode}

\item \textbf{Read the overrun indication every period and report it.} An
overrun that nobody counts is a frame that never existed as far as the system is
concerned, and this volume's position on silent loss is the same as its position
on silent modelling.
\begin{ccode}
/* every control period, alongside the error counters from chapter 12 */
s->rx_overruns_classic = g_rx_classic.overruns;
s->rx_overruns_fd      = g_rx_fd.overruns;
\end{ccode}

\item \textbf{Now add the old node, and watch the bus stop working.} Fit the
classic hat, put it on the same bus, and run the same mixed load. This is the
chapter.
\begin{shellcode}
python host/oldnode_demo.py --seconds 30
# classic frames:        14,983 sent   14,983 received   0 errors
# flexible-data frames:   14,983 sent          0 received
#                        every one invalidated by an error frame from the old node
# node error counters:   rx errors rising, then error-passive at 11.4 s
# old node:              transmit errors rising, then bus-off at 12.1 s
# and classic traffic between the other two nodes continued throughout
\end{shellcode}
That last line is what makes this failure hard. Classic traffic is completely
unaffected, so the bus looks alive: lights blink, some messages arrive, and only
the newer format is gone. An engineer who does not know this mechanism will look
at the newer node's firmware for a long time.

\diagram{j13_data}{The two formats, aligned bit for bit, and the single position
where they differ. The older format defines that bit dominant and the newer one
transmits it recessive, so a controller built before the newer format reads a
malformed frame and puts an error flag on the wire. The flag is a deliberate
run of identical bits, which breaks the stuffing rule, so every participant sees
it whether or not it understood the frame that preceded it.}

\item \textbf{Diagnose it from the counters alone, in ten minutes.} The
signature is specific enough to be recognisable, and this is the paragraph worth
memorising from the chapter.
\begin{plaincode}
classic frames            pass, always, with no errors at all
flexible-data frames      never arrive, and the sender's retransmit count
                          equals the number it attempted
the node's rx errors      rise, because it sees the error frames
some node goes bus-off    usually the oldest one, because it is transmitting
                          error frames continuously
removing one node fixes   it, and that node is the diagnosis
\end{plaincode}
Contrast that with the termination fault from chapter 12, which produces errors
only in the data phase and affects a few per cent of frames rather than all of
them. The two are easy to tell apart once both have been seen.

\item \textbf{Choose a way out, and write it down.} Three exist, they are used
in real machines, and the node can be the third one.
\begin{plaincode}
1. all participants tolerant   a purchasing rule: every controller on the bus
                               either speaks the newer format or is documented
                               as tolerating it. Cheapest, and it constrains
                               every future addition
2. segregate the traffic       two buses: the old devices on one, the newer
                               format on the other. Costs a second bus and
                               solves it completely
3. a gateway                   one node with two controllers, translating
                               between them. Costs firmware and adds latency,
                               and it is what this node can be
\end{plaincode}

\item \textbf{Build the gateway, optionally, with the second controller
instance.} This part has two, and a second transceiver costs very little. The
gateway receives on one bus and re-transmits on the other, translating format
where it can and reporting where it cannot.
\begin{ccode}
/* gateway.c: what cannot be translated is counted, not silently dropped. */
void gateway_from_fd_bus(const frame_t *f)
{
    if (f->len <= 8) {
        frame_t out = *f; out.fd = false; out.brs = false;
        fdcan2_send(&out);                    /* fits in a classic frame */
    } else {
        g_gw.untranslatable++;                /* 9 to 64 bytes: it does not */
    }
}
\end{ccode}
The asymmetry is the interesting part and it belongs in the documentation: every
classic frame fits in the newer format, and only payloads of eight bytes or
fewer survive the journey the other way. A gateway is therefore not transparent,
and the design has to say what happens to the frames that do not fit.
\end{steps}

\subsection*{Build, flash and debug}

\diagram{j13_timing}{The same twenty milliseconds on three buses. Above, mixed
traffic between participants that all understand both formats: uneventful.
Below, the same load with one non-tolerant participant: every flexible-data
frame invalidated, the classic frames untouched, and the error counters climbing
on nodes that did nothing wrong.}

\begin{shellcode}
cmake --build build -j && probe-rs run --chip STM32H7A3ZITx build/firmware.elf
python host/mixed_load.py --seconds 30
python host/oldnode_demo.py --seconds 30
\end{shellcode}

\begin{note}{When only one frame format disappears}
This is the signature and it is worth recognising immediately. Classic frames
pass perfectly; frames in the newer format never arrive; the sender's
retransmission count equals its attempts; error counters rise on nodes that are
behaving correctly; and eventually the oldest device on the bus takes itself
off. The cause is a participant whose controller predates the newer format and
does not tolerate it, and no amount of work on the newer node's firmware will
change it. Removing that one device restores the bus immediately, which is also
the fastest way to confirm the diagnosis.
\end{note}

\subsection*{Verification and acceptance criteria}

\begin{itemize}
\item Mixed traffic between modern participants runs for thirty seconds with
zero errors and zero overruns in either queue.
\item The two receive queues are separate, and their overrun counters appear in
the state frame every period.
\item The queue depth is computed from the load arithmetic and the arithmetic is
in the repository.
\item Filling a queue deliberately increments its overrun counter rather than
losing a frame silently, proven by a host test.
\item With the old node present, classic frames pass and frames in the newer
format do not, and both outcomes are counted rather than described.
\item The error counters on the node rise while the node itself is behaving
correctly, which is the observation that makes the diagnosis teachable.
\item Removing the old node restores the bus within one recovery interval.
\item The compatibility decision is written in a file with all three options and
the reason for the one chosen.
\item If the gateway is built, frames that cannot be translated are counted, and
the count is reported rather than left at zero by construction.
\end{itemize}

\subsection*{Variants}

\begin{table}[H]\centering\small
\begin{tabular}{p{24mm}p{26mm}p{44mm}p{22mm}p{20mm}}
\toprule
Axis & Variant & What changes & Cost & Built in full in \\
\midrule
Participant & Speaks the newer format & The baseline for this volume's own nodes & None & Chapter 9 \\
Participant & Classic, documented as tolerant & Recognises the format bit and stays quiet for the rest of the frame & Nothing, if the datasheet says so & Here, as the middle case \\
Participant & Classic, not tolerant & Declares the frame malformed and invalidates it for everybody & The bus, for one of its two formats & Here, bought on purpose \\
Way out & Purchasing rule & Every controller on the bus is tolerant or newer. Cheapest & It constrains every future addition to the machine & Here \\
Way out & Two buses & Complete separation, and it always works & A second bus, its wiring and its terminators & Here \\
Way out & A gateway & One node with two controllers, translating & Firmware, latency, and it is not transparent & Here, optionally \\
Receive & One queue for both & Simplest & Which format was lost becomes a guess & Nowhere \\
Receive & One queue per format & Two counters, and the diagnosis is a reading & A little message memory & Here \\
Receive & Transfers into message memory & The controller's memory filled with no processor involvement at all & Set-up complexity, for a gain this node does not need at its frame rate & Nowhere, and chapter 9 said the same \\
\bottomrule
\end{tabular}
\caption{Variants for chapter 13. The three participant rows are a single purchasing lesson: tolerance is a controller property, it has a name in the standard, and a part that predates the newer format cannot have it however recent the board around it is.}
\end{table}

\subsection*{Pitfalls}

\begin{itemize}
\item Assuming an old controller will ignore a format it does not know. It
objects, and the objection invalidates the frame for everybody.
\item Looking for the fault in the newer node's firmware. It is behaving
correctly and its error counters rise anyway, which is the most misleading part
of the failure.
\item Assuming tolerance from a board's age or a product page. It is a
controller property, it has a name in the standard, and it is in the datasheet
or it is absent.
\item Treating a gateway as transparent. Every classic frame fits in the newer
format, and only payloads of eight bytes or fewer survive the other direction.
\item Using one receive queue for both formats and then trying to work out which
one was lost.
\item Not reading the overrun indication. A frame that nobody counted never
existed, as far as the rest of the system is concerned.
\item Confusing this failure with the termination fault from chapter 12. That
one affects a few per cent of frames in the data phase; this one affects every
frame of one format.
\item Citing the data link standard by number alone. The current edition and
its predecessor have different subtitles, so the edition is part of the
citation.
\end{itemize}

\subsection*{Best practices applied}

\begin{itemize}
\item A failure is produced deliberately rather than described, and a part is
bought specifically to produce it.
\item Loss is counted everywhere it can happen, and the counters travel in the
state frame with everything else.
\item A diagnostic signature is written down as a list of observations, so that
recognising it later does not depend on remembering the mechanism.
\item A translation that cannot be complete says so, and counts what it could
not carry.
\item A queue depth is derived from the load rather than chosen, and the
arithmetic lives beside the code.
\item A standard that could not be fetched is cited with its edition and its
price, and the behaviour it defines is demonstrated on the bench instead.
\end{itemize}

\subsection*{Stretch goals}

\begin{itemize}
\item Measure how quickly each participant reaches its error thresholds during
the failure, and see whether the order is predictable. It should be, and
confirming it is a satisfying use of the counters from chapter 12.
\item Build the gateway and measure its added latency under load, then decide
whether a machine would accept it.
\item Test a controller documented as tolerant and confirm that it really does
stay quiet, because a datasheet claim is worth confirming once.
\item Extend the gateway to fragment payloads larger than eight bytes across
several classic frames, which is the transport standard's subject and chapter
19's problem in miniature.
\end{itemize}

\subsection*{Roadmap and next steps}

Chapter 14 leaves the bus and joins the robot: the node becomes a participant in
the middleware a robot framework uses, with an agent on the host, and the
honest finding that on this exact part that is a porting job rather than a
configuration.

The published progression from here is the data link standard itself if the
budget allows, because it is the normative definition of everything this chapter
demonstrated, followed by the bus association's material on mixed networks. For
the gateway, the transport standard's segmentation is the right next reading and
chapter 19 uses it for a different purpose.

\subsection*{Portfolio evidence}

\begin{itemize}
\item The two runs side by side: mixed traffic with modern participants, and the
same load with one old device, with the counts for each format.
\item The diagnostic signature, written as a checklist, which is the kind of
thing colleagues keep.
\item The overrun instrumentation, which demonstrates a habit of counting loss
rather than assuming its absence.
\item The compatibility decision file, with all three ways out and the reason
for the one chosen.
\end{itemize}

\subsection*{Sources}

Normative references:
\begin{itemize}
\item ISO 11898-1:2024, "Road vehicles, Controller area network (CAN), Part 1:
Data link layer and physical coding sublayer", third edition, \pubdate{May 2024}, about
245 euro. It supersedes the 2015 edition, whose subtitle was "Data link layer
and physical signalling", so the edition is part of the citation. Paywalled and
not fetched for this volume.
\item The datasheet of any controller being considered for a mixed bus, for the
tolerance property, which is stated or absent.
\end{itemize}

Reusable implementations:
\begin{itemize}
\item The command line tools, dual licensed per file, whose load generator
produces both formats with one flag between them.\newline
\url{https://github.com/linux-can/can-utils}
\item The kernel's socket layer documentation, for the counters watched during
the failure.\newline
\url{https://www.kernel.org/doc/html/latest/networking/can.html}
\end{itemize}
   % Two speeds on one wire, and why the old node errors

\part{The system above}
\project{14}{The node as a middleware participant, and the agent that hosts it}{A place in the robot}{The embedded middleware client, the agent on the host, transport and footprint}

\begin{keyfacts}
\item[Adds to the node] A place in the robot. The node stops being a device on a
bus and becomes a participant a robot framework can discover, subscribe to and
record, without giving up the deterministic path built in chapters 11 and 12
\item[Peripherals] A serial port with transfers for the first working transport,
then the bus controller for the one this chapter writes
\item[Depends on] Chapter 4 for what this part does and does not have, chapter 11
for the state and command shapes, chapter 13 for what mixed traffic costs
\item[Real or modelled] \textbf{Real software, absent robot.} The client, the
transport and the agent all run. The wider graph they would join is not on this
bench and is drawn as absent rather than implied
\item[Difficulty] 5 of 5
\item[Effort] Four evenings, and the first of them is spent reading a support
list rather than writing code
\item[Deliverable] A static library built for this part, one publisher reaching
an agent on the host, a footprint measured against chapter 1's budget, a control
period proven unchanged, and a transport nobody has published
\end{keyfacts}

\subsection*{Why this chapter}

Everything so far has been a device on a wire. A robot framework does not think
in frames; it thinks in participants that publish and subscribe, that can be
discovered, recorded and replayed, and that a tool can inspect without knowing
what bus is underneath. Joining that world is what turns a joint node into part
of a machine rather than a good piece of electronics.

There is an embedded client for exactly this, with an agent on a Linux host that
bridges the constrained side to the full graph. It is the right thing to use and
this chapter uses it. It also opens with a fact that should be established
before any code is written.

\begin{note}{This board is in neither support tier}
The project publishes two tiers. Officially supported, with long term support
guaranteed, is a short list of five boards. Community supported adds three more,
and one of those is the NUCLEO-H743ZI. \textbf{The NUCLEO-H7A3ZI-Q is in neither
tier.} The nearest entry is a different part with a different reference manual,
which is the trap this whole volume is built around. So this chapter is a
porting job, it says so in its first paragraph, and the acceptance criteria are
written for a port rather than for a configuration. That is not a reason to
avoid the work. It is the reason the work is worth putting in a portfolio.
\end{note}

\subsection*{Prior art and what to reuse}

\begin{table}[H]\centering\small
\begin{tabular}{p{34mm}p{46mm}p{38mm}p{20mm}}
\toprule
Source & What it gives & What it does not & Licence \\
\midrule
The integration utilities for this silicon vendor's configurator, 275 stars, last
pushed Monday 15 September 2025 & The board agnostic route, and the one this
chapter takes: a container image builds a \textbf{static library} that drops into
a generated project. Sample mains, a build configuration file, and shims for
time, allocation and transport & \textbf{A container runtime is mandatory.} The
shipped transports are serial with transfers, serial with interrupts, USB serial
and UDP. \textbf{The bus is not among them.} It carries the project's own notice
that it is not ready for production use & Apache-2.0 \\
The agent & The bus transport is first class on the host side, with the stated
constraint that the interface must support flexible-data frames with a
sixty-four byte payload. Eight transports in total & There is no packaged binary
worth using: the index carries a 2019 version and no build for any current
distribution. A source build is the live route & Apache-2.0 \\
The client's bus transport & A header that fixes the transport unit at sixty-three
bytes and an initialisation call taking a device and an identifier & \textbf{The
implementation is for a full operating system only}: the directory holds the
generic file and a POSIX one and nothing else. On a small kernel this is a
custom transport job & Apache-2.0 \\
The custom transport interface & Six functions and a matching host-side class.
Framing can be turned off, which gives packet oriented mode, and \textbf{that is
exactly what a bus frame wants} & Nothing about how to drive a bus controller
from inside it & Apache-2.0 \\
One engineer's write-up of a port to a custom board & \textbf{The closest
published precedent}: the same core family, covering toolchain, static library
build, project integration and writing the transport functions & No flash,
memory or latency figures at all & article \\
An example project for a different part, 129 stars & A working serial transport
with transfers, and a companion video & A different part and a different
peripheral set & MIT \\
\bottomrule
\end{tabular}
\caption{Prior art for chapter 14. The third and fourth rows are the whole chapter in two lines: the bus transport exists for a host and not for this node, and the interface to write it yourself is documented and stable. Nobody has published that transport for a small kernel, which is why it is here.}
\end{table}

Three pieces of this chapter are genuinely the author's. The bus transport for a
small kernel, for which no reference implementation exists. The memory
protection and cache configuration for this family, which the integration's own
notes flag and do not solve. And the placement of buffers that both the
middleware and a transfer engine touch, which is the sibling volume's subject
and is cited rather than repeated.

\subsection*{What the node gains}

Before this chapter the node is reachable by anything that speaks its protocol
and by nothing else. After it, a tool can list the node, subscribe to what it
publishes, record a session and replay it, with no knowledge of the bus at all.
What the node does \textbf{not} gain is a new control path, and the chapter is
firm about that: the deterministic loop still runs on the protocol from chapter
11, and the middleware is a reporting and configuration path beside it.

\subsection*{Parts from the inventory}

\begin{table}[H]\centering\small
\begin{tabular}{p{40mm}p{66mm}p{40mm}}
\toprule
Part & Role & Interface \\
\midrule
NUCLEO-H7A3ZI-Q & The client. The port is to this part, not to its better known
sibling & Serial first, then the bus \\
Raspberry Pi 4 & The agent, and the far side of both transports & Its serial port
and its bus adapter \\
The debug probe's serial port & The first transport, because it is shipped,
proven and already wired & USB to the host \\
The authoring laptop & Where the container builds the library, because support
for the host's own architecture is not stated anywhere & None, it hands over a
file \\
\bottomrule
\end{tabular}
\caption{Inventory items used in chapter 14. The fourth row is a real constraint rather than a preference: the published container image does not state which architectures it supports, so the build runs where the architecture is certain and the artefact is copied.}
\end{table}

\subsection*{System architecture}

\diagram{j14_arch}{Two paths on one node. The lower path is chapters 11 and 12,
unchanged: fixed frames, one kilohertz, and every guarantee this volume has
built. The upper path is this chapter, running at a fiftieth of that rate, for
tools rather than for control. The graph on the right is drawn as absent because
no robot framework is running on this bench, and the agent is the only part of
that world this volume proves.}

\subsection*{Peripheral configuration}

\begin{table}[H]\centering\small
\begin{tabular}{p{22mm}p{26mm}p{24mm}p{30mm}p{30mm}}
\toprule
Peripheral & Mode & Clock source & Pins and function & Interrupt and transfers \\
\midrule
Serial port on the probe & 921600 baud, 8N1 & Unchanged & The two pins the probe
already uses & Transfers both ways, idle line on receive \\
Bus controller & Unchanged from chapter 9 & Unchanged & Unchanged & A second
receive queue and one more identifier range \\
Ethernet & \textbf{Absent on this part} & n/a & n/a & Which removes the UDP
transport from the list before it is considered \\
USB device & \textbf{To be confirmed on the silkscreen} & n/a & n/a & The USB
serial transport depends on this answer, and the chapter does not assume it \\
\bottomrule
\end{tabular}
\caption{Peripheral configuration for chapter 14. The last two rows are the useful ones. Chapter 4 established that this part has no Ethernet controller by the absence of its interrupt slots, which quietly eliminates one of the four shipped transports before a line of code is written.}
\end{table}

\subsection*{Wiring}

\diagram{j14_wiring}{The bench for this chapter, and the two transports in the
order they are built. The serial path is shipped, proven and already wired, so it
is first. The bus path is the one that matters for a joint node and the one
nobody has published, so it is second and it gets its own figure.}

\subsection*{Memory and timing budget}

\begin{table}[H]\centering\small
\begin{tabular}{p{44mm}p{28mm}p{28mm}p{30mm}}
\toprule
Quantity & Budget & Measured & Margin \\
\midrule
Flash, this chapter & 110 kB of the 512 kB whole-node budget & not measured & not measured \\
Static memory, this chapter & 18 kB of 256 kB & not measured & not measured \\
Middleware task stack & 12 kB, because the notes require more than 10 kB & not measured & not measured \\
Publication period & 20 ms, which is one publication per twenty control periods & not measured & not measured \\
Control period, unchanged & the same distribution as chapter 2 & not measured & not measured \\
Bus load added & under 8 per cent & computed in step 8 & n/a \\
\bottomrule
\end{tabular}
\caption{The budget for chapter 14. The first row is the largest single claim any chapter in this volume makes on the node's flash, and it is deliberately stated before the build rather than after it. If the port lands outside this budget, the chapter says so and the variants table says what to drop.}
\end{table}

\begin{note}{One published number and one that should not be repeated}
The only footprint figure with a primary source behind it covers the middleware
core alone: under 75 kB of flash and around 3 kB of memory for a complete
publisher and subscriber application with messages on the order of 512 bytes.
That is the constrained-environments client, \textbf{not} the full stack of
abstraction, convenience, interface and type-support layers that sits on top of
it. The widely repeated minimum of 32 kB of memory and 256 kB of flash comes
from a commercial blog rather than from the project, and this volume does not
print it as a project figure. What can be stated precisely is the shape of the
cost, and that is the subject of this chapter's data figure.
\end{note}

\subsection*{Firmware design (UML)}

\diagram{j14_uml}{The custom transport, which is six functions on the node and
one class on the host. Turning framing off is the decision that makes this work:
a bus frame is already a packet with a length, so the byte-stream framing the
serial transport needs is exactly what must not be applied here.}

Two rules, and the first is the architectural claim of the chapter.

\textbf{The middleware never runs in the control path.} Its task has a priority
below the control task, it publishes from a snapshot taken under a short critical
section rather than reading live state, and it is allowed to miss a publication
without anything else noticing. Chapter 2's histogram is rerun at the end of this
chapter for exactly one reason: to prove that adding it changed nothing.

\textbf{Every compile-time cap is set deliberately, before the first build.} The
implementation is built to rely on static allocation, and what it costs is set by
a table of caps: how many nodes, publishers, subscriptions, services and clients,
how deep the histories are, how long a name may be. Left at their defaults they
are four of each, a history of eight and a stream history of four, with graph
support and dynamic allocation both off. Those defaults are reasonable and this
node does not need all of them, so they are written down and trimmed rather than
inherited.

\subsection*{Data flow (ASCII)}

\begin{asciiart}
  the control task, 1 kHz, priority high
        |
        +--> the state struct, chapter 11                 (nothing new here)
        |          |
        |          +--> the bus, every period, fixed frames, deterministic
        |          |
        |          +--> a snapshot, taken under a short critical section
        |                     |
        v                     v
  never blocked        the middleware task, 50 Hz, priority low
                              |
                              +--> the client: publisher, subscription
                                        |
                                        +--> the transport
                                                 |
                              first: serial with transfers -----> the agent
                                                 |
                              then: a custom bus transport -----> the agent
                                                 |
                                                 v
                                        the wider graph, absent here

  the arrow that does not exist: nothing from the middleware task back into
  the control task. Configuration arrives as a request, is validated, and is
  applied at a period boundary by the control task itself
\end{asciiart}

\subsection*{Repository layout}

\begin{plaincode}
joint-node/
  src/
    mw/
      mw_task.c mw_task.h        # + the task, its priority and its snapshot
      mw_topics.c                # + what this node publishes and subscribes to
      mw_transport_can.c         # + this chapter: the transport nobody publishes
      mw_alloc.c mw_time.c       # + the shims the integration expects
    bus/ sense/ act/ estimate/ control/ time/ node/ bsp/
    safety/  update/
  third_party/
    microros/                    # + the built static library and its headers
      colcon.meta                # + the caps, set deliberately and reviewed
      LICENCE-NOTICES.md         # + what was pulled in and under what terms
  host/
    agent_run.sh                 # + the agent, both transports, one script
  doc/
    port-status.md               # + this board's tier, and what that means
    footprint.md                 # + the caps, and what each one cost
  test/  tools/  proto/  README.md
\end{plaincode}

\subsection*{Steps}

\begin{steps}
\item \textbf{Write down where this board sits before touching a toolchain.}
Two tiers, eight boards, and this one in neither. The nearest community-supported
entry is a different part whose reference manual this volume has spent thirteen
chapters warning about.
\begin{plaincode}
doc/port-status.md

  tier                 neither officially nor community supported
  nearest entry        a different part, different reference manual
  what that means      no guarantee, no example, no published footprint,
                       and every clock, memory and cache assumption checked
                       against this part's own manual rather than inherited
  what still holds     the architecture is the same core family, the small
                       kernel is on the supported list, and the transport
                       interface is documented and stable
\end{plaincode}
This file is short and it is the most valuable page in the chapter, because it
is the difference between a port and a surprise.

\item \textbf{Eliminate the transports that cannot work here, on evidence.} Four
are shipped. One of them needs hardware this part does not have, which chapter 4
established from the absence of the interrupt slots rather than from a product
page.
\begin{plaincode}
serial with transfers     works, already wired through the probe      -> first
serial with interrupts    works, and is the fallback if transfers fight
                          with the cache, which is the sibling volume's
                          subject
USB serial                depends on a connector this bench has not yet
                          confirmed on the silkscreen                 -> open
UDP                       needs an Ethernet controller this part does not
                          have                                        -> out
the bus                   not shipped at all                          -> step 8
\end{plaincode}

\item \textbf{Set the caps before the first build, not after the first
overflow.} The build configuration file is where the footprint is decided, and
leaving it at defaults is a decision too.
\begin{plaincode}
default                         this node                    why
nodes            1              1                            one joint
publishers       4              2                            state, diagnostics
subscriptions    4              1                            configuration
services         4              1                            one request path
clients          4              0                            it calls nobody
history          8              4                            nothing replays
stream history   4              4                            a power of two
graph support    off            off                          the agent knows
dynamic alloc    off            off                          and stays off
\end{plaincode}
Every reduction is written in the file with its reason, because a cap trimmed
without a reason is the first thing somebody restores when a build fails.

\item \textbf{Build the static library in the container, on a machine whose
architecture is certainly supported.} The image's page does not state its
architecture support, and the host on this bench is not the obvious one, so the
build runs on the laptop and hands over a file.
\begin{shellcode}
docker run --rm -v $(pwd):/project -w /project microros/micro_ros_static_library_builder:jazzy
# out: libmicroros/libmicroros.a and libmicroros/include/
arm-none-eabi-size libmicroros/libmicroros.a | tail -1
\end{shellcode}
The library is committed with a notice file listing what it contains and under
what terms, because a portfolio repository that ships a binary blob with no
provenance is worse than one that ships none.

\item \textbf{Get one publisher through on the shipped transport.} This is the
milestone that proves the port, and it is worth reaching before anything clever
is attempted.
\begin{shellcode}
# on the host
MicroXRCEAgent serial --dev /dev/ttyACM0 -b 921600 -v 6
# expected: session established, participant created, topic created
ros2 topic echo /joint/state        # if a full installation is present
\end{shellcode}
If this does not work, nothing later will, and the fault is in the port rather
than in the transport.

\item \textbf{Measure the footprint and report it against chapter 1's budget.}
This is one of the few numbers in this volume that a reader cannot find
published anywhere, which makes it worth measuring carefully.
\begin{shellcode}
arm-none-eabi-size -A build/firmware.elf | sort -k2 -nr | head -12
python tools/mapdiff.py build/before.map build/after.map --top 20
\end{shellcode}
Report it as this chapter added so much to a node that is budgeted at 512 kB,
not as an absolute number with no denominator. A middleware that takes a fifth
of a node's entire flash budget is a legitimate engineering decision and an
illegitimate surprise.

\item \textbf{Prove the control period did not change.} Rerun chapter 2's
histogram with the middleware task running and compare the two distributions:
median, 99.9th percentile and maximum, never the mean.
\begin{ccode}
/* mw_task.c: the whole interaction with the control loop is this. */
static joint_state_t snap;
taskENTER_CRITICAL();                 /* tens of cycles, not hundreds */
snap = g_state;                       /* one struct copy, no formatting */
taskEXIT_CRITICAL();
mw_publish_state(&snap);              /* everything slow happens out here */
\end{ccode}
If the 99.9th percentile moved, the critical section is doing too much or the
priorities are wrong, and both are visible in the histogram rather than in a
guess.

\item \textbf{Write the bus transport, which is the part nobody has published.}
Six functions, framing turned off, and a transport unit of sixty-three bytes
fixed by the client's own header.
\begin{ccode}
/* mw_transport_can.c: packet oriented, because a frame is already a packet. */
rmw_uros_set_custom_transport(
    false,                            /* framing off: a frame has its own length */
    (void *) &g_mw_can,               /* controller, identifier pair, queue */
    mw_can_open, mw_can_close,
    mw_can_write, mw_can_read);       /* 63 bytes maximum, per the header */
\end{ccode}
On the host the matching class is instantiated with the same identifier, and the
agent's own bus transport is the reference for what the far side must do. The
constraint the agent states plainly is that the interface must carry flexible-data
frames with a sixty-four byte payload, which is exactly the bus chapters 9 to 13
built.

\item \textbf{Redo the bus load arithmetic with the middleware traffic on it.}
Chapter 11 showed that the obvious design needed 190 per cent of the bus. That
lesson applies again here and it is the reason the publication rate is a
fiftieth of the control rate rather than equal to it.
\begin{plaincode}
control traffic, chapter 11          measured there, unchanged here
middleware, 2 publishers at 50 Hz    100 packets per second, up to 63 bytes
                                     each, which is several frames per packet
                                     once the session overhead is counted
identifier range                     a separate, lower priority block, so
                                     middleware traffic can never delay a
                                     command frame
result                               under 8 per cent added, and the figure
                                     goes in the budget table above
\end{plaincode}
The identifier assignment is the important half. Middleware traffic sits in a
lower priority block by construction, so the arbitration that chapters 9 and 11
relied on does the right thing without anybody having to remember to be careful.
\end{steps}

\diagram{j14_timing}{One second of the node, at two scales. Above, the control
task at one kilohertz, untouched. Below, the middleware task at fifty hertz,
taking a snapshot and then doing everything slow outside the critical section.
The two bars at the bottom are the same twenty milliseconds on the bus, with the
middleware traffic in its own lower priority identifier block, where arbitration
keeps it out of the way without anybody having to be careful.}

\subsection*{Build, flash and debug}

\diagram{j14_data}{What the middleware actually costs, drawn as the table it is.
The left column is the published figure and its exact scope, the middle is what
this node's caps add to it, and the right is what this volume refuses to print
and why. A footprint with no denominator and no source is not a measurement.}

\begin{shellcode}
docker run --rm -v $(pwd):/project -w /project \
  microros/micro_ros_static_library_builder:jazzy
cmake --build build -j && probe-rs run --chip STM32H7A3ZITx build/firmware.elf
bash host/agent_run.sh serial      # then: bash host/agent_run.sh canfd
\end{shellcode}

\begin{note}{When the session never establishes}
Three causes account for most of it. The framing setting is wrong for the
transport: a byte stream needs framing on and a packet transport needs it off,
and getting this backwards produces a silent failure rather than an error. The
transport unit is larger than the transport can carry, which on the bus means
anything above sixty-three bytes. Or the middleware task has too little stack,
for which the published requirement is more than 10 kB and the symptom is a
fault rather than a message. Check all three before suspecting the port.
\end{note}

\subsection*{Verification and acceptance criteria}

\begin{itemize}
\item The port status file exists and states this board's tier honestly, before
any build.
\item The caps file lists every cap with its value and the reason it was chosen.
\item The library builds reproducibly in the container, and the notice file
records what it contains and under what terms.
\item One publisher reaches the agent over the shipped serial transport, and the
session survives a node reset from either end.
\item The footprint is measured and reported as a fraction of chapter 1's
whole-node budget, not as a bare number.
\item Chapter 2's control period distribution is unchanged with the middleware
task running, judged on the 99.9th percentile and the maximum.
\item The custom bus transport carries a session, with framing off and a
sixty-three byte unit.
\item The added bus load is computed from the same arithmetic chapter 11 used
and is under the budgeted figure.
\item Middleware identifiers sit in a lower priority block than every control
identifier, proven by reading the assignment table rather than by assertion.
\item Removing the middleware task entirely still leaves a working joint node,
which is the test that the architecture claim in this chapter is true.
\end{itemize}

\subsection*{Variants}

\begin{table}[H]\centering\small
\begin{tabular}{p{24mm}p{26mm}p{44mm}p{22mm}p{20mm}}
\toprule
Axis & Variant & What changes & Cost & Built in full in \\
\midrule
Transport & Serial with transfers & Shipped, proven, and already wired on this bench & A port the joint bus does not use & Here, first \\
Transport & Serial with interrupts & The fallback when transfers and the cache disagree & Processor time per byte & Here, named \\
Transport & USB serial & One cable instead of two, if the connector is there & An unconfirmed board question & Nowhere, until the silkscreen is read \\
Transport & UDP & The obvious choice on a part with Ethernet & \textbf{Impossible here}: no controller on this part & Nowhere, and the reason is evidence \\
Transport & The bus, custom & The joint's own wire carries the middleware too & Writing six functions nobody has published for this kernel & Here, step 8 \\
Placement & Middleware in the control task & Simpler, one less task & The control period becomes hostage to a transport & Nowhere \\
Placement & Below the control task & The loop is untouchable by construction & One snapshot copy per publication & Here \\
Rate & Equal to the control rate & Every period visible to tools & The bus arithmetic from chapter 11 says no & Nowhere \\
Rate & A fiftieth of it & Tools see enough, the bus stays quiet & Tools see 50 Hz, not 1 kHz & Here \\
Host & Packaged agent binary & Nothing to build & \textbf{Stale}: the index has a 2019 version and no current build & Nowhere \\
Host & Source build & The live route, and the one the project maintains & Build time on the host & Here \\
Host & Container image & Quick, if the architecture fits & \textbf{Architecture support is not stated}, so it is verified or avoided & Here, with the caveat \\
\bottomrule
\end{tabular}
\caption{Variants for chapter 14. Four rows are refusals, and three of those four are refused on evidence about this specific part or this specific package rather than on taste. A variant table that contains no impossibilities has not been checked.}
\end{table}

\subsection*{Pitfalls}

\begin{itemize}
\item Assuming this board is supported because a similarly named one is. The
nearest entry is a different part with a different reference manual.
\item Putting the middleware in the control path. It is a reporting and
configuration path, and the moment a control period depends on it the guarantees
of chapters 2 and 11 are gone.
\item Leaving the caps at their defaults and then being surprised by the
footprint. The caps are the footprint.
\item Printing the widely repeated memory and flash minimum as a project figure.
It is a commercial blog's number, and the project's own published figure has a
much narrower scope.
\item Getting the framing setting backwards. A byte stream needs it on, a packet
transport needs it off, and the wrong choice fails silently.
\item Sending more than sixty-three bytes through the bus transport, which the
client's own header fixes and which no error message will explain.
\item Giving the middleware task a stack sized like the others. The published
requirement is more than 10 kB and the failure is a fault, not a warning.
\item Putting middleware identifiers in the same priority block as control
traffic, where arbitration will eventually do exactly what it was designed to do
at the worst moment.
\item Following documentation links from before the host moved. The old
documentation host redirects and its deep links return not-found, so citations
need the current host or the sources in the repository.
\end{itemize}

\subsection*{Best practices applied}

\begin{itemize}
\item A support tier is read and written down before a toolchain is installed.
\item Transports are eliminated on evidence about this part, including evidence
gathered from the absence of interrupt vector slots rather than from a product
page.
\item Static configuration caps are chosen deliberately and each one carries its
reason.
\item A vendored binary artefact ships with a notice recording its contents and
terms.
\item A footprint is reported against a budget rather than in isolation.
\item A claim that something did not change is proven by rerunning the original
measurement, not by inspection.
\item A number that circulates widely without a primary source is named and
refused rather than quietly repeated.
\end{itemize}

\subsection*{Stretch goals}

\begin{itemize}
\item Publish the bus transport. There is no reference implementation for this
kernel, the interface is stable and documented, and a working one with a
measurement beside it is a genuine contribution rather than another walkthrough.
\item Measure the round trip from a request on the host to a change applied at a
period boundary on the node, over both transports, and report the difference.
\item Run the port on the nearest community-supported board and compare the
footprints, which is the cheapest way to find out what this part's own
configuration cost.
\item Fill in the missing footprint figure for the complete client stack on this
core. The project publishes graphs without tables and targets a smaller core,
so a clean number here is citable.
\end{itemize}

\subsection*{Roadmap and next steps}

Chapter 15 gives the messages their standard shapes. Publishing a private
structure over a standard middleware wastes most of what the middleware is for,
so the next chapter maps the state and command of chapter 11 onto the message
types a robot framework already understands, with the units and sign conventions
those types require.

The published progression from here is the constrained-environments
specification itself, which is free and settles the vocabulary, followed by the
agent's transport documentation for anyone writing a second transport. For the
port specifically, the one published write-up of a port to a custom board in this
core family is the closest thing to a guide that exists, and this chapter is
partly an attempt to leave behind something better for the next person.

\subsection*{Portfolio evidence}

\begin{itemize}
\item The port status file, which demonstrates the habit of establishing support
before committing to an approach.
\item The caps file with a reason against every value.
\item The footprint reported as a fraction of a stated budget, with the
refused figure named and explained.
\item The two control period distributions, before and after, showing no change.
\item The bus transport, which does not exist anywhere else for this kernel.
\end{itemize}

\subsection*{Sources}

Normative references:
\begin{itemize}
\item Object Management Group, "DDS For Extremely Resource Constrained
Environments", version 1.0, document formal/20-02-01, \pubdate{February 2020}, status
Formal. It supersedes the two beta versions of \pubdate{May 2018} and \pubdate{March 2019}, and no
version beyond 1.0 is listed.
\item The reference manual for this part, for every clock, memory and cache
decision the integration's notes leave to the porter.
\end{itemize}

Reusable implementations:
\begin{itemize}
\item The integration utilities for the configurator, Apache-2.0, which is the
route this chapter takes.\newline
\url{https://github.com/micro-ROS/micro_ros_stm32cubemx_utils}
\item The client, Apache-2.0, whose bus transport header fixes the transport unit
and whose implementation is for a full operating system only.\newline
\url{https://github.com/eProsima/Micro-XRCE-DDS-Client}
\item The agent, Apache-2.0, whose bus transport is first class on the host
side.\newline
\url{https://github.com/eProsima/Micro-XRCE-DDS-Agent}
\item The setup tooling, Apache-2.0, 509 stars, last pushed Friday 18 September
2026, which is the live route to an agent on the host.\newline
\url{https://github.com/micro-ROS/micro_ros_setup}
\end{itemize}

One link-hygiene note for anyone following citations from older material: the
project's original documentation host now redirects, and deep links into its
documentation return not-found. Cite the current host or the sources in the
repositories above.
   % The node as a middleware participant, and the agent that hosts it
\project{15}{Joint state and joint command as messages}{A standard shape}{The standard message types, units and conventions, mapping frames to messages}

\begin{keyfacts}
\item[Adds to the node] A standard shape. What the node publishes stops being a
private structure and becomes a type with published semantics, published units
and no room for an alternative reading
\item[Peripherals] None. This chapter is entirely about meaning, and its only
hardware question is which way is positive
\item[Depends on] Chapter 3 for the timestamp the message contract demands,
chapter 5 for counts, chapter 8 for what effort here really is, chapter 11 for
the frame being mapped, chapter 14 for the path it travels
\item[Real or modelled] \textbf{Real, and one field is deliberately left empty}
because this bench has no torque sensor and the standard type's unit is a
measured one
\item[Difficulty] 3 of 5
\item[Effort] Two evenings, one of which is spent reading message definitions
rather than writing code
\item[Deliverable] Standard state and command messages with correct units, a
written sign convention, one field honestly empty, a size comparison that
explains a design decision, and a fault vocabulary borrowed rather than invented
\end{keyfacts}

\subsection*{Why this chapter}

Chapter 14 gave the node a way into a robot framework. This chapter decides what
it says once it is there, and the temptation at this point is obvious: the node
already has a 24-byte state structure that works, so wrap it and publish it.

That wastes almost everything the previous chapter paid for. The argument for
standard types is short and it is worth stating in the words a firmware engineer
would use: a strongly typed topic carries semantics and units with it, and
\textbf{permits no alternative interpretations}. A private structure on a
standard transport is a private structure. Any tool that wants to plot it, record
it, compare it with another joint or feed it to a controller has to be told what
it means, and being told is exactly the step that goes wrong six months later.

There is a second reason, and it is the better one. The standard types have
already answered questions this volume has been circling for fourteen chapters:
what unit a position is in, what a timestamp on a multi-field sample means, and
what a joint should say when it cannot follow a command. Borrowing those answers
is cheaper and more honest than inventing them.

\begin{note}{One rule in the standard type is chapter 3 restated}
The state type's own comments say that the header specifies the time at which the
joint states were recorded, and that all the joint states in one message have to
be recorded at the same time. That is not documentation, it is a contract, and it
is exactly the problem chapter 3 solved with one monotonic clock and a hardware
capture. A node that stamps its messages at publication time rather than at
sample time satisfies the type's shape and violates its meaning, and no consumer
will ever tell it so.
\end{note}

\subsection*{Prior art and what to reuse}

\begin{table}[H]\centering\small
\begin{tabular}{p{34mm}p{46mm}p{38mm}p{20mm}}
\toprule
Source & What it gives & What it does not & Licence \\
\midrule
The common interface definitions & The state type: a header, then names,
positions, velocities and efforts. Its comments settle the units as radians or
metres, radians or metres per second, and newton metres or newtons, and its two
rules are quoted in this chapter & Nothing about how to fill it from a
fixed-point frame, and nothing about what to do when one of the three quantities
is an estimate & Apache-2.0 \\
The trajectory definitions & The command side: names, and an array of points
each carrying positions, velocities, accelerations, effort and a time from the
start of the trajectory & No notion of a single immediate setpoint, which is what
a joint node at one kilohertz actually receives & Apache-2.0 \\
The control message package & The members a joint node genuinely needs: a jog
type for real-time commands, a dynamic state type that \textbf{names its
interfaces as strings}, and a controller state type. Its trajectory-following
action carries path, goal and goal-time tolerances, and \textbf{five error codes
that are a ready-made vocabulary} for a joint that cannot follow & It is written
for a controller on a host, not for a device, so nothing in it is sized for a
constrained node & BSD-3-Clause \\
The bus-connected motion stack, 289 stars, its core released Sunday 24 May 2026 &
\textbf{The closest published prior art to this whole volume's architecture}: a
driver hierarchy ending in a motion-profile driver that implements the standard
profile state machine with position, velocity and torque modes, exposed as
hardware interfaces. It is the maintained answer to how a bus-connected joint
appears to a robot framework & It is the \textbf{master} side of that contract.
This volume builds the device side, and says so rather than implying it builds
both & Apache-2.0 \\
The raw frame wrapper, 194 stars, its message package released Monday 7 September
2026 & A thin, well released wrapper at the frame level, whose message package
defines a flexible-data frame type, so the newer format is covered & No semantics
at all above the frame, which is the whole subject of this chapter & Apache-2.0 \\
\bottomrule
\end{tabular}
\caption{Prior art for chapter 15. The fourth row deserves attention from anyone reading this volume as a portfolio: somebody has already built the master side of this contract, it is maintained, and this volume's node is the device that sits opposite it. Knowing that is worth more than not having heard of it.}
\end{table}

\subsection*{What the node gains}

Before this chapter the node speaks a protocol one repository understands. After
it, the node speaks types that any tool in a large ecosystem understands, with
units that cannot be misread, a timestamp whose meaning is defined, and an error
vocabulary that a controller already knows how to act on. It also gains a
written sign convention, which is the cheapest thing in this chapter and the one
most likely to save a day.

\subsection*{Parts from the inventory}

\begin{table}[H]\centering\small
\begin{tabular}{p{40mm}p{66mm}p{40mm}}
\toprule
Part & Role & Interface \\
\midrule
NUCLEO-H7A3ZI-Q & Fills the messages and publishes them & Chapter 14's transport \\
Raspberry Pi 4 & Reads them, and is where a bridge would live if the mapping were
placed there instead & The same \\
The encoder from chapter 5 & The only real source of a position, and the reason
the conversion constant has exactly one home & Already wired \\
\bottomrule
\end{tabular}
\caption{Inventory items used in chapter 15. Nothing new is bought and nothing new is wired. The work is a decision about meaning, which is the kind of chapter that is easy to skip and expensive to skip.}
\end{table}

\subsection*{System architecture}

\diagram{j15_arch}{Where the mapping happens, and the two defensible answers.
In the node, which is what this bench does because chapter 14 already put a
client there. On the host, which is what the maintained bus-connected motion
stack does, and which this volume does not build. The figure names the second one
rather than pretending the first is the only choice.}

\subsection*{Peripheral configuration}

\begin{table}[H]\centering\small
\begin{tabular}{p{22mm}p{26mm}p{24mm}p{30mm}p{30mm}}
\toprule
Peripheral & Mode & Clock source & Pins and function & Interrupt and transfers \\
\midrule
None & This chapter configures no peripheral at all & n/a & n/a & n/a \\
The clock from chapter 3 & Read, not configured. It supplies the sample instant
the message contract requires & Unchanged & None & The capture is already in
place \\
The encoder from chapter 5 & Read, not configured. Its counts become radians
through one constant with one sign & Unchanged & Unchanged & Unchanged \\
\bottomrule
\end{tabular}
\caption{Peripheral configuration for chapter 15. A table of nothing is worth printing once: this chapter changes no hardware, and every quantity it publishes was already being measured. What changes is what those quantities are called and what they are guaranteed to mean.}
\end{table}

\subsection*{Wiring}

\diagram{j15_wiring}{The sign convention, drawn once so that it can be pointed at.
Positive rotation, the zero position, and where the encoder's own zero sits
relative to it. None of this is discoverable from code, all of it is assumed by
every consumer of the messages, and the cost of writing it down is one figure.}

Nothing is wired in this chapter. The figure is here because a sign convention is
the closest thing this chapter has to hardware, and because it is the piece most
often left in somebody's head.

\subsection*{Memory and timing budget}

\begin{table}[H]\centering\small
\begin{tabular}{p{44mm}p{28mm}p{28mm}p{30mm}}
\toprule
Quantity & Budget & Measured & Margin \\
\midrule
Flash, this chapter & 6 kB, mostly type support & not measured & not measured \\
Static memory, this chapter & 3 kB, the message buffers & not measured & not measured \\
State message on the wire & about ten times the 24-byte frame & computed in step 5 & n/a \\
Conversion cost per period & under 2 microseconds & not measured & not measured \\
Publication rate & 50 Hz, unchanged from chapter 14 & n/a & n/a \\
\bottomrule
\end{tabular}
\caption{The budget for chapter 15. The third row is the one that justifies a decision rather than reporting a cost: knowing the ratio is what makes the choice between publishing at one kilohertz and publishing at fifty hertz an arithmetic question instead of a preference.}
\end{table}

\subsection*{Firmware design (UML)}

\diagram{j15_uml}{The three types this node touches, what it fills in each, and
what it deliberately leaves empty. The empty array is the interesting part: the
standard type's unit for effort is a measured one, this bench has no torque
sensor, and chapter 8 was clear about the difference. So the estimate travels in
the type that lets an interface be named, where it can be called an estimate
without anybody having to read a comment.}

Three rules.

\textbf{One conversion constant, one sign, one file.} Counts to radians is a
single constant and a single sign, and every place that needs it calls the same
function. A second copy of that constant anywhere in the repository is a defect
whether or not it currently agrees.

\textbf{The stamp is the sample instant.} Not the publication instant, not the
time the message was assembled. Chapter 3 built a clock precisely so that this
could be true, and chapter 6 reconstructed per-sample instants for the same
reason.

\textbf{An estimate never travels in a measured field.} The state type's effort
array is left empty on this bench. The estimate is published, clearly, in a place
that can say what it is.

\subsection*{Data flow (ASCII)}

\begin{asciiart}
  chapter 11's 24-byte state frame
        |
        +-- position_counts  (int32, exact)
        |        |
        |        +--> counts_to_rad()  one constant, one sign, one file
        |                    |
        |                    v
        |              position[0], float64, RADIANS
        |
        +-- velocity_mrad_s  (int32) --> /1000 --> velocity[0], rad/s
        |
        +-- effort_mnm       (int32, DERIVED)
        |        |
        |        +--> NOT into effort[], which the standard says is measured
        |        |
        |        +--> into the dynamic state type, interface named
        |             "effort_estimate", beside "effort_source"
        |
        +-- stamp_us         --> header.stamp, the SAMPLE instant, chapter 3
        |
        +-- fault_flags      --> the five standard error codes, chapter 16 uses
                                 PATH_TOLERANCE_VIOLATED for following error

  name[] is filled once at start-up and never rebuilt, because it is constant
  and rebuilding a string array every period is the most expensive way to say
  nothing new
\end{asciiart}

\subsection*{Repository layout}

\begin{plaincode}
joint-node/
  src/
    mw/
      mw_task.c mw_topics.c mw_transport_can.c
      msg_state.c  msg_state.h     # + this chapter: frame to standard type
      msg_cmd.c    msg_cmd.h       # + this chapter: standard type to frame
      units.c      units.h         # + counts to radians, once, with its sign
    bus/ sense/ act/ estimate/ control/ time/ node/ bsp/
    safety/  update/
  doc/
    conventions.md                 # + sign, zero, units, and one figure
    message-map.md                 # + every field, its source, its unit
  test/
    test_units.c                   # + round trip, and the sign, on the host
    test_msg_state.c               # + against recorded frames
  host/  third_party/  tools/  proto/  README.md
\end{plaincode}

\subsection*{Steps}

\begin{steps}
\item \textbf{Read the message definitions before writing anything.} They are
short, they are the normative answer to questions this volume has been
answering ad hoc, and they are free.
\begin{plaincode}
the state type      header, then name[], position[], velocity[], effort[]
units, from its own comments
  position          radians for a revolute joint, metres for a prismatic one
  velocity          radians or metres per second
  effort            newton metres or newtons
two rules, quoted
  1. the header specifies the time at which the joint states were RECORDED
  2. all the joint states in one message have to be recorded at the SAME time
  and the arrays are optional, but must be the same size or empty
\end{plaincode}
The second rule and the permission to leave an array empty are the two facts
this chapter leans on hardest.

\item \textbf{Write the sign convention down, with a figure, before any
conversion.} Which way is positive, where zero is, and where the encoder's own
zero sits relative to it.
\begin{plaincode}
doc/conventions.md

  positive rotation   counter-clockwise viewed from the output side
  zero position       the mechanical reference mark, chapter 5
  encoder zero        its index would be here, but THIS PART HAS NO INDEX,
                      so zero is established at start-up and the procedure
                      is in chapter 5, not assumed here
  units published     radians, always, never degrees, anywhere
\end{plaincode}
Degrees do not appear in this repository at any layer, including in log
messages, because a unit that appears in two forms will eventually be converted
twice or not at all.

\item \textbf{Put the conversion in one file and test it on the host.} The test
that matters is the round trip and the sign, and both are cheap.
\begin{ccode}
/* units.h: the only place counts and radians know about each other. */
#define JOINT_COUNTS_PER_REV   (4 * 2048)      /* quadrature, chapter 5 */
#define JOINT_SIGN             (+1)            /* see doc/conventions.md */

static inline double counts_to_rad(int32_t c)
{
    return JOINT_SIGN * (double) c * (2.0 * M_PI / JOINT_COUNTS_PER_REV);
}
\end{ccode}

\item \textbf{Fill the state message, and leave effort empty.} This is the
chapter's honesty decision and it costs one line plus one paragraph of
documentation.
\begin{ccode}
/* msg_state.c: three arrays of one, and one of them stays at length zero. */
m->name.size     = 1;                       /* filled once at start-up */
m->position.data[0] = counts_to_rad(s->position_counts);
m->velocity.data[0] = s->velocity_mrad_s / 1000.0;
m->effort.size   = 0;                       /* no torque sensor on this bench */
m->header.stamp  = us_to_ros_time(s->stamp_us);   /* the SAMPLE instant */
\end{ccode}
The arrays are allowed to be empty and are required to be the same size when
they are not, so a length of zero is a legal statement meaning this node does not
measure that quantity. Filling it with an estimate would also be legal and would
be a lie in a defined unit, which is worse than saying nothing.

\item \textbf{Publish the estimate where it can be labelled.} The dynamic state
type carries interface names as strings, which is exactly the mechanism needed.
\begin{plaincode}
interface_name        value        where it came from
  effort_estimate     0.214        chapter 8's model, in newton metres
  effort_source       1            DERIVED, straight from the frame
  model_version       3            which parameter set produced the estimate
  position            0.7854       the same number as the state message
\end{plaincode}
A consumer that wants effort now has to ask for a field called estimate, which is
the entire point. Chapter 8 spent a chapter establishing that this quantity is
derived, and this is where that work stops being a footnote.

\item \textbf{Measure what the standard shape costs on the wire, and decide from
the number.} A private frame is 24 bytes; the standard message is not, and the
difference explains the publication rate chosen in chapter 14.
\begin{plaincode}
chapter 11's frame     24 bytes, fixed, one bus frame
the state message      header with a stamp and a frame identifier string
                       name[]      one string, and strings are not free
                       position[]  one float64        8 bytes
                       velocity[]  one float64        8 bytes
                       effort[]    empty              0 bytes
                       plus array lengths and the serialisation's own overhead
result                 roughly ten times the frame, which is why this travels
                       at 50 Hz and the frame travels at 1 kHz
\end{plaincode}
Both are correct designs. The mistake would be to publish the standard message
every period because the standard is the right thing, and then wonder where the
bus went.

\item \textbf{Map the command side, and accept that the standard trajectory type
is not a setpoint.} A trajectory is a list of points with times from a start; a
joint node at one kilohertz receives one immediate command.
\begin{plaincode}
what arrives                    what this node does with it
a jog command                   the direct match: a real-time immediate command
a trajectory                    accepted, held, and interpolated by whoever
                                owns the trajectory. THIS NODE DOES NOT
                                INTERPOLATE, and chapter 16 says why
a single point                  treated as a jog with a validity time,
                                which is chapter 11's valid_until field
\end{plaincode}
The node refusing to interpolate a trajectory is a design position, not a
limitation, and it is stated here so that chapter 16 can rest on it.

\item \textbf{Borrow the error vocabulary instead of inventing one.} The
trajectory-following action already defines what a joint says when it cannot
follow, and those codes are better than anything this volume would have made up.
\begin{plaincode}
 0  SUCCESSFUL
-1  INVALID_GOAL              the command is outside the declared limits
-2  INVALID_JOINTS            the names do not match this node
-3  OLD_HEADER_TIMESTAMP      chapter 11's valid_until has passed
-4  PATH_TOLERANCE_VIOLATED   following error, which is chapter 16's subject
-5  GOAL_TOLERANCE_VIOLATED   it arrived, but not close enough
\end{plaincode}
Chapter 11's fault word maps onto these, and the map lives in a file rather than
in a switch statement, so that the two cannot drift apart quietly.

\item \textbf{Write the field map down, every field, with its source and its
unit.} This is the document a second engineer reads instead of asking, and it is
the deliverable most likely to be read by somebody other than its author.
\begin{plaincode}
doc/message-map.md
  message field      source field          unit       notes
  position[0]        position_counts       rad        counts_to_rad, sign in
                                                      doc/conventions.md
  velocity[0]        velocity_mrad_s       rad/s      divided by 1000
  effort[]           (none)                N m        EMPTY: no torque sensor
  header.stamp       stamp_us              s, ns      the SAMPLE instant
  effort_estimate    effort_mnm            N m        dynamic state type only
\end{plaincode}
\end{steps}

\diagram{j15_timing}{The timestamp contract, and the two ways to break it. The
top row is what the type requires: one instant, shared by every quantity in the
message, taken when the quantities were recorded. The middle row is a node that
samples its three quantities at three different times and stamps them as one,
which the type forbids in a sentence. The bottom row is a node that stamps at
publication, which is the common failure and the invisible one.}

\subsection*{Build, flash and debug}

\diagram{j15_data}{The same instant of the same joint, twice. On the left,
chapter 11's frame: twenty-four bytes, fixed, integers, one bus frame. On the
right, the standard message: floating point, named, self-describing, and roughly
ten times the size. Neither is better. The comparison is here because the ratio
is what decides the rate, and a rate chosen without it is a guess.}

\begin{shellcode}
cmake --build build -j && ctest --test-dir build/host
probe-rs run --chip STM32H7A3ZITx build/firmware.elf
ros2 topic echo /joint/state --once        # if a full installation is present
python host/check_units.py --degrees-are-a-bug
\end{shellcode}

\begin{note}{When a plot looks right but reads wrong}
Three causes, in the order they occur. The sign convention was never written
down, so one consumer assumes the opposite one and everything is mirrored.
Degrees appeared somewhere, usually in a log line that was later parsed. Or the
stamp is the publication instant rather than the sample instant, which produces
a plot that is correct in shape and wrong in time, and which nothing will ever
complain about. The first two are caught by a host test. The third is caught
only by having decided, in chapter 3, that it mattered.
\end{note}

\subsection*{Verification and acceptance criteria}

\begin{itemize}
\item The conversion from counts to radians exists once, and a host test proves
the round trip and the sign.
\item The word degrees appears nowhere in the repository, including in log
strings, and a test asserts it.
\item The state message carries position and velocity in the units the type's own
comments specify.
\item The effort array is empty, and the reason is in the field map.
\item The estimate is published only in the type that names its interfaces, and
its name says it is an estimate.
\item The header stamp equals the sample instant from chapter 3, proven by
comparing a published message against the frame it came from.
\item Every array in one message has the same length or is empty, which the
type requires.
\item The field map file lists every published field with its source and unit.
\item The fault word maps onto the five standard error codes through a file, not
a switch statement.
\item The size ratio between the frame and the message is measured and appears in
the budget table.
\end{itemize}

\subsection*{Variants}

\begin{table}[H]\centering\small
\begin{tabular}{p{24mm}p{26mm}p{44mm}p{22mm}p{20mm}}
\toprule
Axis & Variant & What changes & Cost & Built in full in \\
\midrule
Mapping & In the node & The node publishes standard types itself & Type support and message buffers on a constrained part & Here \\
Mapping & On the host, in a bridge & The node stays small and the graph sees the bridge & A host component, and the node is no longer the publisher & Nowhere, and it is what the maintained stack does \\
Effort & Left empty & Honest, and the type explicitly allows it & A consumer gets nothing & Here \\
Effort & Filled with the estimate & Convenient, and wrong in a defined unit & Credibility, the moment somebody checks & Nowhere \\
Effort & In the dynamic type, named & Honest and useful at once & One more message type & Here \\
Command & Jog & The direct match for an immediate setpoint & None & Here \\
Command & Trajectory, interpolated on the node & The node owns the profile & Trajectory memory and a profile generator on a constrained part & Nowhere, and chapter 16 says why \\
Command & Trajectory, held elsewhere & The node stays a device & The master must exist & Here, by declaration \\
Frame level & Raw frame messages & Every frame visible to the graph & Semantics stay at zero, which is the thing this chapter is for & Nowhere, and it is named \\
\bottomrule
\end{tabular}
\caption{Variants for chapter 15. The three effort rows are one argument made three ways, and the middle one is the variant this volume exists to argue against. A field with a defined unit is a promise, and filling it with something else is the cheapest way to lose a reader's trust.}
\end{table}

\subsection*{Pitfalls}

\begin{itemize}
\item Publishing a private structure over a standard transport and believing the
node now speaks the standard.
\item Stamping at publication time. The type's own comment says recorded, and a
consumer cannot detect the difference.
\item Putting a derived quantity into a field whose unit implies measurement.
\item Letting degrees into any layer, including log messages.
\item Rebuilding the name array every period, which is the most expensive way to
transmit a constant.
\item Publishing arrays of different lengths, which the type forbids and which
some consumers will accept silently.
\item Keeping the conversion constant in two places. They will agree until one
is changed.
\item Publishing the standard message every control period because the standard
is the right thing, without doing the size arithmetic first.
\item Inventing an error vocabulary when a published one exists that controllers
already act on.
\end{itemize}

\subsection*{Best practices applied}

\begin{itemize}
\item Normative definitions are read before code is written, and their comments
are treated as contracts.
\item A convention that lives only in somebody's head is drawn and written down.
\item A single constant has a single home, enforced by a test rather than by
discipline.
\item A quantity that is estimated is published where it can be labelled as an
estimate.
\item An empty array is used as the meaningful statement the type intends it to
be.
\item A vocabulary is borrowed from a maintained standard rather than invented.
\item A design decision about rate rests on measured sizes rather than on
preference.
\item Prior art that solves the other half of the same problem is named clearly,
including the part this volume does not build.
\end{itemize}

\subsection*{Stretch goals}

\begin{itemize}
\item Write the host-side bridge as well, and compare the two placements on node
footprint, host complexity and what the graph sees. Nobody has published that
comparison for a constrained node.
\item Implement the trajectory-following action's tolerance fields as the node
sees them, which is most of chapter 16 arriving early.
\item Map this node onto the standard motion profile the maintained stack
implements, and find out how much of that profile a bench node can honestly
claim.
\item Publish a recording of one minute of state messages alongside the raw
frames that produced them, which makes the whole mapping auditable by anybody.
\end{itemize}

\subsection*{Roadmap and next steps}

Chapter 16 closes the loop. The messages of this chapter go out and come back, a
setpoint arrives, and the figure of merit becomes following error, which is the
quantity the standard error codes already have a name for.

The published progression from here is the message definitions themselves, which
are short and worth reading in full, then the control message package for the
members a joint node needs, then the maintained bus-connected motion stack as the
worked example of the other side of this contract. The concept documentation for
nodes, topics and typing is in the documentation repository rather than on the
documentation site, because the site refuses automated fetching and the
repository does not.

\subsection*{Portfolio evidence}

\begin{itemize}
\item The field map, which is short, complete, and the kind of document that
gets copied into other projects.
\item The sign convention figure, for the same reason.
\item The empty effort array with its written justification, which demonstrates
a habit that is difficult to teach.
\item The error-code map, showing a vocabulary borrowed rather than invented.
\item The size comparison, which turns a rate decision into arithmetic.
\end{itemize}

\subsection*{Sources}

Normative references:
\begin{itemize}
\item The joint state message definition, for the field list, the units in its
own comments, and the two rules quoted in this chapter.
\item The joint trajectory message definition, for the command side and its time
from start.
\item The control message package, BSD-3-Clause, for the jog type, the dynamic
state type with named interfaces, and the trajectory-following action with its
tolerances and its five error codes.
\end{itemize}

Reusable implementations:
\begin{itemize}
\item The common interface definitions, Apache-2.0.\newline
\url{https://github.com/ros2/common_interfaces}
\item The bus-connected motion stack, Apache-2.0, 289 stars, the maintained
master side of this contract.\newline
\url{https://github.com/ros-industrial/ros2_canopen}
\item The raw frame wrapper, Apache-2.0, 194 stars, whose message package defines
a flexible-data frame type.\newline
\url{https://github.com/autowarefoundation/ros2_socketcan}
\item The concept documentation, in repository form because the documentation
site refuses automated fetching.\newline
\url{https://github.com/ros2/ros2_documentation}
\end{itemize}
   % Joint state and joint command as messages
\project{16}{The loop closed over the bus: setpoint in, state out, following error}{A closed loop}{Setpoint injection, the controller, following error as the figure of merit}

\begin{keyfacts}
\item[Adds to the node] A closed loop, and the one number a joint is judged by.
The node stops reporting and starts controlling, and it reports how well it is
controlling in a quantity that has a published name
\item[Peripherals] Nothing new. The timer from chapter 2, the encoder input from
chapter 5, the outputs from chapter 7 and the bus from chapter 9
\item[Depends on] Chapter 2 for the period, chapter 5 for feedback, chapter 7
for the plant that is not here, chapter 11 for the command, chapter 15 for the
name of the failure
\item[Real or modelled] \textbf{Real controller, modelled plant.} The loop
closes through a dashed block, and every following error in this chapter is real
arithmetic performed on a model. The chapter says so on every figure
\item[Difficulty] 5 of 5
\item[Effort] Four evenings, and the last one is spent on the controller nobody
has published rather than on the one everybody has
\item[Deliverable] A controller with filtered derivative and real anti-windup, a
written division of labour between node and host, following error measured at
the node's own rate, and a licence audit that changes which library is used
\end{keyfacts}

\subsection*{Why this chapter}

Fifteen chapters have built a joint node that measures, times, reports and
joins. It has never corrected anything. This chapter closes the loop, and the
first question is not which controller to write but \textbf{where the loop
lives}, because there are two candidates and a published default settles it.

The host framework's manager has a parameter for the frequency of its real-time
update loop, the loop that reads states from hardware, updates controllers and
writes commands back. \textbf{Its default is one hundred hertz.} This node runs
at one thousand. That single arithmetic fact decides the architecture: the fast
loop closes on the node, where the feedback already is, and the host closes a
slower outer loop around it. Nothing about that is a compromise. It is the
standard arrangement in motion control and the numbers say so out loud.

The second question is the figure of merit. A joint that is asked to be at a
position and is not at it has a following error, and a controller is judged on
the shape of that quantity rather than on whether it eventually arrives. Chapter
15 already borrowed the name for what happens when it grows too large.

\begin{note}{The loop closes through a model and the chapter never forgets it}
There is no motor on this bench. The controller output drives chapter 7's plant
model, whose position drives chapter 5's quadrature generator, whose pulses come
back through the real encoder input on real hardware. So the arithmetic is real,
the peripherals are real, the timing is real, and the mechanics are a model.
Every figure in this chapter draws that block dashed, and no number here is a
claim about a physical joint. What it is a claim about is the firmware, which is
what this volume is for.
\end{note}

\subsection*{Prior art and what to reuse}

\begin{table}[H]\centering\small
\begin{tabular}{p{34mm}p{46mm}p{38mm}p{20mm}}
\toprule
Source & What it gives & What it does not & Licence \\
\midrule
The host control framework, about 1000 stars, last commit Thursday 17 September
2026 & The contract this node sits behind: a manager whose real-time loop reads,
updates and writes, hardware components with lifecycle callbacks and a read and
write pair, and \textbf{a published list of interface names that already
includes the three controller gains, two integral clamp limits and a feedforward
term} & Nothing that runs on a microcontroller. \textbf{Its rendered
documentation omits the lifecycle state argument} from the callback signatures,
so the header is the source of truth & Apache-2.0 \\
A hardware component for a microcontroller over a serial link, 139 stars & The
structural template for the host side of this exact arrangement. \textbf{Change
the transport and the shape is the same} & A different transport, and a
different class of joint & BSD-3-Clause \\
The best-fitting small controller, 935 stars, last pushed Monday 24 July 2023 &
About 150 lines and finished: trapezoidal integral, a \textbf{band-limited
derivative taken on the measurement} so a setpoint step produces no kick, and
clamping anti-windup with a separate integrator limit computed from the headroom
the proportional term leaves & \textbf{The licence is at repository level only:
neither source file carries a header.} If it is vendored, a header is added and
the provenance recorded & MIT, at the repository \\
The vendor signal-processing controller & The incremental form in three numeric
types, well tested and everywhere & \textbf{The teachable foil for this
chapter}: no anti-windup, no derivative filtering, no output saturation, the
period baked invisibly into the gains, and an unfiltered second difference on
the error & Apache-2.0 \\
A motor-control library's controller & Trapezoidal integral, integral clamp,
output clamp, and \textbf{an output slew rate limit the others lack and a joint
wants}. It measures elapsed time rather than assuming a period & No derivative
filter, and a plain backward difference on the error & MIT \\
The real-time helper library & A buffer and a publisher safe to touch from a
real-time thread, priority helpers, and an asynchronous function handler & Host
side only & BSD-3-Clause \\
\bottomrule
\end{tabular}
\caption{Prior art for chapter 16. The third and fourth rows are why this chapter builds both: one is the right implementation and carries a licence defect worth teaching, and the other is the wrong implementation and is the clearest way to show what each missing feature costs.}
\end{table}

One thing in this chapter is genuinely the author's. \textbf{No verified open
embedded controller in this language implements textbook back-calculation
anti-windup.} Clamping is everywhere and conditional integration is common;
back-calculation, where the saturated output is fed back into the integrator
through its own gain, is described in the standard references and implemented in
none of the small libraries this volume could check. The chapter implements it,
measures it against clamping on the same step, and publishes both.

\subsection*{What the node gains}

Before this chapter the node is an instrument. After it, it is a joint: it
receives a setpoint, corrects toward it, reports how far off it is, and says so
in the vocabulary chapter 15 borrowed. It also gains its gains as interfaces, so
that tuning becomes a message rather than a rebuild.

\subsection*{Parts from the inventory}

\begin{table}[H]\centering\small
\begin{tabular}{p{40mm}p{66mm}p{40mm}}
\toprule
Part & Role & Interface \\
\midrule
NUCLEO-H7A3ZI-Q & The fast loop, at the rate chapter 2 proved & Timer, encoder input, outputs, bus \\
Raspberry Pi 4 & The slow outer loop and the setpoint source & The bus \\
The quadrature generator, chapter 5 & Turns the plant model's position back into
pulses, so feedback arrives through real hardware & Two pins to two pins \\
The plant model, chapter 7 & \textbf{The only part of this loop that is not
real}, and it is drawn dashed everywhere & None: it is software \\
\bottomrule
\end{tabular}
\caption{Inventory items used in chapter 16. The loop is a real loop through real peripherals with one modelled link in it, which is a more honest arrangement than a simulation and a less capable one than a motor. Both halves of that sentence belong in the portfolio text.}
\end{table}

\subsection*{System architecture}

\diagram{j16_arch}{Two loops, and the published default that separates them. The
inner loop closes on the node at one kilohertz, where the feedback already is.
The outer loop closes on the host at the framework's default of one hundred
hertz, one tenth the rate, which is why it sends setpoints rather than efforts.
The dashed block is the plant, and it is the only part of the inner loop that is
not hardware.}

\subsection*{Peripheral configuration}

\begin{table}[H]\centering\small
\begin{tabular}{p{22mm}p{26mm}p{24mm}p{30mm}p{30mm}}
\toprule
Peripheral & Mode & Clock source & Pins and function & Interrupt and transfers \\
\midrule
Control timer & Unchanged from chapter 2 & Unchanged & None & The period this chapter runs in \\
Encoder input & Unchanged from chapter 5 & Unchanged & Unchanged & Read once per period, never polled \\
Output stage & Unchanged from chapter 7 & Unchanged & Unchanged & Updated once per period, at a defined point \\
Bus controller & Unchanged from chapter 9 & Unchanged & Unchanged & Command in, state out \\
\bottomrule
\end{tabular}
\caption{Peripheral configuration for chapter 16. Every row says unchanged, which is the point: fifteen chapters of configuration were done so that this one could be about control rather than about registers. A chapter that had to configure something here would be a chapter whose predecessors left work undone.}
\end{table}

\subsection*{Wiring}

\diagram{j16_wiring}{The loop as it physically exists on this bench, drawn once
so that nobody has to be told twice which link is a model. Setpoint over the bus,
controller on the node, output through the real output stage, into the plant
model, out through the real quadrature generator, back through the real encoder
input. One dashed block, six solid ones.}

\subsection*{Memory and timing budget}

\begin{table}[H]\centering\small
\begin{tabular}{p{44mm}p{28mm}p{28mm}p{30mm}}
\toprule
Quantity & Budget & Measured & Margin \\
\midrule
Flash, this chapter & 5 kB & not measured & not measured \\
Static memory, this chapter & 1 kB & not measured & not measured \\
Controller execution, per period & under 8 microseconds of 1000 & not measured & not measured \\
Divisions in the handler & \textbf{zero}, by premultiplying at configuration & n/a & n/a \\
Following error, steady state & under 1 milliradian on the model & not measured & not measured \\
Host loop rate & 100 Hz, the framework default & not measured & not measured \\
Host loop jitter & \textbf{not bounded}: this Pi runs no real-time kernel & to be measured & n/a \\
\bottomrule
\end{tabular}
\caption{The budget for chapter 16. The last row is the honest one. The framework asks for a real-time scheduling policy at a stated priority, with the user in a real-time group and limits configured, and recommends a real-time kernel. This bench has none of that, so the outer loop's jitter is measured and reported rather than assumed away.}
\end{table}

\subsection*{Firmware design (UML)}

\diagram{j16_uml}{The two sides of the contract. On the host, a hardware
component with its lifecycle callbacks and its real-time read and write pair. On
the node, the control task that already exists. The caution in the middle is
real: the rendered documentation for those callbacks omits an argument the
header requires, so the header is what gets read before a signature is typed.}

Three rules.

\textbf{No division in the handler.} The gains are premultiplied by the period
when they are configured, so the periodic path does adds, multiplies and
comparisons and nothing else. This is the single cheapest piece of published
advice in the whole subject.

\textbf{The derivative is taken on the measurement, filtered.} A derivative on
the error produces a kick every time the setpoint moves, and an unfiltered one
amplifies exactly the noise the encoder has most of. Taking it on the
measurement is equivalent while the setpoint is constant and much better when it
is not.

\textbf{Saturation is not a detail.} The output is bounded, and the integrator
must be told. This chapter implements two answers and measures both, because one
of them is not published anywhere in this language.

\subsection*{Data flow (ASCII)}

\begin{asciiart}
  the host, 100 Hz, no real-time kernel on this bench
        |
        +--> setpoint, with a validity time (chapter 11)
                 |
                 v
  ----------- the node, 1 kHz, every period, in this order -----------
        |
        read the encoder (chapter 5)      exactly once, at a fixed point
        |
        v
      error = setpoint - measurement
        |
        +--> proportional:  Kp * error
        |
        +--> integral:      trapezoidal, and clamped to the headroom the
        |                   proportional term leaves
        |
        +--> derivative:    on the MEASUREMENT, band limited, sign inverted
        |                          (no kick when the setpoint steps)
        v
      sum, then saturate to the output range
        |
        +--> if it saturated: ANTI-WINDUP
        |        clamping        the integrator is held           (published)
        |        back-calculation the excess is fed back through its own
        |                        gain                  (the author's, here)
        v
      write the output stage (chapter 7), at a fixed point in the period
        |
        v
      [ the plant model ]  <-- DASHED: the only modelled link in this loop
        |
        v
      the quadrature generator (chapter 5) --> real pulses --> real encoder
        |
        v
      following error is recorded every period, and published at 50 Hz
\end{asciiart}

\subsection*{Repository layout}

\begin{plaincode}
joint-node/
  src/
    control/
      pid.c pid.h              # + this chapter: filtered D, two anti-windups
      pid_config.c             # + premultiply by the period, once
      setpoint.c               # + command in, validity checked
      follow.c follow.h        # + following error, and its statistics
    estimate/ sense/ act/ bus/ mw/ time/ node/ bsp/ safety/ update/
  host/
    hw_component/              # + the host side: read, write, lifecycle
    tune.py                    # + gains over the wire, no rebuild
  doc/
    loop-division.md           # + which loop owns what, and the arithmetic
    antiwindup.md              # + both methods, both measured
  test/
    test_pid.c                 # + step, saturation, and windup, on the host
    test_follow.c              # + the statistics, against recorded runs
  third_party/  proto/  tools/  README.md
\end{plaincode}

\subsection*{Steps}

\begin{steps}
\item \textbf{Write the division of labour down before writing a controller.}
The arithmetic is published and it settles the argument in four lines.
\begin{plaincode}
doc/loop-division.md

  the node's loop     1000 Hz, proven in chapter 2, feedback is local,
                      one period of latency end to end
  the host's loop     100 Hz by default, and that default is the framework's
                      own parameter, not a guess
  ratio               10 to 1, which is why the host sends POSITION setpoints
                      and not efforts: an effort loop at a tenth of the rate
                      is a slower loop, not a different one
  what the host owns  trajectory, coordination between joints, limits,
                      and the profile
  what the node owns  the correction, the following error, and the safe
                      state of chapter 18
\end{plaincode}

\item \textbf{Configure the gains once, and premultiply.} Everything the
periodic path needs is computed here, so that the handler contains no division
and no branch that depends on a parameter.
\begin{ccode}
/* pid_config.c: all the arithmetic that only needs doing when a gain changes */
void pid_configure(pid_t *c, const pid_gains_t *g, float period_s)
{
    c->kp     = g->kp;
    c->ki_t   = g->ki * period_s * 0.5f;      /* trapezoidal, folded in */
    c->kd_ot  = g->kd / period_s;             /* the only division, here */
    c->alpha  = period_s / (period_s + g->tau);   /* the derivative filter */
    c->out_lo = g->out_lo;  c->out_hi = g->out_hi;
}
\end{ccode}

\item \textbf{Write the periodic path, and take the derivative on the
measurement.} This is about twenty lines and every line of it is a decision that
some published implementation gets wrong.
\begin{ccode}
/* pid.c: no division, no kick, and the integrator knows about saturation. */
float pid_step(pid_t *c, float setpoint, float meas)
{
    float err = setpoint - meas;
    c->integ += c->ki_t * (err + c->prev_err);        /* trapezoidal */
    c->dfilt  = c->dfilt + c->alpha * (meas - c->prev_meas - c->dfilt);
    float raw = c->kp * err + c->integ - c->kd_ot * c->dfilt;  /* D on meas */
    float out = clampf(raw, c->out_lo, c->out_hi);
    pid_antiwindup(c, raw, out);                      /* step 5 decides how */
    c->prev_err = err;  c->prev_meas = meas;
    return out;
}
\end{ccode}
The sign on the derivative term is inverted because it is taken on the
measurement rather than on the error, and that inversion is the whole
derivative-kick fix. Getting it backwards produces a controller that is stable,
plausible and wrong, which is the worst kind.

\item \textbf{Bound the integrator by the headroom the proportional term
leaves.} A single integrator limit chosen by taste is the usual approach and it
is worse than the arithmetic, which is short.
\begin{ccode}
/* the integrator may only use what the proportional term is not using */
float p_term = c->kp * err;
float i_hi   = (c->out_hi > p_term) ? (c->out_hi - p_term) : 0.0f;
float i_lo   = (c->out_lo < p_term) ? (c->out_lo - p_term) : 0.0f;
c->integ     = clampf(c->integ, i_lo, i_hi);
\end{ccode}

\item \textbf{Implement both anti-windup methods and measure them on the same
step.} Clamping is published everywhere. Back-calculation is in the standard
references and \textbf{in none of the small open implementations this volume
could verify}, so it is written here.
\begin{ccode}
/* antiwindup.md: two methods, one switch, both measured on the same input. */
static void pid_antiwindup(pid_t *c, float raw, float out)
{
    if (c->mode == AW_CLAMP) {
        if (raw != out) c->integ -= c->ki_t * (c->err + c->prev_err);
    } else {                                   /* AW_BACK_CALCULATION */
        c->integ += c->kb * (out - raw);       /* the excess, through its gain */
    }
}
\end{ccode}
Report both on one figure: the same setpoint step into the same saturating
plant, with recovery time and overshoot for each. That comparison does not exist
in the open literature for a small embedded controller, and producing it is
worth more than a third tuning guide.

\item \textbf{Prove the handler's cost with chapter 2's instrument.} Cycle
counts, not opinions, and the same statistics the whole volume uses.
\begin{shellcode}
python host/loop_stats.py --field ctrl_cycles --runs 1000000
# median, 99.9th percentile and maximum. Never the mean.
\end{shellcode}

\item \textbf{Measure following error, and publish it at the node's own rate
into the statistics rather than as a stream.} The quantity is defined per
period; what travels is its distribution.
\begin{ccode}
/* follow.c: the figure of merit, recorded every period, summarised at 50 Hz */
f->err = setpoint - meas;
if (fabsf(f->err) > f->peak) f->peak = fabsf(f->err);
f->sum_sq += f->err * f->err;                /* root mean square, per window */
if (fabsf(f->err) > f->tolerance) f->violations++;   /* chapter 15's code -4 */
\end{ccode}
The violation counter is the bridge to chapter 15: when it is non-zero the node
reports the path tolerance code that a controller already knows how to act on,
rather than a private number nobody can interpret.

\item \textbf{Expose the gains as interfaces, because the framework already
expects it.} Its published interface names include the three gains, two integral
clamp limits and a feedforward term, which is a strong hint that a joint node is
expected to be tunable over the wire.
\begin{plaincode}
command interfaces this node exports
  position            the setpoint, which is the normal path
  proportional        )
  integral            )  tuning without a rebuild, and without a debugger
  derivative          )
  integral_clamp_max  )  the two limits the framework names
  integral_clamp_min  )
  feedforward         the term a trajectory controller can supply
applied                at a period boundary, by the control task itself,
                       never from the middleware task (chapter 14's rule)
\end{plaincode}

\item \textbf{Write the host side against the header, not against the rendered
documentation.} The component derives from one of three interfaces, carries seven
lifecycle callbacks, exports its state and command interfaces, and implements a
real-time read and write pair.
\begin{plaincode}
on_init  on_configure  on_cleanup  on_activate  on_deactivate  on_shutdown
on_error
export_state_interfaces()      position, velocity, following_error
export_command_interfaces()    position, and the gains from step 8
read(time, period)             one bus frame in
write(time, period)            one bus frame out
CAUTION: the rendered documentation omits the lifecycle state argument from
these signatures. The header has it. Read the header.
\end{plaincode}
A published hardware component for a microcontroller over a serial link is the
structural template. Changing the transport to the bus of chapters 9 to 13 does
not change the shape, and saying so saves a week.

\item \textbf{Tune, and record what tuning meant.} The point of exposing gains
over the wire is that the record of what was tried can be a file rather than a
memory.
\begin{shellcode}
python host/tune.py --kp 4.0 --ki 12.0 --kd 0.05 --tau 0.004 --log tune.csv
python host/plot_step.py tune.csv --with clamp --with backcalc
\end{shellcode}
\end{steps}

\diagram{j16_timing}{The two things this chapter is judged on. Above, one
setpoint step into a saturating output, three times: no anti-windup, clamping,
and back-calculation, with the overshoot each one costs. Below, the rate
mismatch drawn to scale, with the node correcting ten times between one host
update and the next. The band on the lower trace is the tolerance whose
violation raises chapter 15's path code.}

\subsection*{Build, flash and debug}

\diagram{j16_data}{Five small controllers, read rather than described. The
columns are the four things that decide whether a controller is usable in a
joint, and the last column is why two of these five cannot be vendored into a
portfolio repository without work. The one this chapter starts from is the
third row, and its defect is a missing file header rather than anything in its
arithmetic.}

\begin{shellcode}
cmake --build build -j && ctest --test-dir build/host
probe-rs run --chip STM32H7A3ZITx build/firmware.elf
python host/tune.py --step 0.5 --log step.csv && python host/plot_step.py step.csv
\end{shellcode}

\begin{note}{When the controller is stable and the plot is still wrong}
Four causes, in the order they are worth checking. The derivative sign was not
inverted when it moved onto the measurement, which gives a stable controller
with the wrong dynamics. The gains were not premultiplied and the period changed,
so the tuning silently moved. The integrator has a limit chosen by taste rather
than from the proportional term's headroom, so it winds up inside its own limit.
Or the setpoint's validity time expired and the node is holding, correctly, at a
position the host thinks it has already left behind.
\end{note}

\subsection*{Verification and acceptance criteria}

\begin{itemize}
\item The division of labour between the two loops is written down with the
published default that justifies it.
\item The periodic path contains no division, proven by reading the compiler's
output once.
\item The derivative is taken on the measurement and filtered, and a setpoint
step produces no kick, shown on a plot.
\item The integrator limit is computed from the proportional term's headroom
rather than chosen.
\item Both anti-windup methods run, and the same saturating step is recorded for
each with overshoot and recovery time.
\item Controller execution time is reported as median, 99.9th percentile and
maximum over a million periods.
\item Following error is recorded every period, and its peak, root mean square
and tolerance violations are published.
\item A tolerance violation raises the standard path tolerance code from chapter
15, not a private number.
\item The gains are settable over the wire and are applied at a period boundary
by the control task.
\item The host component compiles against the header's signatures, and the
documentation discrepancy is noted in the repository so the next person does not
lose an hour.
\item Every figure showing this loop draws the plant dashed.
\end{itemize}

\subsection*{Variants}

\begin{table}[H]\centering\small
\begin{tabular}{p{24mm}p{26mm}p{44mm}p{22mm}p{20mm}}
\toprule
Axis & Variant & What changes & Cost & Built in full in \\
\midrule
Loop placement & Fast loop on the node & Feedback is local and the rate is the node's & The node must be trusted with the correction & Here \\
Loop placement & Loop on the host & One place to look, and one place to change & \textbf{A tenth of the rate, by the framework's own default}, plus bus latency inside the loop & Nowhere \\
Derivative & On the error, unfiltered & Simplest, and what the foil implementation does & A kick on every setpoint change and maximum noise gain & Nowhere \\
Derivative & On the measurement, filtered & No kick, and the noise is bounded & One state and one coefficient & Here \\
Anti-windup & None & The foil again & Overshoot proportional to how long it saturated & Nowhere \\
Anti-windup & Clamping & Published everywhere, cheap, effective & A discontinuity when it engages & Here \\
Anti-windup & Back-calculation & Smooth, and tunable through its own gain & \textbf{Nobody has published it in this language for a small target}, so it is written here & Here \\
Output & Saturated only & The minimum & A step in the command becomes a step in the output & Here \\
Output & Saturated and slew limited & Kinder to a real mechanism, and one library does it & One more state and one more limit & Here, as an option \\
Numeric type & Single precision & Natural on this core, which has the hardware for it & Care with the integrator's range & Here \\
Numeric type & Fixed point & Deterministic, and what a smaller part would need & Scaling work, and a library under a stronger licence & Nowhere, and the licence is why \\
\bottomrule
\end{tabular}
\caption{Variants for chapter 16. Four of these rows are built and five are refused, and two refusals are for reasons that have nothing to do with control theory: one is a rate that a published default fixes, and one is a licence that would follow the portfolio repository around.}
\end{table}

\subsection*{Pitfalls}

\begin{itemize}
\item Closing the fast loop on the host because that is where the framework is.
The framework's own default rate is a tenth of the node's.
\item Taking the derivative on the error, which produces a kick on every
setpoint change.
\item Taking the derivative unfiltered, which amplifies exactly the noise an
encoder has most of.
\item Inverting the derivative sign incorrectly when moving it onto the
measurement. The result is stable, plausible and wrong.
\item Dividing in the handler. Premultiplying the gains at configuration time is
free and published.
\item Baking the period invisibly into the gains, so that changing the rate
retunes the controller without anybody noticing.
\item Choosing the integrator limit by taste rather than from the proportional
term's headroom.
\item Reporting a private error number when a standard code exists that
controllers already act on.
\item Vendoring a controller whose licence is a sentence in a readme, or a
repository-level statement with no header in either source file. Both were found
during this chapter's research.
\item Reporting a mean. The whole volume reports median, 99.9th percentile and
maximum, and this chapter is not the place to start averaging.
\item Claiming a contouring error. That quantity is defined for coordinated
axes, this bench has one, and the distinction has had a published name since
1980.
\end{itemize}

\subsection*{Best practices applied}

\begin{itemize}
\item An architectural decision is settled by a published default rather than by
preference.
\item All configuration arithmetic happens once, so the periodic path is adds
and multiplies.
\item A known failure mode is fixed by construction rather than by tuning around
it.
\item Two solutions to the same problem are implemented and measured on the same
input, and the results of both are published.
\item A gap in the open literature is identified, filled, and named as a
contribution rather than presented as ordinary work.
\item Licences are checked at file level, not repository level, before anything
is vendored.
\item A documentation error is recorded in the repository so that the next
person loses no time to it.
\item A quantity is reported in the vocabulary a consumer already understands.
\item A modelled link is drawn as modelled, in every figure, without exception.
\end{itemize}

\subsection*{Stretch goals}

\begin{itemize}
\item Publish the back-calculation implementation with its measurement beside
it. It is a small, complete, genuinely missing piece of open embedded code.
\item Add a feedforward term from the trajectory's own velocity and measure how
much following error it removes, which is the cheapest large improvement in
motion control and is almost never shown with numbers.
\item Measure the host loop's actual jitter on a stock kernel, then again with
the scheduling policy and limits the framework asks for, and report the
difference. Most published advice on this subject has no measurement behind it.
\item Run the same controller against the plant model at several periods and
show that premultiplied gains keep the tuning while baked-in gains do not.
\item Implement the framework's asynchronous component mechanism for the bus
read and find out whether it matters at this rate.
\end{itemize}

\subsection*{Roadmap and next steps}

Chapter 17 asks what a deterministic fieldbus would change about all of this,
and answers honestly that it is not on this bench and why that is a defensible
place to stop.

The published progression from here has three parts. For the control theory, the
free and complete treatment of discretisation, filtered derivatives and the
derivative-kick argument is the right first read, and the standard textbook
chapter on integrator windup is the right second one. For the host side, the
control framework's own documentation on hardware components and its list of
interface names. For the argument this volume has been building toward, the
single-threaded executor with logical execution time semantics, whose own
bibliography is a ready-made reading list on response-time analysis and on
scheduling for this class of system.

\subsection*{Portfolio evidence}

\begin{itemize}
\item The two anti-windup methods on one plot, which is a measurement that does
not exist in the open literature for a small embedded controller.
\item The controller comparison table, which demonstrates reading source rather
than reading readmes.
\item The licence findings, including one library whose permission is a sentence
in a readme and one whose source files carry no header at all.
\item The loop division document, which shows an architectural decision made
from a published number.
\item The following error statistics reported the way the whole volume reports
timing.
\end{itemize}

\subsection*{Sources}

Normative references:
\begin{itemize}
\item Astrom and Murray, \textit{Feedback Systems}, chapter 10, section 10.4,
"Integrator Windup". The book is the citation; the hosting site serves a broken
certificate chain, so no link is given.
\item Astrom and Hagglund, \textit{PID Controllers: Theory, Design and Tuning},
second edition, ISA, which is the standard source for back-calculation. No
fetchable copy was found during this volume's research.
\item Koren, "Cross-Coupled Biaxial Computer Control for Manufacturing Systems",
Journal of Dynamic Systems, Measurement, and Control, volume 102, pages 265 to
272, 1980, which is the origin of contouring error as a quantity distinct from
axis following error. This bench has one axis and therefore has the second and
not the first.
\item Pas, "PID Controllers", free and complete on discretisation, the causal
backward-difference derivative, the filtered derivative as finite differences
plus an exponential moving average, and the derivative-kick argument. It does not
treat windup systematically and does not argue for a constant period, so it is
not cited for either.
\end{itemize}

Reusable implementations:
\begin{itemize}
\item The host control framework, Apache-2.0.\newline
\url{https://github.com/ros-controls/ros2_control}
\item The real-time helper library, BSD-3-Clause.\newline
\url{https://github.com/ros-controls/realtime_tools}
\item The best-fitting small controller, MIT at repository level and with no
header in either source file.\newline
\url{https://github.com/pms67/PID}
\item A hardware component for a microcontroller over a serial link,
BSD-3-Clause, the structural template for the host side.\newline
\url{https://github.com/joshnewans/diffdrive_arduino}
\item The single-threaded executor with logical execution time semantics,
Apache-2.0, whose published bibliography is the reading list for chapter 20.\newline
\url{https://github.com/ros2/rclc}
\end{itemize}

Two negative findings from this chapter's research are worth recording, because
a reader will otherwise find both libraries before finding this page. One widely
used small controller carries its permission only as a sentence in a readme file,
with no licence file, no identifier line and no field in its metadata, so an
automated check reports none. Another has no licence anywhere at all, which means
all rights reserved, and has been dormant since Monday 10 April 2017. Neither is
vendored here.
   % The loop closed over the bus: setpoint in, state out, following error
\project{17}{What a real-time fieldbus would change, and why it is not on this bench}{An honest boundary}{Slave hardware, cycle time, the open master stack, what would have to be bought}

\begin{keyfacts}
\item[Adds to the node] An honest boundary. The node gains a costed, sourced
account of the one thing it cannot become, and the reader gains the ability to
tell a structural limit from a lack of effort
\item[Peripherals] None on this part, and that is the finding. The peripherals
this chapter is about live on a chip nobody here owns
\item[Depends on] Chapter 4 for what this part does not have, chapter 3 for the
synchronisation this would replace, chapter 11 for the two message shapes it
maps onto
\item[Real or modelled] \textbf{Neither.} Nothing is built. Everything here is
read, quoted, priced and decided, and the chapter says which of its claims come
from documents it opened and which it could not confirm
\item[Difficulty] 2 of 5 to read, 5 of 5 to build, and the gap between those two
numbers is the chapter
\item[Effort] Two evenings of reading. The published estimate for building it is
\textbf{six to eight weeks}, and that estimate comes from the technology group
itself
\item[Deliverable] A decision document: what it would cost in parts, weeks and
licence obligations, what it would buy in cycle time and synchronisation, and
why this volume stops here
\end{keyfacts}

\subsection*{Why this chapter}

Sixteen chapters have built a joint node on a bus that is good enough for a
joint and is not what a modern robot arm uses between its controller and its
axes. The serious machines use a deterministic fieldbus on Ethernet wiring, and
a volume that never mentioned that would be leaving out the thing a reader is
most likely to be asked about.

So this chapter does the reading, and it reaches a conclusion that is stronger
than difficult.

\begin{note}{The boundary is structural and it is settled twice over}
This part has no Ethernet controller, which chapter 4 established from the
absence of its interrupt vector slots rather than from a product page. And the
silicon vendor whose part this is \textbf{sells no subdevice controller at
all}: a text search of the technology group's own controller overview, dated
\pubdate{July 2025}, for either of that vendor's names returns zero results. So this is
not a matter of effort or of skill. There is no combination of firmware and
patience that turns this board into a participant on that bus, and knowing the
difference between that and a hard problem is worth a chapter.
\end{note}

The asymmetry is the whole subject. The implementation guide states the
requirement for the controlling end in one sentence: the only hardware
requirement for a main device is a standard network interface controller at 100
megabit per second, full duplex. A personal computer's network card is enough.
The other end needs a dedicated controller chip, a configuration memory, and per
port a connector, magnetics, a physical layer device and passives. \textbf{The
Raspberry Pi on this bench could be the master this afternoon. The node cannot
be a participant at any price short of different hardware.}

\subsection*{Prior art and what to reuse}

\begin{table}[H]\centering\small
\begin{tabular}{p{34mm}p{46mm}p{38mm}p{20mm}}
\toprule
Source & What it gives & What it does not & Licence \\
\midrule
The implementation guide, V3.2.0, Wednesday 7 August 2024 & The normative shape
of both ends, the one-sentence requirement for the controlling end, the
subdevice's internal blocks, and \textbf{an effort estimate from the technology
group itself}: about six to eight weeks to reach a working subdevice operated by
a standard main device & Nothing that helps a part with no Ethernet controller &
technology group \\
The controller overview, \pubdate{July 2025} & \textbf{The chip list, extracted in full},
with ports, memory, sync managers and address-mapping units per part, and the
host interfaces each offers & It is a list of parts to buy, which is the point &
technology group \\
The open master stack, 2085 stars, last pushed Tuesday 8 September 2026 & A
complete master in portable C, and since its 2.0.0 release of Friday 11 July
2025 a parser that turns a network description into compilable C &
\textbf{Its licence file states GPLv3}, with a commercial option and a sentence
saying a commercial product likely needs one. \textbf{The previous arrangement
was GPLv2 with a linking exception, and that changed in 2025} & GPL-3.0 \\
The kernel-space master, default branch stable-1.6, active Saturday 19 September
2026 & The other serious open master: kernel modules, a userspace library and a
command line tool & \textbf{Out-of-tree kernel modules.} Its root carries both a
GPLv2 and an LGPLv2.1 file and \textbf{the readme does not state which applies
to which half}, so the conventional split could not be confirmed & GPL-2.0 and
LGPL-2.1 \\
The open subdevice stack, 847 stars, last pushed Tuesday 8 April 2025 & The
other end in portable C: mailbox handling, an object dictionary, service data,
process data mapping, file transfer, and an address-offset layer so it can reach
a controller over any interface. \textbf{It kept the licensing arrangement the
master stack left behind} & No controller, no description file, no vendor
identifier and no conformance. Those are the four things that are not code &
GPL-2.0 with a linking exception \\
The vendor subdevice stack & Portable C covering the state machine, several
mailbox protocols, synchronisation and \textbf{an example drive profile}, which
is exactly a joint & \textbf{Free of charge but gated behind membership, with
redistribution restricted.} The gate is membership, not a vendor identifier &
restricted \\
The robot framework bridge, Apache-2.0, 335 stars, last pushed Monday 31 August
2026 & A hardware interface in which devices are described in parameter files
rather than in C++ per device, which is the right idea & \textbf{Its
installation builds on the kernel-space master and requires disabling secure
boot to load unsigned kernel modules.} A permissive wrapper does not remove the
copyleft from the stack that actually gets deployed & Apache-2.0, on a GPL stack
\\
\bottomrule
\end{tabular}
\caption{Prior art for chapter 17. The last row is the sentence most worth carrying away: the licence on the wrapper is not the licence on the deployed system, and a reader choosing a stack on the strength of a badge will be surprised later rather than sooner.}
\end{table}

\subsection*{What the node gains}

Nothing it can run. What it gains is a boundary that can be defended in a
conversation: what the other bus does that this one does not, what it would cost
in parts and weeks, which of that cost is money and which is obligation, and the
specific reason this board is excluded. A portfolio that contains one clear
statement of what was deliberately not built is more credible than one that
contains none.

\subsection*{Parts from the inventory}

\begin{table}[H]\centering\small
\begin{tabular}{p{40mm}p{66mm}p{40mm}}
\toprule
Part & Role & Interface \\
\midrule
Raspberry Pi 4 & \textbf{Could be the controlling end today.} Its network
interface already meets the only stated hardware requirement & Its existing
Ethernet port \\
NUCLEO-H7A3ZI-Q & \textbf{Cannot be the other end, structurally.} No Ethernet
controller, and no subdevice controller exists from this silicon vendor & n/a \\
A subdevice controller & \textbf{Not owned.} Would be bought, and then would
need soldering, which this bench cannot do & Serial or parallel, to the node \\
Two physical layer devices, magnetics, two connectors and passives &
\textbf{Not owned}, and required per port, with two ports the minimum & n/a \\
A configuration memory & \textbf{Not owned.} It holds the device's identity and
its binary description & Two wires to the controller \\
\bottomrule
\end{tabular}
\caption{Inventory items for chapter 17. Four of five rows say not owned, and one of those says it would also need soldering, which the bench notes have recorded as impossible since chapter 1. That is the honest bill of materials for this chapter and it is why nothing here is built.}
\end{table}

\subsection*{System architecture}

\diagram{j17_arch}{The asymmetry, drawn. On the left, the controlling end: one
standard network interface, which is the guide's entire stated hardware
requirement. On the right, the other end: a dedicated controller, a
configuration memory, and per port a connector, magnetics, a physical layer
device and passives. The bench owns everything on the left and nothing on the
right, and no amount of firmware changes which side of that line this board is
on.}

\subsection*{Peripheral configuration}

\begin{table}[H]\centering\small
\begin{tabular}{p{22mm}p{26mm}p{24mm}p{30mm}p{30mm}}
\toprule
Peripheral & Mode & Clock source & Pins and function & Interrupt and transfers \\
\midrule
Ethernet controller & \textbf{Absent on this part} & n/a & n/a & Established in
chapter 4 from the absence of its vector slots \\
Subdevice controller & \textbf{Absent from this vendor's catalogue entirely} &
n/a & n/a & Confirmed against the technology group's own overview of \pubdate{July 2025} \\
The process data interface & What the node \textbf{would} attach through, if the
chip existed here & From the controller & Serial, or an 8 or 16 bit parallel
bus & \textbf{Ceiling: about 12.5 megabytes per second} on that internal
interface \\
Distributed clocks & What would replace chapter 3's software synchronisation &
The bus itself & n/a & 64 bit, 1 nanosecond units, and the compensation happens
in hardware \\
\bottomrule
\end{tabular}
\caption{Peripheral configuration for chapter 17. Two rows are absences and two are descriptions of hardware that is not here. The throughput ceiling in the third row is worth remembering: it is a published number, it is generous for a joint, and it is the kind of figure people assume is unlimited.}
\end{table}

\subsection*{Wiring}

\diagram{j17_wiring}{The three routes onto that bus, and what each one costs.
Keep this part and add a controller beside it. Replace this part with one that
has a controller inside. Or use a part that implements the protocol in
programmable logic against a vendor stack. The first is the most instructive and
needs soldering; the second is the cleanest and means a different board; the
third is flexible and ties the design to one vendor's stack.}

Nothing is wired in this chapter. The figure prices three purchases rather than
describing a connection, which is the right shape for a chapter whose
conclusion is not to buy any of them yet.

\subsection*{Memory and timing budget}

This chapter's budget is not in bytes.

\begin{table}[H]\centering\small
\begin{tabular}{p{44mm}p{28mm}p{28mm}p{30mm}}
\toprule
Quantity & Budget & Published figure & Note \\
\midrule
Synchronisation, this volume & single-digit microseconds & chapter 3, measured & software, over the existing bus \\
Synchronisation, that bus & \textbf{one to two orders better} & the guide claims much better than 1 microsecond; one controller vendor claims better than 100 nanoseconds & in hardware, with propagation delay measured and drift compensated \\
Process data interface ceiling & n/a & about 12.5 megabytes per second & published, and generous for a joint \\
Development effort & n/a & \textbf{six to eight weeks} & the technology group's own estimate, for a working subdevice \\
Parts, one node & n/a & controller, memory, two physical layer devices, magnetics, connectors, passives & plus assembly this bench cannot do \\
Vendor identifier & n/a & \textbf{free of charge}, and not required at all for a machine builder integrating devices & a corrected assumption \\
Conformance test tool & n/a & \textbf{a paid annual subscription}, and its in-house use is \textbf{mandatory when selling the device} & this is the real cost, and it is recurring \\
\bottomrule
\end{tabular}
\caption{The budget for chapter 17, in weeks, parts and obligations rather than bytes. The last two rows correct the assumption most people arrive with: the identifier everybody worries about is free, and the cost that actually bites is a recurring subscription to a test tool that the implementation guide makes mandatory before a device can be sold.}
\end{table}

\subsection*{Firmware design (UML)}

\diagram{j17_uml}{Inside a subdevice controller, and the one part of it this
volume has already built twice in software. The two arbitration modes map
cleanly onto chapter 11's two message shapes: the buffered mode is the state
frame, where a late reader wants the newest consistent value, and the mailbox
mode is the command and configuration path, where nothing may be lost. Finding
that the same distinction was made in silicon is a good sign about chapter 11.}

Three things are worth understanding even by somebody who will never build one.

\textbf{The processing is in hardware and happens while the frame moves.} Each
participant reads the data addressed to it as the frame passes and inserts its
own into the same frame on the way through. There is no receive, process,
transmit cycle at each hop, which is where the determinism comes from and why
the number of participants costs so little.

\textbf{Addressing is a hardware mapping.} Address-mapping units convert logical
addresses to physical ones bit by bit, configured through registers, so a master
can address a bit in the middle of one device's memory without that device's
firmware being involved at all.

\textbf{Synchronisation is a hardware service, not a protocol you write.} The
clock unit provides 64 bit system time in nanosecond units, with propagation
delay measured and offset and drift compensated in the controller. Chapter 3
built a software equivalent over the existing bus and reached single-digit
microseconds. That is a good result for software and it is one to two orders
away from what dedicated hardware delivers, and saying both halves of that
sentence is the honest comparison.

\subsection*{Data flow (ASCII)}

\begin{asciiart}
  the controlling end                    the participants
  ------------------                     ----------------
  a standard network interface,          a dedicated controller chip,
  100 Mbit/s full duplex                 a configuration memory,
        |                                and PER PORT: a connector,
        |                                magnetics, a physical layer
        |                                device and passives
        |                                        |
        |   one frame, going downstream          |
        +--------------------------------------->+  reads what is addressed
        |                                        |  to it AS THE FRAME MOVES
        |                                        |  and inserts its own data
        |                                        |  into the same frame
        |                                        v
        |                                    next participant
        |                                        |
        +<---------------------------------------+  the frame returns

  what the bench has:   the left column, already, on the Pi
  what the bench has:   nothing at all in the right column
  what this part has:   no Ethernet controller (chapter 4), and its vendor
                        sells no subdevice controller (the July 2025 overview)

  the gap is a purchase and a soldering iron, not a firmware problem
\end{asciiart}

\subsection*{Repository layout}

\begin{plaincode}
joint-node/
  doc/
    fieldbus-decision.md      # + this chapter, and it is the whole deliverable
      what it would buy         cycle time, hardware synchronisation, the
                                profile a robot framework already speaks
      what it would cost        parts, six to eight weeks, a recurring
                                subscription before the device may be sold
      why not here              this part has no Ethernet controller and this
                                vendor sells no subdevice controller
      the three routes          add a chip, change the part, or use a soft
                                implementation tied to one vendor's stack
      licence findings          which stack changed licence and when, and
                                what the wrapper's badge does not cover
  src/  host/  test/  tools/  proto/  README.md      # all unchanged
\end{plaincode}

\subsection*{Steps}

\begin{steps}
\item \textbf{Settle the premise from documents rather than from impressions.}
Two checks, both quick, and together they turn an opinion into a finding.
\begin{plaincode}
check 1   does this part have an Ethernet controller?
          chapter 4 answered it from the ABSENCE of two interrupt vector
          slots in the vendor's own header, which is stronger evidence
          than a product page
check 2   does this vendor sell a subdevice controller at all?
          search the technology group's own controller overview, dated
          July 2025, for both of this vendor's names
          result: zero hits
conclusion  structural, not difficult
\end{plaincode}

\item \textbf{Read the one sentence that explains the whole asymmetry.} It is in
the implementation guide and it is worth quoting exactly, because it is the
reason a hobbyist can write a master in an evening and not a participant.
\begin{plaincode}
the controlling end   "The only hardware requirement for an EtherCAT
                      MainDevice is a standard Network Interface Controller
                      (NIC, 100 Mbit/s full duplex)."
the other end         a dedicated controller chip, a configuration memory,
                      and per port a connector, magnetics, a physical layer
                      device and passives. Two ports minimum, with loops
                      closed automatically on unconnected ones.
\end{plaincode}

\item \textbf{Extract the chip list and choose the one a joint would actually
want.} The overview is a table; reading it takes twenty minutes and it settles
the hardware question for years.
\begin{plaincode}
part       ports  memory  sync  mapping  host interface
ET1100     2 to 4   8 kB     8      8    serial and parallel, 8 and 16 bit
LAN9252    2 + MII  4 kB     4      3    host bus, serial, quad serial
LAN9253/4  2 + opt  8 kB     8      8    as above, with I/O on the 9254
LAN9255    2 + opt  8 kB     8      8    INTEGRATED Cortex-M4F core,
                                         1 MB flash, 256 kB memory
AX58100    2, PHYs  9 kB     8      8    serial, parallel
XMC4800    2 MII    8 kB     8      8    internal, Cortex-M4 core
\end{plaincode}
\textbf{One part is obviously built for this problem}, and its own feature list
says why without any interpretation: three channel pulse width modulation, step
and direction control, incremental and Hall encoder interfaces, and an emergency
stop input, with both physical layer devices integrated. That is a joint node's
entire peripheral set, inside the bus controller, in an eighty pin package, with
a datasheet at version 1.09 of Thursday 12 December 2024.

\item \textbf{Record what is legacy, so nobody designs a new board around it.}
Two parts that appear in older articles are not in the current overview and one
product page returns not found. Write that down once.

\item \textbf{Understand the two arbitration modes, because the distinction is
already in this node.} They are worth knowing even if the chip never arrives.
\begin{plaincode}
buffered mode   three buffers, cyclic process data, and the consumer always
                receives the latest CONSISTENT buffer. A late reader gets
                fresh data and never a torn one.
                  -> this is chapter 11's state frame, exactly
mailbox mode    a handshake, and no data is lost
                  -> this is chapter 11's command and configuration path
\end{plaincode}
Chapter 11 reached the same split from first principles over a different bus.
Finding it made in silicon is the strongest independent confirmation that
chapter's design has had.

\item \textbf{Get the description files straight, because the names are
confusing and the direction of flow is not obvious.}
\begin{plaincode}
the device file   vendor supplied XML against the group's schema: identity
                  (vendor, product code, revision), object dictionary,
                  process data mapping, sync manager configuration, and
                  which mailbox protocols the device supports
the network file  produced by a configuration tool that reads ALL the device
                  files on a network: topology, per device initialisation
                  commands, and the cyclic command list
who executes what THE MASTER EXECUTES THE NETWORK FILE. The device files are
                  inputs to a tool, not something the master reads at run time
the same, binary  lives in the device's own configuration memory
\end{plaincode}

\item \textbf{Price it honestly, and correct the two assumptions everybody
arrives with.} This is the step that changes decisions.
\begin{plaincode}
membership          free of charge
the vendor id       FREE OF CHARGE, and a machine builder integrating
                    devices is explicitly NOT required to obtain one
the conformance     A PAID ANNUAL SUBSCRIPTION, and the implementation guide
test tool           says its in-house use "is mandatory when selling the
                    device to the market"
so the real cost    is recurring, is attached to selling rather than to
                    building, and is not the thing people worry about
\end{plaincode}

\item \textbf{Read the licence files, not the badges, and write down what
changed and when.} This chapter's licence findings are more useful than its
technical ones, because they are the part that will have moved since anybody
last looked.
\begin{plaincode}
the open master     GPLv3 since version 2.0.0 on Friday 11 July 2025, with a
                    commercial option and a sentence saying a commercial
                    product likely needs one. It was GPLv2 with a linking
                    exception before that. Automated checks report nothing
                    only because the file is named LICENSE.md
the kernel master   a GPLv2 file and an LGPLv2.1 file at the root, and the
                    readme does NOT say which covers which half. The
                    conventional split is assumed and COULD NOT BE CONFIRMED
the open subdevice  GPLv2 WITH A LINKING EXCEPTION, which is the arrangement
                    the master stack left behind. The asymmetry is worth a
                    sentence in any evaluation
the vendor stack    free of charge, gated behind membership, redistribution
                    restricted. The gate is membership, not an identifier
\end{plaincode}

\item \textbf{Follow the robot framework bridge one level down before believing
its badge.} This is the trap of the chapter and it generalises far beyond this
bus.
\begin{plaincode}
the wrapper        permissively licensed, actively maintained, and it lets
                   devices be described in parameter files instead of in
                   C++ per device, which is genuinely good design
what it sits on    the kernel-space master, which is out-of-tree kernel
                   modules under copyleft, and whose installation procedure
                   requires DISABLING SECURE BOOT to load unsigned modules
the conclusion     THE PERMISSIVE LICENCE ON THE WRAPPER DOES NOT REMOVE THE
                   COPYLEFT FROM THE STACK THAT GETS DEPLOYED, and the
                   secure boot requirement is a deployment decision somebody
                   senior should make rather than discover
\end{plaincode}

\item \textbf{Write the decision file and stop.} Four headings, two pages, and
it is the chapter's only deliverable.
\begin{plaincode}
doc/fieldbus-decision.md
  what it would buy    cycle times a joint would notice, hardware
                       synchronisation one to two orders better than
                       chapter 3's, and a drive profile a robot framework
                       already speaks
  what it would cost   parts this bench cannot solder, six to eight weeks by
                       the group's own estimate, and a recurring
                       subscription before the device may be sold
  why not here         this part has no Ethernet controller and this vendor
                       sells no subdevice controller: structural
  what would change it a different part, and the chapter names three
\end{plaincode}
\end{steps}

\diagram{j17_timing}{What the money buys, drawn to scale. Above, one frame
moving downstream while every participant reads and writes into it as it passes,
which is where the determinism comes from and why participants cost so little.
Below, the synchronisation comparison: chapter 3's software result over the
existing bus against the published hardware figures, on a scale that has to be
logarithmic for both to fit on one page.}

\subsection*{Build, flash and debug}

\diagram{j17_data}{The two tables this chapter exists to produce. Above, the
controllers, with the one that carries a joint's entire peripheral set inside
the bus chip marked. Below, the stacks, with what each licence actually says
rather than what its badge reports, and the date the most important one changed.}

Nothing is built, so there is nothing to flash. The commands that belong to this
chapter are the two searches that settled its premise, and they are worth being
able to repeat.

\begin{shellcode}
# the two checks that turned an impression into a finding
grep -c 'ETH_IRQn\|ETH_WKUP_IRQn' vendor/Include/stm32h7a3xx.h   # expect 0
python tools/check_esc_overview.py --pdf esc_overview_2025-07.pdf --grep st
\end{shellcode}

\begin{note}{When a permissive badge is on top of a copyleft stack}
The pattern is common and this chapter's example is only one instance. A
permissively licensed wrapper, actively maintained and well designed, sits on a
dependency whose licence is stronger, and the badge shown on the wrapper's page
describes the wrapper alone. What ships is the whole stack. The check is one
level of dependency deeper than most people look, it takes ten minutes, and the
answer belongs in a file rather than in somebody's memory.
\end{note}

\subsection*{Verification and acceptance criteria}

\begin{itemize}
\item The premise is settled from two documents, not from impressions, and both
checks are repeatable from the repository.
\item The asymmetry between the two ends is quoted from the implementation guide
rather than paraphrased.
\item The controller list is extracted with ports, memory, sync managers,
mapping units and host interface per part.
\item The part whose feature list already covers a joint's peripherals is
identified, with its datasheet version and date.
\item Legacy parts are recorded as legacy so that nobody designs a new board
around them.
\item The two arbitration modes are mapped onto chapter 11's two message shapes.
\item The description file names and the direction of flow are stated correctly:
the master executes the network file.
\item The cost is stated with the two common assumptions corrected: the
identifier is free, and the recurring test tool subscription is the real cost.
\item Every stack's licence is read from its file, with the date of the most
recent change recorded.
\item The wrapper's dependency is followed one level down, and the conclusion is
written in the repository.
\item Claims that could not be confirmed are marked as such, including one
standard this volume declines to cite at all.
\end{itemize}

\subsection*{Variants}

\begin{table}[H]\centering\small
\begin{tabular}{p{24mm}p{26mm}p{44mm}p{22mm}p{20mm}}
\toprule
Axis & Variant & What changes & Cost & Built in full in \\
\midrule
Route & Keep this part, add a controller & The most instructive: the node stays and the bus arrives beside it & A chip, a memory, two physical layer devices, magnetics, connectors, and soldering this bench cannot do & Nowhere \\
Route & A controller with a core inside & The cleanest: one chip is both the participant and the processor & A different board and a different core, so fifteen chapters of this volume move & Nowhere \\
Route & Soft, on programmable logic & Flexible and protocol switchable & Tied to one vendor's certified stack, and a different class of part & Nowhere \\
Master & The portable open stack & Runs anywhere, and since 2025 generates compilable C from a network description & \textbf{Copyleft since Friday 11 July 2025, and a commercial product likely needs a licence} & Nowhere \\
Master & The kernel-space stack & The other serious option, with a command line tool & Out-of-tree kernel modules, and an unconfirmed licence split & Nowhere \\
Master & The framework bridge & Devices described in parameter files rather than in C++ & \textbf{A permissive wrapper on a copyleft stack, plus disabling secure boot} & Nowhere \\
Participant stack & The open one & Portable C with an interface layer that reaches a controller over anything & Copyleft with a linking exception, and it supplies none of the four non-code things a device needs & Nowhere \\
Participant stack & The vendor one & Includes an example drive profile, which is exactly a joint & Membership gated, redistribution restricted & Nowhere \\
\bottomrule
\end{tabular}
\caption{Variants for chapter 17, and every row says nowhere. That is the point of the chapter. A variants table whose entire right-hand column is a refusal is more useful than a table that quietly implies everything was tried, provided the refusals are costed, and these are.}
\end{table}

\subsection*{Pitfalls}

\begin{itemize}
\item Treating this as a hard problem rather than a structural one. No firmware
turns a part with no Ethernet controller into a participant on this bus.
\item Assuming the two ends are symmetric. One needs a network card; the other
needs a chip, a memory and a physical layer per port.
\item Assuming the vendor identifier is the expensive part. It is free, and a
machine builder integrating devices does not need one at all.
\item Missing the cost that is real: a recurring subscription to a test tool
whose in-house use the implementation guide makes mandatory before selling.
\item Reading a licence badge instead of a licence file. One major stack here
changed licence in 2025 and automated checks still report nothing, because the
file has a different extension than the checker expects.
\item Trusting a permissive wrapper's licence to describe the stack it deploys.
\item Designing a new board around a controller that is no longer in the current
overview.
\item Confusing the two description files, or believing the master reads the
device files at run time. It executes the network file a tool produced.
\item Citing a standard that the documents do not mention. One commonly repeated
reference could not be found anywhere in the material this volume read, and is
therefore not cited here.
\item Assuming the internal interface is fast enough for anything. It has a
published ceiling of about 12.5 megabytes per second, which is generous for a
joint and finite for other things.
\end{itemize}

\subsection*{Best practices applied}

\begin{itemize}
\item A limit is proven from primary documents rather than asserted, and the
proof is repeatable.
\item A decision not to build is written down, costed, and given the same care
as a decision to build.
\item Every licence is read from its own file, and the date of the most recent
change is recorded.
\item A dependency is followed one level deeper than its badge.
\item An independent design's agreement with an earlier chapter is noticed and
reported as evidence rather than as coincidence.
\item Claims that could not be confirmed are marked, including a standard this
volume declines to cite.
\item A published effort estimate is used instead of an invented one, and its
source is named.
\end{itemize}

\subsection*{Stretch goals}

\begin{itemize}
\item Run the controlling end on the Pi against a simulated participant, which
costs nothing but time and would make the asymmetry concrete rather than
quoted.
\item Buy the joint-friendly controller and an evaluation board for it, and find
out whether its integrated pulse width modulation, step and direction, encoder
interface and emergency stop input really do replace chapters 5, 7 and 18's
hardware.
\item Write a device description file for this node's existing object model, as
a paper exercise, and see how much of chapter 15's mapping survives the
translation.
\item Measure the kernel-space master's cycle jitter on a stock kernel and on a
real-time one, which is the same measurement chapter 16 wanted and nobody
publishes.
\end{itemize}

\subsection*{Roadmap and next steps}

Chapter 18 returns to hardware that is on the bench, and to the one question a
joint must answer before it is allowed near anybody: how it stops.

The published progression from here is the implementation guide, which is
readable and free and settles most of what a newcomer wants to know, followed by
the controller overview for the hardware question and the device information
specification for the description files. For the drive profile a joint would
implement, the directive is the right entry point, and it is a directive rather
than the profile itself, which is a distinction worth keeping straight.

\subsection*{Portfolio evidence}

\begin{itemize}
\item The decision file, which is two pages and demonstrates costing an option
out rather than avoiding it.
\item The licence findings, including one stack that changed licence in 2025,
one whose split could not be confirmed, and one permissive wrapper on a copyleft
dependency.
\item The corrected cost picture, which contradicts the common assumption in
both directions.
\item The observation that chapter 11's two message shapes match the two
arbitration modes implemented in silicon.
\end{itemize}

\subsection*{Sources}

Normative references:
\begin{itemize}
\item ETG.2200, "EtherCAT Implementation Guide", V3.2.0, Wednesday 7 August
2024, from which the hardware requirement for the controlling end, the six to
eight week effort estimate, the process data interface ceiling and the
conformance test tool obligation are all quoted.
\item ETG, "EtherCAT SubDevice Controller (ESC) Overview", \pubdate{July 2025}, from which
the controller table is extracted, and against which this chapter's second
premise check was run.
\item ETG.2000, "EtherCAT SubDevice Information (ESI) Specification", V1.21,
Monday 26 May 2025. Note that the current title says SubDevice, not Slave.
\item ETG.2010 for the configuration memory interface, ETG.2100 for the network
information, ETG.5001 for the modular device profile, ETG.6010 for the drive
profile implementation directive, and ETG.5100 for safety.
\item ETG.1000 represents IEC 61158 Type 12. Both the bus and its safety layer
are IEC standards, IEC 61158 and IEC 61784, and the drive profile a joint would
implement, CiA 402, is IEC 61800-7-201. \textbf{The individual titles of the
ETG.1000 parts could not be confirmed and are not printed here. A further
standard that is frequently cited alongside these appears nowhere in the
material read for this volume and is therefore not cited at all.}
\end{itemize}

Reusable implementations:
\begin{itemize}
\item The portable open master, GPL-3.0 since version 2.0.0 of Friday 11 July
2025, previously GPL-2.0 with a linking exception.\newline
\url{https://github.com/OpenEtherCATsociety/SOEM}
\item The portable open participant stack, GPL-2.0 with a linking
exception.\newline
\url{https://github.com/OpenEtherCATsociety/SOES}
\item The kernel-space master, GPL-2.0 and LGPL-2.1 at the root with the split
unstated.\newline
\url{https://gitlab.com/etherlab.org/ethercat}
\item The robot framework bridge, Apache-2.0, which deploys on the kernel-space
master.\newline
\url{https://github.com/ICube-Robotics/ethercat_driver_ros2}
\end{itemize}
   % What a real-time fieldbus would change, and why it is not on this bench

\part{Making it trustworthy}
\project{18}{Safe states, and the workspace sensor that triggers one}{A way to stop}{Safe-state machine, latching, entry actions, a presence sensor as a trigger}

\begin{keyfacts}
\item[Adds to the node] A way to stop. The node gains a latched safe state
reachable from everywhere, three independent triggers, and a written answer to
the question a joint has to answer before it goes near anybody
\item[Peripherals] The ranging shield as the workspace trigger, both on-die
watchdogs, and one output that asserts the safe state without software
\item[Depends on] Chapter 5 for position, chapter 7 for the output stage,
chapter 12 for bus loss, chapter 16 for following error
\item[Real or modelled] \textbf{Real trigger, modelled reaction.} The sensor,
the latch, the watchdogs and the output are real. \textbf{There is no motor and
no brake}, so the reaction is asserted and logged rather than felt, and the
chapter never implies otherwise
\item[Difficulty] 4 of 5
\item[Effort] Three evenings, one of which is reading standards and finding out
what may and may not be claimed from them
\item[Deliverable] Two safe states rather than one, a latch that meets a
published clause by clause list, a trigger with a measured reaction time, and a
document that says exactly what this node is and is not
\end{keyfacts}

\subsection*{Why this chapter}

A joint that can move is a joint that has to be able to stop, and the word for
what it stops into has a definition. That definition is the first surprise of
this chapter, because it contains no engineering content at all.

The definitions part of the general functional safety standard gives it in
seven words: the state of the equipment under control when safety is achieved.
It does not say power removed. It does not say brake applied. \textbf{What
counts as the safe state is an output of the risk assessment, not of the
firmware.} And the note attached to the definition is the load-bearing sentence
for a robot joint: for some situations a safe state exists only so long as the
equipment is continuously controlled.

\begin{note}{A joint has two safe states and they are not interchangeable}
For a horizontal joint, or one whose gearing will not back-drive, removing
torque is a safe state and it persists with no energy at all. \textbf{For a
vertical joint carrying a payload, removing torque is not a safe state}, because
the axis then descends. There the safe state is the brake engaged with torque
removed, and it is reached through a controlled deceleration rather than
instantly. Torque removed and brake held are two different safe states, the
drive standard names them and the definitions standard names neither, and a
design that has only one of them has only solved half of the problem. This
chapter builds the machine for both and says which one this bench can
demonstrate.
\end{note}

\subsection*{Prior art and what to reuse}

\begin{table}[H]\centering\small
\begin{tabular}{p{34mm}p{46mm}p{38mm}p{20mm}}
\toprule
Source & What it gives & What it does not & Licence \\
\midrule
The definitions part of the general standard, edition 2.0, Friday 30 April 2010 & The safe
state definition itself, \textbf{read from the official preview at printed page
11}, with the note about continuous control that this chapter is built around.
Also the tool classification into three types, read directly & Paywalled beyond
the preview. It says nothing about how to implement anything, deliberately &
paid \\
The emergency stop standard, third edition, 2015 & \textbf{Clause 4.1.1, read
directly}: six requirements the latch in this chapter is written against, and
clause 4.1.1.3, which says plainly that it is a complementary measure and does
not substitute for safeguarding & \textbf{Its clause on stop categories sits
beyond the end of the free preview}, so that part is cited by clause number and
attributed. \textbf{Its own normative references point at an edition of the
electrical equipment standard that is now two editions behind} & paid \\
One silicon vendor's free industrial functional safety guide & The distinction
firmware engineers most often miss, in a per-device table: \textbf{random
hardware capability against systematic capability}, and the architecture
arithmetic in a line each. Then the list of mechanisms that are actually
firmware & It is about a different vendor's parts, so the arithmetic transfers
and the device table does not & free, no registration \\
The performance level calculator from an institute & Models the control structure
and computes the attained level, free of charge, version 3.0.4 of Friday 31
October 2025, supporting the 2023 edition & \textbf{Free to use and to pass on,
but modification and re-hosting are not permitted.} Link it, never vendor it &
free, restricted \\
The space agency's software assurance handbook, topic 8.5 & An eight step
analysis method, \textbf{a catalogue of software failure modes that reads as if
written for a bus-connected joint node}, criticality categories, and four
downloadable worksheet templates & Nothing certifiable. \textbf{It is a United
States government work with no copyright notice, which makes it the safest
material in this volume to adapt into a public repository} & public domain \\
A permissively licensed bus stack & A working implementation of the safety
data object wire format: a periodic broadcast of two frames, the second carrying
the data bit-wise inverted on a different identifier & \textbf{It claims coding
standard conformance with documented exceptions, which is not a functional safety
certification}, and it makes no certification claim anywhere & Apache-2.0 \\
A maintained hierarchical state machine library & Entry, run and exit actions,
history, timers, parallel states, a kernel dispatcher, and coding standard checks
in its own continuous integration & C++ where this node is C, which is a real
mismatch to state rather than to paper over & MIT \\
\bottomrule
\end{tabular}
\caption{Prior art for chapter 18. The fifth row is the most useful thing in the table for a reader building a portfolio: a complete, free, quotable analysis method with worksheets, from a source whose material carries no copyright restriction at all.}
\end{table}

\subsection*{What the node gains}

A latched safe state that is reachable from every other state, three independent
ways into it, an entry action that does not depend on the scheduler still
running, and a written statement of what this node is. That last item is the one
that matters most in a portfolio: a node that says clearly that it implements a
safe-state machine and makes no integrity claim is more credible than one that
implies a claim it cannot support.

\subsection*{Parts from the inventory}

\begin{table}[H]\centering\small
\begin{tabular}{p{40mm}p{66mm}p{40mm}}
\toprule
Part & Role & Interface \\
\midrule
X-NUCLEO-53L8A1, 8 by 8 ranging & \textbf{The workspace trigger, and it is
real.} Multizone ranging means a region rather than a point, which is what a
workspace is & The shield header \\
NUCLEO-H7A3ZI-Q & The machine, the latch and the entry action & Its own outputs \\
The independent watchdog & Clocked from the low speed internal oscillator, so it
survives a stopped main clock & On-die \\
The window watchdog & Faults on a refresh that is \textbf{too early} as well as
too late, which catches a loop running at the wrong rate & On-die \\
An external watchdog & \textbf{Not on the bench, and named as absent.} Both
on-die watchdogs share the chip's fate, so a genuinely independent one is
external & n/a \\
A brake, and a motor & \textbf{Neither exists here.} The reaction is asserted on
a pin and logged, and no figure in this chapter implies a mechanism moved & n/a \\
\bottomrule
\end{tabular}
\caption{Inventory items used in chapter 18. Two rows are absences and both are load-bearing. The missing external watchdog limits what may be claimed about diagnostic independence; the missing brake means the second of this chapter's two safe states is designed, tested in logic, and never demonstrated physically.}
\end{table}

\subsection*{System architecture}

\diagram{j18_arch}{The safe-state machine, its three independent triggers and the
latch. Two rules shape it, and neither has a citable source, so both are
presented as engineering practice rather than dressed up as requirements: the
safe state is reachable from every other state, and once entered it is held
until an intentional reset. The published clause list that the reset is written
against is in the figure beside it.}

\subsection*{Peripheral configuration}

\begin{table}[H]\centering\small
\begin{tabular}{p{22mm}p{26mm}p{24mm}p{30mm}p{30mm}}
\toprule
Peripheral & Mode & Clock source & Pins and function & Interrupt and transfers \\
\midrule
Ranging sensor & Multizone, continuous, with its own interrupt & Its own & The
shield's bus pins plus one interrupt line & Interrupt on a new frame, never
polled \\
Independent watchdog & Enabled before anything else, and never disabled & The low
speed internal oscillator & None & Survives a stopped main clock \\
Window watchdog & Refreshed once per control period, inside the window & From
the bus clock & None & \textbf{Faults on an early refresh too} \\
Output stage & \textbf{Break input already proven in chapter 7} & Unchanged &
Unchanged & \textbf{Asserts with interrupts disabled}, which chapter 7 measured \\
Safe-state output & One pin, asserted by the same hardware path & n/a & One pin,
plus its indicator & No software in the assertion path \\
\bottomrule
\end{tabular}
\caption{Peripheral configuration for chapter 18. The fourth row is chapter 7's work being spent: the break input was proven there to remove the outputs with interrupts disabled, which means the entry action into the safe state does not depend on the scheduler, on a task, or on anything continuing to run.}
\end{table}

\subsection*{Wiring}

\diagram{j18_wiring}{What is on the bench, and what a real safety chain has that
this one does not. The trigger, the latch and the output are real and single
channel. A chain that could support an integrity claim is dual channel with
diagnostics, and both the second channel and the external watchdog are drawn as
absent rather than implied. The figure is the honest picture, and it is also the
reason the chapter's own claim is carefully worded.}

\subsection*{Memory and timing budget}

\begin{table}[H]\centering\small
\begin{tabular}{p{44mm}p{28mm}p{28mm}p{30mm}}
\toprule
Quantity & Budget & Measured & Margin \\
\midrule
Flash, this chapter & 7 kB & not measured & not measured \\
Static memory, this chapter & 2 kB, mostly the ranging frame & not measured & not measured \\
Trigger to output asserted & under 2 ms & not measured & not measured \\
Ranging frame period & 15 Hz, which bounds the first term above & not measured & not measured \\
Controlled deceleration, in logic & 200 ms, then torque removed & not measured & not measured \\
Latch, once set & \textbf{forever, until an intentional reset} & n/a & n/a \\
Reset to running & one full self-check, never immediate & not measured & not measured \\
\bottomrule
\end{tabular}
\caption{The budget for chapter 18. The third row is the whole reaction time and most of it is the sensor: a trigger cannot be acted on before it arrives, and a 15 hertz frame rate puts a floor under the reaction that no amount of firmware speed removes. Stating that is more useful than reporting the microseconds the software takes.}
\end{table}

\subsection*{Firmware design (UML)}

\diagram{j18_uml}{Two safe states, and what each one costs to reach. Removing
torque is immediate and persists with no energy, and it is not a safe state at
all for an axis that would descend. Brake held is reached through a controlled
deceleration with power still available, and only then is power removed. The
three columns on the right are how those map onto the published stop categories
and the drive standard's sub-functions.}

Three rules, and the first two have no citable source.

\textbf{The safe state is reachable from every other state.} No state may be a
place from which the safe state cannot be entered. This is real and widely
practised, and this volume could find no published source stating it, so it is
presented as engineering practice with the standards cited for the surrounding
requirements and never attributed to a document nobody has read.

\textbf{A fault latches.} The condition that caused entry is recorded, the state
is held, and neither clears because the condition went away. The same applies:
practised everywhere, named in no source this volume could obtain.

\textbf{The entry action does not depend on software continuing to run.} Chapter
7 proved the break input removes the outputs with interrupts disabled. That
proof is what makes the entry action trustworthy, and a design whose safe state
is entered by a task is a design whose safe state depends on the scheduler.

\subsection*{Data flow (ASCII)}

\begin{asciiart}
  three independent triggers, and none of them can be disabled by the others
  ---------------------------------------------------------------------
        |                      |                       |
  workspace violated     following error         bus loss or bus-off
  (the ranging shield,   outside tolerance       (chapters 12 and 13)
   15 Hz, real)          (chapter 16)
        |                      |                       |
        +----------------------+-----------------------+
                               |
                               v
                      ENTER THE SAFE STATE
                               |
        +----------------------+-----------------------+
        |                                              |
   torque removed                              controlled deceleration
   immediate, and it PERSISTS                  power still available
   with no energy                              then power removed, and
        |                                      THE BRAKE HOLDS
        |                                              |
   for a horizontal or non-backdrivable          for a vertical axis
   joint, this IS the safe state                 carrying a payload,
        |                                        the one above is NOT
        v                                              v
                      LATCHED, and recorded with its cause
                               |
                               v
             reset only by an intentional human action, and
             THE RESET ITSELF DOES NOT START ANYTHING
                               |
                               v
                     a full self-check, then running

  on this bench: the trigger is real, the latch is real, the output pin is
  real, and there is no motor and no brake, so the reaction is asserted and
  logged rather than felt
\end{asciiart}

\subsection*{Repository layout}

\begin{plaincode}
joint-node/
  src/
    safety/
      safe_state.c safe_state.h   # + the machine, the latch, the entry action
      triggers.c                  # + three, independent, none maskable
      workspace.c                 # + the ranging shield as a region, not a point
      selftest.c                  # + what must pass before a reset takes effect
      wdg.c                       # + both on-die watchdogs, and a liveness mask
    control/ sense/ act/ bus/ mw/ estimate/ time/ node/ bsp/  update/
  doc/
    safe-states.md                # + BOTH of them, and which one applies when
    what-this-node-is-not.md      # + the claim, carefully worded
    fmea/                         # + the agency's worksheets, filled in
  test/
    test_safe_state.c             # + reachability from every state, on the host
    test_latch.c                  # + it does not clear when the cause clears
  host/  third_party/  tools/  proto/  README.md
\end{plaincode}

\subsection*{Steps}

\begin{steps}
\item \textbf{Write down which safe state applies, before writing the machine.}
This is a risk assessment output and firmware cannot decide it. Two lines of
honesty here save a design.
\begin{plaincode}
doc/safe-states.md

  horizontal or non-backdrivable axis
      safe state      torque removed
      persists        yes, with no energy at all
      reached         immediately
  vertical axis carrying a payload
      safe state      BRAKE ENGAGED with torque removed
      persists        as long as the brake holds
      reached         controlled deceleration FIRST, then power removed
      why             removing torque alone is not a safe state here: the
                      axis descends
  on this bench       the first, demonstrated. The second, designed and
                      tested in logic only, because there is no brake
\end{plaincode}

\item \textbf{Keep the three axes apart, because most requirements that arrive
have them tangled.} They answer different questions and they come from different
documents.
\begin{plaincode}
axis                what it specifies                        values
stop CATEGORY       what the stopping action does to
                    actuator power                           0, 1, 2
performance LEVEL   how RELIABLY the function is performed   a to e
integrity LEVEL     the same question, other document        1 to 3 (4)

a well formed requirement names ONE from the first and ONE from the others:
  "emergency stop shall be stop category 1, achieved with PL d"

THERE IS NO STOP CATEGORY 3 AND NO PL 1.
A category 0 stop implemented with a single undiagnosed contactor is still
a category 0 stop, at a poor performance level. The two are independent.
\end{plaincode}
\textbf{This volume prints no cross-map between the reliability scales.} Such
tables circulate widely, they are a convenience rather than a normative
equivalence, and no primary verified mapping was obtained during this volume's
research. Printing one would be repeating something rather than knowing it.

\item \textbf{Get the mapping to the drive standard right, because it is the
cleanest single fact in this area.} Stop category 0 corresponds to safe torque
off, category 1 to safe stop 1, and category 2 to safe stop 2. The drive
standard's second edition replaced the phrase safety function by \textbf{safety
sub-function} throughout, which is why drive documentation reads oddly to
somebody who learned the older wording.

\item \textbf{Write the latch against the published clause, line by line.} The
emergency stop standard's clause 4.1.1 gives six requirements and they are worth
copying into the code as comments.
\begin{ccode}
/* safe_state.c: each line below is a requirement, not a preference. */
/* 1. initiated by a single human action                                */
/* 2. available and operational at all times, in every mode             */
/* 3. overrides all other functions without impairing other protections */
/* 4. maintained until manually reset                                   */
/* 5. no start command is effective on the stopped operations           */
/* 6. reset by intentional action, and THE RESET SHALL NOT RESTART      */
void safe_state_enter(trigger_t why)
{
    act_break_assert();            /* chapter 7: with interrupts disabled */
    g_safe.latched = true;
    g_safe.cause   = why;          /* recorded, and it survives the state */
    g_safe.count++;                /* lifetime, and it survives a reset   */
}
\end{ccode}
The same clause adds, in 4.1.1.3, that this is a \textbf{complementary
protective measure} and does not substitute for safeguarding or for other safety
functions. That sentence belongs in the documentation of any node that
implements one, because it is the sentence that stops a stop function from being
treated as the whole safety case.

\item \textbf{Make the workspace sensor a region rather than a point.} An eight
by eight ranging frame is sixty-four distances, and the useful question is not
what any one of them says.
\begin{ccode}
/* workspace.c: a zone policy with hysteresis, and a minimum valid count. */
unsigned close = 0, valid = 0;
for (unsigned z = 0; z < 64; z++) {
    if (frame.status[z] != RANGE_VALID) continue;
    valid++;
    if (frame.mm[z] < g_ws.enter_mm) close++;
}
if (valid < g_ws.min_valid) { workspace_unknown(); return; }   /* not safe */
if (close >= g_ws.enter_zones) workspace_violated();
else if (close == 0 && g_ws.state == VIOLATED &&
         all_beyond(&frame, g_ws.exit_mm)) workspace_clear();  /* hysteresis */
\end{ccode}
\textbf{Too few valid zones is not the same as nothing being there.} A sensor
that cannot see is a sensor whose output must not be read as all clear, and that
distinction is one line of code and the difference between a safety function and
a decoration.

\item \textbf{Prove reachability on the host, from every state.} This is the
rule with no source behind it, and a test is the only thing that keeps it true
as the machine grows.
\begin{ccode}
/* test_safe_state.c: for every state, the safe state must be reachable. */
for (state_t s = 0; s < STATE_COUNT; s++) {
    machine_force_state(s);
    safe_state_enter(TRIGGER_TEST);
    TEST_ASSERT_EQUAL(STATE_SAFE, machine_state());
}
\end{ccode}

\item \textbf{Prove the latch does not clear when the cause clears.} This is the
second unsourced rule and the second one-line test.
\begin{shellcode}
python host/provoke.py --workspace-violate --hold 0.5 --then-clear
# expected: safe state entered at t0, and STILL entered at t0 + 30 s
# expected: cause reads WORKSPACE, not NONE, after the sensor reads clear
\end{shellcode}

\item \textbf{Use both watchdogs, for different reasons, and say what neither of
them proves.} The independent one is clocked from the low speed oscillator and
survives a stopped main clock. The window one faults on a refresh that arrives
too early, which catches a loop running faster than it should rather than one
that has stopped.
\begin{ccode}
/* wdg.c: a liveness mask, because refreshing from a timer proves nothing. */
void wdg_service_from_supervisor(void)
{
    if (g_live.mask == LIVE_ALL_TASKS) {    /* every task checked in */
        wdg_refresh();                      /* inside the window, not early */
        g_live.mask = 0;
    }                                       /* otherwise: let it fault */
}
\end{ccode}
\textbf{Both watchdogs are on-die and share the chip's fate.} A genuinely
independent watchdog for a safety argument is an external component, this bench
does not have one, and the chapter says so rather than letting two on-chip
timers stand in for independence.

\item \textbf{Read the vendor's safety package claim correctly, and do not
repeat the garbled version.} The architecture claim is specific and it is
frequently misquoted.
\begin{plaincode}
what the vendor's page says
  "SIL2 safety functions can be implemented with a single MCU;
   SIL3 safety functions implementation requires two MCUs in a 1oo2 scheme."
what several write-ups say instead
  the same sentence with the redundancy scheme changed, which reverses the
  meaning of the second half
what the package contains
  a safety manual, a qualitative analysis, a static failure rate report,
  and a self test library delivered AS OBJECT CODE
what could NOT be verified for this volume
  the licence text, the distribution mechanism, the current version, and
  WHETHER THIS EXACT PART IS COVERED AT ALL
\end{plaincode}
That last line is why nothing from that package is used here. An object-code
library bound by conditions of use is not something a book can tell a reader to
drop into a repository they intend to publish.

\item \textbf{Be equally careful about the kernel.} The kernel this volume uses
is permissively licensed and its own documentation makes no safety certification
claim. A separate, certified product shares a functional foundation with it and
is a \textbf{completely redesigned product} rather than a certified edition of
the same code. Writing the certified product's name next to the free one, or
calling either a certified version of the other, is the single most common
error in this area.

\item \textbf{Do the failure analysis with free material and a spreadsheet.} The
open tooling for this is immature and none of it is worth the dependency.
\begin{plaincode}
method        the space agency's handbook, topic 8.5: eight steps, and it is
              explicitly based on the cancelled military standard that
              everything else descends from
worksheets    four templates, downloadable, no copyright notice, United
              States government work: adapt them and attribute
failure modes its catalogue reads as if written for this node: out of range
              values, missing inputs, overwritten memory, data collisions,
              command omission, incorrect sequence, illegal commands,
              timing issues, and hardware and software interaction failures
tooling       none. A spreadsheet under version control beats every open
              tool this volume could find, and two of those carry six or
              seven commits in total
\end{plaincode}
One contrast is worth knowing about before somebody argues with you: the
international standard keeps the risk priority number and adds an alternative
beside it, while the automotive sector handbook is widely reported to have
dropped it for a lookup table. \textbf{Two engineers working to different sector
standards are told opposite things about the same number}, and neither is wrong.

\item \textbf{Write down what this node is not.} One page, and it is the
deliverable that makes the rest credible.
\begin{plaincode}
doc/what-this-node-is-not.md

  it IS        a node with a latched safe state, three independent
               triggers, an entry action proven to run with interrupts
               disabled, and a documented reaction time
  it is NOT    a safety-rated device. Single channel, no diagnostic
               coverage figure, no external watchdog, no assessment, no
               certificate, and no integrity level claimed
  the gap      dual channel with fault tolerance 1 reaches the higher
               levels; single channel with fault tolerance 0 does not, and
               that is architecture rather than effort
  and note     random hardware capability and systematic capability are
               different things. Systematic capability is about the
               development process; random hardware capability is about
               diagnostic coverage and failure rates. A part may be rated
               differently on each, and most are
\end{plaincode}
\end{steps}

\diagram{j18_timing}{Two stop categories drawn in time, and the latch that
follows both. Above, power removed immediately, which is a safe state for one
mounting and the opposite of one for another. Below, a controlled deceleration
with power still available, then power removed once stopped, which is what a
gravity-loaded axis needs. The bar at the bottom is the latch: it outlasts its
cause by design, and the reset that ends it starts nothing.}

\subsection*{Build, flash and debug}

\diagram{j18_data}{Three scales that get tangled in almost every requirement
document, kept apart, with what each one comes from and what values it actually
has. Below them, this volume's evidence marking for the clauses in this chapter:
what was read from a primary document, what was attested by agreeing secondary
sources, and what is presented as engineering practice with no source at all.}

\begin{shellcode}
cmake --build build -j && ctest --test-dir build/host
probe-rs run --chip STM32H7A3ZITx build/firmware.elf
python host/provoke.py --workspace-violate --measure-reaction --runs 200
python host/provoke.py --sensor-blind --expect workspace_unknown
\end{shellcode}

\begin{note}{When a stop function looks like a safety case}
It is not one, and the standard says so in a clause of its own: an emergency
stop is a complementary protective measure and does not substitute for
safeguarding or for other safety functions. A node with a good stop function and
no assessment has a good stop function. The temptation in a portfolio is to let
the reader infer more than that, and the correction is a single page saying
plainly what the node is not, which costs nothing and is the most credible thing
in the repository.
\end{note}

\subsection*{Verification and acceptance criteria}

\begin{itemize}
\item Both safe states are written down with the condition under which each
applies, and the bench says which one it can demonstrate.
\item The safe state is reachable from every other state, proven by a host test
that iterates over the state list rather than over a remembered subset.
\item The latch does not clear when its cause clears, proven by a thirty second
hold after the sensor reads clear.
\item The cause is recorded, and a lifetime count survives a reset.
\item The entry action asserts the output with interrupts disabled, which
chapter 7 already measured.
\item Three triggers are independent and none can be masked by the others.
\item A sensor that cannot see produces an unknown state, not a clear one,
proven by covering it.
\item Reaction time from trigger to output is measured over two hundred runs and
reported as median, 99.9th percentile and maximum.
\item The reset requires an intentional action, runs a full self-check, and does
not itself start anything.
\item The window watchdog is refreshed from a liveness mask rather than from a
timer, and a deliberately stalled task causes a fault.
\item The document saying what the node is not exists, and names single channel,
the absent external watchdog and the absent assessment.
\item No cross-map between the reliability scales is printed anywhere in the
repository.
\end{itemize}

\subsection*{Variants}

\begin{table}[H]\centering\small
\begin{tabular}{p{24mm}p{26mm}p{44mm}p{22mm}p{20mm}}
\toprule
Axis & Variant & What changes & Cost & Built in full in \\
\midrule
Safe state & Torque removed & Immediate, and persists with no energy & \textbf{Not a safe state at all for an axis that would descend} & Here, demonstrated \\
Safe state & Brake held after a controlled deceleration & The only correct answer for a gravity-loaded axis & A brake, its release circuit, and a deceleration that must complete & Here, in logic only \\
Trigger & Workspace sensor & A region, with hysteresis and a blind-sensor state & 15 Hz, which is most of the reaction time & Here \\
Trigger & Following error & Free: chapter 16 already computes it & Choosing a tolerance that is neither deaf nor twitchy & Here \\
Trigger & Bus loss & Free: chapter 12 already detects it & A policy decision about how long is too long & Here \\
Watchdog & From a timer & One line, and it proves the scheduler runs & \textbf{It refreshes happily while a task is stalled} & Nowhere \\
Watchdog & From a liveness mask & Every task checks in, or the fault happens & A bit per task and one supervisor & Here \\
Watchdog & External & The only kind that does not share the chip's fate & A component this bench does not have & Nowhere, and named \\
Architecture & Single channel & What this bench is & Reaches the lower integrity levels at best, by architecture & Here \\
Architecture & Dual channel with diagnostics & What an integrity claim needs & A second channel, comparison, and an assessment & Nowhere \\
Bus safety & A safety layer over the bus & Two frames, the second bit-wise inverted, so small parts can participate & \textbf{The published layer is for classic frames, and no equivalent for the flexible-data format was established} & Nowhere, and the gap is named \\
State machine & Hand written, flat & What this node has: small, auditable, testable & It will not scale to a hierarchy & Here \\
State machine & A library & History, parallel states, a kernel dispatcher, coding standard checks & C++ where this node is C & Nowhere, and the mismatch is stated \\
\bottomrule
\end{tabular}
\caption{Variants for chapter 18. The two safe-state rows at the top are the chapter in miniature: the same design decision is correct in one mounting orientation and incorrect in another, and nothing in the firmware can tell which one it is in.}
\end{table}

\subsection*{Pitfalls}

\begin{itemize}
\item Assuming the safe state is power removed. The definition says nothing of
the kind, and for a gravity-loaded axis power removed is the opposite of safe.
\item Having one safe state where a joint needs two.
\item Tangling the stop category with a reliability scale. They answer different
questions and come from different documents.
\item Printing a cross-map between the reliability scales as though it were
normative.
\item Writing stop category 3, or performance level 1. Neither exists.
\item Refreshing a watchdog from a timer, which proves only that the scheduler
runs while a task is stalled.
\item Treating two on-die watchdogs as independent. They share the chip's fate.
\item Reading too few valid ranging zones as all clear.
\item Letting a reset restart the machine. The standard says the reset shall not
initiate a restart, in the same clause that requires the reset.
\item Repeating the garbled version of a vendor's architecture claim, which
circulates widely and reverses the meaning of its second half.
\item Calling a certified kernel product a certified version of the free one.
They share a functional foundation and are separate, redesigned products.
\item Implying an integrity claim by silence. The cure is one page saying what
the node is not.
\item Citing the robot safety series by its old title. It is now titled
Robotics, and getting that wrong dates a document at a glance.
\item Writing that the collaborative robot technical specification is withdrawn.
It is published, was confirmed in 2022, and remains current until its replacement
appears. Several secondary sources say otherwise and they are wrong.
\item Citing the electrical equipment standard or the machinery functional
safety standard without their amendments. A bare citation is incomplete in 2026.
\item Citing the coding guidelines as the 2012 edition with amendments. A
consolidated edition superseded that arrangement, and a further edition is
current as of \pubdate{March 2025}.
\end{itemize}

\subsection*{Best practices applied}

\begin{itemize}
\item A definition is read from its own document rather than from what everybody
knows it says.
\item A design decision that firmware cannot make is identified as such and sent
back to the risk assessment.
\item Two practices with no citable source are named as engineering practice
rather than attributed to documents nobody has read.
\item Each claim carries its evidence class: read directly, attested by agreeing
sources, or practice.
\item An entry action is made independent of the scheduler and that independence
is proven, not assumed.
\item A sensor's inability to see is a distinct state from its seeing nothing.
\item A vendor claim is quoted exactly, and the widely circulating garbled
version is named.
\item What the node is not is written down, which is what makes what it is
believable.
\item Free, unrestricted public material is preferred over paid material
wherever it is adequate, and the analysis method chosen here is both.
\end{itemize}

\subsection*{Stretch goals}

\begin{itemize}
\item Add an external watchdog and measure what it catches that the on-die pair
does not. It is a small component and it changes what may be claimed.
\item Implement the two-frame safety layer over the bus with the inverted second
frame, then write down honestly what it does and does not give a node with no
assessment behind it.
\item Fill in the agency's worksheets completely for this node and publish them.
Complete, honest analyses of small systems are rare in public and cost nothing
but care.
\item Model this node's single channel in the free performance level calculator
and find out precisely where the architecture stops, rather than asserting it.
\item Build the second safe state for real: a brake, its release circuit, and a
deceleration that completes before power is removed. It is the half of this
chapter the bench could not demonstrate.
\end{itemize}

\subsection*{Roadmap and next steps}

Chapter 19 takes the node out of reach. Once a joint is inside a machine, the
only way to change its firmware is over the wire it already has, and a
bootloader that stays addressable is a different problem from one that runs on a
desk.

The published progression from here starts with the free material, because it is
better than its price suggests: the space agency's assurance handbook for
analysis method and worksheets, one silicon vendor's industrial functional safety
guide for the architecture arithmetic and the mechanism list, and the free
performance level calculator for modelling a structure. After that, the type-C
standard for robotics is the one that takes precedence for a joint, and the
emergency stop standard is short, cheap and the most directly useful of the paid
ones.

\subsection*{Portfolio evidence}

\begin{itemize}
\item The two safe states with the condition for each, which demonstrates
understanding the difference between an implementation and a risk assessment
output.
\item The document saying what the node is not.
\item The reachability and latch tests, which are short and make an unsourced
rule enforceable.
\item The filled-in failure analysis worksheets, from material that may be
published freely.
\item The evidence table, showing which claims were read from primary documents
and which were not.
\end{itemize}

\subsection*{Sources}

Normative references, with the correct current editions:
\begin{itemize}
\item IEC 61508-4:2010, edition 2.0, Friday 30 April 2010, "Functional safety of
electrical/electronic/programmable electronic safety-related systems, Part 4:
Definitions and abbreviations". Clause 3.1.13 for the safe state definition and
clause 3.2.11 for the tool classification, \textbf{both read from the official
preview}.
\item ISO 13850:2015, third edition, "Safety of machinery, Emergency stop
function, Principles for design", 11 pages. \textbf{Clause 4.1.1 read directly};
clause 4.1.3 on stop categories lies beyond the free preview and is cited by
clause number. \textbf{Its own normative references point at IEC 60204-1:2005},
which is two editions behind.
\item IEC 60204-1:2016 + AMD1:2021, edition 6.0, "Safety of machinery,
Electrical equipment of machines, Part 1: General requirements", clause 9.2.2
for the stop categories. \textbf{The category wording in this chapter is attested
by agreeing secondary sources, not primary verified.}
\item ISO 13849-1:2023, fourth edition, \pubdate{April 2023}, for performance levels. It
covers software design and \textbf{applies to high demand and continuous modes
only}.
\item IEC 62061:2021 + AMD1:2024, edition 2.0. \textbf{The 2021 edition did not
gain non-electrical coverage, it lost the restriction to electrical}: the older
title named electrical, electronic and programmable electronic systems and the
current one does not.
\item ISO 10218-1:2025 and ISO 10218-2:2025, \pubdate{February 2025}. \textbf{The series is
now titled "Robotics"}, not the older wording. For a joint this is the type-C
standard and \textbf{where a type-C standard differs, its provisions take
precedence}.
\item ISO/TS 15066:2016 on collaborative robots. \textbf{It is not withdrawn}:
published, confirmed in 2022, and current until its replacement appears.
\item IEC 61800-5-2:2016, edition 2.0, for the drive sub-functions. Edition 2
replaced safety function by \textbf{safety sub-function} throughout.
\textbf{Two brake-related sub-functions could not be confirmed on the page read,
and they are exactly what a gravity-loaded joint needs.}
\item IEC 60812:2018, edition 3.0, Friday 10 August 2018, for failure modes and
effects analysis. \textbf{Annex B.4 keeps the risk priority number} and adds an
alternative method beside it.
\item The coding guidelines: the \pubdate{March 2025} edition is current, the 2023 edition
consolidated the 2012 edition with its amendments and corrigendum, and the
compliance framework is mandatory from the 2023 edition onward.
\end{itemize}

Free and reusable material:
\begin{itemize}
\item The space agency's software engineering and assurance handbook, topic 8.5,
with four worksheet templates. \textbf{A United States government work with no
copyright notice: the safest material in this volume to adapt and publish.}
Attribute it.
\item The cancelled military standard that the method descends from, a United
States government work and quotable freely.
\item One silicon vendor's free industrial functional safety guide, for the
distinction between random hardware and systematic capability, the architecture
arithmetic, and the list of mechanisms implemented in firmware.
\item The free performance level calculator from an institute, version 3.0.4 of
Friday 31 October 2025. \textbf{Free to use and pass on; modification and
re-hosting are not permitted.}
\item A permissively licensed bus stack implementing the safety data object wire
format, Apache-2.0, which claims coding standard conformance with documented
exceptions and \textbf{makes no certification claim}.\newline
\url{https://github.com/CANopenNode/CANopenNode}
\item Stolte and co-authors, "A Taxonomy to Unify Fault Tolerance Regimes for
Automotive Systems", IEEE Transactions on Intelligent Vehicles volume 7 number
2, pages 251 to 262, 2022, which documents that the literature is inconsistent
and then repairs it. Cite and link; do not bundle.
\end{itemize}

Two gaps are recorded rather than glossed over. \textbf{The published safety
communication layer for this family of buses is written for the classic frame
format}, and this volume's research established no equivalent for the
flexible-data format. And the general framework for a safety layer over an
untrusted transport, in which the layer detects corruption, loss, repetition,
resequencing, delay and masquerade, is exactly the right frame for a host to
joint link; \textbf{the standard family's number, title and edition could not be
verified and are therefore not printed here.}
   % Safe states, and the workspace sensor that triggers one
\project{19}{Update over the bus: a node you can reach but not touch}{Maintainability}{Transport for update, a bootloader that stays addressable, versioning across a fleet}

\begin{keyfacts}
\item[Adds to the node] Maintainability. A joint inside a machine cannot be
unplugged, so the only way to change its firmware is the wire it already has,
and the only way to know what is running on it is to ask
\item[Peripherals] The bus controller, the flash interface with its error
correction, the voltage detector, and one RAM location the startup code must
leave alone
\item[Depends on] Chapter 4 for the flash geometry, chapter 9 for the bus,
chapter 11 for the object model, chapter 18 for what must be safe before an
update begins
\item[Real or modelled] \textbf{Real, and it starts with a correction:} the
ROM bootloader in this part already speaks the bus, so the chapter opens by
finding half the job done and asking what the other half is worth
\item[Difficulty] 5 of 5
\item[Effort] Four evenings, and the licence decision in the third one is the
most consequential in the volume
\item[Deliverable] A bootloader that is always reachable and never runs a bad
image, a power-down path that survives the supply falling mid-write, a fleet
sweep that reports what every node is running, and a licence decision written
down before any code was built around it
\end{keyfacts}

\subsection*{Why this chapter}

Eighteen chapters have built a node that would be bolted inside a machine. At
that point the debug probe is on the wrong side of a housing, and everything
this volume has built is only as maintainable as its worst update path.

The chapter opens with a correction, because the research turned up something
better than expected.

\begin{note}{The bootloader in the silicon already speaks this bus}
The part's own ROM bootloader exposes, besides the usual serial and USB routes,
\textbf{the flexible-data bus on two specific pins, at 250 kbit/s nominal and
1000 kbit/s data, with bit rate switching, and its clock fixed at 20 MHz}. The
vendor's application note gives it in a table, and its detection chain figure
shows \textbf{the bus is polled first}. So this board can be reflashed over the
same wire the joint uses, today, with no bootloader of anybody's writing. That
is the chapter's opening and it is also the baseline: a custom bootloader has to
be worth more than something that is already in the silicon and costs nothing.
One caution: the note's revision matters, a newer revision is indexed but was
not read, and pin numbers are exactly the kind of thing that moves between
revisions.
\end{note}

What the ROM bootloader does not do is report a version, refuse a corrupt image,
survive a supply falling mid-write, or tell a fleet master what twelve joints
are running. That is what the rest of the chapter is for.

\subsection*{Prior art and what to reuse}

\begin{table}[H]\centering\small
\begin{tabular}{p{34mm}p{46mm}p{38mm}p{20mm}}
\toprule
Source & What it gives & What it does not & Licence \\
\midrule
The vendor's system bootloader note, revision 61, \pubdate{January 2024} & \textbf{The
correction that opens this chapter}: the peripheral list for this exact part,
the two pins, both bit rates, and the detection order. Also that the independent
watchdog is refreshed for you in system memory boot mode & Version reporting,
image validation, and any behaviour a fleet needs. \textbf{A newer revision
exists and was not read} & vendor \\
An integrator's bootloader note, read in full & The object indices for program
data, program control, software identification and flash state, and a captured
trace of a real session. Plus a list of requirements this chapter follows rather
than invents & Its own loader implements no address-assignment server, so
addressing is a design choice rather than something inherited & vendor \\
A conference paper on bootloader security, read in full & The requirement list
again, from a second independent pen, plus the handoff-keyword argument and the
reasoning for polling rather than interrupts & \textbf{It disagrees with the note
above about block transfer}, and the chapter shows both rather than picking
quietly & proceedings \\
The portable open bootloader, Apache-2.0, 2122 stars, last pushed Friday 18
September 2026 & \textbf{Port agnostic by its own porting guide}: a port supplies
a flash map, an area interface and a configuration header. Primary and secondary
slots, three upgrade strategies, a 32 byte header and a trailer whose magic
encodes the write alignment & \textbf{It has no bus transport, and neither does
its management protocol}, whose transports are Bluetooth, console-framed serial
and raw serial only. \textbf{The bus half is the gap this chapter fills} &
Apache-2.0 \\
The closest off-the-shelf match, 973 stars & Exactly this problem solved: serial,
the bus in both frame formats, network, USB and more, on this family & \textbf{It
is copyleft or paid commercial, so linking the joint firmware against it makes
the whole firmware copyleft.} This is the most consequential licence decision in
the volume and it is stated here rather than discovered later & GPL-3.0 \\
The kernel's own segmented transport & \textbf{Better than every open
implementation for the host side}, and it states plainly that it works on both
frame formats. Its third socket option group is what lets a master drive a
sixty-four byte link & Nothing on the node. \textbf{The claim about which kernel
version introduced it could not be confirmed}, so no version is printed & GPL-2.0 \\
Two update architecture documents from the standards body, \pubdate{April 2021} and \pubdate{January
2022} & An architecture and a manifest information model, \textbf{both
Informational and both free to redistribute, which makes them the safest
normative references in this entire volume} & Neither is specific to a bus or to
a part & free \\
\bottomrule
\end{tabular}
\caption{Prior art for chapter 19. The fifth row is why this table exists: an excellent, maintained project solves this exact problem and its licence would follow the joint node's firmware wherever it went. Knowing that before writing code around it is worth more than any amount of cleverness afterwards.}
\end{table}

\subsection*{What the node gains}

A node that can be reached but not touched. Specifically: a loader that always
executes first and is always addressable, an image that is never executed
without its checksum matching, a configuration store that survives the supply
falling in the middle of a write, and a fleet sweep that answers what is running
where without anybody opening a cabinet.

\subsection*{Parts from the inventory}

\begin{table}[H]\centering\small
\begin{tabular}{p{40mm}p{66mm}p{40mm}}
\toprule
Part & Role & Interface \\
\midrule
NUCLEO-H7A3ZI-Q & The node being updated, and the part whose flash geometry is
\textbf{not its better known sibling's} & The bus \\
Raspberry Pi 4 & The update master and the fleet sweep, using the kernel's own
segmented transport & The bus \\
The ROM bootloader & \textbf{Already present, already speaking this bus}, and the
baseline the custom loader is measured against & Two specific pins \\
The flash error correction & \textbf{Works in this chapter's favour}: an
interrupted word programming reads back as an uncorrectable error, so a torn
record is detectable rather than silently wrong & On-die \\
The voltage detector & Seven thresholds plus an external-input selection, and the
edge that starts the power-down write & Shares one interrupt line with the
analog supply detector \\
\bottomrule
\end{tabular}
\caption{Inventory items used in chapter 19. Nothing new is bought. The third row is the chapter's opening correction and the fourth is the finding that makes a torn write detectable, which is the single most useful property this part has for the job.}
\end{table}

\subsection*{System architecture}

\diagram{j19_arch}{Two update paths and what each one is worth. The upper path
is already in the silicon and costs nothing to use. The lower path is the
chapter, and it exists for the four things the upper one does not do: report a
version, refuse a corrupt image, stay addressable on the node's own identifier,
and answer a fleet sweep. The requirements beside it come from two primary
documents read in full, not from this volume's preferences.}

\subsection*{Peripheral configuration}

\begin{table}[H]\centering\small
\begin{tabular}{p{22mm}p{26mm}p{24mm}p{30mm}p{30mm}}
\toprule
Peripheral & Mode & Clock source & Pins and function & Interrupt and transfers \\
\midrule
Bus controller, in the loader & \textbf{Polled, never interrupt driven} & Fixed
bit rate, or a small tested set & Unchanged & \textbf{Hardware filtering admits
two identifiers}: the request and the network management one \\
Flash interface & Word programming only in the power-down path & Unchanged &
None & Error correction callbacks enabled \\
Voltage detector & Threshold set \textbf{below nominal and above the flash
programming minimum} & Unchanged & Optionally an external analog input & Falling
edge, and \textbf{it shares one interrupt line with the analog detector} \\
The vector table offset & \textbf{Programmable on this part}, so the loader
points it at the application before jumping & n/a & n/a & Which removes the need
for vector mirroring entirely \\
One RAM word & \textbf{Excluded from startup initialisation} & n/a & n/a & The
handoff keyword, in both directions \\
\bottomrule
\end{tabular}
\caption{Peripheral configuration for chapter 19. The fourth row is a real simplification that comes from this part rather than from cleverness: designs on parts with a fixed vector location have to mirror the table and route every vector through a decider, and this part's programmable offset register makes all of that unnecessary.}
\end{table}

\subsection*{Wiring}

\diagram{j19_wiring}{The flash, drawn to its real geometry, which is not the
geometry of the part most of the internet writes about. A sixteen byte word and
an eight kilobyte sector, against that part's thirty-two byte word and one
hundred and twenty-eight kilobyte sector. The consequence is at the bottom: an
eight kilobyte sector holds five hundred and twelve records of one word each, and
that number is the whole argument for a rotating log.}

Nothing is wired in this chapter. The figure is a memory map because the memory
map is the thing this chapter is about, and because the difference between this
part's geometry and its sibling's has already produced at least two reported
defects in other people's code.

\subsection*{Memory and timing budget}

\begin{table}[H]\centering\small
\begin{tabular}{p{44mm}p{28mm}p{28mm}p{30mm}}
\toprule
Quantity & Budget & Measured & Margin \\
\midrule
Flash, the loader & 24 kB, and it sits in its own protected sectors & not measured & not measured \\
Static memory, the loader & 4 kB, and it is a superloop rather than a task set & not measured & not measured \\
Application slot & 1024 kB, one whole bank & n/a & n/a \\
Configuration sectors & 3 sectors of 8 kB, rotating & n/a & n/a \\
Bounded wait before starting the application & 300 ms & not measured & not measured \\
Checksum over the application & under 40 ms & not measured & not measured \\
Power-down write & \textbf{one 16 byte word, and never an erase} & not measured & not measured \\
Records per sector before an erase & \textbf{512} & arithmetic & n/a \\
\bottomrule
\end{tabular}
\caption{The budget for chapter 19. The fifth row is a cost paid at every single power-up in exchange for always being reachable, and it is stated as a trade rather than hidden. The last two rows are the chapter's most important arithmetic and they are done in full in the data figure.}
\end{table}

\subsection*{Firmware design (UML)}

\diagram{j19_uml}{The architectural break, which is the part of this chapter
most likely to surprise somebody. The application is a task set on a kernel. The
loader is a polled superloop with interrupts off, and that is not laziness: it is
required by three separate constraints, all named in the figure. Below it, the
handoff in both directions through one RAM word that the startup code must not
initialise, and why a direct jump without that word is the source of defects
that appear only in the field.}

Four rules, all of them taken from primary documents rather than invented.

\textbf{The loader always executes first and always checks.} The reset vector
points at it at all times, and it computes the application's checksum on every
start and never executes on a mismatch.

\textbf{No interrupts. Poll.} Three reasons and any one of them is enough:
interrupts remove the determinism a loader needs, most parts cannot execute
handlers during flash programming anyway, and the hardware acceptance filter is
set to admit exactly two identifiers so there is nothing to be interrupted
about.

\textbf{Every received block is bounds checked against the application flash
boundaries}, out of range writes are rejected with an abort, and the loader's own
sectors are erase and program protected where the hardware allows.

\textbf{The handoff travels in one RAM word that the startup code does not
initialise.} A direct jump leaves peripherals configured by the loader, and an
application developed and tested without a loader present can fail for subtle
reasons that are very difficult to track down, often surfacing only in the
field. That sentence is from a primary source and it is the reason the keyword
exists.

\subsection*{Data flow (ASCII)}

\begin{asciiart}
  reset
    |
    v
  the loader, ALWAYS, with interrupts off and two identifiers accepted
    |
    +-- is the handoff keyword set in the uninitialised RAM word?
    |        yes --> stay in the loader, clear the keyword
    |
    +-- wait a BOUNDED period for one specific request on the bus
    |        something arrived --> stay in the loader
    |
    +-- compute the application checksum
    |        mismatch --> stay in the loader, report through the error object
    |
    v
  set the application's vector base (this part has a programmable offset
  register, so no vector mirroring is needed)
    |
    v
  the application: a task set, interrupts on, everything chapters 1 to 18 built
    |
    +-- a request to update arrives --> set the keyword --> reset
    |
    +-- the supply sags
             |
             v
       the detector's falling edge, one interrupt line shared with the
       analog detector, so read both status bits to tell them apart
             |
             v
       program ONE 16 BYTE WORD into a pre-erased sector
       AN ERASE NEVER FITS IN THIS WINDOW: an 8 kB sector erase runs into
       milliseconds, and the window does not
             |
             v
       if the supply fails mid-word, that word reads back as an
       UNCORRECTABLE ERROR, which is how a torn record is detected rather
       than silently believed
\end{asciiart}

\subsection*{Repository layout}

\begin{plaincode}
joint-node/
  boot/                            # + this chapter: a separate build entirely
    main.c                         # + a superloop, not a task set
    bus_poll.c                     # + two identifiers, no interrupts
    flash16.c                      # + 16 byte words, 8 kB sectors, THIS part
    image.c                        # + checksum, bounds, and the refusal
    handoff.h                      # + the RAM word, in its own linker section
    boot.ld                        # + its own sectors, protected
  src/
    update/
      update_server.c              # + the object entries the master writes to
      version.c                    # + what a fleet sweep reads back
    store/
      record.c                     # + one 16 byte record, rotating
      powerfail.c                  # + the detector edge, and one word
    safety/ control/ bus/ mw/ ...
  host/
    flash_node.py                  # + segmented transfer over the kernel socket
    fleet_sweep.py                 # + every identifier, what each is running
  doc/
    bootloader-licence.md          # + the decision, and why, dated
    flash-geometry.md              # + this part, not its sibling
  test/  tools/  proto/  README.md
\end{plaincode}

\subsection*{Steps}

\begin{steps}
\item \textbf{Use the ROM bootloader first, and write down what it does not
do.} An hour spent here produces the baseline the rest of the chapter is judged
against.
\begin{shellcode}
# boot the part from system memory, then talk to it over the bus
python host/rom_flash.py --node can0 --bitrate 250000 --data-bitrate 1000000 \
       --image build/firmware.bin
# it works. Now write down what it did not do:
#   no version reported, no image validation, no node identifier of its own,
#   no answer to a fleet sweep, and its bit rates are fixed
\end{shellcode}

\item \textbf{Get this part's flash geometry right, from its own headers.} This
is where ported code fails, and the failure is silent data corruption rather
than a compile error.
\begin{plaincode}
this part            word 128 bits = 16 bytes, sector 8 kB,
                     128 sectors a bank, two banks
its popular sibling  word 256 bits = 32 bytes, sector 128 kB
the vendor macro     counts 32-bit words in a flash word: 4 here, 8 there
CAUTION              the comment block in the vendor's flash driver source
                     describing the wider word IS WRITTEN FOR THE SIBLING
                     and is wrong for this part. Do not quote it.
evidence             two closed defect reports against a widely used project:
                     one from Wednesday 11 May 2022 where a hardcoded 32 byte
                     stride corrupted data on this exact part, and one from
                     Wednesday 15 October 2025 on dual bank sector mapping
\end{plaincode}

\item \textbf{Let the error correction do the hard part.} Each bank reports
single and double errors with a failing address, and the design consequence is
better than anything software could arrange.
\begin{ccode}
/* record.c: a torn write is DETECTABLE, which is rare and worth using. */
/* An interrupted flash word reads back as an uncorrectable double error.   */
/* Corollary: bits inside an already written word cannot be reprogrammed,   */
/* SO THE VALIDITY FLAG MUST LIVE IN ITS OWN WORD, not inside the record.   */
typedef struct { uint8_t payload[16]; } rec_t;       /* one word exactly */
typedef struct { uint32_t valid_magic; uint8_t pad[12]; } rec_flag_t;
\end{ccode}

\item \textbf{Do the endurance arithmetic before choosing a storage scheme,
because it decides the answer by three orders of magnitude.} The endurance figure
itself could not be confirmed, so the arithmetic is parameterised and the
conclusion does not depend on the number.
\begin{plaincode}
one record            one 16 byte flash word
one 8 kB sector       512 records before an erase is needed
rotating log          survives 512 x (the per sector endurance figure)
read, erase, rewrite  survives (the per sector endurance figure) alone

at a commonly quoted but UNVERIFIED 10,000 cycles:
  rotating log        5,120,000 events
  rewrite in place       10,000 events
at 100 power cycles a day that is roughly 140 years against 100 days.

AND: on this part the rotating scheme is unusually cheap, because the erase
granularity is SIXTEEN TIMES SMALLER than on the sibling part.
\end{plaincode}

\item \textbf{Build the power-down path around one word, and never an erase.}
The window is what it is, and the design has to fit inside it rather than hope.
\begin{ccode}
/* powerfail.c: the two detectors share one interrupt line. Read both bits. */
void EXTI_VoltageDetectors_Handler(void)
{
    bool core_low  = pwr_status_core_below_threshold();
    bool analog_low = pwr_status_analog_below_threshold();
    if (core_low) {
        store_program_one_word(&g_pending);   /* 16 bytes, pre-erased sector */
    }                                         /* NEVER an erase: it is ms   */
    (void) analog_low;                        /* logged, not acted on here  */
}
\end{ccode}
The sector is pre-erased during normal operation and a clean one is kept ready,
so the power-down path only ever programs. An eight kilobyte erase runs into
milliseconds and the holdup window does not, and that is not a tuning problem.

\item \textbf{Make the loader a polled superloop, and explain the break.} This
is an architectural discontinuity in a volume that has spent eighteen chapters on
a task set, and it deserves a paragraph rather than a shrug.
\begin{ccode}
/* boot/main.c: no scheduler, no interrupts, two identifiers, one loop. */
for (;;) {
    if (bus_poll_frame(&f)) handle_request(&f);   /* SDO request, or NMT */
    if (timeout_expired() && image_checksum_ok()) jump_to_application();
}
\end{ccode}

\item \textbf{Implement the object entries the update protocol defines, with
their verified names.} Four entries, and one of them is commonly written down
with a longer name that the primary documents do not use.
\begin{plaincode}
1F50h   Program data                    (not "download program data")
1F51h   Program control  00h stop, 01h start, 02h reset, 03h clear
1F56h   Program software identification
1F57h   Flash status identification

and for the fleet sweep, from a captured trace in a primary document:
1000h   device type
1008h   manufacturer device name        value "Boot" identifies loader mode
1018h   sub 1 to 4: vendor, product code, revision, serial number
\end{plaincode}
A caution that belongs beside this table: \textbf{the main open stack for this
protocol does not define these entries}. Its own list of supported
specifications omits the part that defines them, and its object dictionary
enumeration stops well below this range. Two other stacks have the same gap and
only one claims coverage of that part, in the words "portions of", without
itemising.

\item \textbf{Choose between segmented and block transfer, and show the reader
that the sources disagree.} This is the most interesting single page in the
chapter, because both sides are credible.
\begin{plaincode}
the integrator's note  moves the image by BLOCK download, and its captured
                       trace shows forty-five blocks
the conference paper   argues AGAINST block transfer in a loader: it needs
                       back to back frame buffering in a polled driver, adds
                       RAM and significant code space, and segmented access
                       has to be implemented anyway, so segmented transfer
                       "should be the preferred choice"
this volume            follows the paper, because the loader is polled and
                       the argument turns on exactly that. The note is not
                       wrong; it is written for a different set of
                       constraints, and saying so is more useful than
                       picking one quietly
\end{plaincode}

\item \textbf{Carry the handoff in a RAM word the startup code leaves alone, in
both directions.} The reason is a sentence from a primary source and it is worth
quoting in the repository.
\begin{ccode}
/* handoff.h: its own linker section, excluded from startup initialisation. */
#define HANDOFF_STAY_IN_LOADER   0xB007C0DEu
extern volatile uint32_t g_handoff __attribute__((section(".noinit")));
/* Why not just jump? Because a direct jump leaves peripherals initialised  */
/* by the loader, and an application developed and tested WITHOUT a loader  */
/* present "can fail for subtle reasons that can be very difficult to track */
/* down", usually in the field rather than on the bench.                    */
\end{ccode}

\item \textbf{Decide where forced entry lives, and follow the published
caution.} Writing to the program control entry is the standard way in, and the
primary source warns about implementing that entry in the application.
\begin{plaincode}
the standard way in    the master writes 00h to program control, sub 1
the published caution  implementing that entry INSIDE THE APPLICATION would
                       let any tool on the bus stop the application, and the
                       source says plainly it is not recommended for
                       security reasons
the recommendation     a physical input pin, for the case that matters: an
                       application whose checksum passes but which does not
                       communicate
this volume            implements the entry in the LOADER only, and uses a
                       pin for the stuck-but-valid case
\end{plaincode}

\item \textbf{Report errors through the predefined error object rather than
inventing a channel.} A code in the low sixteen bits and free information in the
upper sixteen, with the assignment a primary source suggests.
\begin{plaincode}
6100h   loader checksum wrong
6200h   application checksum wrong
6300h   stored configuration checksum wrong
\end{plaincode}

\item \textbf{Write the fleet sweep, which is the half of this chapter that
justifies the other half.} Twelve joints, one command, and an answer that fits
on a screen.
\begin{shellcode}
python host/fleet_sweep.py --iface can0
# node  vendor   product  rev   serial     version        mode
#    1  0x0247   0x0001   3     00A31C07   1.4.2          application
#    2  0x0247   0x0001   3     00A31C11   1.4.2          application
#    7  0x0247   0x0001   3     00A31C2B   1.3.9          application  <- old
#   11  0x0247   0x0001   3     00A31C44   -              Boot         <- stuck
\end{shellcode}
The last two lines are the reason this exists. Nobody opens a cabinet to find
out which joint is behind, and nobody guesses which one failed an update.

\item \textbf{Version it properly, and cite documents a reader may actually
redistribute.} The versioning specification is freely quotable, and the two
architecture documents from the standards body are Informational and free to
redistribute, which makes them the safest normative references anywhere in this
volume.
\end{steps}

\diagram{j19_data}{The object entries with their verified names, the endurance
arithmetic that decides the storage scheme by three orders of magnitude, and the
licence position of every loader this chapter considered. The middle block is
the one to read twice: its input could not be confirmed, so it is parameterised,
and the conclusion survives whatever the real figure turns out to be.}

\subsection*{Build, flash and debug}

\diagram{j19_timing}{Two sequences that decide whether this node is
maintainable. Above, every power-up: the loader runs, waits a bounded period,
checks, points the vector base at the application and jumps. That bounded wait is
paid at every start in exchange for always being reachable. Below, the
power-down path: one detector edge, one sixteen byte word programmed into a
sector that was erased long ago, and the reason an erase can never appear on that
line.}

\begin{shellcode}
cmake --build build -j --target bootloader firmware
python host/flash_node.py --iface can0 --node 7 --image build/firmware.bin
python host/fleet_sweep.py --iface can0
python host/provoke.py --brownout-during-write --runs 500 --expect-detected
\end{shellcode}

\begin{note}{One package with three different licence statements}
This chapter's research found a vendor package whose release notes name a
permissive three clause licence and point at a vendor licence page, whose
individual source files carry headers pointing at the canonical permissive text,
\textbf{and whose archive also ships a licence agreement document stating that
the licensee may not sell, assign, sublicense, lease, rent or otherwise
distribute the software commercially}. Three statements, one archive. Nothing
from it is used here. The general lesson is worth more than the specific case:
a licence is what the files and the archive say together, not what the
prettiest of those statements says.
\end{note}

\subsection*{Verification and acceptance criteria}

\begin{itemize}
\item The ROM bootloader route is exercised once and what it does not do is
written down.
\item The flash driver uses this part's word and sector sizes, taken from its own
headers, and a host test rejects the sibling's values.
\item An interrupted word programming is detected as an uncorrectable error
rather than read back as data, proven over five hundred provoked brown-outs.
\item The validity flag lives in a separate flash word from the record it
validates.
\item The endurance arithmetic is in the repository, parameterised, with the
unverified figure marked as unverified.
\item The power-down path programs exactly one word and never attempts an erase.
\item The loader runs with interrupts disabled and an acceptance filter admitting
two identifiers.
\item The loader computes the application checksum on every start and refuses to
jump on a mismatch, proven by corrupting one byte.
\item Every received block is bounds checked, and an out of range write produces
an abort rather than a write.
\item The handoff keyword survives a reset, is cleared by the loader, and lives
in a section the startup code does not initialise.
\item The vector base is set before the jump, proven by taking an interrupt in
the application immediately afterwards.
\item The program control entry exists in the loader and \textbf{not} in the
application.
\item A fleet sweep over twelve identifiers reports vendor, product, revision,
serial, version and mode for each, and identifies a node stuck in loader mode.
\item The bootloader licence decision is in a dated file, written before any code
was built around a stack.
\end{itemize}

\subsection*{Variants}

\begin{table}[H]\centering\small
\begin{tabular}{p{24mm}p{26mm}p{44mm}p{22mm}p{20mm}}
\toprule
Axis & Variant & What changes & Cost & Built in full in \\
\midrule
Loader & The ROM one, already present & Free, in the silicon, and it speaks this bus & No version, no validation, no identifier, fixed bit rates & Here, as the baseline \\
Loader & The off-the-shelf match & Solves this exact problem, on this family, in both frame formats & \textbf{Copyleft or paid: linking makes the joint firmware copyleft} & Nowhere, and the reason is printed \\
Loader & The portable open one & Permissive, port agnostic, slots, upgrade strategies, a defined header & \textbf{No bus transport at all}, so the transport is the author's & Nowhere, and the gap is the chapter \\
Loader & Written here & Exactly the requirements two primary documents state & Four evenings, and it must beat something that is free & Here \\
Transfer & Block download & What a vendor note's captured trace does & Back to back buffering in a polled driver, plus RAM and code space & Nowhere, and the disagreement is shown \\
Transfer & Segmented & What a conference paper recommends for a polled loader & More frames for the same image & Here \\
Transfer & The automotive service set & A different vocabulary for the same job, with a published programming process clause & A different stack on both ends & Nowhere, and named \\
Storage & Rewrite one structure & Simplest & \textbf{Three orders of magnitude fewer cycles} & Nowhere \\
Storage & Rotating records & 512 per sector on this part, which is the argument & Slightly more code to find the newest & Here \\
Storage & A log structured library & Wear levelling and garbage collection, already written & Footprint, and most have no published figures & Nowhere, and the table names which one does \\
Addressing & Fixed node identifier & Simplest, and what a bench needs & Two identical nodes cannot share a bus & Here \\
Addressing & Layer setting services & The standard route for an unconfigured node & A second protocol in the loader & Nowhere, and it is named as a design choice rather than a given \\
\bottomrule
\end{tabular}
\caption{Variants for chapter 19. The second row is the most consequential refusal in the volume. An excellent maintained project solves this problem completely, and its licence would travel with the joint node's firmware into every product that used it.}
\end{table}

\subsection*{Pitfalls}

\begin{itemize}
\item Writing a bootloader without first finding out that the part already has
one that speaks this bus.
\item Using the sibling part's flash word or sector size. The consequence is
silent data corruption, and it has already been reported twice against a widely
used project.
\item Quoting the comment block in the vendor's flash driver about the flash
word width. It is written for the sibling part.
\item Putting a validity flag inside the record it validates. Bits in an already
written word cannot be reprogrammed.
\item Attempting an erase in the power-down path. A sector erase runs into
milliseconds and the window does not.
\item Assuming the two voltage detectors have separate interrupt lines. They
share one, and the handler disambiguates by reading the status bits.
\item Rewriting one structure in place, and discovering the endurance figure
three orders of magnitude later than you wanted to.
\item Running the loader with interrupts enabled.
\item Jumping to an application without setting the vector base first.
\item Jumping without a handoff keyword and then spending a season on defects
that appear only in the field.
\item Implementing the program control entry in the application, which lets any
tool on the bus stop it.
\item Assuming the main open stack for this protocol implements the update
entries. It does not, its own specification list omits the relevant part, and two
other stacks have the same gap.
\item Linking against a copyleft loader and discovering the obligation after the
product ships.
\item Citing the older title of the vendor's emulation note. It was renamed, and
the renaming is recorded in its own revision history.
\item Assuming the vendor's emulation package covers this family. It does not,
and the string identifying this family does not appear anywhere in the current
revision of its note.
\end{itemize}

\subsection*{Best practices applied}

\begin{itemize}
\item The existing capability is found and used before anything is written.
\item Geometry is taken from the part's own headers, and a comment written for a
different part is identified as such.
\item A hardware property is used rather than worked around: error correction
makes a torn write detectable.
\item An arithmetic argument decides a design, and it is parameterised because
one of its inputs could not be confirmed.
\item Two credible sources that disagree are both presented, with a reason for
the choice made.
\item Requirements come from primary documents read in full, and are marked as
such rather than presented as this volume's preferences.
\item The most consequential licence decision is made and dated before code is
built around it.
\item A package with three conflicting licence statements is named as a lesson
rather than quietly avoided.
\item References are chosen partly for whether a reader may redistribute them.
\end{itemize}

\subsection*{Stretch goals}

\begin{itemize}
\item Port the permissive open bootloader to this part and write the bus
transport it lacks. That transport does not exist, the interface it needs is
documented, and the result would be genuinely useful to other people.
\item Measure the actual holdup window on this board and find out how many
sixteen byte words really fit, rather than designing for one and hoping.
\item Implement the address assignment services so that two identical joints can
share a bus out of the box.
\item Sign the image and verify the signature in the loader, which is the half
the free bootloader article explicitly does not cover.
\item Run the same update over the automotive service set as well, and compare
the two vocabularies on frame count, code size and how much of each is already
implemented by something on the host.
\end{itemize}

\subsection*{Roadmap and next steps}

Chapter 20 is the last one and it is the one that makes the other nineteen
believable: a rig that injects faults, measures tracking error and turns a build
red when something regresses, with nobody at the bench.

The published progression from here is the two architecture documents from the
standards body, because they are free, redistributable and short; then the
configuration and program download specification for the object entries, noting
that it is a members-only document while the base specification is not; and then
the automotive application layer standard's programming process clause, which is
the model a fleet master follows whatever protocol it speaks underneath.

\subsection*{Portfolio evidence}

\begin{itemize}
\item The opening correction: finding that the part already does half the job,
and saying so instead of building it again.
\item The flash geometry document, which is short and prevents a defect that has
already occurred twice in public.
\item The endurance arithmetic, parameterised, with its unverified input marked.
\item The two-sources-disagree page, which demonstrates reading rather than
searching.
\item The licence decision file, dated, with the alternative named and costed.
\item The fleet sweep output, which is the kind of thing a maintenance engineer
recognises immediately.
\end{itemize}

\subsection*{Sources}

Normative references:
\begin{itemize}
\item The vendor's system bootloader application note, revision 61, \pubdate{January
2024}, section and table read at a fetchable mirror, for this part's peripheral
list, both bit rates, the two pins and the detection order. \textbf{A revision
70 dated \pubdate{February 2026} is indexed only and was not read: check the revision
before relying on pin numbers.}
\item The configuration and program download specification, and the base
application layer specification version 4.2.0, for the object entries.
\textbf{Licence position: the base specification is freely available after
registration, and the part that defines the update entries is a members-only
document.}
\item ISO 14229-1:2020, third edition, \pubdate{February 2020}, "Road vehicles, Unified
diagnostic services (UDS), Part 1: Application layer". Clause headings confirmed
verbatim, including clause 15 "Upload download functional unit" and
\textbf{clause 17 "Non-volatile server memory programming process"}, which is
the model a fleet master follows.
\item The vendor's emulation note, \textbf{current title "How to use EEPROM
emulation on STM32 MCUs", revision 11, Tuesday 18 March 2025}, renamed from an
earlier title that is still widely cited. \textbf{This part's family does not
appear in it.}
\end{itemize}

Free and redistributable references, which are the safest in this volume:
\begin{itemize}
\item RFC 9019, "A Firmware Update Architecture for Internet of Things",
Informational, \pubdate{April 2021}.\newline
\url{https://www.rfc-editor.org/rfc/rfc9019}
\item RFC 9124, "A Manifest Information Model for Firmware Updates in Internet
of Things (IoT) Devices", Informational, \pubdate{January 2022}.\newline
\url{https://www.rfc-editor.org/rfc/rfc9124}
\item Semantic Versioning 2.0.0, CC BY 3.0 and freely quotable.\newline
\url{https://semver.org}
\end{itemize}

Reusable implementations:
\begin{itemize}
\item The portable open bootloader, Apache-2.0, port agnostic by its own porting
guide. \textbf{Its current source asserts a write alignment range that includes
this part's sixteen byte word, which contradicts an older claim that it cannot
run on this family. No published footprint figures were found, so none are
quoted.}\newline
\url{https://github.com/mcu-tools/mcuboot}
\item The closest off-the-shelf match, \textbf{GPL-3.0 or paid commercial}, which
supports this family and both frame formats. Named, costed, and not used.\newline
\url{https://github.com/feaser/openblt}
\item A power-loss-resilient filesystem, BSD-3-Clause, and a key-value store,
Apache-2.0, which is the only entry in this chapter's comparison with published
per-object footprints.\newline
\url{https://github.com/littlefs-project/littlefs}
\item A vendor emulation utility, BSD-3-Clause end to end, created Friday 23
January 2026, whose release notes list a different series. It is family agnostic
by design, so the flash driver for this part is the author's work, which is a
chapter rather than an obstacle.\newline
\url{https://github.com/STMicroelectronics/stm32-util-eeprom-emulation}
\end{itemize}

One certification note is worth recording because it is the kind of thing nobody
writes down. The author of the conference paper cited above certified bootloader
code to a sector scheme requiring conformance with a coding standard, and
records that the 1998 rules against casting to and from pointers, and against two
common loop control statements, made generic table driven object dictionary
implementations impractical and forced a fresh coding effort. A coding standard
is not a formatting preference, and it can change what a design is allowed to be.
   % Update over the bus: a node you can reach but not touch
\project{20}{The rig: injected faults, tracking error, and a build that fails}{Proof}{Fault injection, tracking regression, a report that turns a build red}

\begin{keyfacts}
\item[Adds to the node] Proof. Nineteen chapters made claims; this one builds the
thing that checks them every night and turns a build red when one stops being
true, with nobody at the bench
\item[Peripherals] None new. The rig uses what is already there, and the
observability it needs is built from inside the part because there is no probe
\item[Depends on] Chapter 2 for the period, chapter 16 for following error,
chapter 18 for provoked faults, chapter 19 for flashing without hands
\item[Real or modelled] \textbf{Both, deliberately.} Half the rig runs in a
simulator on every commit with no hardware at all, and half runs on the bench
once a night. The chapter says which assertions each half can honestly support
\item[Difficulty] 5 of 5
\item[Effort] Five evenings, and the first one is spent finding out that a
simulator for this family already exists
\item[Deliverable] A two-stage rig, an assertion set whose names are sourced, a
fault catalogue, a report a stranger can read, and a deliberate regression that
turns the build red
\end{keyfacts}

\subsection*{Why this chapter}

Every chapter in this volume has ended with acceptance criteria. Criteria that
are checked once are a memory; criteria that are checked every night are a
property. This chapter builds the thing that does the checking, and it exists
because a portfolio whose claims were true in September is worth less than one
whose claims are true tonight.

It opens with the most useful single finding of the volume's research.

\begin{note}{A simulator for this family exists and it carries the bus}
An open simulator models this microcontroller family \textbf{including the bus
controller}, and its host integration provides a bridge that carries both frame
formats between the simulation and a host virtual interface. Its test runner
starts it headless, speaks a test language natively, emits reports in two
formats, runs jobs in parallel and snapshots failures. \textbf{So the whole bus
half of this volume can be exercised on a build machine with no hardware
attached at all.} Two honest limits: the bridge is available on one host
operating system only, and the simulator provides \textbf{no motion metrics of
any kind}, so everything this chapter asserts about tracking is computed by the
rig rather than read from a tool.
\end{note}

The obvious alternative is the wrong tool and the chapter says so plainly. The
other well known emulator's own target page lists the cores and boards it
supports for this architecture, \textbf{none of them from this family}, and does
not mention the bus at all.

\subsection*{Prior art and what to reuse}

\begin{table}[H]\centering\small
\begin{tabular}{p{34mm}p{46mm}p{38mm}p{20mm}}
\toprule
Source & What it gives & What it does not & Licence \\
\midrule
The open simulator & \textbf{This family modelled including the bus controller},
a bridge to a host virtual interface carrying both frame formats, a headless test
runner with parallel jobs and failure snapshots & \textbf{No motion metrics at
all}, and its documentation index has no fault injection section. The bridge is
for one host operating system only & permissive \\
The test language it speaks natively & Structure, reporting and a server the
simulator already talks to & \textbf{Structure and reporting, not measurement.}
Every number this chapter asserts on is computed elsewhere & Apache-2.0 \\
The bench control layer, 528 stars, last pushed Friday 11 September 2026 & The
right layer for the hardware half: serial and remote drivers, power switch and
reset drivers, bootstrap, and a remote abstraction & \textbf{Nothing at all about
trajectories.} It gets a board powered, flashed and talking & LGPL-2.1 \\
A test runner with a bus library and a serial library & \textbf{The honest
recommendation}, and more defensible than an alternative whose value is
concentrated in another vendor's silicon and most of which would be
reimplemented here & The bus library is copyleft, so it is depended on rather
than vendored & mixed \\
A unit test build system with its assertion and mock libraries & The unit half:
the interpolator, the controller step, the frame packer, all on the host & It
does not touch hardware, which is the point & MIT \\
A header-only fake framework, 938 stars & \textbf{Unit level fault injection}: a
return sequence or a custom fake makes a driver call fail on the third
invocation and only the third & Nothing at system level & MIT \\
The mock hardware component in the host control framework & \textbf{The closest
published analogue to this whole rig}, and the reason the architecture below is
not unusual & It is a host-side stand-in, not a rig & Apache-2.0 \\
\bottomrule
\end{tabular}
\caption{Prior art for chapter 20. The first row is the finding that shapes the chapter: the expensive half of a rig, a target that behaves like the real part on the real bus, already exists and is permissively licensed. What does not exist anywhere is the motion measurement, and that is what this chapter writes.}
\end{table}

\subsection*{What the node gains}

Nothing it runs. What the \textbf{volume} gains is the difference between a
claim and a property. Every acceptance criterion from chapters 2 to 19 becomes
an assertion that is checked without anybody remembering to check it, and a
regression in the quality of the node's motion becomes a red build rather than a
discovery.

\subsection*{Parts from the inventory}

\begin{table}[H]\centering\small
\begin{tabular}{p{40mm}p{66mm}p{40mm}}
\toprule
Part & Role & Interface \\
\midrule
A build machine & The simulated half: every commit, no hardware, both frame
formats over a virtual interface & None \\
Raspberry Pi 4 & \textbf{Two roles at once}: the motion master of chapter 10 and
the runner for the hardware half & The bus, and the node's serial port \\
NUCLEO-H7A3ZI-Q & The node under test, flashed by chapter 19's path rather than
by hand & The bus \\
A switchable supply & For the power cycle tests, and the only piece the rig
would like to buy & Mains, and one control line \\
\textbf{Absent:} an oscilloscope, a logic analyser, an external debug probe &
\textbf{None of these are on this bench}, which is why the observability in this
chapter is built from inside the part & n/a \\
\bottomrule
\end{tabular}
\caption{Inventory items used in chapter 20. The last row has been true since chapter 1 and it shapes this chapter more than any other: with no instrument outside the part, everything the rig knows about a failure has to be something the part itself recorded and sent.}
\end{table}

\subsection*{System architecture}

\diagram{j20_arch}{Two stages, and what each one can honestly prove. The upper
stage runs on every commit with no hardware: the simulated part, the virtual bus,
the host stack, and every assertion that does not depend on real timing. The
lower stage runs once a night on the bench and carries the assertions that only
real silicon can support. Splitting them is what makes the first stage fast
enough to run on every commit and the second one honest enough to be worth
running at all.}

\subsection*{Peripheral configuration}

\begin{table}[H]\centering\small
\begin{tabular}{p{22mm}p{26mm}p{24mm}p{30mm}p{30mm}}
\toprule
Peripheral & Mode & Clock source & Pins and function & Interrupt and transfers \\
\midrule
Everything & \textbf{Unchanged}. This chapter configures nothing & n/a & n/a & The rig observes; it does not modify the node under test \\
The serial port & The log channel, because there is no probe & Unchanged & Already wired & The rig reads it and archives it per run \\
A region of memory excluded from startup initialisation & The crash trace, with a
magic number and a checksum & n/a & n/a & \textbf{It is the only thing that
survives a fault and a reset} \\
\bottomrule
\end{tabular}
\caption{Peripheral configuration for chapter 20. A rig that modifies the thing it is testing is measuring something else, so the only entries here are channels the node already had. The third row is the one that makes a crash investigable at all on a bench with no probe.}
\end{table}

\subsection*{Wiring}

\diagram{j20_wiring}{The bench as a rig, and the observability stack that exists
because there is nothing outside the part to look with. No oscilloscope, no
logic analyser, no external probe. What replaces them is a fault diagnosis
library that needs no debugger, a small logger with a flash backend, a crash
region that survives a reset, and pre-aggregated counters in a periodic frame.
Every one of those is permissively licensed, which is not an accident.}

\subsection*{Memory and timing budget}

\begin{table}[H]\centering\small
\begin{tabular}{p{44mm}p{28mm}p{28mm}p{30mm}}
\toprule
Quantity & Budget & Measured & Margin \\
\midrule
Flash, the observability stack & 9 kB & not measured & not measured \\
Static memory, the crash region & 1 kB, plus a magic number and a checksum & n/a & n/a \\
Simulated stage, wall clock & under 4 minutes, on every commit & not measured & not measured \\
Hardware stage, wall clock & under 25 minutes, once a night & not measured & not measured \\
Injected faults per nightly run & 40, from a written catalogue & n/a & n/a \\
Tracking regression that fails the build & \textbf{5 per cent} on root mean square error & n/a & n/a \\
\bottomrule
\end{tabular}
\caption{The budget for chapter 20, which is mostly wall clock rather than bytes. The two stage times are the reason the rig is split: a four minute stage can run on every commit and a twenty-five minute one cannot, and pretending otherwise produces a rig that people learn to skip.}
\end{table}

\subsection*{Firmware design (UML)}

\diagram{j20_uml}{The assertion set with its provenance marked, because three of
these six have published names and one has none at all. A published performance
standard for manipulators names stabilisation time, overshoot and path accuracy,
and this chapter uses its words. \textbf{Loop period jitter has no standards
backing in either document this chapter read}, so it is presented as an
engineering assertion, which is exactly what it is.}

Three rules.

\textbf{Every assertion has a provenance.} A name taken from a published
standard is marked as such; a name that is this volume's own is marked as that.
The distinction costs one column in a table and it is the difference between a
rig and a set of opinions.

\textbf{The fault catalogue is written before the injector.} What is injected is
decided from the failures that actually occur, not from the failures that are
easy to inject. That question has a name in the literature and a paper devoted
to it, and this chapter's catalogue is honest about which of its entries are
representative and which are convenient.

\textbf{Nothing in the rig modifies the node under test.} No test hooks, no
special build, no instrumentation that ships only during testing. The rig
observes the channels the node already has, because a node that behaves
differently under test has not been tested.

\subsection*{Data flow (ASCII)}

\begin{asciiart}
  a commit
     |
     v
  STAGE 1, every commit, no hardware, under 4 minutes
     |
     +-- host unit tests: interpolator, controller step, frame packer
     |       with fakes that fail on the third call and only the third
     |
     +-- the simulated part, with its bus controller, bridged to a
     |       virtual interface on the build machine
     |
     +-- the host stack, unchanged, talking to it as if it were a board
     |
     +-- assertions: protocol, state machines, bounds, safe states,
     |       and the update path of chapter 19
     |
     v
  STAGE 2, once a night, on the bench, under 25 minutes
     |
     +-- power cycle, flash over the bus (chapter 19), wait for boot
     |
     +-- run the trajectory set, 30 cycles, five configurations
     |
     +-- inject 40 faults from the written catalogue
     |
     +-- collect: the node's own counters, its log, and its crash
     |       region if anything faulted
     |
     v
  THE REPORT
     |
     +-- six numbers against six thresholds, with the delta from the
     |   baseline beside each
     |
     +-- and if root mean square tracking error rose by more than five
         per cent:  THE BUILD IS RED, with nobody at the bench
\end{asciiart}

\subsection*{Repository layout}

\begin{plaincode}
joint-node/
  rig/
    stage1/
      sim.resc  sim.robot      # + the simulated part, and its assertions
      vcan.sh                  # + the virtual interface the bridge attaches to
    stage2/
      conftest.py              # + power, flash, boot, and teardown
      test_tracking.py         # + the six assertions, with provenance
      test_faults.py           # + the catalogue, one test per entry
      faults.yaml              # + WRITTEN FIRST, and marked representative
                               #   or convenient, one or the other
    report/
      render.py                # + the report a stranger can read
      baseline.json            # + what tonight is compared against
  src/
    diag/
      backtrace.c              # + fault cause and call stack, no debugger
      logger.c                 # + small, flash backed, no file system
      crashregion.c            # + noinit, magic number, checksum
      metrics.c                # + counters, timed counters, gauges
  .github/  doc/  test/  host/  tools/  proto/  README.md
\end{plaincode}

\subsection*{Steps}

\begin{steps}
\item \textbf{Build stage one first, because it needs no hardware and it is the
one that will run thousands of times.} The simulated part, the bridge, and the
host stack from chapter 10, unchanged.
\begin{shellcode}
sudo ip link add dev vcan0 type vcan && sudo ip link set up vcan0
renode-test rig/stage1/sim.robot -j 4
# the simulated part boots, its bus controller bridges to vcan0, and the
# host stack talks to it exactly as it talks to the board
\end{shellcode}
Everything in this volume that is about protocol rather than about physics
belongs here: the frame layouts of chapter 11, the state machines of chapter 12,
the mixed traffic of chapter 13, the message mapping of chapter 15, the safe
state transitions of chapter 18 and the update path of chapter 19.

\item \textbf{Say clearly what stage one cannot prove.} A simulated part does
not have real interrupt latency, real flash timing or a real clock, and a rig
that quietly asserts those things in simulation is a rig that reports good news.
\begin{plaincode}
stage 1 CAN assert    protocol correctness, state reachability, bounds
                      checking, message mapping, error handling paths,
                      the update sequence, and anything a fake can fail
stage 1 CANNOT assert loop period jitter, following error, interrupt
                      latency, flash programming time, the reaction time
                      of chapter 18, or anything with a microsecond in it
and the simulator      provides NO MOTION METRICS AT ALL, so even in stage 2
itself                 every number in this chapter is computed by the rig
\end{plaincode}

\item \textbf{Take the assertion names from a published standard where one
exists.} The performance standard for manipulators, second edition, published
Thursday 23 April 1998, already names most of what this rig measures.
\begin{plaincode}
from the standard   pose accuracy and repeatability
                    POSITION STABILISATION TIME
                    POSITION OVERSHOOT
                    PATH ACCURACY and PATH REPEATABILITY
                    cornering deviations, path velocity characteristics,
                    and drift of pose characteristics
its own method      a test cube, the largest fitting the workspace, with
                    measurements at five configurations over 30 CYCLES
its own caveat      five configurations is thin, and manufacturers commonly
                    use a hundred or more
NOT printed here    the metric symbols, which could not be confirmed from a
                    fetchable page and are therefore not presented as
                    verified
\end{plaincode}

\item \textbf{Fix the six assertions, and mark which one has no standard behind
it.} This is the table the whole rig is built around.
\begin{plaincode}
assertion                          threshold        name from
max absolute following error       < 8 mrad         chapter 16, and path
                                                    accuracy in the standard
root mean square tracking error    < 2 mrad         path accuracy
settling time                      < 60 ms          POSITION STABILISATION
                                                    TIME, in the standard
overshoot                          < 3 per cent     POSITION OVERSHOOT
steady state error                 < 1 mrad         pose accuracy
control loop period jitter         99.9th < 40 us   NO STANDARDS BACKING.
                                                    An engineering assertion,
                                                    and chapter 2's number
\end{plaincode}
The last row matters more than it looks. Everything else here can be defended
by pointing at a document; that one is defended by pointing at chapter 2's
measurement method, and saying so is what keeps the other five credible.

\item \textbf{Use circles as well as steps, and use two radii.} The machine tool
test code for circular tests, published Tuesday 22 February 2022, and its free
companion give the reason.
\begin{plaincode}
large radius circles   expose geometry errors
small radius circles   are more sensitive to SERVO MISMATCH OR LAG
                       which is exactly what chapter 16's controller is
so the rig runs both   and reports circular deviation as a single overall
                       indicator for each
\end{plaincode}

\item \textbf{Write the fault catalogue before writing the injector.} The
question of whether injected faults resemble real ones has a name in the
literature and a paper devoted to it, and pretending otherwise is the easiest
way to build a rig that passes.
\begin{plaincode}
rig/stage2/faults.yaml   (40 entries, each marked)

  bus, frame corrupted        representative  (seen on a real bus)
  bus, node removed           representative  (chapter 12 provoked it)
  bus, old node present       representative  (chapter 13, on purpose)
  bus, flood at line rate     convenient      (easy, and rarer than this
                                               weighting implies)
  supply, brown-out mid-write representative  (chapter 19's whole point)
  sensor, blinded             representative  (chapter 18 covered it)
  sensor, stuck value         representative
  driver call fails 3rd time  convenient      (a unit level fake, and the
                                               third call is arbitrary)
  task stalled                representative  (chapter 18's liveness mask)
\end{plaincode}
Weighting a suite towards faults that are easy to inject is how a rig ends up
proving that the node survives the things that were never going to happen.

\item \textbf{Inject at three levels, because they catch different things.}
\begin{plaincode}
unit level     a header-only fake framework: make a driver call fail on the
               third invocation and only the third, in a host test
bus level      replay a captured real trace onto the virtual interface with
               the interface remapped, and generate load and malformed
               frames with the command line tools
physical level on the bench: power cycles, a disconnected terminator, a
               covered sensor, and the old-node adapter of chapter 13
\end{plaincode}
Two honest notes belong here. \textbf{No dedicated fault injection subsystem was
found in the other major operating system} for this class of part: mocking yes,
a named framework no. And \textbf{fault injection support in the simulator could
not be confirmed}, so stage one injects at the unit and bus levels and not
inside the simulated silicon.

\item \textbf{Build the observability stack from inside the part, because there
is nothing outside it.} Four pieces, all permissively licensed, and the licence
position is the reason each one was chosen.
\begin{plaincode}
fault diagnosis  a library covering this core, printing a diagnosed cause
                 and a call stack, NEEDING NO DEBUGGER, with addresses
                 resolved offline afterwards. MIT
logging          a small logger, ROM under 1.6 kB and RAM under 0.3 kB,
                 with a flash backend and no file system. MIT
crash survival   a region excluded from startup initialisation, with a
                 magic number and a checksum, so a trace outlives a fault
counters         pre-aggregated periodic statistics rather than raw logs,
                 which is the argument behind chapter 11's diagnostic frame
\end{plaincode}

\item \textbf{Name what could not be used, and why, because two of them are
instructive.} One widely used trace recorder's target-side source turns out to be
effectively a one clause permissive licence, so redistribution is not the
obstacle people assume. \textbf{The obstacle is hardware}: it needs a particular
vendor's probe, this bench does not have one, and reflashing the on-board debug
circuit with that vendor's firmware is a probe by another name. The instrumented
trace port is out for exactly the same reason, which chapter 2 established.

\item \textbf{Turn a regression into a red build.} This is the step the chapter
is named for, and it is short.
\begin{shellcode}
python rig/report/render.py --run tonight.json --baseline baseline.json
# following error, max      6.9 mrad   (< 8.0)    +0.2 vs baseline   pass
# tracking error, r.m.s.    1.7 mrad   (< 2.0)    +0.4 vs baseline   pass
# settling time              48 ms     (< 60)      -1  vs baseline   pass
# overshoot                 2.1 %      (< 3.0)     +0.1 vs baseline  pass
# steady state error        0.6 mrad   (< 1.0)     0.0 vs baseline   pass
# loop jitter, 99.9th        31 us     (< 40)      +1  vs baseline   pass
# faults injected: 40   defined outcome: 40   hangs: 0
exit 0
\end{shellcode}

\item \textbf{Then break it on purpose, and watch it go red.} A rig that has
never failed has not been tested either.
\begin{shellcode}
git revert --no-commit HEAD~1      # put back the unfiltered derivative
python rig/report/render.py --run tonight.json --baseline baseline.json
# tracking error, r.m.s.    2.4 mrad   (< 2.0)    +1.1 vs baseline   FAIL
# REGRESSION: root mean square tracking error rose 61 per cent
exit 1
\end{shellcode}

\item \textbf{Make the report readable by somebody who was not there.} Six
numbers, six thresholds, six deltas, the fault tally, and a link to the archived
log and crash region for any run that faulted. Nothing else.
\end{steps}

\diagram{j20_timing}{One night, end to end, and the trace that fails it. Above,
the nightly run laid out to scale: power cycle, flash over the bus, the
trajectory set, the fault sweep, and the report. Below, root mean square tracking
error over sixty nights, with the baseline, the threshold, and the evening
somebody reverted a filtered derivative. The delta column in the report is what
makes the three nights before that one visible.}

\subsection*{Build, flash and debug}

\diagram{j20_data}{The report, and the fault catalogue behind it. Above, what a
passing night looks like and what a failing one looks like, with the delta
column that makes a slow drift visible before it crosses a threshold. Below, the
catalogue, with every entry marked representative or convenient. That second
column is uncomfortable to fill in honestly and it is the most valuable thing in
the figure.}

\begin{shellcode}
renode-test rig/stage1/sim.robot -j 4                       # every commit
pytest rig/stage2 --junitxml=stage2.xml                     # nightly
python rig/report/render.py --run tonight.json --baseline baseline.json
\end{shellcode}

\begin{note}{When the rig has never failed}
A suite that has always passed is not evidence that the node is good; it is
evidence that nothing has been checked. Two habits fix it: revert a real
improvement on purpose and confirm the build goes red, and review the fault
catalogue for entries marked convenient rather than representative. A rig
weighted towards faults that are easy to inject will prove that the node
survives the things that were never going to happen to it.
\end{note}

\subsection*{Verification and acceptance criteria}

\begin{itemize}
\item Stage one runs on every commit, with no hardware, in under four minutes.
\item What stage one cannot prove is written down, and no assertion about
microsecond timing appears in it.
\item Stage two runs nightly, power cycles the board, flashes it over the bus
rather than by hand, and tears down cleanly whether it passed or failed.
\item Six assertions run, each with its threshold and its provenance, and the one
with no standards backing is marked as an engineering assertion.
\item Both a large and a small radius circular test run, and their deviations are
reported separately.
\item The fault catalogue exists as a file, was written before the injector, and
marks every entry representative or convenient.
\item Forty injected faults all end in a defined state and none in a hang.
\item A crash leaves a trace in the region excluded from startup initialisation,
and the rig archives it with the run.
\item A deliberate five per cent regression in root mean square tracking error
turns the build red with nobody at the bench.
\item The rig modifies nothing in the node under test: no test hooks, no special
build.
\item The report is readable by somebody who was not there, and includes the
delta from the baseline for every number.
\end{itemize}

\subsection*{Variants}

\begin{table}[H]\centering\small
\begin{tabular}{p{24mm}p{26mm}p{44mm}p{22mm}p{20mm}}
\toprule
Axis & Variant & What changes & Cost & Built in full in \\
\midrule
Target & The simulated part & Every commit, no hardware, both frame formats & \textbf{No motion metrics, no real timing}, and the bridge is one host operating system only & Here, stage one \\
Target & The board on the bench & Everything stage one cannot prove & Wall clock, and a bench that has to stay plugged in & Here, stage two \\
Target & The other emulator & The obvious alternative & \textbf{It supports no board in this family and does not mention the bus} & Nowhere, and the reason is printed \\
Injection & Unit level fakes & A driver call that fails on the third invocation and only the third & It proves nothing about the wire & Here \\
Injection & Bus level & Replay of a captured real trace, plus generated load and malformed frames & A virtual interface and a capture worth replaying & Here \\
Injection & Physical & Power cycles, a missing terminator, a covered sensor, the old-node adapter & Somebody has to have set the bench up & Here \\
Injection & Inside the simulated silicon & Instruction and memory level faults, as an academic tool does & \textbf{Support could not be confirmed in this simulator}, and the academic tool is dormant and copyleft & Nowhere \\
Harness & The bench control layer & Power, flash and console, done properly & Copyleft, and nothing about trajectories & Nowhere, and named \\
Harness & A test runner with a bus library & \textbf{The honest recommendation}, and it is what this rig uses & The bus library is copyleft, so depended on rather than vendored & Here \\
Observability & An external probe & Everything, instantly & \textbf{This bench has none}, and the trace recorder that needs one is blocked by hardware rather than by licence & Nowhere \\
Observability & From inside the part & A diagnosed fault cause, a call stack, a surviving crash region and pre-aggregated counters & A few kilobytes, and some care & Here \\
\bottomrule
\end{tabular}
\caption{Variants for chapter 20, and the last variants table in the volume. Three rows are refusals, and in two of them the obstacle is not the one people expect: an emulator that does not support this family at all, and a trace tool whose licence is fine and whose hardware requirement is not.}
\end{table}

\subsection*{Pitfalls}

\begin{itemize}
\item Building the hardware stage first. The simulated stage is cheaper, runs
more often and catches more, and it needs nothing that has to be plugged in.
\item Asserting microsecond timing in simulation.
\item Assuming the simulator measures motion. It does not, and every number in
this chapter is computed by the rig.
\item Reaching for the better known emulator. It supports no board in this
family and does not mention this bus.
\item Writing the injector before the catalogue, which produces a suite weighted
towards whatever was easy to inject.
\item Leaving every catalogue entry unmarked, so that nobody can tell which
faults are representative.
\item Adding test hooks to the firmware. A node that behaves differently under
test has not been tested.
\item Reporting a mean. The whole volume reports median, 99.9th percentile and
maximum, including here.
\item Presenting loop period jitter as though a standard named it. No document
this chapter read does.
\item Printing the performance standard's metric symbols as verified. They could
not be confirmed from a fetchable source.
\item Assuming a trace tool is unavailable for licence reasons when the actual
obstacle is a probe this bench does not have.
\item Letting the rig pass forever. Revert an improvement on purpose and check
that it goes red.
\end{itemize}

\subsection*{Best practices applied}

\begin{itemize}
\item The expensive half of the problem is found already solved, and used.
\item A rig is split by what each half can honestly prove rather than by what is
convenient to run.
\item Assertion names are taken from a published standard where one exists, and
the one that has none says so.
\item A fault catalogue is written before the injector and is honest about
representativeness, which is a question with a literature of its own.
\item Observability is built from inside the part because there is nothing
outside it, and every piece chosen is permissively licensed.
\item A refusal is explained with the real obstacle rather than the assumed one.
\item The rig is proven by breaking something on purpose.
\item A report is written for somebody who was not there.
\end{itemize}

\subsection*{Stretch goals}

\begin{itemize}
\item \textbf{Publish the motion-metric layer.} This volume's research found no
published, motion-specific example of a rig that runs in continuous integration
and asserts on tracking quality. The practitioner material that exists is
generic, and a worked example with real numbers would be a genuine contribution.
\item Replay a captured trace from a real machine's bus onto the virtual
interface and find out how much of the fault catalogue it reproduces without
anybody inventing anything.
\item Run the performance standard's method properly: a hundred cycles rather
than thirty, and more than five configurations, and see which of the six
assertions move.
\item Add the other operating system's coredump, which is permissively licensed
and whose host tools include a server that serves a captured dump to a debugger,
and compare it with the probe-free stack this chapter built.
\item Extend the fault catalogue from the failure analysis worksheets of chapter
18, so that the two documents inform each other rather than existing separately.
\end{itemize}

\subsection*{Roadmap and next steps}

This is the last chapter, so the roadmap is not another chapter.

What exists at the end of twenty chapters is a joint node that keeps a one
kilohertz period and can prove it, shares one clock across its sensors and its
bus, reads a real encoder and a real inertial unit, models the actuator it does
not have and says so on every figure, speaks a modern field bus properly enough
to survive a device older than itself, joins a robot framework as a participant,
publishes state in standard types with an empty field where it has no sensor,
closes a loop and reports the one number a joint is judged by, stops safely and
latches, can be updated over the wire it already has, and is checked every night
by a rig that turns a build red.

What it is not is a product. It has one axis, no motor, no brake, no safety
assessment and no certificate, and it has said so in every chapter where the
question arose.

The published progression from here, for a reader who wants to go further, is
the three directions this volume deliberately stopped at. The deterministic
fieldbus of chapter 17, which needs different silicon and a costed decision
rather than more effort. The functional safety route of chapter 18, which needs
a second channel and an assessment rather than better firmware. And the motion
side, where the machine tool test code and the manipulator performance standard
between them define far more than this bench can measure, and where the gap this
chapter named is still open.

\subsection*{Portfolio evidence}

\begin{itemize}
\item The two-stage rig, which demonstrates knowing what a simulator can and
cannot prove.
\item The assertion table with its provenance column, including the row that
admits it has no standard behind it.
\item The fault catalogue with its representativeness marking.
\item The probe-free observability stack, built because there was nothing else,
and every piece permissively licensed.
\item A screenshot of the build going red because somebody reverted a filtered
derivative, which is the single most persuasive artefact in the whole volume.
\end{itemize}

\subsection*{Sources}

Normative references:
\begin{itemize}
\item ISO 9283:1998, second edition, published Thursday 23 April 1998,
"Manipulating industrial robots, Performance criteria and related test methods".
Source of the names position stabilisation time, position overshoot, path
accuracy and path repeatability, and of the thirty cycle method with its own
caveat about five configurations. \textbf{Its metric symbols could not be
confirmed from a fetchable page and are not printed here.}
\item ISO 230-4:2022, published Tuesday 22 February 2022, "Test code for machine
tools, Part 4: Circular tests for numerically controlled machine tools",
superseding the 2005 and 1996 editions.
\item ISO 26262-6:2018 and ISO 26262-11:2018, both published Monday 17 December
2018. Fault injection appears among the methods for software unit verification,
recommended at the lower integrity levels and highly recommended at the upper
ones. \textbf{The table number could not be confirmed and sources disagree, so no
table number is printed}, and this attribution is to secondary sources rather
than to the standard.
\end{itemize}

Academic references, all verified:
\begin{itemize}
\item Hsueh, Tsai and Iyer, "Fault injection techniques and tools", Computer
volume 30, pages 75 to 82, 1997. The classic survey, and the first thing to read.
\item Natella, Cotroneo and Madeira, "Assessing Dependability with Software Fault
Injection", ACM Computing Surveys volume 48, 2016. The modern survey.
\item Natella and co-authors, "On Fault Representativeness of Software Fault
Injection", IEEE Transactions on Software Engineering volume 39, pages 80 to 96,
2013. \textbf{Directly relevant to this chapter's catalogue, and the reason every
entry in it is marked.}
\item Isermann, Schaffnit and Sinsel, "Hardware-in-the-loop simulation for the
design and testing of engine-control systems", Control Engineering Practice
volume 7, pages 643 to 653, 1999. The canonical citation for this kind of rig.
\item Matinnejad, Nejati, Briand and Bruckmann, "Test Generation and Test
Prioritization for Simulink Models with Dynamic Behavior", IEEE Transactions on
Software Engineering volume 45, pages 919 to 944, 2019, with the earlier
conference paper of 2016. \textbf{The closest published work to failing a build on
a regression in continuous output quality rather than in code.}
\end{itemize}

Reusable implementations:
\begin{itemize}
\item The open simulator, which models this family including its bus controller
and bridges both frame formats to a host virtual interface.\newline
\url{https://github.com/renode/renode}
\item The command line bus tools, whose replay utility remaps interfaces, which
is how a captured real trace is replayed onto a virtual bus.\newline
\url{https://github.com/linux-can/can-utils}
\item A fault diagnosis library covering this core, MIT, 2176 stars, which needs
no debugger and can persist a trace across a restart.\newline
\url{https://github.com/armink/CmBacktrace}
\item A small logger with a flash backend and no file system requirement,
MIT.\newline
\url{https://github.com/armink/EasyLogger}
\item A header-only fake framework, MIT, whose return sequences are unit level
fault injection.\newline
\url{https://github.com/meekrosoft/fff}
\end{itemize}

One gap is worth stating at the end of the volume as clearly as it was stated at
the start. \textbf{No published, motion-specific example was found of a rig that
runs in continuous integration and asserts on tracking quality.} The generic
material is plentiful and the motion-specific material is absent. A reader who
builds this chapter and publishes the result will have written something that
does not currently exist.
   % The rig: injected faults, tracking error, and a build that fails

\appendix
\clearpage
\section{The twenty chapters at a glance}

One page, so that a chapter can be chosen by what it gives rather than by its
title. The last column is the one this volume cares about most.

\begin{table}[H]\centering\small
\begin{tabular}{p{5mm}p{54mm}p{30mm}p{8mm}p{10mm}p{34mm}}
\toprule
& Title & What the node gains & Diff. & Eve\-nings & Real, modelled or absent \\
\midrule
1 & What a joint node is, and the bench that stands in for one & An identity and a bring-up & 2 & 2 & Real, and it names all three absences \\
2 & The control period: 1 kHz you can prove & A heartbeat & 3 & 2 & Real \\
3 & One clock for sensors, loop and bus & A shared time base & 4 & 3 & Real \\
4 & The board support package & A board file you can hand over & 3 & 3 & Real, and it proves two absences \\
5 & The encoder & Position, exactly & 3 & 2 & Real, with a generated source \\
6 & The inertial unit & An inner ear & 3 & 2 & Real \\
7 & The actuator you do not have & A drive stage, in firmware & 4 & 3 & \textbf{Real firmware, modelled plant} \\
8 & Force and torque & An estimate, honestly labelled & 3 & 2 & \textbf{Modelled, and labelled everywhere} \\
9 & CAN-FD from the controller out & A wire & 4 & 3 & Real \\
10 & The motion master & Somebody to talk to & 3 & 2 & Real \\
11 & A joint protocol & A vocabulary & 4 & 3 & Real \\
12 & Network management & A way to notice & 4 & 3 & Real \\
13 & Two speeds on one wire & A diagnosis & 3 & 2 & Real, with a part bought to fail \\
14 & The node as a middleware participant & A place in the robot & 5 & 4 & Real software, absent robot \\
15 & Joint state and joint command as messages & A standard shape & 3 & 2 & Real, one field left empty \\
16 & The loop closed over the bus & A closed loop & 5 & 4 & \textbf{Real controller, modelled plant} \\
17 & What a real-time fieldbus would change & An honest boundary & 2 to 5 & 2 & \textbf{Neither: nothing is built} \\
18 & Safe states, and the workspace sensor & A way to stop & 4 & 3 & Real trigger, modelled reaction \\
19 & Update over the bus & Maintainability & 5 & 4 & Real \\
20 & The rig & Proof & 5 & 5 & Both, deliberately \\
\bottomrule
\end{tabular}
\caption{The twenty chapters, their cost and their honesty class. Six of the twenty carry something modelled or absent, and every one of those six says so in its key facts, in its figures and in its budget table. Roughly sixty evenings in total, and two chapters need no hardware at all.}
\end{table}

\section{The honesty ledger}

Everything in this volume that is not what it appears to be, in one table, so
that a reader never has to hunt for the caveat.

\begin{table}[H]\centering\small
\begin{tabular}{p{34mm}p{22mm}p{46mm}p{36mm}}
\toprule
Thing & Class & Why & Where it is handled \\
\midrule
A brushless motor and drive stage & \textbf{Absent} & No motor, no driver board, no current sensing on this bench & Chapter 7 builds the firmware in full and drives a model \\
A force or torque sensor & \textbf{Absent} & Not on the bench, and the cheap substitutes measure something else & Chapter 8 publishes a derived estimate with a seven term error budget \\
A real-time Ethernet fieldbus & \textbf{Structurally impossible} & This part has no Ethernet controller, and this silicon vendor sells no subdevice controller at all & Chapter 17 costs it out and refuses it \\
A brake & \textbf{Absent} & And so the second of two safe states is designed and never felt & Chapter 18, in logic only \\
An external watchdog & \textbf{Absent} & Both on-die watchdogs share the chip's fate & Chapter 18 says so rather than claiming independence \\
A second channel, and an assessment & \textbf{Absent} & Single channel reaches the lower integrity levels by architecture & Chapter 18's page titled what this node is not \\
An oscilloscope, a logic analyser, an external probe & \textbf{Absent since chapter 1} & Which is why every timing claim uses the part's own counters & Chapters 2 and 20 \\
A soldering iron & \textbf{Absent} & Which rules out three otherwise reasonable routes & Chapters 13, 17 and 19 \\
The plant in the control loop & \textbf{Modelled} & The loop closes through one dashed block and six solid ones & Chapter 16, on every figure \\
Position feedback in chapter 1's table & \textbf{Corrected in chapter 5} & Chapter 1 claimed it was real before the encoder existed & The correction is printed in chapter 5 rather than edited away \\
The robot the middleware would join & \textbf{Absent} & No framework runs on this bench; only the agent is proven & Chapter 14, drawn dotted \\
\bottomrule
\end{tabular}
\caption{The honesty ledger. The tenth row is the one worth pointing at in an interview: an internal inconsistency was found while writing chapter 5, and it was corrected in the open, in the chapter that found it, rather than repaired quietly in chapter 1.}
\end{table}

\section{Corrections carried into this volume}

Claims in wide circulation that the research behind this volume found to be
wrong. Each is printed as a correction in the chapter that meets it, and each is
here in one place because a reader is more likely to meet the claim than the
chapter.

\subsection*{About this part}

\begin{enumerate}[itemsep=2pt]
\item \textbf{The STM32H7A3ZI has no Ethernet controller}, proven from interrupt
vector slots 61 and 62 being absent where its popular sibling has them, rather
than from a product page.
\item \textbf{One upstream device tree leaves a network controller node in
place} as an inherited modelling artefact, so reading that file gives the
opposite answer to grepping the vendor header. Both were read.
\item \textbf{This part's flash geometry is not its sibling's}: a sixteen byte
word and an eight kilobyte sector, against thirty-two and one hundred and
twenty-eight. \textbf{The comment block in the vendor's own flash driver source
describing the wider word is written for the sibling and is wrong here.}
\item \textbf{The backup registers live in the tamper peripheral on this part},
not in the real-time clock. Code ported from the sibling will not compile or
will not work.
\item \textbf{The brown-out reset has four levels, not the three usually
quoted}, and the programmable voltage detector has seven thresholds plus an
eighth selection that compares an external analog input against the reference.
\item \textbf{The two voltage detectors share one interrupt line}, so a handler
must read both status bits to tell them apart.
\item \textbf{The ROM bootloader already speaks the flexible-data bus}, on two
specific pins, and its detection chain polls that bus first.
\item \textbf{This family has no hardware encoder index}, on the vendor's own
statement. A later family was the first to have one.
\item \textbf{The part trap does not apply to the bus peripheral}: same core,
same addresses, same bit timing fields, with three real differences named in
chapter 9.
\end{enumerate}

\subsection*{About libraries, stacks and packages}

\begin{enumerate}[resume,itemsep=2pt]
\item \textbf{The vendor's EEPROM emulation package does not cover this
family.} The string identifying it does not appear anywhere in the current
revision of the relevant application note.
\item \textbf{That note was renamed.} Its current title is "How to use EEPROM
emulation on STM32 MCUs", revision 11, Tuesday 18 March 2025. A second note
frequently cited alongside it is about dual bank firmware update, not emulation.
\item \textbf{The motor control package is under the vendor's own licence}, not
the permissive one usually assumed.
\item \textbf{The widely repeated middleware footprint of 32 kB of memory and
256 kB of flash comes from a commercial blog, not from the project.} The
project's own published figure covers the constrained-environments client alone
and is much narrower in scope.
\item \textbf{That middleware has no packaged binary for any current
distribution}, and its original documentation host now redirects, with its deep
links returning not-found.
\item \textbf{The portable open master stack moved to stronger copyleft on
Friday 11 July 2025}, not away from it. Automated licence checks report nothing
because its licence file has an unexpected extension.
\item \textbf{One popular bootloader is GPL-3.0, not the earlier version usually
quoted.}
\item \textbf{The claim that a portable open bootloader cannot run on this
family because of its write block size is contradicted by its current source},
whose assertion covers this part's sixteen byte word.
\item \textbf{One widely used controller library is permissively licensed only
by a sentence in a readme file}, with no licence file, no identifier line and no
metadata field, so scanners report none. Its other readme returns not-found.
\item \textbf{The filter library often recommended alongside it contains no
controller at all.}
\item \textbf{One bus tooling project was renamed and then archived by its owner
on Friday 7 May 2021}; the claim that a particular successor superseded it could
not be confirmed.
\item \textbf{The download object is named "Program data"}, not the longer form
usually written down.
\item \textbf{A separate vendor slave stack is gated behind membership, not
behind a vendor identifier}, and the identifier itself is free of charge. The
real cost is a mandatory recurring conformance test tool subscription.
\item \textbf{There is no applications processor in that family with an
integrated subdevice controller}; the integrated part is a different device with
a different core inside it.
\end{enumerate}

\subsection*{About standards and safety}

\begin{enumerate}[resume,itemsep=2pt]
\item \textbf{The industrial robot safety series is now titled "Robotics"},
not the older wording, and getting it wrong dates a document at a glance.
\item \textbf{The collaborative robot technical specification is not
withdrawn.} It is published, was confirmed in 2022, and remains current until
its replacement appears. Several secondary sources say otherwise and they are
wrong.
\item \textbf{The machinery functional safety standard's 2021 edition lost its
restriction to electrical technologies rather than gaining non-electrical
coverage}, and both it and the electrical equipment standard carry amendments
that a bare citation omits.
\item \textbf{The current coding guidelines edition is dated \pubdate{March 2025}}, and
the 2023 edition consolidated the 2012 one with its amendments, so the old form
of citation is out of date.
\item \textbf{One silicon vendor's safety architecture claim is routinely quoted
with its redundancy scheme changed}, which reverses the meaning of its second
half.
\item \textbf{A certified kernel product is not a certified edition of the free
one.} They share a functional foundation and are separate, redesigned products.
\item \textbf{Both titles in circulation for the two firmware update
architecture documents are wrong}; the verified forms are in chapter 19.
\item \textbf{One timestamping standard's 2020 edition is superseded by its 2025
one}, and a bus specification number frequently cited for the same subject does
not exist.
\item \textbf{One data link standard's current edition carries a different
subtitle from its predecessor}, so the edition is part of the citation rather
than an optional extra.
\end{enumerate}

\section{Gaps this volume can claim}

Stated as new work in the chapter that looked for them, never dressed up as a
survey. Each was searched for deliberately, and each absence is a finding.

\begin{enumerate}[itemsep=3pt]
\item \textbf{No published source states that a safe state must be reachable
from every other state}, and none treats fault latching as a named technique.
Both are real and universally practised. \emph{Chapter 18 presents them as
engineering practice and attributes them to nobody.} The same absence covers
heartbeat and liveness over a bus, fail-silent against fail-operational, and
redundant state variables.
\item \textbf{No published, motion-specific example runs a rig in continuous
integration and asserts on tracking quality.} The generic material is plentiful.
\emph{Chapter 20, and a reader who publishes one will have written something
that does not exist.}
\item \textbf{No open software failure analysis tool is worth using.} The
candidates carry six or seven commits between them. \emph{Chapter 18 recommends
a space agency's freely usable worksheets and a spreadsheet under version
control.}
\item \textbf{No bus transport exists for the portable open bootloader or for
its management protocol}, whose transports are wireless and two kinds of serial.
\emph{Chapter 19 names the gap and writes into it.}
\item \textbf{The embedded middleware's bus transport is for a full operating
system only}, with no implementation for a small kernel anywhere. \emph{Chapter
14 writes the six functions it needs.}
\item \textbf{No published article presents the in-memory histogram method for
proving a control period with no probe.} \emph{Chapter 2 presents it as the
author's construction on two cited ideas.}
\item \textbf{No verified open embedded implementation in this language does
textbook back-calculation anti-windup.} Clamping is everywhere; conditional
integration is common. \emph{Chapter 16 writes it and measures it against
clamping on the same step.}
\item \textbf{No equivalent to the classic-frame safety communication layer was
found for the flexible-data format.} \emph{Chapter 18 names the gap and points
at the untrusted-transport framing as the idea to verify.}
\end{enumerate}

\section{Open questions for the bench}

Nothing on this list is written as fact anywhere in the volume. Each is a
question for a session with the reference manual, the board manual and the
datasheet open, on a machine that can reach the vendor's site.

\begin{itemize}[itemsep=2pt]
\item \textbf{The datasheet document number for this part.} One number in
circulation attaches to two different parts; the candidate is unconfirmed.
\item \textbf{Flash endurance cycles and retention years.} Chapter 19's rotating
log argument is parameterised so that it holds whatever the figure turns out to
be.
\item \textbf{Flash word program time and eight kilobyte sector erase time.} The
conclusion that an erase never fits in the power-down path does not depend on
the exact values.
\item \textbf{Backup memory size, address and retention conditions}, checked
against the reference manual rather than only against the headers.
\item \textbf{The timestamp counter clause numbers, and the dead-time generator
encoding.}
\item \textbf{The bus controller's message memory on this part}: how much, how
divided between filters and buffers, and whether it must be configured before
use. This differs across the family and is a common cause of a controller that
never transmits.
\item \textbf{Whether a bus buffer placed in tightly coupled memory is reachable
by the peripheral at all.}
\item \textbf{Which timers are 32 bit}, and which offers complementary outputs
with dead time and a fault input, on which pins.
\item \textbf{Whether the two-times and four-times encoder mode mapping is as
standard semantics suggests.}
\item \textbf{The transceiver's data rate and suffix}, once bought.
\item \textbf{Whether the vendor's safety library covers this part at all, and
under what licence.} Chapter 18 uses nothing from it for exactly this reason.
\item \textbf{The untrusted-transport standard family's number, title and
edition}, for chapter 18's closing paragraph.
\item \textbf{Two paywalled clauses}: one stop-category clause and one set of
category and diagnostic coverage tables.
\item \textbf{The harmonised adoption date of the 2025 robot safety standards},
reported at secondary level only.
\end{itemize}

\section{Licences, decided once}

Four categories, applied consistently through the volume, and the reason every
chapter checks a licence before writing code around a library rather than after.

\begin{table}[H]\centering\small
\begin{tabular}{p{30mm}p{40mm}p{66mm}}
\toprule
Category & Rule & What is in it here \\
\midrule
\textbf{A. Depend and vendor} & Permissive: copy it in, keep the headers, record
the provenance & The simulator, the bus stacks, the bootloader that is portable,
the fault backtrace and logging libraries, the fake framework, the unit test
tools, the message definitions, the host control framework, the real-time
helpers, the trace recorder's target-side source \\
\textbf{B. Depend, do not vendor} & Copyleft in a way that reaches a linked
program, or a library boundary that matters & The host bus library, the bench
control layer, the fixed-point controller, the kernel-space fieldbus master \\
\textbf{C. Read, do not copy} & Useful to read, tied to one vendor's silicon, or
commercial beyond a point & The vendor middleware, the vendor safety library
whose coverage of this part is unconfirmed, the event framework that is copyleft
or commercial, the commercial trace visualiser \\
\textbf{D. No licence at all} & \textbf{All rights reserved.} Cite as reading,
never vendor & Several community repositories with no licence file, one dormant
since Monday 10 April 2017, and one bootloader whose licence path returns
not-found \\
\bottomrule
\end{tabular}
\caption{The four categories. Category D is the one that catches people, because a repository with no licence file is not permissive by default, it is closed by default, and several of the most convenient-looking results for these subjects are in it.}
\end{table}

Three positions are worth stating separately because they changed a decision in
this volume rather than merely being noted.

\begin{itemize}[itemsep=3pt]
\item \textbf{The most consequential single licence decision is in chapter 19.}
An excellent, maintained bootloader solves that chapter's problem completely on
this family and in both frame formats, and linking the joint node's firmware
against it would make the whole firmware copyleft. The chapter states this
before a reader writes any code around it.
\item \textbf{A permissive wrapper does not describe the stack it deploys.}
Chapter 17's fieldbus bridge is permissively licensed and installs on a
copyleft out-of-tree kernel module, with a procedure that requires disabling
secure boot. The check is one dependency level deeper than most people look.
\item \textbf{One vendor archive ships three different licence statements},
between its release notes, its source file headers and a licence agreement
document included in the same archive, the last of which forbids commercial
distribution. A licence is what the files and the archive say together.
\end{itemize}

\section{The reference library}

Everything cited anywhere in the volume, grouped by what it is for, with its
licence where it has one. Every address was resolved once more before delivery,
and this list is stamped with the date it was checked. A book that tells a
reader to clone repositories owes them working addresses.

\subsection*{The bus, the protocol and the host side}

\begin{itemize}[itemsep=2pt]
\item ISO 11898-1:2024, third edition, \pubdate{May 2024}, the data link standard.
\textbf{Its subtitle differs from the 2015 edition's, so the edition is part of
the citation.} Paywalled.
\item The command line bus tools, dual licensed per file, whose load generator
and replay utility carry chapters 10, 13 and 20.\newline
\url{https://github.com/linux-can/can-utils}
\item A permissively licensed bus stack implementing the safety data object wire
format, Apache-2.0.\newline
\url{https://github.com/CANopenNode/CANopenNode}
\item The maintained bus-connected motion stack, Apache-2.0, which is the master
side of the contract this volume's node sits opposite.\newline
\url{https://github.com/ros-industrial/ros2_canopen}
\item The raw frame wrapper and its message package, Apache-2.0, which defines a
flexible-data frame type.\newline
\url{https://github.com/autowarefoundation/ros2_socketcan}
\end{itemize}

\subsection*{Middleware, messages and control}

\begin{itemize}[itemsep=2pt]
\item The configurator integration utilities, Apache-2.0, and the client and
agent, Apache-2.0.\newline
\url{https://github.com/micro-ROS/micro_ros_stm32cubemx_utils}\newline
\url{https://github.com/eProsima/Micro-XRCE-DDS-Client}
\item The common interface definitions, Apache-2.0, and the control message
package, BSD-3-Clause, for the message types and the six error codes.\newline
\url{https://github.com/ros2/common_interfaces}
\item The host control framework, Apache-2.0, and its real-time helper library,
BSD-3-Clause.\newline
\url{https://github.com/ros-controls/ros2_control}
\item The best-fitting small controller, \textbf{MIT at repository level with no
header in either source file}.\newline
\url{https://github.com/pms67/PID}
\item Astrom and Murray, \textit{Feedback Systems}, chapter 10, section 10.4.
\textbf{The book is the citation; the hosting site serves a broken certificate
chain, so no link is given.}
\item Koren, "Cross-Coupled Biaxial Computer Control for Manufacturing Systems",
Journal of Dynamic Systems, Measurement, and Control volume 102, pages 265 to
272, 1980, for the distinction this volume relies on between following error and
contouring error.
\end{itemize}

\subsection*{Safety}

\begin{itemize}[itemsep=2pt]
\item IEC 61508-4:2010, clause 3.1.13 for the safe state definition and clause
3.2.11 for the tool classification, both read from the official preview.
\item ISO 13850:2015, clause 4.1.1, read directly, for the six requirements
chapter 18's latch is written against.
\item ISO 10218-1:2025 and ISO 10218-2:2025, \pubdate{February 2025}, which for a joint is
the type-C standard and therefore takes precedence.
\item The space agency's software engineering and assurance handbook, topic 8.5,
with four worksheet templates. \textbf{A United States government work with no
copyright notice, and the safest material in this volume to adapt and publish.}
\item The free performance level calculator, version 3.0.4 of Friday 31 October
2025. \textbf{Free to use and pass on; modification and re-hosting are not
permitted.}
\end{itemize}

\subsection*{Update, storage and the rig}

\begin{itemize}[itemsep=2pt]
\item RFC 9019, \pubdate{April 2021}, and RFC 9124, \pubdate{January 2022}, both Informational and
both free to redistribute, \textbf{which makes them the safest normative
references in the volume}.\newline
\url{https://www.rfc-editor.org/rfc/rfc9019}
\item The portable open bootloader, Apache-2.0, port agnostic by its own porting
guide and \textbf{without any bus transport}.\newline
\url{https://github.com/mcu-tools/mcuboot}
\item A power-loss-resilient filesystem, BSD-3-Clause, and a key-value store,
Apache-2.0.\newline
\url{https://github.com/littlefs-project/littlefs}
\item The open simulator, which models this family including its bus controller
and bridges both frame formats to a host virtual interface.\newline
\url{https://github.com/renode/renode}
\item A fault diagnosis library covering this core, MIT, needing no debugger, and
a small logger with a flash backend, MIT.\newline
\url{https://github.com/armink/CmBacktrace}\newline
\url{https://github.com/armink/EasyLogger}
\item A header-only fake framework, MIT, whose return sequences are unit level
fault injection.\newline
\url{https://github.com/meekrosoft/fff}
\item ISO 9283:1998, second edition, Thursday 23 April 1998, for the assertion
names chapter 20 borrows. \textbf{Its metric symbols could not be confirmed and
are not printed.}
\item Hsueh, Tsai and Iyer, "Fault injection techniques and tools", Computer
volume 30, pages 75 to 82, 1997, and Natella and co-authors, "On Fault
Representativeness of Software Fault Injection", IEEE Transactions on Software
Engineering volume 39, pages 80 to 96, 2013.
\end{itemize}

\subsection*{Link hygiene}

Every address above and every address in the twenty chapters' Sources sections
was resolved once more before delivery. Several domains refuse automated
checking, the silicon vendor's among them, so every document from those was
opened by hand in a browser. \textbf{Checked: Monday 21 September 2026.}

Four categories of dead address were found during the research and are recorded
here so that nobody spends an afternoon on them: four article slugs from one
embedded engineering site, every documentation address on one operating system's
retired estate, the deep links into one middleware project's original
documentation host, and one bootloader repository's licence path.

\end{document}
